\documentclass[10pt,a4paper]{article}
\usepackage[margin=22mm]{geometry}
\usepackage{fontspec}
\setmainfont{DejaVu Serif}
\setsansfont{DejaVu Sans}
\setmonofont[Scale=0.90]{DejaVu Sans Mono}
\usepackage{amsmath,amssymb,amsthm,mathtools}
\usepackage{booktabs,longtable,array,multirow}
\usepackage{enumitem}
\usepackage{etoolbox}
\usepackage{listings}
\lstset{basicstyle=\ttfamily\small,breaklines=true,columns=fullflexible,
keepspaces=true,showstringspaces=false,frame=none}
\usepackage{xcolor}
\usepackage{hyperref}
\usepackage{xurl}
\hypersetup{colorlinks=true,linkcolor=blue!45!black,urlcolor=blue!45!black}
\theoremstyle{definition}
\newtheorem{theorem}{Theorem}
\newtheorem{lemma}[theorem]{Lemma}
\providecommand{\tightlist}{\setlength{\itemsep}{0pt}\setlength{\parskip}{0pt}}
\newcommand{\passthrough}[1]{#1}
\BeforeBeginEnvironment{longtable}{\begingroup\footnotesize}
\AfterEndEnvironment{longtable}{\endgroup}
\setlength{\emergencystretch}{3em}
\setlength{\parskip}{0.3em}
\setlength{\parindent}{0pt}
\setcounter{tocdepth}{1}
\title{Subset sums and four summands\\A cardinality--85 exclusion and the upper bound 84}
\author{Proof supplement to the supplied upper--85 report}
\date{5 October 2026}
\begin{document}
\maketitle
\begingroup\footnotesize\tableofcontents\endgroup
\clearpage
\section{The new upper bound and its proof boundary}
For a finite real set $A$, let $FS(A)$ include the empty subset sum, and let
$4A$ allow repeated summands. We continue the supplied report on the condition
$FS(A)\subseteq4A$. Its established upper bound and its explicitly stated
general mathematical inputs are retained.

\begin{theorem}
Under the mathematical inputs of the supplied upper--85 report, every feasible
finite real set has at most $84$ elements. The supplied $71$--element
construction remains feasible. Consequently
\[
                       71\le N\le84.
\]
\end{theorem}

The previously established $N\le86$ theorem is an explicit input in the
original report. The supplied attachments do not contain its complete proof
or an identifiable citation for that global theorem. It is not reproved here;
the improved global conclusion retains that dependency. The new contribution
is the complete exclusion of cardinality $85$, in addition to the original
report's exclusion of cardinality $86$.

The new work excludes cardinality $85$. It does not assume that arbitrary
subsets of a feasible set are feasible. If a feasible $85$--element set omitted
zero, adjoining zero would produce a feasible $86$--element set, contrary to
the retained bound. Thus zero is present. The rationalization lemma preserves
both cardinality and this zero coordinate. The fresh structural argument below
then reduces to an integral source $D$ and integral core $E$ with parameters
\[
 |D|=77,\quad |E|=69,\quad r_s+r_t=68,\quad
 I=(r_s-34)^2,\quad \Lambda=1191+I,\quad\Omega=600-I.
\]
The complete $85$--element argument is given in the following supplements.
All references to original covering, peeling, defect and inverse lemmas are to
the supplied report and its bundled mathematical sources; they are not new
claims certified merely by a program run.

\subsection{What makes the improvement possible}
There are three changes. First, all structural boundary cases are recomputed
at the new cardinality, including the $13:55$ and $14:54$ boundaries. Second,
the original fixed-label gap inventory is sharpened to $11$, $8$, and $15$
labels for the first gap, beta-active second gap, and their union. Third, the
two source envelopes are compared using the actual distinct fringe sums.

For the four offsets $1\le a_v<b_v<c_v<e_v$ beyond each core maximum, let
$P_v$ count the distinct pair sums of $\{a_v,b_v,c_v\}$ above $c_v$, and
$B_v$ the distinct pair sums of $\{a_v,b_v,c_v,e_v\}$ above $e_v$. The
elementary inequality
\[
 P_v+B_v-a_v-b_v\le5
\]
improves the source mass parameter to
\[
 T=77j-280-w+I+\tau+\delta\le165,
 \qquad I+\tau+\delta\le445-77j+w.
\]
The one-maximum envelope also yields $e_v\ge\max(4,T-89)$ on each physical
sign. These are bounds on actual source sets, not formal endpoint descriptors.

\subsection{Exhaustive case assembly}
Write $z$ for the core hull-hole count and $j$ for the number of genuinely empty
ordinary whole units. Paired units retain a single capacity throughout. The
physical signs are never reordered after an owner is selected.

\begin{center}
\begin{tabular}{lll}
\toprule
Actual family & Indices & Supporting proof below\\
\midrule
Zero core holes & $Z_0,Z_1,Z_2$ & Adaptive forests and raw gaps\\
Zero core holes & $Z_3$ & Endpoint cards and physical side union\\
Zero core holes & $Z_4$ & Retention cards and ownerless certificate\\
Zero core holes & $Z_5$ & Generic rank cards and final tail arguments\\
Positive low holes & $P_1,P_5$ & Common localization and terminal/source bounds\\
Positive low holes & $P_2$ & Fresh activity and rooted forest proof\\
Positive low holes & $P_3$ & Dense seed, two actual forms, endpoint proofs\\
Positive low holes & $P_4$ & Fresh source profiles and owner certificates\\
High core holes & $H$, $z\ge15$ & Paired profile floors and full localization\\
\bottomrule
\end{tabular}
\end{center}
The no-empty-unit obstruction covers the zero-index positive-hole alternatives
arising in the structural reduction; the fresh no-hole reduction covers $Z_0$.
The scalar budgets exclude indices beyond those shown. These statements are
proved in the supplements, including the actual activity classifications.

\subsection{Verification and dependencies}
The new verifier checks explicit finite arithmetic, profile minima, the dense
seed lemma, small offset inventories, and the final scalar certificates. Its
coverage manifest identifies every new written proof and every checked script.
The inherited verification program and mathematical sources are bundled intact
under \texttt{audit/verification\_program}. The new driver can also run their
full verification suite.

The verifier does not purport to formalize the infinite inverse theorems, the
passage from actual sets to forest inequalities, or the complete prose proof.
Those remain written mathematical arguments, with their hypotheses and
original input boundary retained. The package therefore supplies both the
proof and a reproducible check of its specified finite components.

\clearpage

\clearpage
\hypertarget{complete-size85-structure-integration-and-dependency-check}{%
\section{Complete SIZE85 structure: integration and dependency
check}\label{complete-size85-structure-integration-and-dependency-check}}

The fresh structural argument is complete relative to the general
mathematical inputs already stated and audited in the supplied upper85
foundation sources. It does not import a SIZE86 inverse row or finite
exception census. The two companion proofs are
\texttt{structure85\_\allowbreak{}\allowbreak{}ap\_\allowbreak{}\allowbreak{}bootstrap.md} and
\texttt{../\allowbreak{}inverse/\allowbreak{}proof85\_\allowbreak{}\allowbreak{}inverse.md}. The additional
strict-inverse a13 closure is now included in the latter and its exact
checker.

\hypertarget{entry-a-prescribed-zero-and-rationalization}{%
\subsection{Entry: a prescribed zero and
rationalization}\label{entry-a-prescribed-zero-and-rationalization}}

Assume the already established upper bound
\texttt{|A|<=\allowbreak{}85}. If a hypothetical 85-element feasible
real set did not contain zero,\allowbreak{} adjoining zero would keep every subset
sum feasible: its subset sums are unchanged and its fourfold sumset only
grows. This would give a feasible 86-element set,\allowbreak{} a contradiction. Thus
a hypothetical size85 set contains zero.

To pass from its existence to an integral feasible example,\allowbreak{} enumerate
the 85 coordinates and choose one fourfold representation of every
subset sum. These finitely many identities form a homogeneous linear
system with integer coefficients. Include the equation setting the
distinguished zero coordinate equal to zero. Its real kernel has a basis
over the rationals; rational points are dense in that kernel. Take a
rational kernel point sufficiently close to the original point to
preserve all strict coordinate orders and nonzero signs. The chosen
identities preserve feasibility,\allowbreak{} and the added zero equation preserves
the prescribed zero exactly. A positive common-denominator scale then
gives 85 distinct integers. This argument transfers existence; it does
not claim that the original real coordinates themselves lie on a
lattice.

\hypertarget{exact-resulting-theorem}{%
\subsection{Exact resulting theorem}\label{exact-resulting-theorem}}

After the outer four-on-each-sign deletion and two further
two-on-each-sign peels,\allowbreak{} a positive normalization gives integral source
sets of sizes

\[
 |D|=\allowbreak{}77,\allowbreak{}\quad |D_1|=\allowbreak{}73,\allowbreak{}\quad
 E=\allowbreak{}P\cup\{0\}\cup(-Q),\allowbreak{}\quad |E|=\allowbreak{}69,\allowbreak{}
 \quad a+\allowbreak{}b=\allowbreak{}68,\allowbreak{}\quad15\le a,\allowbreak{}b\le53.
\]

Both positive core sides are individually primitive. With
\texttt{T\_\allowbreak{}\allowbreak{}r=\allowbreak{}r(r+\allowbreak{}1)/\allowbreak{}2},\allowbreak{}
\texttt{k=\allowbreak{}34-min(a,\allowbreak{}b)} and
\texttt{I=\allowbreak{}k\^2},\allowbreak{}

\[
 \Lambda=\allowbreak{}1191+\allowbreak{}I,\allowbreak{}\quad \Omega=\allowbreak{}600-I,\allowbreak{}
 \quad0\le k\le19.
\]

On this common integer lattice the hull-hole counts satisfy

\[
 h_P,\allowbreak{}h_Q\le\kappa,\allowbreak{}\quad z=\allowbreak{}h_P+\allowbreak{}h_Q-\kappa,\allowbreak{}
 \quad0\le z\le\min(h_P,\allowbreak{}h_Q)\le\kappa,\allowbreak{}
 \quad0\le w\le4z.
\]

The actual extreme subset-sum translates in both peels intersect. If
\texttt{calT=\allowbreak{}kappa+\allowbreak{}ell1+\allowbreak{}ell2+\allowbreak{}g},\allowbreak{} then

\[
 c+\allowbreak{}\rho+\allowbreak{}\mathcal T\le600-I,\allowbreak{}
 \quad\mathcal T\ge\kappa+\allowbreak{}144-3z,\allowbreak{}
 \quad\mathcal T\ge\max(48,\allowbreak{}z,\allowbreak{}144-2z).
\]

Each \texttt{rho\_\allowbreak{}\allowbreak{}X>=\allowbreak{}0}. A non-AP side support obeys
\texttt{mu\_\allowbreak{}\allowbreak{}X>=\allowbreak{}r-3} and
\texttt{rho\_\allowbreak{}\allowbreak{}X>=\allowbreak{}r-14} for
\texttt{15<=\allowbreak{}r<=\allowbreak{}20},\allowbreak{} or
\texttt{rho\_\allowbreak{}\allowbreak{}X>=\allowbreak{}r-5} for
\texttt{r>=\allowbreak{}21}. Maximum removal gives

\[
 R_X\le r+\allowbreak{}\rho_X+\allowbreak{}h_X+\allowbreak{}2,\allowbreak{}
 \qquad R_P+\allowbreak{}R_Q\le72+\allowbreak{}\rho+\allowbreak{}\kappa+\allowbreak{}z.
\]

In particular,\allowbreak{} if \texttt{z>0},\allowbreak{} both supports are
non-AP and

\[
 \mu\ge62,\allowbreak{}\quad\tau=\allowbreak{}\mu+\allowbreak{}h_P+\allowbreak{}h_Q\ge62+\allowbreak{}2z,\allowbreak{}
 \quad\rho\ge\begin{cases}58,\allowbreak{}&\min(a,\allowbreak{}b)\ge21,\allowbreak{}\\49,\allowbreak{}&15\le\min(a,\allowbreak{}b)\le20.\end{cases}
\]

No outer-integrality or leading-activity assertion is part of this
theorem.

\hypertarget{exhaustive-branch-checklist}{%
\subsection{Exhaustive branch
checklist}\label{exhaustive-branch-checklist}}

\begin{longtable}[]{@{}
  >{\raggedright\arraybackslash}p{(\columnwidth - 4\tabcolsep) * \real{0.3333}}
  >{\raggedright\arraybackslash}p{(\columnwidth - 4\tabcolsep) * \real{0.3333}}
  >{\raggedright\arraybackslash}p{(\columnwidth - 4\tabcolsep) * \real{0.3333}}@{}}
\toprule\noalign{}
\begin{minipage}[b]{\linewidth}\raggedright
Core branch
\end{minipage} & \begin{minipage}[b]{\linewidth}\raggedright
Alignment and pre-lattice argument
\end{minipage} & \begin{minipage}[b]{\linewidth}\raggedright
Boundary or primitivity closure
\end{minipage} \\
\midrule\noalign{}
\endhead
\bottomrule\noalign{}
\endlastfoot
Initial deficient sign & Provisional available-sign peel; use
\texttt{24*85+\allowbreak{}77=\allowbreak{}2117} before fixing
\texttt{d=\allowbreak{}77} & One-sided 70-point lower count
\texttt{2416} excludes it \\
Sign disappearing in later peel & Fixed \texttt{d=\allowbreak{}77},\allowbreak{}
one-sided 69-point core & Lower \texttt{2347} exceeds
upper \texttt{1925} \\
\texttt{a=\allowbreak{}10} & Larger support must be AP,\allowbreak{} residue
count and strict layer rows & \texttt{G144} contradicts
budget \texttt{Omega+\allowbreak{}8=\allowbreak{}32} \\
\texttt{a=\allowbreak{}11} & AP branch or strict simultaneous
inverse,\allowbreak{} with small-side \texttt{+\allowbreak{}16} correction &
\texttt{G144},\allowbreak{} or
\texttt{T48+\allowbreak{}rho52-8=\allowbreak{}92>71} \\
\texttt{a=\allowbreak{}12} & AP branch or strict inverse,\allowbreak{} with
\texttt{+\allowbreak{}16} correction &
\texttt{z>=\allowbreak{}36}; effective
\texttt{(56,\allowbreak{}5,\allowbreak{}132)} forces
\texttt{rhoQ125>124} \\
Larger support AP,\allowbreak{} \texttt{a>=\allowbreak{}13} & AP-plus-set
formula,\allowbreak{} one residue,\allowbreak{} unclipped gaps,\allowbreak{} \texttt{z=\allowbreak{}0};
layers then \texttt{G144} & Any nonprimitive smaller
side has \texttt{c<=\allowbreak{}59<77},\allowbreak{}
\texttt{w=\allowbreak{}0}; \texttt{j0} obstruction;
boundaries \texttt{a13,\allowbreak{}a14} also have
\texttt{c<77} \\
Only smaller support AP,\allowbreak{} one residue,\allowbreak{} clipped gap or unique off-scale
coefficient & Final disjoint supports cost
\texttt{T\_\allowbreak{}\allowbreak{}a+\allowbreak{}T\_\allowbreak{}\allowbreak{}(b-2)+\allowbreak{}2-2b+\allowbreak{}rhoP} & Exact budget
contradiction throughout \\
Only smaller support AP,\allowbreak{} one residue,\allowbreak{} no gap &
\texttt{z=\allowbreak{}0},\allowbreak{} \texttt{hQ=\allowbreak{}kappa}; both
primitive & \texttt{a13}:
\texttt{192>159}; \texttt{a14}:
\texttt{c<=\allowbreak{}7<77,\allowbreak{}w=\allowbreak{}0},\allowbreak{} \texttt{j0}
obstruction \\
Smaller AP,\allowbreak{} two residues & Order-two residue,\allowbreak{} exact parity-fiber
identity before hole bound; odd coefficient
\texttt{W<=\allowbreak{}15}; strict AP-specific layers &
\texttt{G144} gives
\texttt{c<=\allowbreak{}61<77,\allowbreak{}w<=\allowbreak{}56}; \texttt{j0}
obstruction \\
Both supports non-AP,\allowbreak{} sufficient inverse test fails,\allowbreak{}
\texttt{k3..20} & Fresh 18 counted-endpoint repairs,\allowbreak{}
including strengthened non-AP prefix bounds; short-diameter theorem &
\texttt{z<=\allowbreak{}16}; strict layers before
\texttt{G144}; nonprimitive \texttt{P}
gives \texttt{c<=\allowbreak{}28<77,\allowbreak{}w<=\allowbreak{}64},\allowbreak{} hence
\texttt{j0} \\
Both supports non-AP,\allowbreak{} failed-test \texttt{a13} &
Intrinsic primitive Q effective \texttt{(55,\allowbreak{}2,\allowbreak{}52)}
gives \texttt{hQ<=\allowbreak{}4}; residue and short diameter &
\texttt{G144}:
\texttt{T+\allowbreak{}rho>=\allowbreak{}185>159} \\
Both supports non-AP,\allowbreak{} sufficient inverse test holds & Strict common
lattice and hole bounds; all32 large-hole final-overlap orientations
repaired before \texttt{G144} & Strict hole bounds
imply individual primitivity \\
Both supports non-AP,\allowbreak{} strict-test \texttt{a13} &
\texttt{kappa<=\allowbreak{}110}; corrected layer bound
\texttt{-811<0}; then \texttt{G144},\allowbreak{}
\texttt{z>=\allowbreak{}17} & Effective
\texttt{(55,\allowbreak{}4,\allowbreak{}119)} gives
\texttt{rhoQ>=\allowbreak{}113};
\texttt{48-1+\allowbreak{}113=\allowbreak{}160>159} \\
Both supports non-AP,\allowbreak{} repaired \texttt{a14} & Endpoint
\texttt{H<=\allowbreak{}8}; corrected strict pre-layer &
\texttt{c<=\allowbreak{}23<77,\allowbreak{}w<=\allowbreak{}32},\allowbreak{}
\texttt{j0} \\
Both supports non-AP,\allowbreak{} strict-test \texttt{a14} &
Corrected strict pre-layer; \texttt{G144} & If
\texttt{z>=\allowbreak{}13},\allowbreak{} effective
\texttt{(54,\allowbreak{}6,\allowbreak{}168)} gives
\texttt{rhoQ153}; otherwise
\texttt{c<=\allowbreak{}31,\allowbreak{}w<=\allowbreak{}48}; both impossible \\
\end{longtable}

Reflection is used only to order the scalar core sizes; both physical
sign orientations are included in every actual-owner statement. AP
support type,\allowbreak{} actual strict-test success/failure,\allowbreak{} and the boundary rows
partition the entire bootstrap range \texttt{a>=\allowbreak{}10}; no
actual test is replaced by its sufficient floor.

\hypertarget{dependency-ordering-no-use-of-integer-holes-before-its-lattice}{%
\subsection{Dependency ordering: no use of integer holes before its
lattice}\label{dependency-ordering-no-use-of-integer-holes-before-its-lattice}}

\begin{enumerate}
\def\labelenumi{\arabic{enumi}.}
\tightlist
\item
  The covering,\allowbreak{} peeling,\allowbreak{} \texttt{mu/\allowbreak{}rho} estimates,\allowbreak{}
  exact budget,\allowbreak{} inverse tests,\allowbreak{} residue counts,\allowbreak{} and counted-endpoint
  repairs are cardinality statements and valid before common-source
  integrality.
\item
  The short-diameter theorem is used on the proved common core lattice,\allowbreak{}
  and does not require unproved individual primitivity of the other
  side.
\item
  The pre-G comparisons use \texttt{C+\allowbreak{}Omega=\allowbreak{}-165-2I},\allowbreak{}
  plus \texttt{16} exactly when the small-side
  maximum-deletion prefix has fewer than14 elements. Large-hole overlap
  is supplied independently of \texttt{G144}.
\item
  A strict comparison proves lattice alignment of the next layer and the
  intersection of its two extreme support translates. This is applied to
  both layers before counting their integer restricted-sum hull holes.
\item
  Only then use the sharp integral four-extreme inequality
  \texttt{G=\allowbreak{}144},\allowbreak{} derive
  \texttt{w<=\allowbreak{}4z} and the structural T bounds,\allowbreak{} and
  invoke the conditional \texttt{j0} obstruction where
  the already proved caps supply \texttt{c<77,\allowbreak{}w<=\allowbreak{}104}.
\item
  Individual primitivity follows after these steps in the repaired
  branches; it is available immediately from strict hole bounds in
  strict-inverse branches. Thus neither residue repair nor
  \texttt{G144} assumes primitivity of P.
\end{enumerate}

The remaining retained inputs are the general inverse/stability/defect,\allowbreak{}
restricted-sum and cover lemmas,\allowbreak{} not fresh finite cardinality-specific
rows. The exact checkers \texttt{ap\_\allowbreak{}\allowbreak{}bootstrap85.py}
and \texttt{../\allowbreak{}inverse/\allowbreak{}check85\_\allowbreak{}\allowbreak{}inverse.py} both pass;
their arithmetic checks supplement the actual-set applicability
arguments in the written companion proofs.

\clearpage
\hypertarget{fresh-structural-deductions-at-cardinality-85}{%
\section{Fresh structural deductions at cardinality
85}\label{fresh-structural-deductions-at-cardinality-85}}

These deductions concern a hypothetical zero-containing \textbf{integral
feasible} set A of size 85. Zero follows only after the preceding
exclusion of all sizes at least 86 has been accepted: a zero-free
size-85 example could otherwise be extended by adjoining zero. The
general inverse/stability/defect theorems are the imported mathematical
inputs identified in the audited upper-85 package. This is an
intermediate proof supplement,\allowbreak{} not yet the complete STRUCTURE85 theorem;
the both-non-AP inverse and overlap repairs are supplied separately.

\hypertarget{bootstrap-derived-from-the-new-cardinality}{%
\subsection{1. Bootstrap,\allowbreak{} derived from the new
cardinality}\label{bootstrap-derived-from-the-new-cardinality}}

Four-on-each-sign outer deletion and two four-element peels give d=\allowbreak{}77,\allowbreak{}
e=\allowbreak{}69,\allowbreak{} m=\allowbreak{}68. If an outer sign had fewer than four elements,\allowbreak{} provisional
peeling from the available sign gives a one-sided core with at least 70
points,\allowbreak{} whose subset-sum count is at least 1+\allowbreak{}T\_\allowbreak{}69=\allowbreak{}2416. Before the
outer deletion size has been fixed,\allowbreak{} use the safe upper bound
s\textless=\allowbreak{}24\textbar D\textbar+\allowbreak{}77\textless=\allowbreak{}24*85+\allowbreak{}77=\allowbreak{}2117,\allowbreak{} not a
premature substitution \textbar D\textbar=\allowbreak{}77. If a sign disappears under
the further peels,\allowbreak{} a one-sided 69-point core gives
s\textgreater=\allowbreak{}1+\allowbreak{}T\_\allowbreak{}68=\allowbreak{}2347,\allowbreak{} whereas s\textless=\allowbreak{}24*77+\allowbreak{}77=\allowbreak{}1925. Thus both
physical signs actually support all the required extreme deletions.

The unpaired-absolute-value estimate gives
q\^{}2\textless=\allowbreak{}96e+\allowbreak{}1073-e\^{}2=\allowbreak{}2936. Here q is even,\allowbreak{} so q\textless=\allowbreak{}54;
at least (68-54)/2=\allowbreak{}7 opposite pairs remain in E. In particular zero is
not a diagonal exception and both core sides have at least 7 elements.

Reflect only in scalar analysis and write a=\allowbreak{}34-k,\allowbreak{} b=\allowbreak{}34+\allowbreak{}k. Then
Lambda=\allowbreak{}1+\allowbreak{}T\_\allowbreak{}a+\allowbreak{}T\_\allowbreak{}b=\allowbreak{}1191+\allowbreak{}k\^{}2 and Omega=\allowbreak{}22*77+\allowbreak{}97-Lambda=\allowbreak{}600-k\^{}2. The
larger-side adjusted defect is nonnegative; the smaller-side uniform
bound is rho\_\allowbreak{}P\textgreater=\allowbreak{}-8. The exact slack identity therefore
implies k\^{}2\textless=\allowbreak{}608,\allowbreak{} so k\textless=\allowbreak{}24 and a\textgreater=\allowbreak{}10. The
initial small rows are:

\begin{longtable}[]{@{}lllll@{}}
\toprule\noalign{}
a & b & k & Omega & Omega+\allowbreak{}8 \\
\midrule\noalign{}
\endhead
\bottomrule\noalign{}
\endlastfoot
10 & 58 & 24 & 24 & 32 \\
11 & 57 & 23 & 71 & 79 \\
12 & 56 & 22 & 116 & 124 \\
13 & 55 & 21 & 159 & 167 \\
14 & 54 & 20 & 200 & 208 \\
\end{longtable}

At a=\allowbreak{}14 the threshold-14 theorem improves the last slack bound to
Omega=\allowbreak{}200. At a=\allowbreak{}13,\allowbreak{} the actual threshold-14 induction table gives
Phi\_\allowbreak{}13\textless=\allowbreak{}4,\allowbreak{} hence rho\_\allowbreak{}P\textgreater=\allowbreak{}-2; if its support is
non-AP,\allowbreak{} the general non-AP theorem gives the stronger
rho\_\allowbreak{}P\textgreater=\allowbreak{}-1.

\hypertarget{exact-layer-constants}{%
\subsection{2. Exact layer constants}\label{exact-layer-constants}}

On a common core lattice,\allowbreak{} maximum removal gives
R\_\allowbreak{}s+\allowbreak{}R\_\allowbreak{}t\textless=\allowbreak{}72+\allowbreak{}rho+\allowbreak{}kappa+\allowbreak{}z when both core sizes are at least 15.
A smaller side requires adding 8 to this maximum bound,\allowbreak{} or 16 to a
two-span comparison. The conservative pre-lattice D1 comparison has
constant

\begin{lstlisting}
C=2m+4e+14-Lambda=-765-k^2,
C+Omega=-165-2k^2.
\end{lstlisting}

Thus F1-2F0\textless=\allowbreak{}-165-2k\^{}2+\allowbreak{}2z,\allowbreak{} with an additional +\allowbreak{}16 if
a\textless=\allowbreak{}14. After the source layers and their actual overlaps have
been established,\allowbreak{} the sharp restricted-sum inequality is G=\allowbreak{}2d-10=\allowbreak{}144. It
gives

\begin{lstlisting}
T>=kappa+144-3z, T>=144-2z, T>=z, T>=48,
c+rho+T<=600-k^2, w<=4z.
\end{lstlisting}

The quantitative layer rows,\allowbreak{} using \textbar R(E)\textbar\textless=\allowbreak{}2R-1
instead of the conservative 2R+\allowbreak{}1 bound,\allowbreak{} are freshly derived as follows:

\begin{lstlisting}
F2-2F1=-293-2k^2-c-rho-mu-2kappa-2theta-2ell2-2g;
V0-C1+delta+2z<=-292-2k^2-77j-mu-rho;
F1-2F0<=-167-2k^2+rho-mu+2z-c-2theta-ell1-g;
U0-C0+delta-2z<=-166-2k^2-77j, U0-C0<=0.
\end{lstlisting}

The last two rows assume the maximum-removal bound with core sizes
\textgreater=\allowbreak{}15; they are not silently applied on small boundaries.

\hypertarget{boundaries-10-11-and-12}{%
\subsection{3. Boundaries 10,\allowbreak{} 11 and 12}\label{boundaries-10-11-and-12}}

At a=\allowbreak{}10 a non-AP larger Q would give
rho\textgreater=\allowbreak{}-8+\allowbreak{}(58-5)=\allowbreak{}45\textgreater Omega=\allowbreak{}24. Hence Q is AP. Its
intrinsic step forces P onto the same lattice by residue counting,\allowbreak{}
because kappa\textless=\allowbreak{}32\textless T\_\allowbreak{}58. The AP-plus-set formula
precludes a clipped gap and gives h\_\allowbreak{}Q=\allowbreak{}z=\allowbreak{}0,\allowbreak{}h\_\allowbreak{}P=\allowbreak{}kappa. The conservative
D1 comparison with +\allowbreak{}16 is strict,\allowbreak{} as is the outer comparison
Lambda\textgreater4d-10+\allowbreak{}(Omega+\allowbreak{}8). Consequently D is integral on this
lattice and G144 applies. But T\textgreater=\allowbreak{}144\textgreater Omega+\allowbreak{}8,\allowbreak{}
impossible.

At a=\allowbreak{}11 and non-AP Q,\allowbreak{} rho\_\allowbreak{}Q\textgreater=\allowbreak{}52,\allowbreak{} so
kappa\textless=\allowbreak{}71+\allowbreak{}8-52=\allowbreak{}27\textless T\_\allowbreak{}11-1=\allowbreak{}65. The strict simultaneous
inverse supplies the common lattice and both hole bounds. The layer
comparisons are strict even with z\textless=\allowbreak{}79 and +\allowbreak{}16; G144 is
therefore available. Then T+\allowbreak{}rho\textgreater=\allowbreak{}48+\allowbreak{}52-8=\allowbreak{}92\textgreater71. If
Q is AP,\allowbreak{} the preceding z=\allowbreak{}0 argument gives
T\textgreater=\allowbreak{}144\textgreater79.

At a=\allowbreak{}12 and non-AP Q,\allowbreak{} rho\_\allowbreak{}Q\textgreater=\allowbreak{}51 and
kappa\textless=\allowbreak{}116+\allowbreak{}8-51=\allowbreak{}73\textless T\_\allowbreak{}12-1=\allowbreak{}77. The strict inverse again
supplies the common lattice and strict Q upper-hole bound; both layer
comparisons are strict. From G144 and rho\textgreater=\allowbreak{}43 we obtain
144-2z+\allowbreak{}43\textless=\allowbreak{}116,\allowbreak{} hence z\textgreater=\allowbreak{}36. Therefore
h\_\allowbreak{}Q\textgreater=\allowbreak{}36. The parameterized effective hole witness
(r,\allowbreak{}beta,\allowbreak{}C)=\allowbreak{}(56,\allowbreak{}5,\allowbreak{}132) has

\begin{lstlisting}
K=17,D=149,E=150,M=20,J=20,
(r-J+1)beta-D=36>0,
T_(J-1)-1-D=40>0.
\end{lstlisting}

It gives mu\_\allowbreak{}Q\textgreater=\allowbreak{}189; the retained flexible defect inverse
gives rho\_\allowbreak{}Q\textgreater=\allowbreak{}125. But the budget gives
rho\_\allowbreak{}Q\textless=\allowbreak{}116+\allowbreak{}8=\allowbreak{}124. The AP-Q branch has z=\allowbreak{}0 and is already
excluded by T\textgreater=\allowbreak{}144\textgreater124. Thus every surviving core
has a\textgreater=\allowbreak{}13.

\hypertarget{a-general-no-empty-unit-obstruction}{%
\subsection{4. A general no-empty-unit
obstruction}\label{a-general-no-empty-unit-obstruction}}

This uses only the original 55-fixed-label raw gap bounds,\allowbreak{} not any newly
strengthened gap count. Once D is integral,\allowbreak{} suppose j=\allowbreak{}0 and
w\textless=\allowbreak{}104. All four own centers on both signs are active,\allowbreak{} so all
outer magnitudes are integral. The source maxima satisfy
M\textgreater=\allowbreak{}m+\allowbreak{}8=\allowbreak{}76; active C7 gives t\_\allowbreak{}v=\allowbreak{}d\_\allowbreak{}v-M\textgreater=\allowbreak{}5,\allowbreak{} hence
d\_\allowbreak{}v\textgreater=\allowbreak{}81. Choose a genuinely contributing beta sign s (the
beta pair supplies one in branch A,\allowbreak{} and both beta singletons contribute
in branch B). Then

\begin{lstlisting}
t_s<=(57+eta_s)/2<=(57+w)/2<=80.5,
t_t<=56+eta_t<=56+w<=160.
\end{lstlisting}

An active C4\_\allowbreak{}s forces t\_\allowbreak{}s\textgreater d\_\allowbreak{}t\textgreater=\allowbreak{}81. An active
C4\_\allowbreak{}t forces
t\_\allowbreak{}t\textgreater h\_\allowbreak{}s+\allowbreak{}d\_\allowbreak{}s\textgreater2d\_\allowbreak{}s\textgreater=\allowbreak{}162. Both
alternatives are impossible,\allowbreak{} although pair1 must contribute when j=\allowbreak{}0.
Thus every aligned configuration with j=\allowbreak{}0,\allowbreak{}w\textless=\allowbreak{}104 is excluded. In
particular this applies whenever c\textless77 and w\textless=\allowbreak{}104.

\hypertarget{complete-ap-branch-simplification-for-a13}{%
\subsection{5. Complete AP-branch simplification for
a\textgreater=\allowbreak{}13}\label{complete-ap-branch-simplification-for-a13}}

The larger-AP case forces one common residue and gives z=\allowbreak{}0 as before. If
its smaller side is nonprimitive,\allowbreak{} h\_\allowbreak{}P\textgreater=\allowbreak{}T\_\allowbreak{}a+\allowbreak{}mu\_\allowbreak{}P,\allowbreak{} and hence
c\textless=\allowbreak{}Omega-144-T\_\allowbreak{}a-rho\_\allowbreak{}P. The maximum cap over all
a\textgreater=\allowbreak{}13 is 59; the a=\allowbreak{}13 universal floor rho\_\allowbreak{}P\textgreater=\allowbreak{}-2
is included. Therefore c\textless77 and the no-empty-unit obstruction
excludes every such branch,\allowbreak{} without 152 auxiliary certificate systems.
The larger unit-step AP side is intrinsically primitive. A smaller
primitive side is the remaining legitimate z=\allowbreak{}0 structural alternative.
At a=\allowbreak{}13,\allowbreak{} Q AP also gives c\textless=\allowbreak{}159+\allowbreak{}2-144=\allowbreak{}17,\allowbreak{} and at a=\allowbreak{}14,\allowbreak{} Q AP
gives c\textless=\allowbreak{}200-144=\allowbreak{}56; both boundaries are excluded by the same
obstruction.

Suppose only the smaller support is AP. The residue-count inequality
kappa\textgreater=\allowbreak{} (\#residues-1)(u-1),\allowbreak{} with u\textgreater=\allowbreak{}T\_\allowbreak{}a+\allowbreak{}1,\allowbreak{} and
Omega\textless2T\_\allowbreak{}a over all the present rows permits at most two
residue classes. A one-residue clipped gap or a unique off-scale
generator forces final disjoint supports of Q and
mu\_\allowbreak{}Q\textgreater=\allowbreak{}T\_\allowbreak{}(b-2). The necessary budget cost

\begin{lstlisting}
T_a+T_(b-2)+2-2b+rho_P
\end{lstlisting}

strictly exceeds Omega in every row; the executable checks these
margins. If there is one residue and no clipped gap,\allowbreak{}
z=\allowbreak{}0,\allowbreak{}h\_\allowbreak{}Q=\allowbreak{}kappa,\allowbreak{}h\_\allowbreak{}P=\allowbreak{}0. The larger side is primitive because
h\_\allowbreak{}Q\textless=\allowbreak{}Omega-rho\textless T\_\allowbreak{}b\textless=\allowbreak{}\textbar FS(Q)\textbar-1,\allowbreak{}
and the smaller unit-step AP support is primitive. This is another
legitimate z=\allowbreak{}0 alternative. Its a=\allowbreak{}13 boundary is impossible because
T+\allowbreak{}rho\textgreater=\allowbreak{}144+\allowbreak{}50-2=\allowbreak{}192\textgreater159. At a=\allowbreak{}14,\allowbreak{}
c\textless=\allowbreak{}200-144-49=\allowbreak{}7\textless77 and w=\allowbreak{}0,\allowbreak{} so the no-empty-unit
obstruction excludes it.

In the remaining two-residue case,\allowbreak{} at least two nonzero-residue
generators force that residue to have order two. Normalize
U=\allowbreak{}2{[}0,\allowbreak{}u-1{]}. Let W be the least odd Q coefficient and t\textgreater=\allowbreak{}2
the number of odd Q generators. The parity-fiber identity,\allowbreak{} before using
a hole bound,\allowbreak{} is

\begin{lstlisting}
kappa=u-1+eta, h_P=u-1,
h_Q=kappa-u+W, z=W-1,
mu_Q+kappa=sum(Q)+u-W-T_b.
\end{lstlisting}

Distinct even and odd generators give

\begin{lstlisting}
sum(Q)>=(b-t)(b-t+1)+tW+t(t-1).
\end{lstlisting}

The exact finite minimization in ap\_\allowbreak{}bootstrap85.py proves that
W\textgreater=\allowbreak{}17 contradicts mu\_\allowbreak{}Q+\allowbreak{}kappa\textless=\allowbreak{}Omega+\allowbreak{}2b-2-rho\_\allowbreak{}P,\allowbreak{}
with minimum contradiction margin 22. It uses the floor
rho\_\allowbreak{}P\textgreater=\allowbreak{}0 for a\textgreater=\allowbreak{}14 and rho\_\allowbreak{}P\textgreater=\allowbreak{}-2 for
a=\allowbreak{}13. Since W is odd,\allowbreak{} W\textless=\allowbreak{}15,\allowbreak{} hence z\textless=\allowbreak{}14; also
W\textless u,\allowbreak{} so h\_\allowbreak{}P,\allowbreak{}h\_\allowbreak{}Q\textless=\allowbreak{}kappa.

For this half-step branch,\allowbreak{} distinct-generator bounds give

\begin{lstlisting}
R_P+R_Q<=m+a+mu+kappa-T_a+14.
\end{lstlisting}

The conservative D1 comparison therefore gives

\begin{lstlisting}
F1-2F0<=2367+2a-4T_a-2T_b<=-761<0.
\end{lstlisting}

The outer-layer strict comparison likewise holds. D is integral and its
extreme support translates intersect; w\textless=\allowbreak{}4z\textless=\allowbreak{}56 and G144
apply. Using rho\_\allowbreak{}Q\textgreater=\allowbreak{}b-5 then gives

\begin{lstlisting}
c<=Omega-T_a-144+42-(b-5)-rho_P<=61<77.
\end{lstlisting}

The no-empty-unit obstruction excludes every half-step branch. This
includes the smaller-AP a=\allowbreak{}13 boundary. The only remaining a=\allowbreak{}13 or a=\allowbreak{}14
possibilities have both supports non-AP; their inverse repairs are
treated separately.

\hypertarget{executable-scope}{%
\subsection{6. Executable scope}\label{executable-scope}}

ap\_\allowbreak{}bootstrap85.py checks every scalar row k=\allowbreak{}0,\allowbreak{}...,\allowbreak{}21,\allowbreak{} the exact
odd-coefficient mass minimum,\allowbreak{} half-step and nonprimitive caps,\allowbreak{}
lattice-comparison constants,\allowbreak{} and one-residue contradiction margins. It
supplies no computational proof of the general inverse theorem or the
actual-set parity-fiber identities. The result file is
ap\_\allowbreak{}bootstrap85\_\allowbreak{}results.json.

\clearpage
\hypertarget{fresh-size85-inverse-alignment-and-final-support-overlap}{%
\section{Fresh SIZE85 inverse alignment and final-support
overlap}\label{fresh-size85-inverse-alignment-and-final-support-overlap}}

This companion proves the inverse portion of a SIZE85 structural
argument. Its general mathematical inputs are the previously established
covering and peeling identities,\allowbreak{} prefix defect recurrence,\allowbreak{} counted-hole
lemma,\allowbreak{} effective hole lemma,\allowbreak{} flexible defect inverse,\allowbreak{} and
Stanchescu\textquotesingle s two distinct-set theorems. It additionally
invokes the SIZE85 no-empty-unit obstruction after lattice forcing and
G144; that obstruction must be proved in the common source. The checker
\texttt{check85\_\allowbreak{}\allowbreak{}inverse.py} checks only the explicit
finite arithmetic.

\hypertarget{parameters-and-exact-pre-lattice-budget}{%
\subsection{1. Parameters and exact pre-lattice
budget}\label{parameters-and-exact-pre-lattice-budget}}

Use m=\allowbreak{}68,\allowbreak{}d=\allowbreak{}77,\allowbreak{}e=\allowbreak{}69,\allowbreak{}G=\allowbreak{}144. After reflection a=\allowbreak{}34-k,\allowbreak{}b=\allowbreak{}34+\allowbreak{}k,\allowbreak{}I=\allowbreak{}k². Put
T\_\allowbreak{}r=\allowbreak{}r(r+\allowbreak{}1)/2,\allowbreak{} Lambda=\allowbreak{}1+\allowbreak{}T\_\allowbreak{}a+\allowbreak{}T\_\allowbreak{}b=\allowbreak{}1191+\allowbreak{}I,\allowbreak{} Omega=\allowbreak{}600-I,\allowbreak{} and

\begin{lstlisting}
C=2m+4e+14-Lambda=-765-I, C+Omega=-165-2I.
\end{lstlisting}

For U=\allowbreak{}FS(P),\allowbreak{}V=\allowbreak{}FS(Q),\allowbreak{} write u=\allowbreak{}1+\allowbreak{}T\_\allowbreak{}a+\allowbreak{}muP,\allowbreak{}v=\allowbreak{}1+\allowbreak{}T\_\allowbreak{}b+\allowbreak{}muQ and
kappa=\allowbreak{}\textbar U+\allowbreak{}V\textbar-u-v+\allowbreak{}1. The exact budget is

\begin{lstlisting}
rhoP+rhoQ+kappa+2theta+c+ell1+ell2+g=Omega.
\end{lstlisting}

The non-AP support floors used below are n(r)=\allowbreak{}r-5 for r\textgreater=\allowbreak{}21
and n(r)=\allowbreak{}r-14 for 14\textless=\allowbreak{}r\textless=\allowbreak{}20. At r=\allowbreak{}13,\allowbreak{} the sharper non-AP
floor is -1. All mu values are nonnegative,\allowbreak{} and non-AP
mu\_\allowbreak{}X\textgreater=\allowbreak{}r-3. Also rho\textless=\allowbreak{}mu. The main range is
a\textgreater=\allowbreak{}15,\allowbreak{} k0..19. The a14 and a13 boundaries are included and
closed explicitly below; smaller boundaries are treated in the
bootstrap.

The larger-side sufficient strict inverse condition is automatic because
Omega-n(a)-n(b)\textless T\_\allowbreak{}b+\allowbreak{}b-4=\allowbreak{}v\_\allowbreak{}min-2 for every k0..21. Do not
replace this with Omega\textless T\_\allowbreak{}b-1: that unnecessarily stronger
inequality is false at balanced k0,\allowbreak{} where Omega=\allowbreak{}600,\allowbreak{}T34-1=\allowbreak{}594.

The sufficient small-side condition is

\begin{lstlisting}
muP+rhoP+rhoQ>Omega-T_a+2.
\end{lstlisting}

Its floor margin is positive exactly at k0,\allowbreak{}1,\allowbreak{}2 and nonpositive at
k3..21. Whenever the actual condition fails,\allowbreak{} the following fresh
cardinality bounds hold:

\begin{lstlisting}
rhoQ<=R=Omega-T_a+2-(a-3)-n(a),
kappa<=Kmax=Omega-n(a)-n(b), u<=Umax=Kmax+3,
v>=Vmin=T_b+b-2, u>=Umin=T_a+a-2.
\end{lstlisting}

\hypertarget{two-counted-endpoint-repairs}{%
\subsection{2. Two counted-endpoint
repairs}\label{two-counted-endpoint-repairs}}

The original repair starts with rho14\textgreater=\allowbreak{}0. Given
rhoQ\textless=\allowbreak{}R,\allowbreak{} fix beta\textgreater=\allowbreak{}3 with
(beta-1)(b-14)\textgreater R. Some index j\textgreater=\allowbreak{}15 has
y\_\allowbreak{}j\textless=\allowbreak{}beta; take the last one. The recurrence
rho\_\allowbreak{}i\textgreater=\allowbreak{}rho\_\allowbreak{}(i-1)+\allowbreak{}y\_\allowbreak{}i-2 and nonnegative prefix defects give

\begin{lstlisting}
j>=j0=b-floor(R/(beta-1)).
\end{lstlisting}

The counted-hole lemma applies if j0\textgreater=\allowbreak{}2beta+\allowbreak{}5. It supplies,\allowbreak{}
on an intrinsic prefix lattice,\allowbreak{} a deleted interval with at most beta
lower holes in {[}1,\allowbreak{}2beta{]} and their reflections. That prefix is
primitive on its own lattice because j elements among j+\allowbreak{}beta positive
positions cannot all be divisible by 2.

Assume T\_\allowbreak{}j0-b-1\textgreater R. A first later coefficient off this
lattice has disjoint support translates and raises rho above R; the
subsequent increments,\allowbreak{} all at least beta+\allowbreak{}1,\allowbreak{} cannot undo this. Thus every
coefficient is on that lattice. Every later increment is at most R+\allowbreak{}2. If
a current m-prefix has sum L and the next coefficient x satisfies
x\textgreater=\allowbreak{}L-4beta,\allowbreak{} the new translate adds above L at least
min(x,\allowbreak{}L+\allowbreak{}1)-2beta\textgreater=\allowbreak{}L-6beta\textgreater=\allowbreak{}T\_\allowbreak{}m-6beta points. When
T\_\allowbreak{}m-6beta\textgreater m+\allowbreak{}R+\allowbreak{}3,\allowbreak{} this exceeds its allowed increment.
Therefore x\textless=\allowbreak{}L-4beta-1. Also
x\textgreater max(prefix)\textgreater=\allowbreak{}m\textgreater2beta. The central
intervals join,\allowbreak{} the lower holes remain unchanged,\allowbreak{} and reflection
supplies the upper holes. Induction gives final hQ\textless=\allowbreak{}H=\allowbreak{}2beta and
a central interval {[}2beta+\allowbreak{}1,\allowbreak{}L\_\allowbreak{}Q-2beta-1{]}.

The strengthened repair uses a non-AP prefix floor,\allowbreak{} which is essential
in several SIZE85 rows. Suppose the final Q support is non-AP,\allowbreak{}
rhoQ\textless=\allowbreak{}R,\allowbreak{} and q\textgreater=\allowbreak{}20 satisfies T\_\allowbreak{}q-b-1\textgreater R.
Then every prefix of size at least q is non-AP. Indeed,\allowbreak{} if a last
AP-support prefix had size l\textgreater=\allowbreak{}q,\allowbreak{} the next support must break
the progression by disjoint translates. At that step
y\textgreater=\allowbreak{}T\_\allowbreak{}l-l. Every later increment is positive by the
zero-increment equality theorem. The prefix recurrence,\allowbreak{} starting from
rho\_\allowbreak{}l\textgreater=\allowbreak{}0,\allowbreak{} gives

\begin{lstlisting}
rhoQ>=T_l-b-1>=T_q-b-1>R,
\end{lstlisting}

a contradiction. Hence rho\_\allowbreak{}q\textgreater=\allowbreak{}n(q),\allowbreak{} and
rho\_\allowbreak{}j\textgreater=\allowbreak{}j-5 at j\textgreater=\allowbreak{}max(q,\allowbreak{}21).

If n(q)+\allowbreak{}(beta-1)(b-q)\textgreater R,\allowbreak{} there is a last j\textgreater=\allowbreak{}q+\allowbreak{}1
with y\_\allowbreak{}j\textless=\allowbreak{}beta. The later increments yield

\begin{lstlisting}
R>=rho_j+(beta-1)(b-j)
  >=(j-5)+(beta-1)(b-j),
\end{lstlisting}

and thus

\begin{lstlisting}
j>=j0=max(q+1,21,ceil(((beta-1)b-R-5)/(beta-2))).
\end{lstlisting}

Now apply exactly the same counted-hole,\allowbreak{} off-scale,\allowbreak{} and bridge
arguments. This strengthens the index bound without imposing a new
hypothesis on Q.

\hypertarget{all-failed-inverse-rows-acquire-a-common-lattice}{%
\subsection{3. All failed-inverse rows acquire a common
lattice}\label{all-failed-inverse-rows-acquire-a-common-lattice}}

The explicit choices (k,\allowbreak{}b,\allowbreak{}R,\allowbreak{}beta,\allowbreak{}q,\allowbreak{}j0,\allowbreak{}H) are:

\begin{longtable}[]{@{}rrrrrrr@{}}
\toprule\noalign{}
k & b & R & beta & q & j0 & H \\
\midrule\noalign{}
\endhead
\bottomrule\noalign{}
\endlastfoot
3 & 37 & 43 & 3 & 14 & 16 & 6 \\
4 & 38 & 69 & 4 & 14 & 15 & 8 \\
5 & 39 & 92 & 6 & 14 & 21 & 12 \\
6 & 40 & 112 & 6 & 14 & 18 & 12 \\
7 & 41 & 129 & 7 & 14 & 20 & 14 \\
8 & 42 & 143 & 8 & 14 & 22 & 16 \\
9 & 43 & 154 & 8 & 14 & 21 & 16 \\
10 & 44 & 162 & 8 & 20 & 24 & 16 \\
11 & 45 & 167 & 8 & 14 & 22 & 16 \\
12 & 46 & 169 & 8 & 14 & 22 & 16 \\
13 & 47 & 168 & 7 & 21 & 22 & 14 \\
14 & 48 & 173 & 7 & 21 & 22 & 14 \\
15 & 49 & 166 & 7 & 14 & 22 & 14 \\
16 & 50 & 156 & 7 & 14 & 24 & 14 \\
17 & 51 & 143 & 6 & 14 & 23 & 12 \\
18 & 52 & 127 & 5 & 14 & 21 & 10 \\
19 & 53 & 108 & 5 & 14 & 26 & 10 \\
20 & 54 & 86 & 4 & 14 & 26 & 8 \\
\end{longtable}

Every strict tail/off-scale/bridge margin and every intermediate m in
j0..b-1 is checked with integers. For k10,\allowbreak{} T20-b-1-R=\allowbreak{}3; for k13,\allowbreak{}k14,\allowbreak{}
T21-b-1-R=\allowbreak{}15,\allowbreak{}9. These are the three uses of the strengthened repair.

Residue counting gives kappa\textgreater=\allowbreak{}(t-1)(v-1) if U occupies t
residue classes modulo Q\textquotesingle s prefix lattice. Since
Kmax\textless Vmin-1,\allowbreak{} U has only one class,\allowbreak{} and 0 in U places P on the
same lattice. Set alpha=\allowbreak{}4beta+\allowbreak{}2. Q\textquotesingle s central interval
has length l=\allowbreak{}L\_\allowbreak{}Q-4beta-1=\allowbreak{}v+\allowbreak{}hQ-alpha\textgreater=\allowbreak{}v-alpha. Apply the
AP-plus-set formula to that interval and U. Since it is a subset of V,\allowbreak{}
its clipped gap sum is at most kappa+\allowbreak{}alpha. The checked inequality

\begin{lstlisting}
Vmin-Kmax-2alpha-1>0
\end{lstlisting}

makes kappa+\allowbreak{}alpha\textless l-1,\allowbreak{} so no gap can be clipped; hence
hP\textless=\allowbreak{}kappa+\allowbreak{}alpha. The further positive inequalities

\begin{lstlisting}
Vmin-Kmax-alpha-2>0, Umin-H-2>0
\end{lstlisting}

give both diameters at most u+\allowbreak{}v-3. Stanchescu\textquotesingle s
short-diameter theorem therefore gives the exact hP,\allowbreak{}hQ\textless=\allowbreak{}kappa,\allowbreak{}
and z=\allowbreak{}hP+\allowbreak{}hQ-kappa\textless=\allowbreak{}H. No individual primitivity of P has been
assumed.

The pre-G D1 comparison is bounded by -165-2I+\allowbreak{}2H,\allowbreak{} with the required +\allowbreak{}16
correction at a14; every row is strict. The outer comparison is strict
since 1191\textgreater4*77-10+\allowbreak{}600. This proves D1,\allowbreak{}D integral and the
extreme support intersections before any integer-hole G144 estimate is
used.

At a13/b55,\allowbreak{} the actual failed sufficient test gives rhoQ\textless=\allowbreak{}61.
Normalize Q by its intrinsic gcd. The primitive effective-hole witness
(r,\allowbreak{}beta,\allowbreak{}C\_\allowbreak{}h)=\allowbreak{}(55,\allowbreak{}2,\allowbreak{}52) has (K,\allowbreak{}D,\allowbreak{}M,\allowbreak{}J)=\allowbreak{}(11,\allowbreak{}63,\allowbreak{}13,\allowbreak{}13),\allowbreak{} strict
margins23,\allowbreak{}14,\allowbreak{} and forces muQ\textgreater=\allowbreak{}108,\allowbreak{} rhoQ\textgreater=\allowbreak{}62
whenever hQ\textgreater4. Hence hQ\textless=\allowbreak{}4. Residue counting,\allowbreak{}
endpoint hulls and the same short-diameter theorem give
hP,\allowbreak{}hQ\textless=\allowbreak{}kappa and z\textless=\allowbreak{}4; all required dense-diameter
margins are positive. The D1 comparison with +\allowbreak{}16 is -1023. After G144,\allowbreak{}
rhoP+\allowbreak{}rhoQ+\allowbreak{}T\textgreater=\allowbreak{}-1+\allowbreak{}50+\allowbreak{}144-8=\allowbreak{}185\textgreater159,\allowbreak{} so this boundary
is impossible.

The remaining a13/b55 branch is the actual sufficient strict inverse
branch. Its strict hole bounds give individual primitivity on the joint
lattice. Here kappa\textless=\allowbreak{}159-(-1)-50=\allowbreak{}110,\allowbreak{} so the rough D1 comparison
is at most -165-2*441+\allowbreak{}2*110+\allowbreak{}16=\allowbreak{}-811\textless0; the outer comparison is
strict as above. Thus lattice forcing and G144 are established without
using any later hole estimate. The post-G budget gives

\begin{lstlisting}
-1+50+144-2z<=159,
\end{lstlisting}

so z\textgreater=\allowbreak{}17 and hQ\textgreater=\allowbreak{}z\textgreater8. The fresh
primitive effective-hole witness (r,\allowbreak{}beta,\allowbreak{}C\_\allowbreak{}h)=\allowbreak{}(55,\allowbreak{}4,\allowbreak{}119) has
(K,\allowbreak{}D,\allowbreak{}M,\allowbreak{}J)=\allowbreak{}(17,\allowbreak{}136,\allowbreak{}19,\allowbreak{}19),\allowbreak{} strict margins 12,\allowbreak{}34,\allowbreak{} and forces
muQ\textgreater=\allowbreak{}175 whenever hQ\textgreater8. The flexible defect
inverse then gives rhoQ\textgreater=\allowbreak{}113 (the g7,\allowbreak{}8,\allowbreak{}9,\allowbreak{}10 rows each give
113). Since T\textgreater=\allowbreak{}48,\allowbreak{}

\begin{lstlisting}
rhoP+rhoQ+T>=-1+113+48=160>159.
\end{lstlisting}

Consequently both actual inverse branches at a13 are excluded.

\hypertarget{large-hole-final-support-overlap-all-32-orientations}{%
\subsection{4. Large-hole final-support overlap: all 32
orientations}\label{large-hole-final-support-overlap-all-32-orientations}}

For unrepaired strict-inverse main branches,\allowbreak{} D1 is immediately strict
when z\textless=\allowbreak{}82+\allowbreak{}I. Therefore only z\textgreater=\allowbreak{}Z=\allowbreak{}83+\allowbreak{}I needs an extra
argument. Every repaired branch has z\textless=\allowbreak{}16\textless Z and cannot
enter this large-hole case. The strict inverse upper-hole bounds
hX\textless\textbar FS(X)\textbar-1 are retained here. The budget
kappa\textgreater=\allowbreak{}z and both non-AP floors exclude k16..19 directly. The
remaining k0..15 have 32 oriented possible disjoint-final-support sides.

If the side X has r coefficients and its final prefix translates are
disjoint,\allowbreak{} muX\textgreater=\allowbreak{}T\_\allowbreak{}(r-2),\allowbreak{} so rhoX\textgreater=\allowbreak{}T\_\allowbreak{}(r-2)+\allowbreak{}2-2r.
The elementary contradiction margin is

\begin{lstlisting}
T_(r-2)+2-2r+n(s)+Z-Omega.
\end{lstlisting}

It is positive in 18 orientations. The remaining 14 comprise P at k0..12
and also Q at balanced k0. Put R=\allowbreak{}Omega-(T\_\allowbreak{}(r-2)+\allowbreak{}2-2r)-Z; the budget
forces the other side\textquotesingle s defect at most R. Apply the
strengthened repair with q20 using the following exact (k,\allowbreak{}s,\allowbreak{}R,\allowbreak{}beta,\allowbreak{}j0)
choices:

\begin{longtable}[]{@{}rrrrr@{}}
\toprule\noalign{}
k & s & R & beta & j0 \\
\midrule\noalign{}
\endhead
\bottomrule\noalign{}
\endlastfoot
0 & 34 & 55 & 5 & 26 \\
1 & 35 & 83 & 7 & 25 \\
2 & 36 & 106 & 8 & 24 \\
3 & 37 & 124 & 8 & 22 \\
4 & 38 & 137 & 9 & 24 \\
5 & 39 & 145 & 9 & 24 \\
6 & 40 & 148 & 9 & 24 \\
7 & 41 & 146 & 8 & 23 \\
8 & 42 & 139 & 8 & 25 \\
9 & 43 & 127 & 7 & 26 \\
10 & 44 & 110 & 6 & 27 \\
11 & 45 & 88 & 5 & 29 \\
12 & 46 & 61 & 4 & 36 \\
\end{longtable}

All strict margins hold,\allowbreak{} including T20-s-1\textgreater R; the final
support has at most 2beta\textless=\allowbreak{}18 holes,\allowbreak{} whereas
h\_\allowbreak{}other\textgreater=\allowbreak{}z\textgreater=\allowbreak{}Z\textgreater=\allowbreak{}83,\allowbreak{} a contradiction.
Thus both sides\textquotesingle{} actual final prefix support translates
intersect. This is a fresh endpoint argument; the weaker effective-hole
tests alone do not close k2..8 at SIZE85.

Let Y remove the largest coefficient from an r-side X. Intersecting
final support hulls cover the full side hull,\allowbreak{} so hX\textless=\allowbreak{}2hY. The
exact maximum-deletion identity,\allowbreak{} with rhoY\textgreater=\allowbreak{}0 and
destroyed-defect count nu\textgreater=\allowbreak{}0,\allowbreak{} gives

\begin{lstlisting}
max X<=r+rhoX+hX/2+2.
\end{lstlisting}

The common D1 comparison becomes

\begin{lstlisting}
|FS(D1)|-2s<=C+2rho-mu+z+ell2
           <=C+Omega+rho-mu<=-165-2I<0.
\end{lstlisting}

Here the one budget is used via rho+\allowbreak{}kappa+\allowbreak{}ell2\textless=\allowbreak{}Omega and
z\textless=\allowbreak{}kappa. Together with the outer comparison,\allowbreak{} this finishes
pre-G non-AP lattice forcing without circular use of G144.

\hypertarget{post-g-individual-primitivity-and-a14-closure}{%
\subsection{5. Post-G individual primitivity and a14
closure}\label{post-g-individual-primitivity-and-a14-closure}}

Only after the proved lattice and overlap steps use

\begin{lstlisting}
T>=kappa+144-3z, c+rhoP+rhoQ+T<=Omega, w<=4z.
\end{lstlisting}

In a strict simultaneous-inverse branch,\allowbreak{} an individual gcd\textgreater=\allowbreak{}2
would imply hX\textgreater=\allowbreak{}\textbar FS(X)\textbar-1,\allowbreak{} contradicting the
strict inverse upper-hole bound. In an endpoint-repaired branch,\allowbreak{} Q is
primitive on the constructed prefix lattice. If P had gcd
g\textgreater=\allowbreak{}2,\allowbreak{} its intrinsic support is non-AP and has
h0\textgreater=\allowbreak{}1. Rescaling gives

\begin{lstlisting}
hP=(g-1)(T_a+muP)+g*h0>=T_a+a-1.
\end{lstlisting}

Thus

\begin{lstlisting}
c<=Omega-n(a)-n(b)-T_a-a+1-144+3H.
\end{lstlisting}

The exact caps for k3..20 are

\begin{lstlisting}
-119,-88,-54,-35,-13,6,16,23,27,28,20,24,16,5,-15,-38,-58,-87.
\end{lstlisting}

Negative caps exclude immediately. In the remaining rows
c\textless=\allowbreak{}28\textless77,\allowbreak{} so all 18 ordinary covering units contribute
(j0),\allowbreak{} and w\textless=\allowbreak{}4H\textless=\allowbreak{}64. The SIZE85 no-empty-unit
obstruction,\allowbreak{} valid for w\textless=\allowbreak{}104,\allowbreak{} excludes these. Hence P is
individually primitive in every surviving repaired main row.

For a14/b54,\allowbreak{} the repaired branch has z\textless=\allowbreak{}8,\allowbreak{} so
c\textless=\allowbreak{}200-49-(144-16)=\allowbreak{}23\textless77 and w\textless=\allowbreak{}32; j0 excludes
it entirely. In the strict-inverse branch,\allowbreak{} if z\textgreater=\allowbreak{}13,\allowbreak{} the
witness (r,\allowbreak{}beta,\allowbreak{}C\_\allowbreak{}h)=\allowbreak{}(54,\allowbreak{}6,\allowbreak{}168) has (K,\allowbreak{}D,\allowbreak{}M,\allowbreak{}J)=\allowbreak{}(19,\allowbreak{}187,\allowbreak{}22,\allowbreak{}22),\allowbreak{} strict
margins11,\allowbreak{}43. It yields muQ\textgreater=\allowbreak{}223,\allowbreak{} rhoQ\textgreater=\allowbreak{}153. Since
T\textgreater=\allowbreak{}48,\allowbreak{} the budget would require 200\textgreater=\allowbreak{}153+\allowbreak{}48=\allowbreak{}201.
Therefore z\textless=\allowbreak{}12. Then c\textless=\allowbreak{}31\textless77,\allowbreak{}w\textless=\allowbreak{}48,\allowbreak{}
and j0 again excludes the branch. This closes a14 without using P
primitivity first.

\hypertarget{verification-boundary}{%
\subsection{6. Verification boundary}\label{verification-boundary}}

Run \texttt{python -I -S -B check85\_\allowbreak{}\allowbreak{}inverse.py}. It
reconstructs all SIZE85 parameters,\allowbreak{} all 18 endpoint inverse rows,\allowbreak{} both
a13 effective boundary,\allowbreak{} all 32 overlap orientations and 14 endpoint
repairs,\allowbreak{} nonprimitive caps,\allowbreak{} and the a14 effective closure. The
mathematical passage from actual sets to the displayed inequalities,\allowbreak{} the
retained published general theorems,\allowbreak{} the bootstrap/AP branches and the
no-empty-unit obstruction remain written proofs. No SIZE86 row
combination or old exception census is reused.

\clearpage
\hypertarget{a-sharper-elementary-source-envelope-at-size-85}{%
\section{A sharper elementary source envelope at size
85}\label{a-sharper-elementary-source-envelope-at-size-85}}

This deduction concerns the \textbf{77-point source} in the fresh
size-85 reduction. It does not use the computed fixed-label tables. The
69-point inner source has zero,\allowbreak{} positive maximum
\texttt{R\_\allowbreak{}\allowbreak{}s},\allowbreak{} negative absolute maximum
\texttt{R\_\allowbreak{}\allowbreak{}t},\allowbreak{} and four further entries on each sign

\[
 R_v+\allowbreak{}a_v,\allowbreak{}\quad R_v+\allowbreak{}b_v,\allowbreak{}\quad R_v+\allowbreak{}c_v,\allowbreak{}\quad R_v+\allowbreak{}e_v,\allowbreak{}
 \qquad 1\le a_v<b_v<c_v<e_v.
\]

Write

\[
 R=\allowbreak{}R_s+\allowbreak{}R_t,\allowbreak{}\quad U=\allowbreak{}a_s+\allowbreak{}b_s+\allowbreak{}a_t+\allowbreak{}b_t,\allowbreak{}\quad C=\allowbreak{}c_s+\allowbreak{}c_t,\allowbreak{}
 \quad J=\allowbreak{}(e_s-c_s)+\allowbreak{}(e_t-c_t).
\]

Let \texttt{K} be the number of distinct pair sums of
this 77-point source,\allowbreak{} with repetition. The previously derived mass
identity,\allowbreak{} with \texttt{T=\allowbreak{}77j-280-w+\allowbreak{}I+\allowbreak{}tau+\allowbreak{}delta},\allowbreak{} is

\[
 2K=\allowbreak{}T+\allowbreak{}4R+\allowbreak{}U+\allowbreak{}2C+\allowbreak{}J.                                      \tag{1}
\]

For a positive integer quartet \texttt{a<b<c<e},\allowbreak{} define
the distinct fringes

\[
 P=\allowbreak{}|\{x+\allowbreak{}y>c:x,\allowbreak{}y\in\{a,\allowbreak{}b,\allowbreak{}c\}\}|,\allowbreak{}
 \qquad B=\allowbreak{}|\{x+\allowbreak{}y>e:x,\allowbreak{}y\in\{a,\allowbreak{}b,\allowbreak{}c,\allowbreak{}e\}\}|.
\]

The elementary quartet fringe lemma already used in the source reduction
says

\[
 2B\le a+\allowbreak{}b+\allowbreak{}8-p,\allowbreak{}\qquad p=\allowbreak{}(a+\allowbreak{}b)\bmod2.                 \tag{2}
\]

\hypertarget{the-single-sign-improvement}{%
\subsection{The single-sign
improvement}\label{the-single-sign-improvement}}

For every such quartet,\allowbreak{}

\[
                         P+\allowbreak{}B-a-b\le5.                \tag{3}
\]

Here is a short proof,\allowbreak{} including its two exceptional small cases. Put
\texttt{u=\allowbreak{}a+\allowbreak{}b}. There are at most six unordered pair
sums of the triple,\allowbreak{} so \texttt{P<=\allowbreak{}6}. For
\texttt{u>=\allowbreak{}9},\allowbreak{} (2) gives

\[
 P+\allowbreak{}B-u\le6+\allowbreak{}\lfloor(u+\allowbreak{}8)/2\rfloor-u\le5.
\]

For \texttt{u=\allowbreak{}8},\allowbreak{} necessarily
\texttt{a<=\allowbreak{}3},\allowbreak{} hence \texttt{2a<=\allowbreak{}c};
at least this triple pair lies outside its upper fringe. Thus
\texttt{P<=\allowbreak{}5},\allowbreak{} and (2) gives
\texttt{B<=\allowbreak{}8}. For \texttt{u=\allowbreak{}7},\allowbreak{}
\texttt{P<=\allowbreak{}5} and \texttt{B<=\allowbreak{}7}; for
\texttt{u=\allowbreak{}5},\allowbreak{} \texttt{P<=\allowbreak{}4} and
\texttt{B<=\allowbreak{}6}; for \texttt{u=\allowbreak{}3},\allowbreak{}
\texttt{P=\allowbreak{}3} and \texttt{B<=\allowbreak{}5}. These
stated triple counts are the complete small cases tabulated below.

For \texttt{u=\allowbreak{}4},\allowbreak{} the only way
\texttt{P=\allowbreak{}4} is
\texttt{(a,\allowbreak{}b,\allowbreak{}c)=\allowbreak{}(1,\allowbreak{}3,\allowbreak{}4)}. Its possible values of
\texttt{B} at \texttt{e=\allowbreak{}5,\allowbreak{}6,\allowbreak{}7} are
\texttt{5,\allowbreak{}5,\allowbreak{}4}; at all
\texttt{e>=\allowbreak{}8=\allowbreak{}2c},\allowbreak{} \texttt{B=\allowbreak{}4}. Thus
\texttt{P+\allowbreak{}B-u<=\allowbreak{}5}. For \texttt{u=\allowbreak{}6},\allowbreak{}
the only way \texttt{P=\allowbreak{}5} is
\texttt{(a,\allowbreak{}b,\allowbreak{}c)=\allowbreak{}(2,\allowbreak{}4,\allowbreak{}5)}. Its values of
\texttt{B} at \texttt{e=\allowbreak{}6,\allowbreak{}7,\allowbreak{}8,\allowbreak{}9} are
\texttt{6,\allowbreak{}6,\allowbreak{}5,\allowbreak{}5}; at all
\texttt{e>=\allowbreak{}10=\allowbreak{}2c},\allowbreak{} \texttt{B=\allowbreak{}4}. This
also gives (3). If these respective triple maxima do not occur,\allowbreak{} (2)
already gives (3).

The complete triple counts at \texttt{u<=\allowbreak{}7} are finite:
when \texttt{c>=\allowbreak{}2b},\allowbreak{} exactly
\texttt{a+\allowbreak{}c,\allowbreak{}b+\allowbreak{}c,\allowbreak{}2c} exceed \texttt{c},\allowbreak{}
so \texttt{P=\allowbreak{}3}; the remaining values are:

\begin{longtable}[]{@{}ll@{}}
\toprule\noalign{}
\texttt{(a,\allowbreak{}b)} & \texttt{c} below
\texttt{2b} and corresponding
\texttt{P} \\
\midrule\noalign{}
\endhead
\bottomrule\noalign{}
\endlastfoot
\texttt{(1,\allowbreak{}2)} & \texttt{3:3} \\
\texttt{(1,\allowbreak{}3)} & \texttt{4:4,\allowbreak{} 5:3} \\
\texttt{(1,\allowbreak{}4)} &
\texttt{5:4,\allowbreak{} 6:4,\allowbreak{} 7:3} \\
\texttt{(2,\allowbreak{}3)} & \texttt{4:4,\allowbreak{} 5:4} \\
\texttt{(1,\allowbreak{}5)} &
\texttt{6:4,\allowbreak{} 7:4,\allowbreak{} 8:4,\allowbreak{} 9:3} \\
\texttt{(2,\allowbreak{}4)} &
\texttt{5:5,\allowbreak{} 6:3,\allowbreak{} 7:4} \\
\texttt{(1,\allowbreak{}6)} &
\texttt{7:4,\allowbreak{} 8:4,\allowbreak{} 9:4,\allowbreak{} 10:4,\allowbreak{} 11:3} \\
\texttt{(2,\allowbreak{}5)} &
\texttt{6:5,\allowbreak{} 7:4,\allowbreak{} 8:3,\allowbreak{} 9:4} \\
\texttt{(3,\allowbreak{}4)} &
\texttt{5:5,\allowbreak{} 6:5,\allowbreak{} 7:4} \\
\end{longtable}

In particular the weaker triple fringe inequality
\texttt{2P<=\allowbreak{}u+\allowbreak{}4-p} also holds universally: for
\texttt{u>=\allowbreak{}8} use \texttt{P<=\allowbreak{}6},\allowbreak{} and
otherwise use this table,\allowbreak{} followed by rounding to an even integer.

\hypertarget{two-exact-source-envelopes}{%
\subsection{Two exact source
envelopes}\label{two-exact-source-envelopes}}

Remove the largest entry on each sign from the 77-point source. The
remaining source has 75 points. Its pair sums lie in

\[
 [-2R_t-c_t,\allowbreak{}\ 2R_s+\allowbreak{}c_s]
\]

except for distinct same-sign sums between the three respective tails
whose offset sum exceeds \texttt{c\_\allowbreak{}\allowbreak{}v}. The interval
has \texttt{2R+\allowbreak{}C+\allowbreak{}1} integer positions,\allowbreak{} and there are at
most \texttt{P\_\allowbreak{}\allowbreak{}s+\allowbreak{}P\_\allowbreak{}\allowbreak{}t} distinct exceptions.
Reinserting two entries creates at most

\[
 2\cdot75+\allowbreak{}3=\allowbreak{}153
\]

unordered pair labels,\allowbreak{} including repetition. Hence

\[
 K\le2R+\allowbreak{}C+\allowbreak{}154+\allowbreak{}P_s+\allowbreak{}P_t.                               \tag{4}
\]

The analogous interval for the full source is

\[
 [-2R_t-e_t,\allowbreak{}\ 2R_s+\allowbreak{}e_s],\allowbreak{}
\]

with \texttt{2R+\allowbreak{}C+\allowbreak{}J+\allowbreak{}1} integer positions. Only pairs
between the four same-sign tails can lie outside it,\allowbreak{} yielding

\[
 K\le2R+\allowbreak{}C+\allowbreak{}J+\allowbreak{}1+\allowbreak{}B_s+\allowbreak{}B_t.                               \tag{5}
\]

Combining (1) with (4) and (5) gives

\[
 T\le308+\allowbreak{}2(P_s+\allowbreak{}P_t)-U-J,\allowbreak{}
 \qquad T\le2+\allowbreak{}2(B_s+\allowbreak{}B_t)-U+\allowbreak{}J.
\]

Add these inequalities and divide by two. Using (3) separately on both
signs,\allowbreak{}

\[
 T\le155+\allowbreak{}(P_s+\allowbreak{}B_s-a_s-b_s)+\allowbreak{}(P_t+\allowbreak{}B_t-a_t-b_t)
   \le\boxed{165}.                                  \tag{6}
\]

Consequently the source budget improves to

\[
               I+\allowbreak{}\tau+\allowbreak{}\delta\le445-77j+\allowbreak{}w.             \tag{7}
\]

This deduction is uniform in the source core,\allowbreak{} all offsets,\allowbreak{} and the outer
coefficients. No integrality of the outer coefficients is used beyond
the existing integer source reduction.

\hypertarget{a-sharper-one-maximum-envelope}{%
\subsection{A sharper one-maximum
envelope}\label{a-sharper-one-maximum-envelope}}

Remove only the maximum on sign \texttt{v}; write
\texttt{o} for the opposite sign. Its 76-point
remainder has a pair-sum central interval with
\texttt{2R+\allowbreak{}c\_\allowbreak{}\allowbreak{}v+\allowbreak{}e\_\allowbreak{}\allowbreak{}o+\allowbreak{}1} integer positions,\allowbreak{} and its
exceptions have at most \texttt{P\_\allowbreak{}\allowbreak{}v+\allowbreak{}B\_\allowbreak{}\allowbreak{}o} distinct
values. Reinsertion creates at most77 unordered pair labels. Therefore

\[
 K\le2R+\allowbreak{}c_v+\allowbreak{}e_o+\allowbreak{}78+\allowbreak{}P_v+\allowbreak{}B_o.                          \tag{9}
\]

Put \texttt{j\_\allowbreak{}\allowbreak{}v=\allowbreak{}e\_\allowbreak{}\allowbreak{}v-c\_\allowbreak{}\allowbreak{}v},\allowbreak{} so
\texttt{J=\allowbreak{}j\_\allowbreak{}\allowbreak{}s+\allowbreak{}j\_\allowbreak{}\allowbreak{}t},\allowbreak{} and
\texttt{p=\allowbreak{}p\_\allowbreak{}\allowbreak{}s+\allowbreak{}p\_\allowbreak{}\allowbreak{}t}. Combining (9) with (1),\allowbreak{}
\texttt{2P\_\allowbreak{}\allowbreak{}v<=\allowbreak{}a\_\allowbreak{}\allowbreak{}v+\allowbreak{}b\_\allowbreak{}\allowbreak{}v+\allowbreak{}4-p\_\allowbreak{}\allowbreak{}v},\allowbreak{} and (2) on sign
\texttt{o} gives

\[
 T+\allowbreak{}j_v-j_o\le168-p.
\]

The reflected inequality and \texttt{J>=\allowbreak{}T-18+\allowbreak{}p}
consequently give

\[
 j_v\ge T-93+\allowbreak{}p.
\]

For every positive strictly increasing triple,\allowbreak{}
\texttt{c\_\allowbreak{}\allowbreak{}v+\allowbreak{}p\_\allowbreak{}\allowbreak{}v>=\allowbreak{}4}: the sole possibility with
\texttt{c\_\allowbreak{}\allowbreak{}v=\allowbreak{}3} is
\texttt{(a\_\allowbreak{}\allowbreak{}v,\allowbreak{}b\_\allowbreak{}\allowbreak{}v)=\allowbreak{}(1,\allowbreak{}2)},\allowbreak{} for which
\texttt{p\_\allowbreak{}\allowbreak{}v=\allowbreak{}1}. Since
\texttt{p>=\allowbreak{}p\_\allowbreak{}\allowbreak{}v},\allowbreak{} this proves the uniform endpoint row

\[
                         e_v\ge\max(4,\allowbreak{}T-89).         \tag{10}
\]

For \texttt{Z5},\allowbreak{} \texttt{w=\allowbreak{}0,\allowbreak{}j=\allowbreak{}5} and
\texttt{T=\allowbreak{}105+\allowbreak{}I+\allowbreak{}tau+\allowbreak{}delta=\allowbreak{}105+\allowbreak{}A},\allowbreak{} so this is
\texttt{e\_\allowbreak{}\allowbreak{}v>=\allowbreak{}A+\allowbreak{}16} on \textbf{both physical signs}.
This improves the earlier crude one-maximum floor
\texttt{e\_\allowbreak{}\allowbreak{}v>=\allowbreak{}A+\allowbreak{}8} without an activity,\allowbreak{} normality,\allowbreak{} or
ownership assumption.

\hypertarget{fringe-bound-at-the-possible-endpoint-t165}{%
\subsection{\texorpdfstring{Fringe bound at the possible endpoint
\texttt{T=\allowbreak{}165}}{Fringe bound at the possible endpoint T=\allowbreak{}165}}\label{fringe-bound-at-the-possible-endpoint-t165}}

The original weaker two-maximum envelope
\texttt{K<=\allowbreak{}2R+\allowbreak{}C+\allowbreak{}166},\allowbreak{} combined with (1),\allowbreak{} gives
\texttt{T+\allowbreak{}U+\allowbreak{}J<=\allowbreak{}332}. Summing (2) over both signs,\allowbreak{} with
\texttt{p=\allowbreak{}p\_\allowbreak{}\allowbreak{}s+\allowbreak{}p\_\allowbreak{}\allowbreak{}t},\allowbreak{} and using (5) gives
\texttt{J>=\allowbreak{}T-18+\allowbreak{}p}. Therefore

\[
 2T+\allowbreak{}U+\allowbreak{}p\le350.
\]

If \texttt{T=\allowbreak{}165},\allowbreak{} this says
\texttt{U+\allowbreak{}p<=\allowbreak{}20},\allowbreak{} so

\[
                 B_s+\allowbreak{}B_t\le(U+\allowbreak{}16-p)/2\le18.          \tag{8}
\]

Thus the endpoint fringe saves four in a scalar estimate which had used
\texttt{B\_\allowbreak{}\allowbreak{}s+\allowbreak{}B\_\allowbreak{}\allowbreak{}t<=\allowbreak{}20}. The local scalar computation
supplied by the separate audit agent has raw margin at least
\texttt{1} for \texttt{T<=\allowbreak{}164} and
margin \texttt{-2} at \texttt{T=\allowbreak{}165};
(8) increases the latter margin to at least \texttt{2}.

\hypertarget{exact-verification}{%
\subsection{Exact verification}\label{exact-verification}}

\texttt{source\_\allowbreak{}\allowbreak{}triple\_\allowbreak{}\allowbreak{}fringe85.py} checks every
finite small case needed by (3),\allowbreak{} the two exceptional quartet tables,\allowbreak{} and
the displayed algebra. The unbounded domains are reduced by the explicit
\texttt{c>=\allowbreak{}2b} and \texttt{e>=\allowbreak{}2c}
observations above. It also performs a larger independent numerical
sanity check; that numerical check is supplementary,\allowbreak{} not the
justification for the unbounded statement.

\clearpage
\hypertarget{sharp-fixed-label-inventories-in-the-raw-gaps}{%
\section{Sharp fixed-label inventories in the raw
gaps}\label{sharp-fixed-label-inventories-in-the-raw-gaps}}

This proof does not depend on cardinality 85. It applies to the original
four-extreme cover with positive outer magnitudes
\(v_1>v_2>v_3>v_4>M_v>0\),\allowbreak{} opposite outer magnitudes \(o_i>M_o>0\),\allowbreak{} and
the same 55 raw fixed labels on target v.

Write \[
d=\allowbreak{}v_1-v_2,\allowbreak{}\quad h=\allowbreak{}v_2-v_3,\allowbreak{}\quad x=\allowbreak{}v_3,\allowbreak{}\quad y=\allowbreak{}v_4,\allowbreak{}
\quad M=\allowbreak{}M_v+\allowbreak{}M_o.
\] The ordinary centre identities are \[
\lambda=\allowbreak{}d-x-y,\allowbreak{}\quad\alpha=\allowbreak{}h-y,\allowbreak{}\quad\beta=\allowbreak{}h-d-y,\allowbreak{}
\quad H=\allowbreak{}d+\allowbreak{}h-y,\allowbreak{}\quad O=\allowbreak{}H+\allowbreak{}d.
\]

The raw fixed labels are C labels \(\lambda+\allowbreak{}b_i+\allowbreak{}b_j\) for unordered
pairs in \(\{v_1,\allowbreak{}v_2,\allowbreak{}v_3,\allowbreak{}v_4,\allowbreak{}-o_1,\allowbreak{}-o_2,\allowbreak{}-o_3,\allowbreak{}-o_4\}\); U labels
\(\alpha+\allowbreak{}b_i\) and V labels \(\beta+\allowbreak{}b_i\),\allowbreak{} excluding \(b_i=\allowbreak{}v_1\); and
the five F labels of the original cover.

All statements below count formal labels. Colliding or nonintegral
values only reduce the ability to cover integer target points.

\hypertarget{first-gap-at-most-eleven-labels}{%
\subsection{First gap: at most eleven
labels}\label{first-gap-at-most-eleven-labels}}

Assume the first gap \[
(H+\allowbreak{}M_v,\allowbreak{}O-M_o)
\] is guarded inside the positive target. Its length is \(t=\allowbreak{}d-M\),\allowbreak{} and
the lower guard gives \(H+\allowbreak{}M_v>0\).

Subtract \(\lambda\). A C label in the interval must be either one of
the six positive pairs among \(v_2,\allowbreak{}v_3,\allowbreak{}v_4\),\allowbreak{} or one of the four cross
pairs \(v_1-o_i\). All other pairs are below or above the interval. A U
or V label can use only \(v_2,\allowbreak{}v_3,\allowbreak{}v_4\),\allowbreak{} giving three candidates of each
kind.

The seventeen internal candidates split into the following five chains,\allowbreak{}
each with consecutive difference d: \[
\{C22,\allowbreak{}U2,\allowbreak{}V2\},\allowbreak{}\quad\{C23,\allowbreak{}U3,\allowbreak{}V3\},\allowbreak{}\quad\{C24,\allowbreak{}U4,\allowbreak{}V4\},\allowbreak{}
\quad\{C33,\allowbreak{}F1,\allowbreak{}F3\},\allowbreak{}\quad\{F2,\allowbreak{}F4\},\allowbreak{}
\] and three remaining labels \(C34,\allowbreak{}C44,\allowbreak{}F5\). Since \(t<d\),\allowbreak{} each chain
meets the gap at most once.

After subtracting H,\allowbreak{} the last two labels have values \[
C44-H=\allowbreak{}y-h-(x-y),\allowbreak{}\qquad F5-H=\allowbreak{}y+\allowbreak{}(x-y)-3h-2d.
\] The guard is \(y<d+\allowbreak{}h+\allowbreak{}M_v\). If C44 is above the lower endpoint,\allowbreak{} then
\(y>M_v+\allowbreak{}h+\allowbreak{}(x-y)\),\allowbreak{} hence \(x-y<d\). Consequently \[
F5-H<M_v+\allowbreak{}(x-y)-2h-d<M_v.
\] Thus C44 and F5 cannot both meet the gap. There are at most seven
internal labels and four cross labels: eleven in total. The original
covering gap argument therefore gives \[
t\le12+\allowbreak{}\eta_v.\tag{1}
\]

\hypertarget{beta-active-second-gap-at-most-eight-labels}{%
\subsection{Beta-active second gap: at most eight
labels}\label{beta-active-second-gap-at-most-eight-labels}}

Assume beta is active. Since \(\beta+\allowbreak{}M_v\ge1\),\allowbreak{} \[
h\ge d+\allowbreak{}y-M_v+\allowbreak{}1>d.
\] After subtracting alpha,\allowbreak{} the second gap is \[
J=\allowbreak{}(M_v,\allowbreak{}d-M_o).
\]

Its internal candidate labels,\allowbreak{} expressed relative to alpha,\allowbreak{} are:

\begin{longtable}[]{@{}ll@{}}
\toprule\noalign{}
Chain or singleton & Relative values \\
\midrule\noalign{}
\endhead
\bottomrule\noalign{}
\endlastfoot
D: C33,\allowbreak{}F1,\allowbreak{}F3 & \(d+\allowbreak{}x-h,\allowbreak{}\ x-h,\allowbreak{}\ x-h-d\) \\
G: F2,\allowbreak{}F4 & \(x-2h,\allowbreak{}\ x-2h-d\) \\
B: U3,\allowbreak{}V3 & \(x,\allowbreak{}\ x-d\) \\
C: U4,\allowbreak{}V4 & \(y,\allowbreak{}\ y-d\) \\
V2 & \(x+\allowbreak{}h-d\) \\
F5 & \(x-3h-d\) \\
\end{longtable}

The label U2 has value \(x+\allowbreak{}h>d\),\allowbreak{} so it cannot occur. C34 and C44 are
below the lower endpoint because \[
C34-\alpha=\allowbreak{}d+\allowbreak{}y-h<M_v,\allowbreak{}\qquad C44-\alpha<C34-\alpha.
\] The only cross candidates are C labels using \(v_1-o_i\) or
\(v_2-o_i\). For each i the two values differ by d,\allowbreak{} so their total
contribution to J is at most four.

Each of the four internal chains contributes at most one. If V2 lies in
J,\allowbreak{} then \(x+\allowbreak{}h<2d\),\allowbreak{} hence \(x<2d-h\); G and F5 are below zero. This
leaves V2 and the three other chains,\allowbreak{} at most four labels. If F5 lies in
J,\allowbreak{} then \(x>3h+\allowbreak{}d+\allowbreak{}M_v\); B,\allowbreak{}D,\allowbreak{}G and V2 are above J,\allowbreak{} leaving only chain C
and F5,\allowbreak{} at most two. If neither singleton occurs,\allowbreak{} the four chains give
at most four. Thus the second gap has at most four internal and four
cross labels,\allowbreak{} eight in total.

When H is active its upper guard is valid and the second gap therefore
gives \[
t\le9+\allowbreak{}\eta_v.\tag{2}
\]

\hypertarget{beta-active-union-at-most-fifteen-labels}{%
\subsection{Beta-active union: at most fifteen
labels}\label{beta-active-union-at-most-fifteen-labels}}

After subtracting alpha,\allowbreak{} the union of the two raw gaps is
\(J\cup(J+\allowbreak{}d)\). Its whole hull has length \(2d-M<2d\). The candidate
internal labels are the five d-spaced chains \[
\begin{array}{ll}
A:& C22,\allowbreak{}U2,\allowbreak{}V2=\allowbreak{}(d+\allowbreak{}x+\allowbreak{}h,\allowbreak{}\ x+\allowbreak{}h,\allowbreak{}\ x+\allowbreak{}h-d),\allowbreak{}\\
B:& C23,\allowbreak{}U3,\allowbreak{}V3=\allowbreak{}(d+\allowbreak{}x,\allowbreak{}\ x,\allowbreak{}\ x-d),\allowbreak{}\\
C:& C24,\allowbreak{}U4,\allowbreak{}V4=\allowbreak{}(d+\allowbreak{}y,\allowbreak{}\ y,\allowbreak{}\ y-d),\allowbreak{}\\
D:& C33,\allowbreak{}F1,\allowbreak{}F3=\allowbreak{}(d+\allowbreak{}x-h,\allowbreak{}\ x-h,\allowbreak{}\ x-h-d),\allowbreak{}\\
G:& F2,\allowbreak{}F4=\allowbreak{}(x-2h,\allowbreak{}\ x-2h-d),\allowbreak{}
\end{array}
\] and singleton F5,\allowbreak{} with value \(x-3h-d\). Each chain has at most two
values in the union.

\begin{itemize}
\tightlist
\item
  If \(x<d\),\allowbreak{} A,\allowbreak{}B,\allowbreak{}C have at most two each; D has at most one; G and F5
  are below zero. Total at most seven.
\item
  If \(d\le x\le2d\),\allowbreak{} A has at most one; B,\allowbreak{}C,\allowbreak{}D at most two each; G and
  F5 are below zero. Total at most seven.
\item
  If \(x>2d\),\allowbreak{} A is absent; B has at most one; C,\allowbreak{}D,\allowbreak{}G at most two each.
  If F5 is absent,\allowbreak{} the total is seven. If F5 occurs,\allowbreak{} then
  \(x>3h+\allowbreak{}d+\allowbreak{}M_v>4d\),\allowbreak{} which makes B absent; the total is again seven.
\end{itemize}

The cross pairs \(v_1-o_i,\allowbreak{}v_2-o_i\) give at most eight labels across the
union. All other candidates are excluded exactly as above. Thus both
disjoint gaps together have at most fifteen labels.

When O and beta are active,\allowbreak{} both intervals are guarded inside the
target,\allowbreak{} regardless of H activity. They must share these fifteen labels
and the same \(\eta_v\) holes. Therefore \[
2(t-1)\le15+\allowbreak{}\eta_v,\allowbreak{}\qquad
2t\le17+\allowbreak{}\eta_v.\tag{3}
\]

Nonpositive gap lengths make all three bounds automatic. Leading
activity,\allowbreak{} or the other established lower guards from the original common
lemma,\allowbreak{} licenses (1) whenever O is active. These inventories imply the
uniform cap \(t\le12+\allowbreak{}\eta_v\) for the predicate ``O active or beta,\allowbreak{}H
active.''

\clearpage
\hypertarget{cardinality-85-localization-of-the-entire-high-hole-family}{%
\section{Cardinality-85 localization of the entire high-hole
family}\label{cardinality-85-localization-of-the-entire-high-hole-family}}

Assume the fresh size-85 structural theorem: d=\allowbreak{}77,\allowbreak{} e=\allowbreak{}69,\allowbreak{} a+\allowbreak{}b=\allowbreak{}68,\allowbreak{}
15\textless=\allowbreak{}a,\allowbreak{}b\textless=\allowbreak{}53,\allowbreak{} I=\allowbreak{}(34-min(a,\allowbreak{}b))\^{}2,\allowbreak{} c\textgreater=\allowbreak{}0,\allowbreak{}
individually primitive integer cores,\allowbreak{} and

\begin{lstlisting}
kappa+rho+c+I<=600,
T>=kappa+144-3z, T>=z, T>=144-2z, w<=4z.
\end{lstlisting}

Put H=\allowbreak{}h(P)+\allowbreak{}h(Q)=\allowbreak{}kappa+\allowbreak{}z,\allowbreak{} mu=\allowbreak{}mu(P)+\allowbreak{}mu(Q). The retained verified profile
certificates give the following \textbf{two} lower bounds,\allowbreak{} not just
adjusted defect floors:

\begin{longtable}[]{@{}lrrl@{}}
\toprule\noalign{}
z range & mu floor & rho floor & scalar cap \\
\midrule\noalign{}
\endhead
\bottomrule\noalign{}
\endlastfoot
15-\/-16 & 426 & 294 & c+\allowbreak{}I\textless=\allowbreak{}194 \\
17-\/-20 & 470 & 338 & c+\allowbreak{}I\textless=\allowbreak{}158 \\
21-\/-48 & 514 & 382 & c+\allowbreak{}I\textless=\allowbreak{}170 \\
\textgreater=\allowbreak{}49 & 554 & 422 & c+\allowbreak{}I+\allowbreak{}z\textless=\allowbreak{}178 \\
\end{longtable}

These are minima over every oriented pair a+\allowbreak{}b=\allowbreak{}68. The mu values are
direct profile outputs; they do not follow merely from rho\textless=\allowbreak{}mu.
Keeping the mu information fixes the failure of the weaker
single-maximum comparison in the last band.

\hypertarget{source-width-versus-every-peeled-coefficient}{%
\subsection{1. Source width versus every peeled
coefficient}\label{source-width-versus-every-peeled-coefficient}}

The core\textquotesingle s absolute-multiset support has mass

\begin{lstlisting}
C0=1190+I+tau=1190+I+mu+H.
\end{lstlisting}

If each core mass exceeds H,\allowbreak{} the elementary sumset-hole lemma fills
{[}H,\allowbreak{}C0-H{]}. Write W=\allowbreak{}C0-2H=\allowbreak{}1190+\allowbreak{}I+\allowbreak{}mu-H for its endpoint displacement.

Use source offsets
alpha\_\allowbreak{}v\textless beta\_\allowbreak{}v\textless gamma\_\allowbreak{}v\textless epsilon\_\allowbreak{}v,\allowbreak{} one
quadruple on each physical sign. Write

\begin{lstlisting}
u=sum(alpha_v+beta_v), C=gamma_s+gamma_t,
J=sum(epsilon_v-gamma_v), L=C+J=M-R.
\end{lstlisting}

The two-maximum source envelope,\allowbreak{} freshly at d=\allowbreak{}77,\allowbreak{} is

\begin{lstlisting}
K<=2R+C+166,
J+u<=612+w-I-tau-c.
\end{lstlisting}

Here the central interval,\allowbreak{} twelve exceptional peeled-pair labels,\allowbreak{} and
reinsertion of two maxima contribute 1+\allowbreak{}12+\allowbreak{}(2*75+\allowbreak{}3)=\allowbreak{}166. The freshly
derived outer-layer inequality is

\begin{lstlisting}
2C+J<=C0+u-292-2I-c-mu-rho-2z.
\end{lstlisting}

Adding the J upper bound gives

\begin{lstlisting}
2L<=1510+w-2I-2c-mu-rho-2z.
\end{lstlisting}

Maximum removal gives R\textless=\allowbreak{}72+\allowbreak{}rho+\allowbreak{}H. Thus

\begin{lstlisting}
M<=827+H+(rho-mu)/2+w/2-I-c-z
  <=827+H+(rho-mu)/2+z-I-c,
\end{lstlisting}

using only the original w\textless=\allowbreak{}4z,\allowbreak{} before any localization.

Subtract from W and use H=\allowbreak{}kappa+\allowbreak{}z and the original single budget:

\begin{lstlisting}
W-M >= 363+2I+c+3mu/2-rho/2-2kappa-3z
     >= -837+4I+3c+3(mu+rho)/2-3z.                 (1)
\end{lstlisting}

Every coefficient has the correct sign; no slack is spent twice.
Substituting the paired mu and rho floors gives minimum margins

\begin{lstlisting}
195, 315, 363
\end{lstlisting}

in the first three bands. In the tail,\allowbreak{} (1) gives

\begin{lstlisting}
W-M>=627+4I+3c-3z
     >=93+7I+6c>0,
\end{lstlisting}

because z\textless=\allowbreak{}178-I-c. Therefore every one of the eight absolute
peeled coefficients is at most M\textless W. The enlarged interval has
even greater width after adjoining each entry,\allowbreak{} so all eight entries obey
the propagation criterion x\textless=\allowbreak{}current\_\allowbreak{}mass-2H+\allowbreak{}1.

\hypertarget{each-core-mass-exceeds-h}{%
\subsection{2. Each core mass exceeds
H}\label{each-core-mass-exceeds-h}}

The original budget and T\textgreater=\allowbreak{}kappa+\allowbreak{}144-3z give

\begin{lstlisting}
H<=456-rho-I-c+4z.
\end{lstlisting}

In the first three bands this yields respectively

\begin{lstlisting}
H<=226-I-c, 198-I-c, 266-I-c.
\end{lstlisting}

Their scalar I caps give k\textless=\allowbreak{}13,\allowbreak{}12,\allowbreak{}13 respectively,\allowbreak{} where a=\allowbreak{}34-k
is the smaller side. Exact minimization of T\_\allowbreak{}a-H over those integer
ranges gives lower bounds 171,\allowbreak{}199,\allowbreak{}131. No use of side relabeling occurs
in the covering argument; these are symmetric mass bounds.

In the tail,\allowbreak{} rho\textgreater=\allowbreak{}422 gives kappa\textless=\allowbreak{}178-I-c; also
z\textless=\allowbreak{}kappa. Hence

\begin{lstlisting}
H<=356-2I-2c.
\end{lstlisting}

Since z\textgreater=\allowbreak{}49,\allowbreak{} I+\allowbreak{}c\textless=\allowbreak{}129,\allowbreak{} so k\textless=\allowbreak{}11. The exact
minimum of T\_\allowbreak{}(34-k)-(356-2k\^{}2) for 0\textless=\allowbreak{}k\textless=\allowbreak{}11 is 120,\allowbreak{}
attained at k=\allowbreak{}7. Thus both individual core masses exceed H in every
high-hole band.

\hypertarget{localization-conclusion}{%
\subsection{3. Localization conclusion}\label{localization-conclusion}}

The core sumset-hole lemma supplies {[}H,\allowbreak{}C0-H{]}. Section 1 propagates
this same boundary strip through all eight absolute entries. Complement
symmetry and inclusion of the old absolute support in the new one then
give

\begin{lstlisting}
z and w are even,
0<=w<=z,
eta_s=eta_t=w/2.
\end{lstlisting}

This proves localization on the \textbf{entire} size-85 high-hole
domain. It uses only original structural bounds before establishing
these stronger conclusions. It in fact does not need the covering index
to be zero; the full slack c is kept throughout.

\hypertarget{verification-scope}{%
\subsection{4. Verification scope}\label{verification-scope}}

highhole\_\allowbreak{}localization85.py reads the already verified retained profile
table and recomputes every paired mu/rho minimum over all 39 oriented
sizes. It independently checks the new scalar and core-mass minima and
the exact integer source constants. This arithmetic accompanies the
written interval propagation proof; it does not replace it.

\clearpage
\hypertarget{local-elementary-exclusions-at-size-85}{%
\section{Local elementary exclusions at size
85}\label{local-elementary-exclusions-at-size-85}}

All statements below assume a valid fresh structural reduction for a
size-85 configuration. Put d=\allowbreak{}77,\allowbreak{} \textbar E\textbar=\allowbreak{}69,\allowbreak{} r\_\allowbreak{}s+\allowbreak{}r\_\allowbreak{}t=\allowbreak{}68,\allowbreak{}
Lambda=\allowbreak{}1191+\allowbreak{}I,\allowbreak{} I=\allowbreak{}(34-min(r\_\allowbreak{}s,\allowbreak{}r\_\allowbreak{}t))\^{}2. The structural scalar budget
is Omega=\allowbreak{}600-I and the restricted-sum constant is 144. These are
conditional local deductions,\allowbreak{} not a proof of the imported reduction.

\hypertarget{general-source-formulas}{%
\subsection{General source formulas}\label{general-source-formulas}}

S=\allowbreak{}2K+\allowbreak{}1470-77j+\allowbreak{}w-delta. With A=\allowbreak{}I+\allowbreak{}tau+\allowbreak{}delta and m0=\allowbreak{}77j-280-w,\allowbreak{}

2K=\allowbreak{}m0+\allowbreak{}A+\allowbreak{}4R+\allowbreak{}U+\allowbreak{}2C+\allowbreak{}J.

Deleting the two maxima leaves 75 points. The core interval and 12
exceptional three-peel pair labels,\allowbreak{} followed by 153 reinserted pair
labels,\allowbreak{} give

K \textless=\allowbreak{} 2R+\allowbreak{}C+\allowbreak{}166,\allowbreak{} J+\allowbreak{}U \textless=\allowbreak{} 332-m0-A,\allowbreak{} A \textless=\allowbreak{} 445-77j+\allowbreak{}w,\allowbreak{}
K-M \textgreater=\allowbreak{} R+\allowbreak{}77j-440-w+\allowbreak{}A,\allowbreak{} K-delta \textgreater=\allowbreak{} 77j-142-w+\allowbreak{}I+\allowbreak{}tau.

When z\textgreater0,\allowbreak{} primitive side supports are non-AP,\allowbreak{} so
R\_\allowbreak{}v\textgreater=\allowbreak{}r\_\allowbreak{}v+\allowbreak{}1 and R\textgreater=\allowbreak{}70. The last bound improves to

K-delta \textgreater=\allowbreak{}77j-138-w+\allowbreak{}I+\allowbreak{}tau.

Quantitative layers:

V0-C1+\allowbreak{}delta+\allowbreak{}2z \textless=\allowbreak{} -292-2I-77j-mu-rho,\allowbreak{} M0+\allowbreak{}delta-2z \textless=\allowbreak{}
-166-2I-77j.

The final source cap follows from the triple and quartet fringes. For
one sign let P count sums of the first three peel offsets above the
third offset and B count sums of all four offsets above the fourth. The
proved elementary inequality P+\allowbreak{}B-a-b\textless=\allowbreak{}5 combines the two-max
envelope K\textless=\allowbreak{}2R+\allowbreak{}C+\allowbreak{}154+\allowbreak{}P\_\allowbreak{}s+\allowbreak{}P\_\allowbreak{}t with the full-source envelope.
Their average gives m0+\allowbreak{}A\textless=\allowbreak{}165 and therefore
A\textless=\allowbreak{}445-77j+\allowbreak{}w. The complete unbounded proof is in
foundations/source\_\allowbreak{}triple\_\allowbreak{}fringe85.md.

\hypertarget{sharp-raw-gap-counts}{%
\subsection{Sharp raw-gap counts}\label{sharp-raw-gap-counts}}

The exact inventories are eleven labels in the first gap,\allowbreak{} eight in the
beta-active second gap,\allowbreak{} and fifteen in the beta-active union. Thus

O\_\allowbreak{}v active =\allowbreak{}\textgreater{} t\_\allowbreak{}v=\allowbreak{}d\_\allowbreak{}v-M\textless=\allowbreak{}12+\allowbreak{}eta\_\allowbreak{}v,\allowbreak{} beta\_\allowbreak{}v,\allowbreak{}H\_\allowbreak{}v
active =\allowbreak{}\textgreater{} t\_\allowbreak{}v\textless=\allowbreak{}9+\allowbreak{}eta\_\allowbreak{}v,\allowbreak{} beta\_\allowbreak{}v,\allowbreak{}O\_\allowbreak{}v active
=\allowbreak{}\textgreater{} 2t\_\allowbreak{}v\textless=\allowbreak{}17+\allowbreak{}eta\_\allowbreak{}v.

The first implication uses the original lower guard; leading activity,\allowbreak{} H
activity,\allowbreak{} or own2 activity licenses it in the corresponding
applications. The other two guards are supplied by the stated active
resources. The proof by d-spaced fixed-label chains and three ranges for
v3/d is written in full in common/raw\_\allowbreak{}gap\_\allowbreak{}inventory.md. These counts
imply the uniform cap t\_\allowbreak{}v\textless=\allowbreak{}12+\allowbreak{}eta\_\allowbreak{}v for the predicate ``O
active or beta,\allowbreak{}H active.''

\hypertarget{star-and-two-cap-constants}{%
\subsection{Star and two-cap
constants}\label{star-and-two-cap-constants}}

Four full-core star edges give 2X+\allowbreak{}4R+\allowbreak{}delta\textgreater=\allowbreak{}548. The
four-top-tail count has 4*83+\allowbreak{}5=\allowbreak{}337 labels per target,\allowbreak{} so

d\_\allowbreak{}s+\allowbreak{}d\_\allowbreak{}t+\allowbreak{}G+\allowbreak{}delta-2K\textgreater=\allowbreak{}796-77j.

With both sharp t\_\allowbreak{}v\textless=\allowbreak{}12+\allowbreak{}eta\_\allowbreak{}v caps,\allowbreak{} the two-cap proof,\allowbreak{} using
Y=\allowbreak{}2M-2K+\allowbreak{}24+\allowbreak{}w+\allowbreak{}delta,\allowbreak{} gives

4I+\allowbreak{}4tau+\allowbreak{}delta+\allowbreak{}231j \textless=\allowbreak{}464+\allowbreak{}6w.

At j=\allowbreak{}1,\allowbreak{} with tau\textgreater=\allowbreak{}62+\allowbreak{}2z and w\textless=\allowbreak{}z,\allowbreak{} the two-cap
inequality would require

15+\allowbreak{}2z+\allowbreak{}4I+\allowbreak{}delta\textless=\allowbreak{}0,\allowbreak{}

which is impossible. In particular the P1 two-cap exclusion needs no
additional four-edge estimate.

\hypertarget{general-first-gap-and-terminal-lemma}{%
\subsection{General first-gap and terminal
lemma}\label{general-first-gap-and-terminal-lemma}}

If the absolute core has a hull hole,\allowbreak{} both side sizes\textgreater=\allowbreak{}16,\allowbreak{}
and I+\allowbreak{}tau\textless=\allowbreak{}395,\allowbreak{} its first missing integer f is \textless=\allowbreak{}8.
Write k=\allowbreak{}k\_\allowbreak{}s+\allowbreak{}k\_\allowbreak{}t for the low coefficients. Then
k\textless=\allowbreak{}4,\allowbreak{}k\_\allowbreak{}v\textless=\allowbreak{}3. The high baseline maxima are
B\_\allowbreak{}v=\allowbreak{}r\_\allowbreak{}v+\allowbreak{}f-k\_\allowbreak{}v,\allowbreak{} B=\allowbreak{}68+\allowbreak{}2f-k,\allowbreak{} with baseline excess b,\allowbreak{} residual
sigma=\allowbreak{}tau-b\textgreater=\allowbreak{}0. Direct summation gives

I+\allowbreak{}b=\allowbreak{}69f-1-(f+\allowbreak{}34.5)k+\allowbreak{}k\^{}2/4+\allowbreak{}(r\_\allowbreak{}s-34-(k\_\allowbreak{}s-k\_\allowbreak{}t)/2)\^{}2,\allowbreak{}
I+\allowbreak{}b-5B\textgreater=\allowbreak{}-282.

Proof of f\textless=\allowbreak{}8 follows exactly the old argument: f\textless max
R\_\allowbreak{}v\textless=\allowbreak{}52+\allowbreak{}tau\textless=\allowbreak{}447,\allowbreak{} so k\textless=\allowbreak{}41; for
f\textgreater=\allowbreak{}9,\allowbreak{}k\textless=\allowbreak{}4 the lower bound is
\textgreater=\allowbreak{}65f-139\textgreater=\allowbreak{}446. For k\textgreater=\allowbreak{}5,\allowbreak{} use
f\textgreater=\allowbreak{}k(k+\allowbreak{}2)/4+\allowbreak{}1 to get
{[}-k\^{}3+\allowbreak{}68k\^{}2-4k+\allowbreak{}272{]}/4\textgreater=\allowbreak{}1827/4\textgreater395 (the
polynomial increases over k=\allowbreak{}5,\allowbreak{}...,\allowbreak{}41). The baseline-minus-5B minimum is
attained at f=\allowbreak{}1,\allowbreak{}k=\allowbreak{}0; f=\allowbreak{}2,\allowbreak{}3 and f=\allowbreak{}4,\allowbreak{}...,\allowbreak{}8 give larger lower bounds.

Delete six core entries per sign: the retained set has 57 points,\allowbreak{} and
its certified overlap parameter satisfies

6q=\allowbreak{}750-7B+\allowbreak{}R-sigma.

The old rooted terminal forest and the source estimate at j=\allowbreak{}1 give

2t\_\allowbreak{}u \textgreater=\allowbreak{}389-7B+\allowbreak{}2R+\allowbreak{}I+\allowbreak{}b-w \textgreater=\allowbreak{}107-w.

The sharp double gap gives 2t\_\allowbreak{}u\textless=\allowbreak{}17+\allowbreak{}w/2. Thus the terminal form
forces

90\textless=\allowbreak{}3w/2.

\hypertarget{p1-localization}{%
\subsection{P1 localization}\label{p1-localization}}

The scalar range is 16\textless=\allowbreak{}r\_\allowbreak{}v\textless=\allowbreak{}52: otherwise
I\textgreater=\allowbreak{}361 while delta\textless=\allowbreak{}407-I-rho and rho\textgreater=\allowbreak{}49.
H=\allowbreak{}h\_\allowbreak{}s+\allowbreak{}h\_\allowbreak{}t obeys H\textless=\allowbreak{}306+\allowbreak{}w-I-delta\textless=\allowbreak{}362-I-delta. For
every physical size,\allowbreak{} T\_\allowbreak{}r+\allowbreak{}I\textgreater=\allowbreak{}397,\allowbreak{} ensuring each core
mass\textgreater H. Initial width \textgreater=\allowbreak{}890+\allowbreak{}2I+\allowbreak{}delta. Every inner
entry is \textless=\allowbreak{}637,\allowbreak{} and U0\textgreater=\allowbreak{}142,\allowbreak{} so the next width is
\textgreater=\allowbreak{}1032. The outer-layer row and the same H bound show
C1+\allowbreak{}Q0-4H\textgreater=\allowbreak{}788+\allowbreak{}6I+\allowbreak{}4delta+\allowbreak{}2z\textgreater0,\allowbreak{} where
Q0=\allowbreak{}369+\allowbreak{}mu+\allowbreak{}rho+\allowbreak{}2I+\allowbreak{}delta+\allowbreak{}2z. The last-shift gap J\textless=\allowbreak{}585 is smaller
than the width. Thus the original common-strip argument propagates and
gives z,\allowbreak{}w even,\allowbreak{} w\textless=\allowbreak{}z,\allowbreak{} eta\_\allowbreak{}v=\allowbreak{}w/2.

Leading activity follows from
K-delta\textgreater=\allowbreak{}-61-w+\allowbreak{}I+\allowbreak{}tau\textgreater=\allowbreak{}1+\allowbreak{}I+\allowbreak{}2z-w\textgreater0. At
most one missing own singleton gives full outer integrality. The
sharpened two-cap and terminal contradictions exclude P1.

\hypertarget{p5-complete-source-only-contradiction}{%
\subsection{P5: complete source-only
contradiction}\label{p5-complete-source-only-contradiction}}

The scalar range gives rho\textgreater=\allowbreak{}58,\allowbreak{} delta\textless=\allowbreak{}41-I,\allowbreak{}
28\textless=\allowbreak{}r\_\allowbreak{}v\textless=\allowbreak{}40. Also H\textless=\allowbreak{}13+\allowbreak{}4z-I-delta\textless=\allowbreak{}69.
Quantitative layers give C0-U0\textgreater=\allowbreak{}551+\allowbreak{}2I+\allowbreak{}delta-2z and
C1-V0\textgreater=\allowbreak{}797+\allowbreak{}2I+\allowbreak{}delta+\allowbreak{}2z. Subtracting 2H leaves
\textgreater=\allowbreak{}385 and \textgreater=\allowbreak{}687,\allowbreak{} so common-strip localization
applies: w\textless=\allowbreak{}z\textless=\allowbreak{}14 and parity holds.

The final source envelope gives A\textless=\allowbreak{}60+\allowbreak{}w\textless=\allowbreak{}60+\allowbreak{}z,\allowbreak{} whereas
positive core holes imply tau\textgreater=\allowbreak{}62+\allowbreak{}2z. Thus

2+\allowbreak{}z+\allowbreak{}I+\allowbreak{}delta\textless=\allowbreak{}0,\allowbreak{}

an immediate contradiction. This excludes P5 without any
outer-coordinate integrality or paired-owner assumption.

\hypertarget{the-entire-high-hole-family}{%
\subsection{The entire high-hole
family}\label{the-entire-high-hole-family}}

Use the retained paired (mu,\allowbreak{}rho) floors (426,\allowbreak{}294),\allowbreak{} (470,\allowbreak{}338),\allowbreak{} (514,\allowbreak{}382),\allowbreak{}
(554,\allowbreak{}422),\allowbreak{} respectively for z bands15-16,\allowbreak{}17-20,\allowbreak{}21-48,\allowbreak{} and
z\textgreater=\allowbreak{}49. The whole-family localization proof is in
foundations/highhole\_\allowbreak{}localization85.md. It keeps the independently
certified mu information as well as rho and requires no extension of the
density threshold.

Writing W=\allowbreak{}1190+\allowbreak{}I+\allowbreak{}mu-H for the original core central-strip width,\allowbreak{} the
fresh source and layer envelopes give

M\textless=\allowbreak{}827+\allowbreak{}H+\allowbreak{}(rho-mu)/2+\allowbreak{}z-I-c,\allowbreak{}
W-M\textgreater=\allowbreak{}-837+\allowbreak{}4I+\allowbreak{}3c+\allowbreak{}3(mu+\allowbreak{}rho)/2-3z.

The retained paired floors give W-M margins195,\allowbreak{}315,\allowbreak{}363 in the first
three bands. In the tail,\allowbreak{} c+\allowbreak{}I+\allowbreak{}z\textless=\allowbreak{}178 gives margin at
least93+\allowbreak{}7I+\allowbreak{}6c. The two individual core masses exceed H with
margins171,\allowbreak{}199,\allowbreak{}131,\allowbreak{}120 in the four bands. Thus the same central strip
propagates through every absolute source entry. This proves parity,\allowbreak{}
eta\_\allowbreak{}v=\allowbreak{}w/2 and w\textless=\allowbreak{}z on the entire high-hole domain,\allowbreak{} including
H0.

At j\textgreater=\allowbreak{}2 the source envelope gives I+\allowbreak{}tau+\allowbreak{}delta\textless=\allowbreak{}291+\allowbreak{}w;
with tau\textgreater=\allowbreak{}rho+\allowbreak{}2z,\allowbreak{} this would require
rho+\allowbreak{}z+\allowbreak{}I+\allowbreak{}delta\textless=\allowbreak{}291,\allowbreak{} already impossible at rho294,\allowbreak{}z15. All H\_\allowbreak{}j
with j\textgreater=\allowbreak{}2 are excluded (larger j only improves the bound).

At j=\allowbreak{}1,\allowbreak{} the source envelope leaves only even z=\allowbreak{}16,\allowbreak{}18,\allowbreak{}20 and
I+\allowbreak{}tau\textless=\allowbreak{}395. Both-leading activity follows directly after
localization from
K-delta\textgreater=\allowbreak{}-61-w+\allowbreak{}I+\allowbreak{}tau\textgreater=\allowbreak{}233+\allowbreak{}z+\allowbreak{}I\textgreater0. At
least three own centers per sign imply outer integrality. The two-cap
contradiction forces the terminal form,\allowbreak{} and 90\textless=\allowbreak{}3w/2\textless=\allowbreak{}30
is impossible.

After the same whole-domain localization,\allowbreak{} at j=\allowbreak{}0 the lower bound
K-delta\textgreater=\allowbreak{}-142-w+\allowbreak{}I+\allowbreak{}tau\textgreater=\allowbreak{}152+\allowbreak{}z+\allowbreak{}I\textgreater0 makes
both leadings active. All own centres are active,\allowbreak{} so outer integrality
follows. The two-cap inequality gives 4I+\allowbreak{}4tau+\allowbreak{}c\textless=\allowbreak{}464+\allowbreak{}6w,\allowbreak{} while
tau\textgreater=\allowbreak{}294+\allowbreak{}2z and w\textless=\allowbreak{}z would require
712+\allowbreak{}2z+\allowbreak{}4I+\allowbreak{}c\textless=\allowbreak{}0. Thus H0 is excluded as well.

\clearpage
\hypertarget{conditional-cardinality-85-zero-hole-reduction}{%
\section{Conditional cardinality-85 zero-hole
reduction}\label{conditional-cardinality-85-zero-hole-reduction}}

These notes assume the common integral,\allowbreak{} primitive-core structural
conclusion at cardinality 85. They do not claim to prove that
conclusion. The original outer eight entries may initially be real. The
actual-cover grouping rule and the raw 55 fixed-label inventories are
unchanged.

Write \texttt{d=\allowbreak{}77},\allowbreak{} \texttt{|E|=\allowbreak{}69},\allowbreak{}
\texttt{r\_\allowbreak{}\allowbreak{}s+\allowbreak{}r\_\allowbreak{}\allowbreak{}t=\allowbreak{}68},\allowbreak{}
\texttt{I=\allowbreak{}(r\_\allowbreak{}\allowbreak{}s-34)\^2},\allowbreak{}
\texttt{T\_\allowbreak{}\allowbreak{}r=\allowbreak{}r(r+\allowbreak{}1)/\allowbreak{}2},\allowbreak{} and
\texttt{tau=\allowbreak{}sum(P\_\allowbreak{}\allowbreak{}s)+\allowbreak{}sum(P\_\allowbreak{}\allowbreak{}t)-T\_\allowbreak{}\allowbreak{}\{r\_\allowbreak{}\allowbreak{}s\}-T\_\allowbreak{}\allowbreak{}\{r\_\allowbreak{}\allowbreak{}t\}}.
Put \texttt{R=\allowbreak{}R\_\allowbreak{}\allowbreak{}s+\allowbreak{}R\_\allowbreak{}\allowbreak{}t},\allowbreak{}
\texttt{xi=\allowbreak{}R-68},\allowbreak{} so
\texttt{0<=\allowbreak{}xi<=\allowbreak{}tau}. The four inner peels on side v
have magnitudes \texttt{R\_\allowbreak{}\allowbreak{}v+\allowbreak{}a\_\allowbreak{}\allowbreak{}v},\allowbreak{}
\texttt{R\_\allowbreak{}\allowbreak{}v+\allowbreak{}b\_\allowbreak{}\allowbreak{}v},\allowbreak{}
\texttt{R\_\allowbreak{}\allowbreak{}v+\allowbreak{}c\_\allowbreak{}\allowbreak{}v},\allowbreak{}
\texttt{R\_\allowbreak{}\allowbreak{}v+\allowbreak{}e\_\allowbreak{}\allowbreak{}v=\allowbreak{}M\_\allowbreak{}\allowbreak{}v},\allowbreak{} with integral positive
strictly increasing offsets. Define
\texttt{L=\allowbreak{}e\_\allowbreak{}\allowbreak{}s+\allowbreak{}e\_\allowbreak{}\allowbreak{}t=\allowbreak{}M-R},\allowbreak{}
\texttt{U=\allowbreak{}a\_\allowbreak{}\allowbreak{}s+\allowbreak{}b\_\allowbreak{}\allowbreak{}s+\allowbreak{}a\_\allowbreak{}\allowbreak{}t+\allowbreak{}b\_\allowbreak{}\allowbreak{}t},\allowbreak{}
\texttt{C=\allowbreak{}c\_\allowbreak{}\allowbreak{}s+\allowbreak{}c\_\allowbreak{}\allowbreak{}t},\allowbreak{} \texttt{J=\allowbreak{}L-C},\allowbreak{}
\texttt{A=\allowbreak{}I+\allowbreak{}tau+\allowbreak{}delta}.

\hypertarget{fresh-source-envelope-all-zero-hole-j}{%
\subsection{Fresh source envelope (all zero-hole
j)}\label{fresh-source-envelope-all-zero-hole-j}}

The ordinary whole-unit capacity is 77,\allowbreak{} not 78. The fixed allowance
remains 84. Hence

\begin{lstlisting}
S=2K+1470-77j-delta,
2K=77j-280+A+4R+U+C+L.
\end{lstlisting}

The second equality follows from the actual source mass
\texttt{S=\allowbreak{}1190+\allowbreak{}I+\allowbreak{}tau+\allowbreak{}4R+\allowbreak{}U+\allowbreak{}C+\allowbreak{}L}.

Deleting both maxima leaves 75 values. Its pair-sum hull has at most
\texttt{2R+\allowbreak{}C+\allowbreak{}1} positions and 12 exceptional same-sign
peeled-pair labels. Reinsertion costs
\texttt{2*75+\allowbreak{}3=\allowbreak{}153} labels. Therefore
\texttt{K<=\allowbreak{}2R+\allowbreak{}C+\allowbreak{}166} and
\texttt{J+\allowbreak{}U<=\allowbreak{}612-77j-A}.

The full source envelope is
\texttt{K<=\allowbreak{}2R+\allowbreak{}L+\allowbreak{}1+\allowbreak{}B\_\allowbreak{}\allowbreak{}s+\allowbreak{}B\_\allowbreak{}\allowbreak{}t}. The four-offset
inequalities \texttt{2B\_\allowbreak{}\allowbreak{}v<=\allowbreak{}a\_\allowbreak{}\allowbreak{}v+\allowbreak{}b\_\allowbreak{}\allowbreak{}v+\allowbreak{}8} and
\texttt{2B\_\allowbreak{}\allowbreak{}v-a\_\allowbreak{}\allowbreak{}v-b\_\allowbreak{}\allowbreak{}v-c\_\allowbreak{}\allowbreak{}v<=\allowbreak{}3} consequently give

\begin{lstlisting}
J>=A+77j-298,
L>=A+77j-288,
A<=452-77j,
K-delta>=77j-142+I+tau,
K-M>=R+77j-440+A.
\end{lstlisting}

For the cap on A use \texttt{U>=\allowbreak{}6} and the two bounds
on J. For \texttt{K-delta} combine its exact mass with
the lower bound on L and \texttt{4R+\allowbreak{}U+\allowbreak{}C>=\allowbreak{}272+\allowbreak{}12}. For
\texttt{K-M} use the upper bound on
\texttt{L=\allowbreak{}C+\allowbreak{}J} and \texttt{U>=\allowbreak{}6}. In
particular j\textgreater=\allowbreak{}6 is already impossible.

\hypertarget{adaptive-core-supports}{%
\subsection{Adaptive core supports}\label{adaptive-core-supports}}

The full core has 69 points and span R,\allowbreak{} so its overlap at integral
displacement q is at least \texttt{137-R-q=\allowbreak{}69-xi-q}.

After deleting the largest k positive core values on each sign,\allowbreak{} the
source has \texttt{69-2k} points. The sum of deviations
excluding each maximum is \texttt{tau-xi}. Monotonicity
gives span at most \texttt{68-2k+\allowbreak{}(tau-xi)/\allowbreak{}k},\allowbreak{} so its
overlap is at least \texttt{f\_\allowbreak{}\allowbreak{}k-q},\allowbreak{} where

\begin{lstlisting}
f_k=69-2k-(tau-xi)/k.
\end{lstlisting}

Every use of \texttt{max(f\_\allowbreak{}\allowbreak{}0,\allowbreak{}f\_\allowbreak{}\allowbreak{}k)} means choosing the
actual one of these two subsets that attains the maximum,\allowbreak{} before
constructing the forest. It does not add two overlap bounds or spend two
deficit budgets.

\hypertarget{z0-is-excluded}{%
\subsection{Z0 is excluded}\label{z0-is-excluded}}

All own centres are active,\allowbreak{} making all outer coordinates integral.
\texttt{M>=\allowbreak{}76}; active C7 gives
\texttt{d\_\allowbreak{}\allowbreak{}v>=\allowbreak{}M+\allowbreak{}5>=\allowbreak{}81}. Choose an actual active beta
contributor on side s. The raw two-gap bound gives
\texttt{t\_\allowbreak{}\allowbreak{}s=\allowbreak{}d\_\allowbreak{}\allowbreak{}s-M<=\allowbreak{}57/\allowbreak{}2}; the opposite O gives
\texttt{t\_\allowbreak{}\allowbreak{}t<=\allowbreak{}56}. Active C4\_\allowbreak{}s would imply
\texttt{t\_\allowbreak{}\allowbreak{}s>d\_\allowbreak{}\allowbreak{}t>=\allowbreak{}81}. Active C4\_\allowbreak{}t would imply
\texttt{t\_\allowbreak{}\allowbreak{}t>h\_\allowbreak{}\allowbreak{}s+\allowbreak{}d\_\allowbreak{}\allowbreak{}s>2d\_\allowbreak{}\allowbreak{}s>=\allowbreak{}162}. Both contradict
their caps,\allowbreak{} so pair1 has no actual contributor despite being retained.

\hypertarget{z1-is-excluded-with-a-stronger-adaptive-endpoint}{%
\subsection{Z1 is excluded with a stronger adaptive
endpoint}\label{z1-is-excluded-with-a-stronger-adaptive-endpoint}}

Outer integrality is proved from the same actual one-hole argument as
the report: a side with all four own centres determines its outer values
and lambda; opposite C6/C7 determine the other lower values; its
remaining own centres determine its top values. At least one of C6/C7 is
active on each side,\allowbreak{} so \texttt{d\_\allowbreak{}\allowbreak{}v>=\allowbreak{}81}.

If pair1 were missing,\allowbreak{} an actual beta contributor opposite an actual C5
contributor would contradict the raw O cap. Therefore pair1 is retained.
Its actual contributor forces its own O to be the unique whole hole; its
beta is inactive. Branch B is impossible. In A the opposite side u has
beta,\allowbreak{} H,\allowbreak{} O all active,\allowbreak{} so \texttt{2t\_\allowbreak{}\allowbreak{}u<=\allowbreak{}57}. All
alpha,\allowbreak{} C5,\allowbreak{} U,\allowbreak{} V,\allowbreak{} C7 resources below are retained.

For the six-edge trees \texttt{L\_\allowbreak{}\allowbreak{}v-V\_\allowbreak{}\allowbreak{}v-U\_\allowbreak{}\allowbreak{}v} and
\texttt{L\_\allowbreak{}\allowbreak{}v-C7\_\allowbreak{}\allowbreak{}v},\allowbreak{} across both signs,\allowbreak{} choose

\begin{lstlisting}
f=max(137-R,57-(tau-xi)/6).
\end{lstlisting}

Designate the full leading on u as root. Its native loss is ell\_\allowbreak{}u. The
actual-cover forest inequality yields
\texttt{6f-2X-G+\allowbreak{}ell\_\allowbreak{}\allowbreak{}u<=\allowbreak{}delta}. Active C7\_\allowbreak{}u supplies a
positive leading maximum,\allowbreak{} so
\texttt{d\_\allowbreak{}\allowbreak{}u>=\allowbreak{}K+\allowbreak{}2X\_\allowbreak{}\allowbreak{}u+\allowbreak{}g\_\allowbreak{}\allowbreak{}u-ell\_\allowbreak{}\allowbreak{}u}. Active alpha\_\allowbreak{}v and
C5\_\allowbreak{}u also give \texttt{t\_\allowbreak{}\allowbreak{}u>=\allowbreak{}2X+\allowbreak{}G+\allowbreak{}2}. Adding produces

\begin{lstlisting}
2t_u>=K-M-delta+6f+2
     >=I+max(121-5xi+tau,49+2xi)
     >=I+max(121-4xi,49+2xi)
     >=73+I.
\end{lstlisting}

The last minimum is attained at xi=\allowbreak{}12. This contradicts
\texttt{2t\_\allowbreak{}\allowbreak{}u<=\allowbreak{}57}. The six-rank support exists because
the structural/source size bounds give at least 15 entries on each core
side (and the source cap alone gives at least 15 for j=\allowbreak{}1).

\hypertarget{z2-common-activity-reductions}{%
\subsection{Z2 common activity
reductions}\label{z2-common-activity-reductions}}

Here \texttt{A<=\allowbreak{}298} and
\texttt{K-delta>=\allowbreak{}12+\allowbreak{}I+\allowbreak{}tau},\allowbreak{} so both leading resources
are active. The same two-witness residue argument proves outer
integrality.

The four-edge physical arm bound is freshly
\texttt{2X\_\allowbreak{}\allowbreak{}v+\allowbreak{}X\_\allowbreak{}\allowbreak{}o+\allowbreak{}g\_\allowbreak{}\allowbreak{}v>=\allowbreak{}252-tau-delta}: trim three
ranks per side,\allowbreak{} whose span uses the fourth terminal deviation,\allowbreak{} giving
each edge the lower overlap \texttt{63-tau/\allowbreak{}4}.

Suppose both cap predicates T\_\allowbreak{}s,\allowbreak{}T\_\allowbreak{}t hold. The four-top tail contains
at most \texttt{4*(77+\allowbreak{}6)+\allowbreak{}5=\allowbreak{}337} values on each side,\allowbreak{} so
\texttt{d\_\allowbreak{}\allowbreak{}s+\allowbreak{}d\_\allowbreak{}\allowbreak{}t+\allowbreak{}G+\allowbreak{}delta-2K>=\allowbreak{}796-77j}. Four full-core
stars give \texttt{2X+\allowbreak{}4R+\allowbreak{}delta>=\allowbreak{}548}. With the exact
mass and source envelope these two bounds give,\allowbreak{} at j=\allowbreak{}2,\allowbreak{}

\begin{lstlisting}
4I+4tau+delta<=178.
\end{lstlisting}

If any C7\_\allowbreak{}v is active,\allowbreak{} the four-edge activity implication gives
\texttt{t\_\allowbreak{}\allowbreak{}v>=\allowbreak{}253-tau-delta>=\allowbreak{}75},\allowbreak{} contradicting the cap
\texttt{t\_\allowbreak{}\allowbreak{}v<=\allowbreak{}56}. If neither C7 is active,\allowbreak{} those are
the exact two holes; both lower C6 are active,\allowbreak{} making each missing C7
HIGH. Each corresponding full leading has at least
\texttt{r\_\allowbreak{}\allowbreak{}v+\allowbreak{}81} off-target points: the entire shifted
set \texttt{D\_\allowbreak{}\allowbreak{}v union (M\_\allowbreak{}\allowbreak{}v+\allowbreak{}D\_\allowbreak{}\allowbreak{}v)} is HIGH and has at
least \texttt{77+\allowbreak{}(r\_\allowbreak{}\allowbreak{}v+\allowbreak{}4)} distinct points. Thus
\texttt{delta>=\allowbreak{}r\_\allowbreak{}\allowbreak{}s+\allowbreak{}r\_\allowbreak{}\allowbreak{}t+\allowbreak{}162=\allowbreak{}230},\allowbreak{} contradicting the
bound 178. Hence T\_\allowbreak{}s,\allowbreak{}T\_\allowbreak{}t cannot both hold.

An active C4\_\allowbreak{}v forces C5\_\allowbreak{}v activity in branch A. Otherwise that
missing C5 is HIGH above C4; its leading and at least five retained
higher singletons lose at least \texttt{r\_\allowbreak{}\allowbreak{}v+\allowbreak{}81} and
five copies of \texttt{r\_\allowbreak{}\allowbreak{}v+\allowbreak{}5},\allowbreak{} respectively. This
gives \texttt{delta>=\allowbreak{}6r\_\allowbreak{}\allowbreak{}v+\allowbreak{}106}; but
\texttt{I+\allowbreak{}6r\_\allowbreak{}\allowbreak{}v>=\allowbreak{}195} for integral r\_\allowbreak{}v,\allowbreak{} so
\texttt{A>=\allowbreak{}301>298}. The actual-owner grammar follows
from the unchanged literal centre identities and the active lower bound
\texttt{u>-M\_\allowbreak{}\allowbreak{}v}.

These statements reproduce the eight-row actual Z2 cover classification,\allowbreak{}
with no new owner selected for empty or split paired resources.

\hypertarget{z2-betao-simultaneous-hit-cards}{%
\subsection{Z2 beta/O simultaneous-hit
cards}\label{z2-betao-simultaneous-hit-cards}}

Assume beta\_\allowbreak{}v and O\_\allowbreak{}v are active. The double raw gap gives
\texttt{t\_\allowbreak{}\allowbreak{}v<=\allowbreak{}57/\allowbreak{}2},\allowbreak{} and opposite O\_\allowbreak{}o is the first
hole. Let W be the other whole hole. Use the report\textquotesingle s
four retained-resource cards,\allowbreak{} but every forest now selects
\texttt{f=\allowbreak{}max(137-R,\allowbreak{}61-(tau-xi)/\allowbreak{}4)}.

\begin{enumerate}
\def\labelenumi{\arabic{enumi}.}
\item
  If W is not V\_\allowbreak{}v,\allowbreak{}C7\_\allowbreak{}v,\allowbreak{}C7\_\allowbreak{}o,\allowbreak{} the retained unused C7\_\allowbreak{}v has at most 20
  core target hits and all four negative peel images miss. Its native
  deficiency is therefore at least \texttt{69-20+\allowbreak{}4=\allowbreak{}53}.
  The rooted three-edge arm gives

  \texttt{t\_\allowbreak{}\allowbreak{}v>=\allowbreak{}I+\allowbreak{}max(42-2xi+\allowbreak{}tau,\allowbreak{}18+\allowbreak{}7xi/\allowbreak{}4+\allowbreak{}tau/\allowbreak{}4)>=\allowbreak{}I+\allowbreak{}max(42-xi,\allowbreak{}18+\allowbreak{}2xi)>=\allowbreak{}34+\allowbreak{}I}.
\item
  If W=\allowbreak{}V\_\allowbreak{}v,\allowbreak{} its HIGH status makes alpha\_\allowbreak{}v,\allowbreak{}H\_\allowbreak{}v,\allowbreak{}O\_\allowbreak{}v each lose at least
  r\_\allowbreak{}v+\allowbreak{}5 native points. Two retained shortcut edges and the full-core
  support give \texttt{t\_\allowbreak{}\allowbreak{}v>=\allowbreak{}-65+\allowbreak{}I+\allowbreak{}3r\_\allowbreak{}\allowbreak{}v>=\allowbreak{}35}; the
  integer minimum of \texttt{I+\allowbreak{}3r\_\allowbreak{}\allowbreak{}v} is 100.
\item
  If W=\allowbreak{}C7\_\allowbreak{}v,\allowbreak{} it is HIGH above retained C6\_\allowbreak{}v. The three unused
  U\_\allowbreak{}v,\allowbreak{}H\_\allowbreak{}v,\allowbreak{}O\_\allowbreak{}v losses give
  \texttt{t\_\allowbreak{}\allowbreak{}v>=\allowbreak{}max(104-xi,\allowbreak{}80+\allowbreak{}2xi)>=\allowbreak{}96}.
\item
  If W=\allowbreak{}C7\_\allowbreak{}o,\allowbreak{} its HIGH leading loses at least r\_\allowbreak{}o+\allowbreak{}81 points and
  V\_\allowbreak{}o,\allowbreak{}U\_\allowbreak{}o,\allowbreak{}H\_\allowbreak{}o lose at least r\_\allowbreak{}o+\allowbreak{}5 each. Two edges with both leading
  roots yield \texttt{t\_\allowbreak{}\allowbreak{}v>=\allowbreak{}16+\allowbreak{}I+\allowbreak{}4r\_\allowbreak{}\allowbreak{}o>=\allowbreak{}148}.
\end{enumerate}

Every lower bound exceeds 57/2. All roots and unused native losses
belong to distinct original retained resources; no paired capacity is
added in these cards.

\hypertarget{z2-final-six-edge-contradiction}{%
\subsection{Z2 final six-edge
contradiction}\label{z2-final-six-edge-contradiction}}

After that exclusion the same three literal hole families remain:
\texttt{\{O\_\allowbreak{}\allowbreak{}v,\allowbreak{}B\}},\allowbreak{}
\texttt{\{O\_\allowbreak{}\allowbreak{}s,\allowbreak{}O\_\allowbreak{}\allowbreak{}t\}},\allowbreak{} and
\texttt{\{O\_\allowbreak{}\allowbreak{}v,\allowbreak{}H\_\allowbreak{}\allowbreak{}v\}} with their actual forced owners.
Every alpha,\allowbreak{}C5,\allowbreak{}U,\allowbreak{}V,\allowbreak{}C7 singleton is retained. At least one T\_\allowbreak{}u is true.
The unchanged identity
\texttt{alpha\_\allowbreak{}\allowbreak{}o+\allowbreak{}C5\_\allowbreak{}\allowbreak{}u=\allowbreak{}d\_\allowbreak{}\allowbreak{}u-2(M+\allowbreak{}X)-G} gives
\texttt{2X+\allowbreak{}G<=\allowbreak{}54}.

Use the retained six-edge trees and the adaptive support
\texttt{f=\allowbreak{}max(137-R,\allowbreak{}57-(tau-xi)/\allowbreak{}6)}. The actual cover
gives \texttt{2X+\allowbreak{}G>=\allowbreak{}6f-delta}. Hence

\begin{lstlisting}
A=I+tau+delta>=I+tau+6f-54
 >=I+max(414-5xi,342+xi)-54
 >=300+I,
\end{lstlisting}

because the minimum of that maximum is 354,\allowbreak{} at xi=\allowbreak{}12. This contradicts
\texttt{A<=\allowbreak{}298}. Thus every actual Z2 configuration is
excluded,\allowbreak{} with margin 2 at the balanced worst case.

\hypertarget{sharper-source-interface-and-the-later-families}{%
\subsection{Sharper source interface and the later
families}\label{sharper-source-interface-and-the-later-families}}

The independently proved joint triple/fringe source theorem in
\texttt{../\allowbreak{}foundations/\allowbreak{}source\_\allowbreak{}\allowbreak{}triple\_\allowbreak{}\allowbreak{}fringe85.md}
gives \texttt{T<=\allowbreak{}165} and
\texttt{A<=\allowbreak{}445-77j} when there are zero holes. Thus the
current bounds are \texttt{A<=\allowbreak{}214,\allowbreak{}137,\allowbreak{}60} for Z3,\allowbreak{}Z4,\allowbreak{}Z5.
The Z3 argument retains the harmless weaker cap221 for its arithmetic.
Fresh complete-target rank constants are73-r and S+\allowbreak{}r\_\allowbreak{}other-72.

A BAD ordinary centre has loss at least
\texttt{73-max(r\_\allowbreak{}\allowbreak{}s,\allowbreak{}r\_\allowbreak{}\allowbreak{}t)}. The old Z5 initial fixed
ledgers with credits13 or14 do not by themselves justify normality at
cardinality85. The retained-P cards and two empty-P cards below handle
these endpoints directly. Only the alpha-hole card obtains normality,\allowbreak{}
after establishing its actual credit27 independently of normality.

Z4 and Z5 calculations and literal endpoint lists are in
\texttt{zero85\_\allowbreak{}\allowbreak{}z4\_\allowbreak{}\allowbreak{}z5\_\allowbreak{}\allowbreak{}cards.md}. The common Z3 A1
calculation is summarized at the end of this file.

\hypertarget{further-result-the-z3-a0-family-is-excluded}{%
\subsection{Further result: the Z3 A0 family is
excluded}\label{further-result-the-z3-a0-family-is-excluded}}

At j=\allowbreak{}3 the source cap is A\textless=\allowbreak{}221 and the source-only rank rows,\allowbreak{}
freshly derived from deletion of one maximum or two same-sign maxima,\allowbreak{}
are

\begin{lstlisting}
L>=A-57,
e_v>=max(4,A-146),
e_v<=378-A.
\end{lstlisting}

The one-maximum envelope is
\texttt{K<=\allowbreak{}2R+\allowbreak{}c\_\allowbreak{}\allowbreak{}v+\allowbreak{}e\_\allowbreak{}\allowbreak{}o+\allowbreak{}94},\allowbreak{} because its 76-point
remainder has 16 exceptional pairs and reinsertion cost 77. It yields
\texttt{|j\_\allowbreak{}\allowbreak{}s-j\_\allowbreak{}\allowbreak{}t|<=\allowbreak{}237-A-U}. Combining with
\texttt{J>=\allowbreak{}A-67} and \texttt{U>=\allowbreak{}6}
gives the lower e bound. The two-same-sign-maximum envelope is
\texttt{K<=\allowbreak{}2R+\allowbreak{}b\_\allowbreak{}\allowbreak{}v+\allowbreak{}e\_\allowbreak{}\allowbreak{}o+\allowbreak{}167}: 75 remaining points,\allowbreak{} 13
exceptional labels,\allowbreak{} reinsertion 153. Combining with
\texttt{J+\allowbreak{}U<=\allowbreak{}381-A} gives
\texttt{e\_\allowbreak{}\allowbreak{}v<=\allowbreak{}382-A-a\_\allowbreak{}\allowbreak{}v-a\_\allowbreak{}\allowbreak{}o-b\_\allowbreak{}\allowbreak{}o<=\allowbreak{}378-A}.

The common j3 activity grammar still follows without any new raw-gap
strengthening. Two cap predicates violate
\texttt{4I+\allowbreak{}4tau+\allowbreak{}delta<=\allowbreak{}640-231j=\allowbreak{}-53}. The five-edge
physical-arm bound is now
\texttt{X\_\allowbreak{}\allowbreak{}o+\allowbreak{}2X\_\allowbreak{}\allowbreak{}v+\allowbreak{}G>=\allowbreak{}305-tau-delta}; active C6 gives
\texttt{t\_\allowbreak{}\allowbreak{}v>=\allowbreak{}85+\allowbreak{}I>56}. The four-edge bound gives
active C7 \texttt{t\_\allowbreak{}\allowbreak{}v>=\allowbreak{}32+\allowbreak{}I>57/\allowbreak{}2}. A HIGH C6 above
retained C5 has native and higher-singleton losses at least
\texttt{4r\_\allowbreak{}\allowbreak{}v+\allowbreak{}96}; its sum with I has minimum 228,\allowbreak{}
exceeding the source cap 221. These prove the same exhaustive two A
families: A0 with holes B,\allowbreak{}O\_\allowbreak{}s,\allowbreak{}O\_\allowbreak{}t and pair1 owner s; A1 with holes
O\_\allowbreak{}s,\allowbreak{}O\_\allowbreak{}t,\allowbreak{}H\_\allowbreak{}s and beta owner s,\allowbreak{} pair1 owner t.

For A0 use the five endpoint resources
C4\_\allowbreak{}s,\allowbreak{}alpha\_\allowbreak{}s,\allowbreak{}alpha\_\allowbreak{}t,\allowbreak{}C5\_\allowbreak{}t,\allowbreak{}H\_\allowbreak{}s with weights 1,\allowbreak{}2,\allowbreak{}1,\allowbreak{}1,\allowbreak{}1. The
complete-target lower endpoint is now
\texttt{Z>=\allowbreak{}73-r\_\allowbreak{}\allowbreak{}v-k-b}; the upper is
\texttt{Z<=\allowbreak{}S\_\allowbreak{}\allowbreak{}v+\allowbreak{}r\_\allowbreak{}\allowbreak{}o-72+\allowbreak{}k+\allowbreak{}b}. Thus the endpoint
identity has constant 233. Use six retained edges with
\texttt{f=\allowbreak{}61-(tau-xi)/\allowbreak{}4},\allowbreak{} giving gap 2X+\allowbreak{}G. The source
mass constant 1239 and the forest baseline 12f=\allowbreak{}732-3tau+\allowbreak{}3xi cancel to
give the same master constant:

\begin{lstlisting}
T_r+4r+tau_t+2L-3tau+3xi_s+7xi_t <= 290+delta+Bstar,
Q3=T_r+4r+3I >=626 (r=28,29).
\end{lstlisting}

Use the common eight O HIGH fixed labels,\allowbreak{} so epsilon=\allowbreak{}delta-8. Define
heavy selected clips as k\textgreater=\allowbreak{}72. For k\textless=\allowbreak{}71 the error is
at most both xi\_\allowbreak{}source and k,\allowbreak{} by
\texttt{(73-k)b<=\allowbreak{}tau\_\allowbreak{}\allowbreak{}source} and
\texttt{tau+\allowbreak{}epsilon<=\allowbreak{}213-I}; at k=\allowbreak{}71 equality in
\texttt{k(74-k)>=\allowbreak{}213} is allowed. Heavy tail error is
at most xi\_\allowbreak{}source+\allowbreak{}e\_\allowbreak{}source-4,\allowbreak{} including k=\allowbreak{}72 where the stronger
core-maximum error xi\_\allowbreak{}source is available.

The five nonexceptional scalar finishes (the r-polynomial is the new Q3)
are

\begin{longtable}[]{@{}llr@{}}
\toprule\noalign{}
Endpoint pattern & Upper inequality & Lower bound \\
\midrule\noalign{}
\endhead
\bottomrule\noalign{}
\endlastfoot
alpha\_\allowbreak{}s light,\allowbreak{} none heavy & Q3+\allowbreak{}tau\_\allowbreak{}t+\allowbreak{}2delta\textless=\allowbreak{}625 & 642 \\
alpha\_\allowbreak{}s light,\allowbreak{} one heavy c & Q3+\allowbreak{}tau\_\allowbreak{}t+\allowbreak{}2delta+\allowbreak{}c\textless=\allowbreak{}781 & 858 \\
alpha\_\allowbreak{}s light,\allowbreak{} two heavy H & Q3+\allowbreak{}tau\_\allowbreak{}t+\allowbreak{}2delta+\allowbreak{}H\textless=\allowbreak{}945 & 1074 \\
alpha\_\allowbreak{}s heavy,\allowbreak{} other source t & Q3+\allowbreak{}tau\_\allowbreak{}t+\allowbreak{}2delta+\allowbreak{}c\textless=\allowbreak{}866 &
1002 \\
alpha\_\allowbreak{}s heavy,\allowbreak{} other C4\_\allowbreak{}s & Q3+\allowbreak{}tau\_\allowbreak{}t+\allowbreak{}2delta+\allowbreak{}c\textless=\allowbreak{}948 & 1002 \\
\end{longtable}

For the last two lower bounds it is enough to allow clips 72,\allowbreak{} so
delta\textgreater=\allowbreak{}152 and c\textgreater=\allowbreak{}72. Each table row follows from
the displayed master and,\allowbreak{} respectively,\allowbreak{}
\texttt{2L>=\allowbreak{}2A-114},\allowbreak{}
\texttt{2L-e\_\allowbreak{}\allowbreak{}v>=\allowbreak{}A-53},\allowbreak{}
\texttt{2L-e\_\allowbreak{}\allowbreak{}v-e\_\allowbreak{}\allowbreak{}w>=\allowbreak{}0},\allowbreak{}
\texttt{e\_\allowbreak{}\allowbreak{}t>=\allowbreak{}A-146},\allowbreak{} and
\texttt{2e\_\allowbreak{}\allowbreak{}t-e\_\allowbreak{}\allowbreak{}s>=\allowbreak{}3A-670}.

The sole-heavy-alpha row needs a new fixed ledger. If its clip is
exactly 72,\allowbreak{} the endpoint is the core maximum and has error at most
xi\_\allowbreak{}s; absorb it radially and use the first row,\allowbreak{} which is already
contradictory. If its clip is at least 73,\allowbreak{} alpha\_\allowbreak{}s\textless=\allowbreak{}0. The
eleven fixed labels

\begin{lstlisting}
U5_s,U6_s,U7_s,U8_s,V5_s,V6_s,V7_s,V8_s,F3_s,F4_s,F5_s
\end{lstlisting}

are LOW. The U labels are negative because alpha\_\allowbreak{}s\textless=\allowbreak{}0. The V
labels are negative because beta\_\allowbreak{}s=\allowbreak{}alpha\_\allowbreak{}s-d\_\allowbreak{}s\textless0. The F
labels are negative because the leading bound and
\texttt{K-delta>=\allowbreak{}89+\allowbreak{}I+\allowbreak{}tau} imply
\texttt{d\_\allowbreak{}\allowbreak{}s>g\_\allowbreak{}\allowbreak{}s}. Those 11 labels touch exactly 9
baseline groups: U7 has no baseline group,\allowbreak{} and U8/F5 share one. Together
with the eight O HIGH labels,\allowbreak{} the compatible credit is 10. Replacing
delta-8 by delta-10 lowers the master right-hand constant from 290 to
286. Cancel the heavy alpha error and use
\texttt{2e\_\allowbreak{}\allowbreak{}t>=\allowbreak{}2A-292}; the resulting exact row is

\begin{lstlisting}
Q3+tau_t+2delta<=791.
\end{lstlisting}

But the native clip requires \texttt{delta>=\allowbreak{}10+\allowbreak{}73=\allowbreak{}83},\allowbreak{}
so its left-hand side is at least \texttt{626+\allowbreak{}166=\allowbreak{}792}.
This strictly excludes the last A0 pattern. It is a new argument,\allowbreak{} not an
automatic translation of the old report.

\hypertarget{z3-a1-is-excluded-by-the-actual-physical-side-union}{%
\subsection{Z3 A1 is excluded by the actual physical-side
union}\label{z3-a1-is-excluded-by-the-actual-physical-side-union}}

The full derivation and exact four-pattern checker are
\texttt{../\allowbreak{}common/\allowbreak{}z3a1\_\allowbreak{}\allowbreak{}side\_\allowbreak{}\allowbreak{}union.md} and
\texttt{../\allowbreak{}common/\allowbreak{}check\_\allowbreak{}\allowbreak{}z3a1\_\allowbreak{}\allowbreak{}side\_\allowbreak{}\allowbreak{}union.py}. The
actual holes are O\_\allowbreak{}s,\allowbreak{}O\_\allowbreak{}t,\allowbreak{}H\_\allowbreak{}s; beta is sole on s and pair1 sole on t.
The missing H\_\allowbreak{}s is HIGH because H\_\allowbreak{}s=\allowbreak{}beta\_\allowbreak{}s+\allowbreak{}2d\_\allowbreak{}s and beta\_\allowbreak{}s is
active. Twelve compatible HIGH labels give epsilon=\allowbreak{}delta-12. Grouping
the leading and its seven actual native resources on s gives
\texttt{S\_\allowbreak{}\allowbreak{}s>=\allowbreak{}K+\allowbreak{}539-epsilon}.

Let H be the native deficiency of active H\_\allowbreak{}t,\allowbreak{} so H\textless=\allowbreak{}76. Lower
P\_\allowbreak{}t and upper H\_\allowbreak{}t are the only endpoints. Their actual clips and the
two separate forest rebates,\allowbreak{} on one chosen actual source with
\texttt{2tau+\allowbreak{}6f>=\allowbreak{}398},\allowbreak{} give

\begin{lstlisting}
M+M_s+6f-H/2<=47+2delta+E,
E<=xi_s+xi_t+p(e_t-4)+h(e_s-4), p,h in {0,1}.
\end{lstlisting}

Use L\textgreater=\allowbreak{}A-57 and e\_\allowbreak{}s\textgreater=\allowbreak{}A-146. For
\texttt{Z=\allowbreak{}r\_\allowbreak{}\allowbreak{}s+\allowbreak{}2I+\allowbreak{}2tau+\allowbreak{}xi\_\allowbreak{}\allowbreak{}s+\allowbreak{}6f-H/\allowbreak{}2},\allowbreak{} the four endpoint
positions yield upper bounds182,\allowbreak{}267,\allowbreak{}A+\allowbreak{}32,\allowbreak{}A+\allowbreak{}117. Even with A\textless=\allowbreak{}221
their largest bound is338,\allowbreak{} whereas
\texttt{Z>=\allowbreak{}398+\allowbreak{}34-38=\allowbreak{}394}. This excludes all actual
light/heavy positions of both endpoints; no endpoint or error assumption
is made on the beta translate. Together with the preceding A0 proof it
completes Z3.

\clearpage
\hypertarget{excluding-the-fresh-zero-hole-z3-a1-form}{%
\section{Excluding the fresh zero-hole Z3 A1
form}\label{excluding-the-fresh-zero-hole-z3-a1-form}}

This is a conditional local proof for a cardinality-85 configuration
after a valid integral primitive-core structural reduction. It uses only
the actual Z3 A1 retention form and the common source/rank/forest
lemmas. No scalar cell is assumed to be an actual set.

\hypertarget{actual-form-and-source-quantities}{%
\subsection{Actual form and source
quantities}\label{actual-form-and-source-quantities}}

The original source has 77 points and the core has 69 points. Write

\[
r_s+\allowbreak{}r_t=\allowbreak{}68,\allowbreak{}\quad I=\allowbreak{}(r_s-34)^2,\allowbreak{}\quad
\tau=\allowbreak{}\tau_s+\allowbreak{}\tau_t,\allowbreak{}\quad R_v=\allowbreak{}r_v+\allowbreak{}\xi_v,\allowbreak{}\quad R=\allowbreak{}68+\allowbreak{}\xi_s+\allowbreak{}\xi_t.
\]

The three missing ordinary whole units are exactly \(O_s,\allowbreak{}O_t,\allowbreak{}H_s\). The
beta pair is retained and has sole actual owner \(s\); pair1
\(\{C4_s,\allowbreak{}C4_t\}\) has sole actual owner \(t\). Every unlisted ordinary
singleton is retained. Both leading resources are active and all outer
coordinates are integral,\allowbreak{} by the preceding actual-activity reduction.

Let \(e_s,\allowbreak{}e_t\) be the fourth peel offsets,\allowbreak{} \(L=\allowbreak{}e_s+\allowbreak{}e_t\),\allowbreak{} \(M=\allowbreak{}R+\allowbreak{}L\),\allowbreak{}
and \(M_s=\allowbreak{}r_s+\allowbreak{}\xi_s+\allowbreak{}e_s\). Put \(A=\allowbreak{}I+\allowbreak{}\tau+\allowbreak{}\delta\). The following
already proved source bounds suffice:

\[
A\le221,\allowbreak{}\qquad L\ge A-57,\allowbreak{}\qquad e_s\ge A-146.
\]

The stronger new cap \(A\le216\) is unnecessary here.

\hypertarget{fixed-credit-and-a-whole-physical-side}{%
\subsection{Fixed credit and a whole physical
side}\label{fixed-credit-and-a-whole-physical-side}}

An empty integral ordinary centre on a complete target is either at most
\(-M_v\) or at least \(S_v+\allowbreak{}1\). Since \(\beta_s\) is active and

\[
H_s=\allowbreak{}\beta_s+\allowbreak{}2d_s,\allowbreak{}\qquad d_s>0,\allowbreak{}
\]

the missing \(H_s\) satisfies \(H_s\ge S_s+\allowbreak{}1\). The missing \(O_s,\allowbreak{}O_t\)
are likewise HIGH. The eight \(O_v+\allowbreak{}v_i\) labels and the three
\(H_s+\allowbreak{}s_i\),\allowbreak{} \(i=\allowbreak{}2,\allowbreak{}3,\allowbreak{}4\),\allowbreak{} are HIGH. A twelfth is \(\alpha_s+\allowbreak{}s_2\): beta
activity gives \(s_2>d_s\),\allowbreak{} so \(\alpha_s+\allowbreak{}s_2=\allowbreak{}H_s-d_s+\allowbreak{}s_2>H_s\). Thus
the compatible fixed credit is 12 and \(\epsilon=\allowbreak{}\delta-12\ge0\).

On target \(s\),\allowbreak{} group the full leading resource with these seven actual
ordinary native resources:

\[
\beta_s,\allowbreak{}\ \alpha_s,\allowbreak{}\ U_s,\allowbreak{}\ V_s,\allowbreak{}\ C5_s,\allowbreak{}\ C6_s,\allowbreak{}\ C7_s.
\]

The beta capacity is licensed because its opposite paired contribution
is empty. The group has original capacity \(K+\allowbreak{}7\cdot77=\allowbreak{}K+\allowbreak{}539\); its
capacity deficiency is at most \(\epsilon\). Its actual union lies in
the complete target,\allowbreak{} hence

\[
S_s\ge K+\allowbreak{}539-\epsilon.\tag{1}
\]

This is an actual union bound,\allowbreak{} not an assumption that the seven native
resources are disjoint.

\hypertarget{two-endpoints-and-two-forest-inequalities}{%
\subsection{Two endpoints and two forest
inequalities}\label{two-endpoints-and-two-forest-inequalities}}

Let \(P=\allowbreak{}C4_t\),\allowbreak{} let \(\Gamma_P,\allowbreak{}\Gamma_{\beta_s}\) be the whole-unit
native deficiencies,\allowbreak{} and put

\[
B=\allowbreak{}\Gamma_P+\allowbreak{}\Gamma_{\beta_s},\allowbreak{}\qquad H=\allowbreak{}\Gamma_{H_t}.
\]

The active 77-point singleton \(H_t\) has \(0\le H\le76\). On complete
targets its native loss is exactly its lower plus upper clips. Use the
lower endpoint of \(P\) and the upper endpoint of \(H_t\):

\[
P\ge73-r_t-l_P-b_P,\allowbreak{}\qquad
H_t\le S_t+\allowbreak{}r_s-72+\allowbreak{}h_H+\allowbreak{}b_H.
\]

Their error sources are respectively \(t,\allowbreak{}s\). Define
\(E_c=\allowbreak{}l_P+\allowbreak{}h_H\le B+\allowbreak{}H\) and \(E=\allowbreak{}b_P+\allowbreak{}b_H\).

Choose one actual full or trimmed core subset attaining

\[
f=\allowbreak{}\max\left(69-\xi,\allowbreak{}\ 65-\frac{\tau-\xi}{2},\allowbreak{}\ 57-\frac{\tau-\xi}{6}\right),\allowbreak{}
\qquad\xi=\allowbreak{}\xi_s+\allowbreak{}\xi_t.
\]

In particular,\allowbreak{}

\[
f\ge\frac{199-\tau}{3},\allowbreak{}\qquad 2\tau+\allowbreak{}6f\ge398.\tag{2}
\]

The ten retained-resource edges are

\[
\begin{array}{ll}
s:& L_s-V_s-U_s,\allowbreak{}\quad L_s-C7_s-C6_s-C5_s,\allowbreak{}\\
t:& L_t-V_t-U_t-H_t,\allowbreak{}\quad L_t-C7_t-C6_t.
\end{array}
\]

Their gaps sum to \(2Q\),\allowbreak{} with \(Q=\allowbreak{}X_s+\allowbreak{}X_t+\allowbreak{}g_s+\allowbreak{}g_t+\allowbreak{}h_t\). The forest
omits both paired endpoint resources,\allowbreak{} so the actual grouping budget
gives

\[
Q\ge5f-\frac{\epsilon-B}{2}.
\]

The two \(X_s\) star edges \(L_s-V_s\) and \(L_t-C7_t\) omit these two
paired resources and also \(H_t\),\allowbreak{} giving

\[
X_s\ge f-\frac{\epsilon-B-H}{2}.
\]

Both are valid statements about the same unchanged deficit. Adding their
inequalities gives

\[
Q+\allowbreak{}X_s\ge6f-\epsilon+\allowbreak{}B+\allowbreak{}\frac H2.\tag{3}
\]

\hypertarget{use-the-missing-h-directly}{%
\subsection{Use the missing H
directly}\label{use-the-missing-h-directly}}

The literal centre identity is

\[
H_t=\allowbreak{}2d_s+\allowbreak{}\beta_s+\allowbreak{}P+\allowbreak{}2s_4+\allowbreak{}t_4+\allowbreak{}g_s+\allowbreak{}g_t+\allowbreak{}h_t.
\]

Since \(\beta_s+\allowbreak{}2d_s=\allowbreak{}H_s\ge S_s+\allowbreak{}1\),\allowbreak{} the two endpoint inequalities imply

\[
S_t\ge S_s+\allowbreak{}M+\allowbreak{}M_s+\allowbreak{}(Q+\allowbreak{}X_s)+\allowbreak{}78-E_c-E.
\]

Use (1) and the exact source mass \(S_s+\allowbreak{}S_t=\allowbreak{}2K+\allowbreak{}1239-\delta\). The result
is

\[
M+\allowbreak{}M_s+\allowbreak{}(Q+\allowbreak{}X_s)\le59+\allowbreak{}\delta+\allowbreak{}E_c+\allowbreak{}E.
\]

Apply (3) and \(E_c\le B+\allowbreak{}H\). The paired-loss variable \(B\) cancels and
the unchanged residual deficit gives the crucial inequality

\[
M+\allowbreak{}M_s+\allowbreak{}6f-\frac H2\le47+\allowbreak{}2\delta+\allowbreak{}E.\tag{4}
\]

No endpoint for beta is used,\allowbreak{} so all positions of its actual native
translate are included automatically.

\hypertarget{all-four-endpoint-positions}{%
\subsection{All four endpoint
positions}\label{all-four-endpoint-positions}}

For a selected clip at most 72,\allowbreak{} the core rank bound gives
\(b_P\le\xi_t\) or \(b_H\le\xi_s\). A selected clip at least 73 is a
peeled rank and gives respectively

\[
b_P\le\xi_t+\allowbreak{}e_t-4,\allowbreak{}\qquad b_H\le\xi_s+\allowbreak{}e_s-4.
\]

Let \(p,\allowbreak{}h\in\{0,\allowbreak{}1\}\) record the latter two alternatives. Thus

\[
E\le\xi_s+\allowbreak{}\xi_t+\allowbreak{}p(e_t-4)+\allowbreak{}h(e_s-4).
\]

Put

\[
Z=\allowbreak{}r_s+\allowbreak{}2I+\allowbreak{}2\tau+\allowbreak{}\xi_s+\allowbreak{}6f-\frac H2.
\]

Substitute \(M+\allowbreak{}M_s=\allowbreak{}68+\allowbreak{}r_s+\allowbreak{}2\xi_s+\allowbreak{}\xi_t+\allowbreak{}2e_s+\allowbreak{}e_t\) into (4). The offset
rows \(L\ge A-57\),\allowbreak{} \(e_s\ge A-146\) and \(A=\allowbreak{}I+\allowbreak{}\tau+\allowbreak{}\delta\) give:

\begin{longtable}[]{@{}lll@{}}
\toprule\noalign{}
Lower P endpoint & Upper Ht endpoint & Consequence \\
\midrule\noalign{}
\endhead
\bottomrule\noalign{}
\endlastfoot
core rank & core rank & \(Z\le182\) \\
peeled rank & core rank & \(Z\le267\) \\
core rank & peeled rank & \(Z\le A+\allowbreak{}32\le253\) \\
peeled rank & peeled rank & \(Z\le A+\allowbreak{}117\le338\) \\
\end{longtable}

All these cancellations are exact. For the four rows the residual offset
expressions after subtracting the endpoint error are respectively
\(L+\allowbreak{}e_s\),\allowbreak{} \(2e_s+\allowbreak{}4\),\allowbreak{} \(L+\allowbreak{}4\),\allowbreak{} \(e_s+\allowbreak{}8\).

But (2),\allowbreak{} \(H\le76\),\allowbreak{} \(\xi_s\ge0\),\allowbreak{} and the integer identity
\(r_s+\allowbreak{}2(r_s-34)^2\ge34\) give

\[
Z\ge398+\allowbreak{}34-38=\allowbreak{}394.
\]

Every upper bound is strictly smaller. Therefore every actual Z3 A1
configuration is impossible,\allowbreak{} in either physical owner orientation,\allowbreak{}
including every heavy or light position of all three original endpoint
resources.

\clearpage
\hypertarget{fresh-z4z5-reductions-at-cardinality-85}{%
\section{Fresh Z4/Z5 reductions at cardinality
85}\label{fresh-z4z5-reductions-at-cardinality-85}}

Conditional structural notation and basic cover are in
zero85\_\allowbreak{}reduction.md. The stronger source theorem now supplies
A\textless=\allowbreak{}137 for Z4 and A\textless=\allowbreak{}60 for Z5. Below proofs sometimes
keep the weaker Z4 cap143,\allowbreak{} so they remain valid.

All forests/groupings are original full resources. Repeated inequalities
may use separate copies of the same deficit; endpoint rebates are made
only when the endpoint resources are disjoint from that particular
forest/group.

\hypertarget{z4-source-and-classification}{%
\subsection{Z4 source and
classification}\label{z4-source-and-classification}}

Fresh rows: S=\allowbreak{}2K+\allowbreak{}1162-delta,\allowbreak{} L\textgreater=\allowbreak{}A+\allowbreak{}20,\allowbreak{}
e\_\allowbreak{}v\textgreater=\allowbreak{}max(4,\allowbreak{}A-69),\allowbreak{} e\_\allowbreak{}v\textless=\allowbreak{}301-A; f=\allowbreak{}61-(tau-xi)/4. A
complete-target rank endpoint has lower constant73-r\_\allowbreak{}v and upper
constant S\_\allowbreak{}v+\allowbreak{}r\_\allowbreak{}o-72. Light clip\textless73 has
error\textless=\allowbreak{}xi\_\allowbreak{}source; a heavy clip\textgreater=\allowbreak{}73 has
error\textless=\allowbreak{}xi\_\allowbreak{}source+\allowbreak{}e\_\allowbreak{}source-4. At most one heavy endpoint is
possible in any selected native loss
budget\textless=\allowbreak{}delta-8\textless=\allowbreak{}135. For light outside-hull clips k,\allowbreak{}
b\textless=\allowbreak{}k follows from tau\_\allowbreak{}v\textless=\allowbreak{}135-k and
k(74-k)\textgreater=\allowbreak{}144 on 28\textless=\allowbreak{}k\textless=\allowbreak{}72.

Both O are absent. Two cap predicates already contradict the
five-/four-top tail and source depth. If beta\_\allowbreak{}v,\allowbreak{}O\_\allowbreak{}v were
simultaneously active,\allowbreak{} their opposite O and their own C7,\allowbreak{}C6 are holes,\allowbreak{}
since active C7 gives t\_\allowbreak{}v\textgreater=\allowbreak{}110+\allowbreak{}I while O\textquotesingle s
raw cap is12. The negative physical-v arm L\_\allowbreak{}o-C7\_\allowbreak{}o-C6\_\allowbreak{}o is fully
retained. A retained positive star L\_\allowbreak{}v-V\_\allowbreak{}v plus that arm,\allowbreak{} rooted at
L\_\allowbreak{}v,\allowbreak{} gives t\_\allowbreak{}v\textgreater=\allowbreak{}K-M+\allowbreak{}3f-delta+\allowbreak{}4. With full-core support this
is \textgreater=\allowbreak{}147-xi+\allowbreak{}I; with the four-rank support it is
\textgreater=\allowbreak{}123+\allowbreak{}2xi+\allowbreak{}I,\allowbreak{} so it exceeds the double-gap cap27/2. If V\_\allowbreak{}v is
the fourth hole,\allowbreak{} it is HIGH; alpha\_\allowbreak{}v,\allowbreak{}H\_\allowbreak{}v,\allowbreak{}O\_\allowbreak{}v are three retained HIGH
resources. Use the two shortcuts L\_\allowbreak{}v-U\_\allowbreak{}v,\allowbreak{}L\_\allowbreak{}o-C7\_\allowbreak{}o,\allowbreak{} root L\_\allowbreak{}v,\allowbreak{} and
their three unused losses. Full-core overlap gives
t\_\allowbreak{}v\textgreater=\allowbreak{}93+\allowbreak{}I+\allowbreak{}3r\_\allowbreak{}v\textgreater=\allowbreak{}193. Thus beta/O coexistence is
impossible. The remaining one-active-O grammar forces a retained
beta/H/C7 triple on the opposite sign,\allowbreak{} again impossible. Both O are
holes and eight HIGH fixed labels give epsilon=\allowbreak{}delta-8.

The exact A-E/A-H/B retention table follows from the actual cross chains
and beta/H incompatibilities. Put
\texttt{N=\allowbreak{}\{U\_\allowbreak{}\allowbreak{}s,\allowbreak{}U\_\allowbreak{}\allowbreak{}t,\allowbreak{}V\_\allowbreak{}\allowbreak{}s,\allowbreak{}V\_\allowbreak{}\allowbreak{}t,\allowbreak{}alpha\_\allowbreak{}\allowbreak{}s,\allowbreak{}alpha\_\allowbreak{}\allowbreak{}t\}}
and \texttt{W=\allowbreak{}N union \{P,\allowbreak{}H\_\allowbreak{}\allowbreak{}s,\allowbreak{}H\_\allowbreak{}\allowbreak{}t,\allowbreak{}C5\_\allowbreak{}\allowbreak{}s,\allowbreak{}C5\_\allowbreak{}\allowbreak{}t\}}.
Every unlisted unit is retained.

\begin{longtable}[]{@{}
  >{\raggedright\arraybackslash}p{(\columnwidth - 4\tabcolsep) * \real{0.3333}}
  >{\raggedright\arraybackslash}p{(\columnwidth - 4\tabcolsep) * \real{0.3333}}
  >{\raggedright\arraybackslash}p{(\columnwidth - 4\tabcolsep) * \real{0.3333}}@{}}
\toprule\noalign{}
\begin{minipage}[b]{\linewidth}\raggedright
Type
\end{minipage} & \begin{minipage}[b]{\linewidth}\raggedright
Exact whole holes
\end{minipage} & \begin{minipage}[b]{\linewidth}\raggedright
Actual paired status
\end{minipage} \\
\midrule\noalign{}
\endhead
\bottomrule\noalign{}
\endlastfoot
A-E & B,\allowbreak{}O\_\allowbreak{}s,\allowbreak{}O\_\allowbreak{}t,\allowbreak{}Z; Z in W & P empty when Z=\allowbreak{}P; owner opposite v when
Z=\allowbreak{}C5\_\allowbreak{}v; either actual owner otherwise \\
A-H & O\_\allowbreak{}s,\allowbreak{}O\_\allowbreak{}t,\allowbreak{}H\_\allowbreak{}v,\allowbreak{}Z; Z in W except H\_\allowbreak{}v & B sole on v; P empty when
Z=\allowbreak{}P,\allowbreak{} owner v when Z=\allowbreak{}C5\_\allowbreak{}opposite,\allowbreak{} owner opposite otherwise \\
B & O\_\allowbreak{}s,\allowbreak{}O\_\allowbreak{}t,\allowbreak{}X\_\allowbreak{}r,\allowbreak{}Y\_\allowbreak{}opposite; X\_\allowbreak{}r in \{beta\_\allowbreak{}r,\allowbreak{}H\_\allowbreak{}r\},\allowbreak{} Y\_\allowbreak{}opposite in
\{beta\_\allowbreak{}opposite,\allowbreak{}H\_\allowbreak{}opposite\} & Both paired units retained,\allowbreak{} P sole on
r \\
\end{longtable}

For completeness the two-beta-hole B exclusion has a fresh scalar check.
Suppose the beta opposite the active pair1 owner were HIGH. The exact
centre/leading bounds give
\texttt{2S\_\allowbreak{}\allowbreak{}v+\allowbreak{}K+\allowbreak{}3g\_\allowbreak{}\allowbreak{}v+\allowbreak{}6X\_\allowbreak{}\allowbreak{}v<=\allowbreak{}M\_\allowbreak{}\allowbreak{}o+\allowbreak{}1161+\allowbreak{}delta}. All the
relevant ordinary vertices are retained,\allowbreak{} so the actual three-edge arm
gives \texttt{2X\_\allowbreak{}\allowbreak{}v+\allowbreak{}g\_\allowbreak{}\allowbreak{}v>=\allowbreak{}189-3tau/\allowbreak{}4-delta}. Use
\texttt{K-M>=\allowbreak{}R-132+\allowbreak{}A},\allowbreak{}
\texttt{S\_\allowbreak{}\allowbreak{}v>=\allowbreak{}T\_\allowbreak{}\allowbreak{}r+\allowbreak{}4r+\allowbreak{}10+\allowbreak{}tau\_\allowbreak{}\allowbreak{}v},\allowbreak{}
\texttt{M\_\allowbreak{}\allowbreak{}v>=\allowbreak{}r+\allowbreak{}4},\allowbreak{} the conservative
\texttt{A<=\allowbreak{}143},\allowbreak{} and R\textgreater=\allowbreak{}68. The necessary
consequence is
\texttt{2T\_\allowbreak{}\allowbreak{}r+\allowbreak{}9r+\allowbreak{}4I+\allowbreak{}7tau/\allowbreak{}4+\allowbreak{}2tau\_\allowbreak{}\allowbreak{}v<=\allowbreak{}1063}. But the
integral minimum of \texttt{2T\_\allowbreak{}\allowbreak{}r+\allowbreak{}9r+\allowbreak{}4(r-34)\^2}
is1192,\allowbreak{} at r=\allowbreak{}26. Hence that beta is LOW. On the pair1-owner side,\allowbreak{} HIGH
beta and retained H would give2d\textless=\allowbreak{}M\_\allowbreak{}o-1,\allowbreak{} whereas actual pair1
activity gives d\textgreater M. That beta is LOW too. Their sum cannot
then have branch-B\textquotesingle s positive sign,\allowbreak{} excluding both beta
holes. The literal dispatch is: A-E retained P uses F for the four U,\allowbreak{}V
holes or H\_\allowbreak{}t,\allowbreak{} E-H for H\_\allowbreak{}s,\allowbreak{} E-alpha for either alpha,\allowbreak{} E-C5 for C5\_\allowbreak{}t
after naming its owner s. A-H uses F for U,\allowbreak{}V or empty P,\allowbreak{} O for
opposite-owner choices H\_\allowbreak{}t,\allowbreak{}C5\_\allowbreak{}s,\allowbreak{}alpha\_\allowbreak{}s,\allowbreak{}alpha\_\allowbreak{}t,\allowbreak{} and S for C5\_\allowbreak{}t.
The three B cases HH,\allowbreak{}beta\_\allowbreak{}rH\_\allowbreak{}opposite,\allowbreak{}H\_\allowbreak{}rbeta\_\allowbreak{}opposite use O,\allowbreak{}O,\allowbreak{}S
respectively. The only ownerless residual is A-E with Z=\allowbreak{}P.

No empty pair is assigned an owner. Exact endpoint/forest predicates are
in Appendix D of the retained report,\allowbreak{} labels sec:z4-cardF,\allowbreak{}
sec:z4-cardO,\allowbreak{}sec:z4-cardS,\allowbreak{}sec:z4-cardEH,\allowbreak{}sec:z4-cardEA,\allowbreak{} sec:z4-cardEC and
the immediately following dispatch (source lines 1962--2158). The
calculations below use these literal lists with fresh capacity77 and
source mass1162; their formulas are derived again from actual target
union bounds.

\hypertarget{z4-card-f}{%
\subsection{Z4 Card F}\label{z4-card-f}}

The five weighted endpoint identity has constant233. Its six-edge forest
gives the fresh master

\texttt{T\_\allowbreak{}\allowbreak{}r+\allowbreak{}4r+\allowbreak{}tau\_\allowbreak{}\allowbreak{}res+\allowbreak{}2L-3tau+\allowbreak{}2e <=\allowbreak{}229+\allowbreak{}delta+\allowbreak{}B}.

For all light endpoints B\textless=\allowbreak{}2epsilon. For one heavy endpoint of
weight q=\allowbreak{}1,\allowbreak{}2 and clip c\textgreater=\allowbreak{}73,\allowbreak{}
B\textless=\allowbreak{}2epsilon-qc+\allowbreak{}q*tau+\allowbreak{}q*e\_\allowbreak{}v. Since e\_\allowbreak{}v\textless=\allowbreak{}L,\allowbreak{} source
L\textgreater=\allowbreak{}A+\allowbreak{}20 and A\textless=\allowbreak{}143 give respectively

\begin{lstlisting}
Q3+tau_res+4e <=332,
Q_(3+q)+tau_res+q*delta+4e+q*c <=332+306q.
\end{lstlisting}

The integer minima are
Q3\textgreater=\allowbreak{}626,\allowbreak{}Q4\textgreater=\allowbreak{}649,\allowbreak{}Q5\textgreater=\allowbreak{}665. At e=\allowbreak{}8 the
light lower bound658 exceeds332; q1 lower835 exceeds638; q2 lower1005
exceeds944. All Card F predicates close.

\hypertarget{z4-card-o}{%
\subsection{Z4 Card O}\label{z4-card-o}}

Use the common actual opposite-owner predicate; e=\allowbreak{}12,\allowbreak{} epsilon=\allowbreak{}delta-12.
Source-side s has at least six native ordinary resources. The
independent common-agent route uses H\_\allowbreak{}s HIGH and
V\_\allowbreak{}t=\allowbreak{}H\_\allowbreak{}s+\allowbreak{}P\_\allowbreak{}t+\allowbreak{}2a+\allowbreak{}b+\allowbreak{}g\_\allowbreak{}s. Together with the physical side union and
four-star/double-arm forest it proves all four endpoint-heavy flags; see
the common-agent proof file. A direct new endpoint master is also

\texttt{T\_\allowbreak{}\allowbreak{}r+\allowbreak{}4r+\allowbreak{}5xi\_\allowbreak{}\allowbreak{}s+\allowbreak{}tau\_\allowbreak{}\allowbreak{}s+\allowbreak{}R+\allowbreak{}L+\allowbreak{}2e\_\allowbreak{}\allowbreak{}s+\allowbreak{}10f+\allowbreak{}60+\allowbreak{}4B+\allowbreak{}(7/\allowbreak{}2)H <=\allowbreak{}1006+\allowbreak{}4delta+\allowbreak{}Eerr}.

With f=\allowbreak{}max(137-R,\allowbreak{}65-(tau-xi)/2,\allowbreak{}57-(tau-xi)/6),\allowbreak{}
4tau+\allowbreak{}10f\textgreater=\allowbreak{}1990/3. Heavy source-s gives upper1091.5,\allowbreak{} heavy
source-t gives1142,\allowbreak{} all-light A\textless=\allowbreak{}73 gives1137,\allowbreak{} all-light
A\textgreater=\allowbreak{}73 gives1207,\allowbreak{} against respective minima1380+\allowbreak{}1/3
or1312+\allowbreak{}1/3. Thus Card O is closed independently.

\hypertarget{z4-card-s-stronger-physical-union-proof}{%
\subsection{Z4 Card S: stronger physical union
proof}\label{z4-card-s-stronger-physical-union-proof}}

Same owner of beta\_\allowbreak{}s and P\_\allowbreak{}s forces C5\_\allowbreak{}t inactive. Target t has six
actual ordinary native resources H,\allowbreak{}U,\allowbreak{}V,\allowbreak{}C6,\allowbreak{}C7,\allowbreak{}alpha,\allowbreak{} in either branch.
Grouping them with its leading gives S\_\allowbreak{}s\textless=\allowbreak{}K+\allowbreak{}700-e,\allowbreak{} with e=\allowbreak{}11.

The endpoint identity alpha\_\allowbreak{}s=\allowbreak{}P\_\allowbreak{}s+\allowbreak{}beta\_\allowbreak{}s+\allowbreak{}alpha\_\allowbreak{}t+\allowbreak{}d\_\allowbreak{}t+\allowbreak{}2(M+\allowbreak{}X)+\allowbreak{}G has
constant155. Leading t and the disjoint six- and three-edge forests
therefore give

\texttt{2M+\allowbreak{}9f+\allowbreak{}4e+\allowbreak{}2N <=\allowbreak{}545+\allowbreak{}3delta+\allowbreak{}B}.

Baseline errors B\textless=\allowbreak{}2xi\_\allowbreak{}s+\allowbreak{}2xi\_\allowbreak{}t. In the all-light case the
radial residual is nonnegative,\allowbreak{} and L\textgreater=\allowbreak{}A+\allowbreak{}20 implies

\texttt{(9/\allowbreak{}4)I-(3/\allowbreak{}4)delta+\allowbreak{}4e+\allowbreak{}2N <=\allowbreak{}-577/\allowbreak{}4}.

Since delta\textless=\allowbreak{}143-I-tau and e=\allowbreak{}11,\allowbreak{} its left side is at least
\texttt{-253/\allowbreak{}4},\allowbreak{} larger by81. In a heavy case cancel
e\_\allowbreak{}v-4 from 2L,\allowbreak{} retaining L+\allowbreak{}e\_\allowbreak{}o\textgreater=\allowbreak{}A+\allowbreak{}24. The resulting row is

\texttt{(9/\allowbreak{}4)I-(3/\allowbreak{}4)delta+\allowbreak{}4e+\allowbreak{}2N <=\allowbreak{}43/\allowbreak{}4}.

Since N\textgreater=\allowbreak{}73,\allowbreak{} delta\textless=\allowbreak{}143-I-tau and e=\allowbreak{}11,\allowbreak{} its left side
is at least331/4,\allowbreak{} exceeding43/4.

\hypertarget{z4-empty-beta-cards-e-h-e-alpha-e-c5}{%
\subsection{Z4 empty-beta Cards E-H,\allowbreak{} E-alpha,\allowbreak{}
E-C5}\label{z4-empty-beta-cards-e-h-e-alpha-e-c5}}

All side caps come from the actual opposite physical target union,\allowbreak{}
rather than a syntactically chosen set of centres.

\hypertarget{e-h}{%
\subsubsection{E-H}\label{e-h}}

Physical t has seven native resources,\allowbreak{} so S\_\allowbreak{}s\textless=\allowbreak{}K+\allowbreak{}623-e. Five
endpoints and the two forests give

\texttt{R+\allowbreak{}2L+\allowbreak{}xi\_\allowbreak{}\allowbreak{}t+\allowbreak{}e\_\allowbreak{}\allowbreak{}t+\allowbreak{}10f+\allowbreak{}3e <=\allowbreak{}416+\allowbreak{}2delta+\allowbreak{}Jstar}.

If no lower endpoint is heavy,\allowbreak{} Jstar\textless=\allowbreak{}xi\_\allowbreak{}t+\allowbreak{}e\_\allowbreak{}t-4. Source
L\textgreater=\allowbreak{}A+\allowbreak{}20 gives
\texttt{(5/\allowbreak{}2)I+\allowbreak{}delta/\allowbreak{}2+\allowbreak{}3e<=\allowbreak{}-469/\allowbreak{}2},\allowbreak{} impossible. If a
lower endpoint has heavy clip c,\allowbreak{} Jstar\textless=\allowbreak{}xi\_\allowbreak{}t+\allowbreak{}xi\_\allowbreak{}v+\allowbreak{}e\_\allowbreak{}v-4-c;
2L+\allowbreak{}e\_\allowbreak{}t-e\_\allowbreak{}v\textgreater=\allowbreak{}A+\allowbreak{}28 gives
\texttt{(5/\allowbreak{}2)I+\allowbreak{}delta/\allowbreak{}2+\allowbreak{}3e+\allowbreak{}c<=\allowbreak{}-159/\allowbreak{}2},\allowbreak{} also impossible.

\hypertarget{owner-alpha_s-missing}{%
\subsubsection{Owner alpha\_\allowbreak{}s missing}\label{owner-alpha_s-missing}}

Physical t again has seven native resources,\allowbreak{} so S\_\allowbreak{}s\textless=\allowbreak{}K+\allowbreak{}623-e.
The raw-order identity with alpha\_\allowbreak{}s\textgreater=\allowbreak{}1-a and the
endpoint/ten-edge forests gives

\texttt{L+\allowbreak{}2e\_\allowbreak{}\allowbreak{}t+\allowbreak{}2xi\_\allowbreak{}\allowbreak{}t+\allowbreak{}10f+\allowbreak{}3e+\allowbreak{}N <=\allowbreak{}419+\allowbreak{}2delta+\allowbreak{}B}.

All-light baseline B\textless=\allowbreak{}xi\_\allowbreak{}s+\allowbreak{}3xi\_\allowbreak{}t leaves 610-(5/2)tau radially;
L+\allowbreak{}2e\_\allowbreak{}t\textgreater=\allowbreak{}A+\allowbreak{}28 yields
\texttt{(5/\allowbreak{}2)I+\allowbreak{}delta/\allowbreak{}2+\allowbreak{}3e+\allowbreak{}N<=\allowbreak{}-9/\allowbreak{}2}.

A heavy error on source t cancels e\_\allowbreak{}t-4; L+\allowbreak{}e\_\allowbreak{}t\textgreater=\allowbreak{}A+\allowbreak{}24 yields
the same row with RHS -1/2. A heavy error on source s can only be P\_\allowbreak{}s;
cancelling e\_\allowbreak{}s leaves3e\_\allowbreak{}t\textgreater=\allowbreak{}3A-207 and yields
\texttt{3I+\allowbreak{}tau/\allowbreak{}2+\allowbreak{}delta+\allowbreak{}3e+\allowbreak{}N<=\allowbreak{}12}. Here
delta\textgreater=\allowbreak{}e+\allowbreak{}73 and N\textgreater=\allowbreak{}73,\allowbreak{} so it is impossible.

\hypertarget{opposite-alpha_t-missing}{%
\subsubsection{Opposite alpha\_\allowbreak{}t
missing}\label{opposite-alpha_t-missing}}

Physical t has six native resources,\allowbreak{} so S\_\allowbreak{}s\textless=\allowbreak{}K+\allowbreak{}700-e. The
weighted endpoint identity and eleven edges give

\texttt{R+\allowbreak{}2L+\allowbreak{}xi\_\allowbreak{}\allowbreak{}s+\allowbreak{}e\_\allowbreak{}\allowbreak{}s+\allowbreak{}11f+\allowbreak{}3e <=\allowbreak{}492+\allowbreak{}2delta+\allowbreak{}B}.

Baseline errors are3xi\_\allowbreak{}s+\allowbreak{}2xi\_\allowbreak{}t,\allowbreak{} leaving739-(11/4)tau. All light
offsets satisfy2L+\allowbreak{}e\_\allowbreak{}s\textgreater=\allowbreak{}2A+\allowbreak{}44 and give a negative upper
bound. A heavy endpoint of weightq=\allowbreak{}1,\allowbreak{}2 satisfies2L+\allowbreak{}e\_\allowbreak{}s-q
e\_\allowbreak{}v\textgreater=\allowbreak{}A+\allowbreak{}24; its remaining inequality is
\texttt{(11/\allowbreak{}4)I+\allowbreak{}(3/\allowbreak{}4)delta+\allowbreak{}3e+\allowbreak{}4q<=\allowbreak{}-83/\allowbreak{}4},\allowbreak{} impossible.

\hypertarget{e-c5}{%
\subsubsection{E-C5}\label{e-c5}}

Physical t has six native resources H,\allowbreak{}U,\allowbreak{}V,\allowbreak{}C6,\allowbreak{}C7,\allowbreak{}alpha,\allowbreak{} so
S\_\allowbreak{}s\textless=\allowbreak{}K+\allowbreak{}700-e. Five endpoints and eleven edges give

\texttt{R+\allowbreak{}2L+\allowbreak{}xi\_\allowbreak{}\allowbreak{}t+\allowbreak{}e\_\allowbreak{}\allowbreak{}t+\allowbreak{}11f+\allowbreak{}3e+\allowbreak{}N <=\allowbreak{}493+\allowbreak{}2delta+\allowbreak{}B}.

Baseline errors2xi\_\allowbreak{}s+\allowbreak{}3xi\_\allowbreak{}t leave739-(11/4)tau. Light
offsets2L+\allowbreak{}e\_\allowbreak{}t\textgreater=\allowbreak{}2A+\allowbreak{}44 yield a negative upper bound. A unique
heavy error leaves2L+\allowbreak{}e\_\allowbreak{}t-e\_\allowbreak{}v\textgreater=\allowbreak{}A+\allowbreak{}28 and gives
\texttt{(11/\allowbreak{}4)I+\allowbreak{}(3/\allowbreak{}4)delta+\allowbreak{}3e+\allowbreak{}N<=\allowbreak{}-111/\allowbreak{}4}. Thus this
card is closed.

All displayed Z4 retention rows are excluded. The final ownerless family
\texttt{\{P,\allowbreak{}B,\allowbreak{}O\_\allowbreak{}\allowbreak{}s,\allowbreak{}O\_\allowbreak{}\allowbreak{}t\}} has seven actual native
ordinary resources on each physical sign. Its complete independent
geometric and scalar certificate is
\texttt{../\allowbreak{}scalar/\allowbreak{}z4\_\allowbreak{}\allowbreak{}ownerless\_\allowbreak{}\allowbreak{}certificate.tex},\allowbreak{}
checked by
\texttt{../\allowbreak{}scalar/\allowbreak{}check\_\allowbreak{}\allowbreak{}z4\_\allowbreak{}\allowbreak{}ownerless\_\allowbreak{}\allowbreak{}certificate.py}.
The fresh four-edge arm gives115+\allowbreak{}I rather than the old131+\allowbreak{}I; this still
yields X+\allowbreak{}g\textgreater=\allowbreak{}58 and d-h-X\textgreater=\allowbreak{}3\textgreater0,\allowbreak{} enough
for the literal ten-label gap. The original eighteen-credit ledger is
therefore valid. Deleting the two largest core values gives the
actual67-point support with f=\allowbreak{}67-(tau-xi); this stronger support closes
all four H core-hit statuses. Both empty pairs remain ownerless
throughout that proof.

\hypertarget{z5-common-activity-and-endpoints}{%
\subsection{Z5 common activity and
endpoints}\label{z5-common-activity-and-endpoints}}

Here S=\allowbreak{}2K+\allowbreak{}1085-delta,\allowbreak{} A\textless=\allowbreak{}60,\allowbreak{} K-M\textgreater=\allowbreak{}R-55+\allowbreak{}A,\allowbreak{}
K\textgreater=\allowbreak{}153. Every retained singleton has at least17 hits,\allowbreak{}
exceeding all eight source peels together; thus it has an actual core
hit. Once sole pair ownership is proved,\allowbreak{} the same assertion holds for
that whole paired capacity. All actual retained endpoint errors
are\textless=\allowbreak{}xi\_\allowbreak{}source because their selected clips are\textless73.

Both O are empty before outer integrality: an active O has raw
t\textless=\allowbreak{}12,\allowbreak{} whereas leading activity gives
t\textgreater=\allowbreak{}R-55+\allowbreak{}I+\allowbreak{}tau+\allowbreak{}2X+\allowbreak{}g\textgreater13.

If beta\_\allowbreak{}v,\allowbreak{}H\_\allowbreak{}v were both active,\allowbreak{} its raw cap21 would
give2X\_\allowbreak{}v+\allowbreak{}g\_\allowbreak{}v\textless=\allowbreak{}8-I-tau,\allowbreak{} so tau\textless8. If either U\_\allowbreak{}v or
C6\_\allowbreak{}o were active then q=\allowbreak{}X\_\allowbreak{}v+\allowbreak{}g\_\allowbreak{}v would be integral. The two opposite
leading/core shortcut intersections
give2q\textgreater=\allowbreak{}2(137-R)-delta\textgreater=\allowbreak{}78+\allowbreak{}I-tau\textgreater70,\allowbreak{}
contradiction q\textless8. If either V\_\allowbreak{}v or C7\_\allowbreak{}o were active the same
argument with q=\allowbreak{}X\_\allowbreak{}v contradicts the cap. Therefore these four ordinary
singleton units are all holes; together with the two O holes they exceed
j=\allowbreak{}5. So beta/H cannot both be active.

The report\textquotesingle s literal one- or two-coordinate residue
classifier then proves outer integrality; the actual cross chains and
HIGH own-chain rules prove extreme-leaf retention and the A0/A1/B hole
lists. The four outer-coordinate identities used there are unchanged. An
empty cross centre is LOW because a HIGH cross loses at least
r\_\allowbreak{}v+\allowbreak{}81\textgreater delta leading points.

An empty alpha below a retained H is LOW without a normality bootstrap.
A retained H has a core hit,\allowbreak{} so H\textless=\allowbreak{}S\_\allowbreak{}v+\allowbreak{}R\_\allowbreak{}o. The leading bound
gives d\_\allowbreak{}v\textgreater=\allowbreak{}153-delta+\allowbreak{}g\_\allowbreak{}v+\allowbreak{}2X\_\allowbreak{}v\textgreater R\_\allowbreak{}o (since
R\_\allowbreak{}o\textless=\allowbreak{}r\_\allowbreak{}o+\allowbreak{}tau\_\allowbreak{}o,\allowbreak{} r\_\allowbreak{}o\textless=\allowbreak{}41 and delta+\allowbreak{}tau\textless=\allowbreak{}60).
Therefore alpha=\allowbreak{}H-d\textless S\_\allowbreak{}v. Its emptiness forces
alpha\textless=\allowbreak{}-M\_\allowbreak{}v,\allowbreak{} giving h\_\allowbreak{}v\textless=\allowbreak{}X\_\allowbreak{}v.

Eight O HIGH fixed labels give epsilon=\allowbreak{}delta-8. After the
outer-integrality classifier every extreme star leaf is retained. On
every chosen core support with overlap f,\allowbreak{} weighted stars give
ca*X\_\allowbreak{}s+\allowbreak{}cb*X\_\allowbreak{}t\textgreater=\allowbreak{}(ca+\allowbreak{}cb)f-max(ca,\allowbreak{}cb)*epsilon/2. The four
disjoint ordinary groups,\allowbreak{} including the one-empty form when U or C6 is
missing,\allowbreak{} give cg*g\_\allowbreak{}s+\allowbreak{}dg*g\_\allowbreak{}t\textgreater=\allowbreak{}(cg+\allowbreak{}dg)f-max(cg,\allowbreak{}dg)*epsilon/2.
These are separate valid scalar inequalities,\allowbreak{} not a claim that stars and
pairs are disjoint.

Source-only sharpenings needed below are

\begin{lstlisting}
M>=R+99+A,
e_v>=A+16.
\end{lstlisting}

The first uses the quartet Q\textless=\allowbreak{}3 equality case and central-gap
correction,\allowbreak{} with fresh mass coefficient105:
L\textgreater=\allowbreak{}103+\allowbreak{}A-Q\_\allowbreak{}s-Q\_\allowbreak{}t; both Q=\allowbreak{}3 would force L\textless=\allowbreak{}12,\allowbreak{} so
L\textgreater=\allowbreak{}98+\allowbreak{}A. Equality implies a huge offset with c\textless=\allowbreak{}15,\allowbreak{}
leaving a central gap of at least11+\allowbreak{}I+\allowbreak{}delta; hence L\textgreater=\allowbreak{}99+\allowbreak{}A.
These are the same actual exceptional-pair counts and equality
conditions of Appendix E,\allowbreak{} with the new source mass.

For the second,\allowbreak{} the sharpened one-maximum envelope is
\texttt{K<=\allowbreak{}2R+\allowbreak{}c\_\allowbreak{}\allowbreak{}v+\allowbreak{}e\_\allowbreak{}\allowbreak{}o+\allowbreak{}78+\allowbreak{}P\_\allowbreak{}\allowbreak{}v+\allowbreak{}B\_\allowbreak{}\allowbreak{}o}. The actual
distinct exceptional triples and pairs satisfy
\texttt{2P\_\allowbreak{}\allowbreak{}v<=\allowbreak{}a\_\allowbreak{}\allowbreak{}v+\allowbreak{}b\_\allowbreak{}\allowbreak{}v+\allowbreak{}4-p\_\allowbreak{}\allowbreak{}v} and
\texttt{2B\_\allowbreak{}\allowbreak{}o<=\allowbreak{}a\_\allowbreak{}\allowbreak{}o+\allowbreak{}b\_\allowbreak{}\allowbreak{}o+\allowbreak{}8-p\_\allowbreak{}\allowbreak{}o},\allowbreak{} where p\_\allowbreak{}v records
quartet equality. Consequently
\texttt{T+\allowbreak{}J\_\allowbreak{}\allowbreak{}v-J\_\allowbreak{}\allowbreak{}o<=\allowbreak{}168-p},\allowbreak{} with
\texttt{T=\allowbreak{}105+\allowbreak{}A}. Combined with
\texttt{J>=\allowbreak{}A+\allowbreak{}87+\allowbreak{}p},\allowbreak{} this gives
\texttt{J\_\allowbreak{}\allowbreak{}v>=\allowbreak{}A+\allowbreak{}12+\allowbreak{}p}. Since
\texttt{c\_\allowbreak{}\allowbreak{}v+\allowbreak{}p>=\allowbreak{}4},\allowbreak{} it follows that
\texttt{e\_\allowbreak{}\allowbreak{}v>=\allowbreak{}A+\allowbreak{}16}. The endpoint tables below
deliberately use the weaker \texttt{e\_\allowbreak{}\allowbreak{}v>=\allowbreak{}A+\allowbreak{}8} to
preserve their independently audited residuals. The alpha-hole terminal
proof uses the stronger floor16.

\hypertarget{all-thirteen-z5-a0-retained-p-templates}{%
\subsection{All thirteen Z5 A0 retained-P
templates}\label{all-thirteen-z5-a0-retained-p-templates}}

Keep the report\textquotesingle s literal endpoint lists,\allowbreak{} which are
exhaustive after its now freshly proved retention grammar. The fresh
constants and baseline rank-error weights are:

Here P=\allowbreak{}C4\_\allowbreak{}s is the actual sole pair1 owner,\allowbreak{} a=\allowbreak{}s4,\allowbreak{} b=\allowbreak{}t4,\allowbreak{} and
G=\allowbreak{}g\_\allowbreak{}s+\allowbreak{}g\_\allowbreak{}t. These are the literal endpoint identities; each upper list
is summed with unit weight and each number in the lower list is its
coefficient. No empty pair appears as a resource.

\begin{longtable}[]{@{}
  >{\raggedright\arraybackslash}p{(\columnwidth - 6\tabcolsep) * \real{0.2500}}
  >{\raggedright\arraybackslash}p{(\columnwidth - 6\tabcolsep) * \real{0.2500}}
  >{\raggedright\arraybackslash}p{(\columnwidth - 6\tabcolsep) * \real{0.2500}}
  >{\raggedright\arraybackslash}p{(\columnwidth - 6\tabcolsep) * \real{0.2500}}@{}}
\toprule\noalign{}
\begin{minipage}[b]{\linewidth}\raggedright
Card
\end{minipage} & \begin{minipage}[b]{\linewidth}\raggedright
Upper endpoints
\end{minipage} & \begin{minipage}[b]{\linewidth}\raggedright
Lower endpoints with weights
\end{minipage} & \begin{minipage}[b]{\linewidth}\raggedright
Centre identity body
\end{minipage} \\
\midrule\noalign{}
\endhead
\bottomrule\noalign{}
\endlastfoot
A & V\_\allowbreak{}s,\allowbreak{}V\_\allowbreak{}t & P,\allowbreak{}2C5\_\allowbreak{}t,\allowbreak{}2alpha\_\allowbreak{}s,\allowbreak{}alpha\_\allowbreak{}t & 5a+\allowbreak{}5b+\allowbreak{}2G \\
B & V\_\allowbreak{}s,\allowbreak{}H\_\allowbreak{}t & P,\allowbreak{}2C5\_\allowbreak{}t,\allowbreak{}2alpha\_\allowbreak{}s & 5a+\allowbreak{}4b+\allowbreak{}2g\_\allowbreak{}s+\allowbreak{}3g\_\allowbreak{}t+\allowbreak{}2h\_\allowbreak{}t \\
C & H\_\allowbreak{}s,\allowbreak{}U\_\allowbreak{}t & P,\allowbreak{}2C5\_\allowbreak{}t,\allowbreak{}3alpha\_\allowbreak{}s & 6a+\allowbreak{}4b+\allowbreak{}3G+\allowbreak{}h\_\allowbreak{}t \\
D & H\_\allowbreak{}s,\allowbreak{}H\_\allowbreak{}t & P,\allowbreak{}2C5\_\allowbreak{}t & 3a+\allowbreak{}4b+\allowbreak{}3G+\allowbreak{}3h\_\allowbreak{}s+\allowbreak{}2h\_\allowbreak{}t \\
E & H\_\allowbreak{}s,\allowbreak{}U\_\allowbreak{}t & P,\allowbreak{}2C5\_\allowbreak{}t,\allowbreak{}alpha\_\allowbreak{}t & 3a+\allowbreak{}5b+\allowbreak{}3G+\allowbreak{}3h\_\allowbreak{}s \\
F & U\_\allowbreak{}s,\allowbreak{}H\_\allowbreak{}t & P,\allowbreak{}2C5\_\allowbreak{}t,\allowbreak{}2alpha\_\allowbreak{}t & 3a+\allowbreak{}6b+\allowbreak{}3G+\allowbreak{}2h\_\allowbreak{}s \\
G & H\_\allowbreak{}s,\allowbreak{}H\_\allowbreak{}t & P,\allowbreak{}2C6\_\allowbreak{}t,\allowbreak{}2alpha\_\allowbreak{}t & 3a+\allowbreak{}6b+\allowbreak{}3G+\allowbreak{}h\_\allowbreak{}s \\
H & H\_\allowbreak{}s,\allowbreak{}H\_\allowbreak{}t & P,\allowbreak{}2C6\_\allowbreak{}t,\allowbreak{}alpha\_\allowbreak{}s & 4a+\allowbreak{}4b+\allowbreak{}3G+\allowbreak{}2h\_\allowbreak{}t \\
I & H\_\allowbreak{}s,\allowbreak{}U\_\allowbreak{}t & P,\allowbreak{}2C6\_\allowbreak{}t,\allowbreak{}alpha\_\allowbreak{}s,\allowbreak{}alpha\_\allowbreak{}t & 4a+\allowbreak{}5b+\allowbreak{}3G \\
J & H\_\allowbreak{}s,\allowbreak{}H\_\allowbreak{}t & P,\allowbreak{}2C7\_\allowbreak{}t,\allowbreak{}alpha\_\allowbreak{}s,\allowbreak{}2alpha\_\allowbreak{}t & 4a+\allowbreak{}6b+\allowbreak{}g\_\allowbreak{}s+\allowbreak{}3g\_\allowbreak{}t \\
K & U\_\allowbreak{}s,\allowbreak{}H\_\allowbreak{}t & P,\allowbreak{}2C6\_\allowbreak{}t,\allowbreak{}2alpha\_\allowbreak{}t & 3a+\allowbreak{}6b+\allowbreak{}3G \\
\end{longtable}

All templates have base holes B,\allowbreak{}O\_\allowbreak{}s,\allowbreak{}O\_\allowbreak{}t. The two extra holes are
respectively: D(alpha\_\allowbreak{}s,\allowbreak{}alpha\_\allowbreak{}t); A(U\_\allowbreak{}s,\allowbreak{}H\_\allowbreak{}s),\allowbreak{} A(U\_\allowbreak{}t,\allowbreak{}H\_\allowbreak{}t),\allowbreak{}
A(H\_\allowbreak{}s,\allowbreak{}H\_\allowbreak{}t); B(alpha\_\allowbreak{}t,\allowbreak{}H\_\allowbreak{}s); C(alpha\_\allowbreak{}t,\allowbreak{}H\_\allowbreak{}t); E(alpha\_\allowbreak{}s,\allowbreak{}H\_\allowbreak{}t);
F(alpha\_\allowbreak{}s,\allowbreak{}H\_\allowbreak{}s); G(C5\_\allowbreak{}t,\allowbreak{}alpha\_\allowbreak{}s); H(C5\_\allowbreak{}t,\allowbreak{}alpha\_\allowbreak{}t); I(C5\_\allowbreak{}t,\allowbreak{}H\_\allowbreak{}t);
J(C5\_\allowbreak{}t,\allowbreak{}C6\_\allowbreak{}t); K(C5\_\allowbreak{}t,\allowbreak{}H\_\allowbreak{}s). This is all thirteen templates,\allowbreak{} including
both physical owner orientations under reflection.

\begin{longtable}[]{@{}lllr@{}}
\toprule\noalign{}
Card & Constant & Error weights (s,\allowbreak{}t) & Maximum native clip weight \\
\midrule\noalign{}
\endhead
\bottomrule\noalign{}
\endlastfoot
A & 310 & 4,\allowbreak{}4 & 2 \\
B & 305-r & 4,\allowbreak{}3 & 2 \\
C & 378-2r & 5,\allowbreak{}3 & 3 \\
D & 159+\allowbreak{}r & 2,\allowbreak{}3 & 2 \\
E & 164+\allowbreak{}2r & 2,\allowbreak{}4 & 2 \\
F,\allowbreak{}G,\allowbreak{}K & 169+\allowbreak{}3r & 2,\allowbreak{}5 & 2 \\
H & 232 & 3,\allowbreak{}3 & 2 \\
I & 237+\allowbreak{}r & 3,\allowbreak{}4 & 2 \\
J & 242+\allowbreak{}2r & 3,\allowbreak{}5 & 2 \\
\end{longtable}

Each constant is constructed as76+\allowbreak{}sum(weight*(73-r\_\allowbreak{}side)),\allowbreak{} with
r\_\allowbreak{}t=\allowbreak{}68-r. Use f=\allowbreak{}61-(tau-xi)/4 and the mass lower bounds
M\textgreater=\allowbreak{}R+\allowbreak{}99+\allowbreak{}A,\allowbreak{} M\_\allowbreak{}v\textgreater=\allowbreak{}r\_\allowbreak{}v+\allowbreak{}xi\_\allowbreak{}v+\allowbreak{}A+\allowbreak{}8. Drop positive h
terms,\allowbreak{} except in D. Use the independent weighted star/pair rows and
Eclip\textless=\allowbreak{}maxweight*epsilon. The residual lower-minus-upper
polynomial,\allowbreak{} after applying delta\textless=\allowbreak{}60-I-tau,\allowbreak{} is bounded below as
follows:

\begin{longtable}[]{@{}ll@{}}
\toprule\noalign{}
Card & Strict positive lower residual \\
\midrule\noalign{}
\endhead
\bottomrule\noalign{}
\endlastfoot
A & 356+\allowbreak{}(9/2)I+\allowbreak{}tau+\allowbreak{}xi/2 \\
B & 166+\allowbreak{}5I+\allowbreak{}(3/2)tau+\allowbreak{}xi/2 \\
C & 351+\allowbreak{}(13/2)I+\allowbreak{}(5/2)tau+\allowbreak{}xi \\
E & 2+\allowbreak{}5I+\allowbreak{}(3/2)tau+\allowbreak{}xi/2 \\
F,\allowbreak{}G,\allowbreak{}K & 178+\allowbreak{}(11/2)I+\allowbreak{}(7/4)tau+\allowbreak{}(3/4)xi \\
H & 51+\allowbreak{}(9/2)I+\allowbreak{}tau+\allowbreak{}xi/2 \\
I & 227+\allowbreak{}5I+\allowbreak{}(5/4)tau+\allowbreak{}(3/4)xi \\
J & 281+\allowbreak{}(11/2)I+\allowbreak{}2tau+\allowbreak{}xi/2 \\
\end{longtable}

For D both alpha resources are missing and LOW. The two retained HUV
triples and cross g\_\allowbreak{}t pair are distinct,\allowbreak{} with native group deficiencies
Psi\_\allowbreak{}s,\allowbreak{}Psi\_\allowbreak{}t,\allowbreak{}Omega. Their common f support gives

\texttt{3G+\allowbreak{}3h\_\allowbreak{}\allowbreak{}s+\allowbreak{}2h\_\allowbreak{}\allowbreak{}t>=\allowbreak{}305+\allowbreak{}6f-3Psi\_\allowbreak{}\allowbreak{}s-2Psi\_\allowbreak{}\allowbreak{}t-Omega>=\allowbreak{}305+\allowbreak{}6f-3epsilon}.

This replaces the weaker3G bound. The final residual is at least
\texttt{53+\allowbreak{}6I+\allowbreak{}(11/\allowbreak{}4)tau+\allowbreak{}xi/\allowbreak{}4>0}. Thus all thirteen
literal A0 retained-P templates are excluded without assuming any
retained endpoint is normal.

\hypertarget{a0-both-pairs-empty-h_s-or-c5_s-extra-hole}{%
\subsection{A0 both pairs empty,\allowbreak{} H\_\allowbreak{}s or C5\_\allowbreak{}s extra
hole}\label{a0-both-pairs-empty-h_s-or-c5_s-extra-hole}}

Physical side s has six native resources; the other side has seven.
Grouping side s gives S\_\allowbreak{}t\textless=\allowbreak{}K+\allowbreak{}623-e. Let A0=\allowbreak{}alpha\_\allowbreak{}s,\allowbreak{}
B=\allowbreak{}alpha\_\allowbreak{}t,\allowbreak{} D=\allowbreak{}C5\_\allowbreak{}t. The identity H\_\allowbreak{}t=\allowbreak{}2(a+\allowbreak{}b)+\allowbreak{}G+\allowbreak{}A0+\allowbreak{}B+\allowbreak{}D,\allowbreak{} the generic
upper endpoint,\allowbreak{} and K\textless=\allowbreak{}M+\allowbreak{}R+\allowbreak{}21 give

\texttt{M+\allowbreak{}2X+\allowbreak{}G+\allowbreak{}A0+\allowbreak{}B+\allowbreak{}D<=\allowbreak{}R+\allowbreak{}572+\allowbreak{}r\_\allowbreak{}\allowbreak{}s-e+\allowbreak{}u\_\allowbreak{}\allowbreak{}Ht+\allowbreak{}xi\_\allowbreak{}\allowbreak{}s}.

Alpha\_\allowbreak{}s is retained. Lower it generically; pack B+\allowbreak{}D in its actual
two-resource group Gamma,\allowbreak{} giving
B+\allowbreak{}D\textgreater=\allowbreak{}2(73-r\_\allowbreak{}t)+\allowbreak{}f-2Gamma-2xi\_\allowbreak{}t. Hence

\texttt{M+\allowbreak{}2X+\allowbreak{}G+\allowbreak{}f<=\allowbreak{}R+\allowbreak{}489-e+\allowbreak{}u\_\allowbreak{}\allowbreak{}Ht+\allowbreak{}l\_\allowbreak{}\allowbreak{}alpha\_\allowbreak{}\allowbreak{}s+\allowbreak{}2Gamma+\allowbreak{}2xi}.

Let N=\allowbreak{}Gamma+\allowbreak{}Gamma\_\allowbreak{}Ht+\allowbreak{}Gamma\_\allowbreak{}alpha\_\allowbreak{}s. All these endpoints avoid the
stars and the four ordinary g groups. Their independent rebates
give2X\textgreater=\allowbreak{}4f-(epsilon-N),\allowbreak{} G\textgreater=\allowbreak{}2f-(epsilon-N)/2. Since
the endpoint weighted loss is\textless=\allowbreak{}2N,\allowbreak{} combination yields

\texttt{M+\allowbreak{}7f<=\allowbreak{}R+\allowbreak{}489-e+\allowbreak{}2epsilon+\allowbreak{}2xi}.

Substitute M\textgreater=\allowbreak{}R+\allowbreak{}99+\allowbreak{}A,\allowbreak{} f=\allowbreak{}61-(tau-xi)/4,\allowbreak{} xi\textless=\allowbreak{}tau and
A\textless=\allowbreak{}60 to obtain \texttt{2I+\allowbreak{}3e<=\allowbreak{}23}. Already e=\allowbreak{}8
gives24\textless=\allowbreak{}23,\allowbreak{} impossible. Thus both cases close without the old
normality assumption or extra fixed ledgers.

\hypertarget{a0-both-pairs-empty-alpha_s-extra-hole}{%
\subsection{A0 both pairs empty,\allowbreak{} alpha\_\allowbreak{}s extra
hole}\label{a0-both-pairs-empty-alpha_s-extra-hole}}

Initial fixed ledger e=\allowbreak{}14 remains valid without normality because both
paired components are LOW,\allowbreak{} the alpha hole is LOW,\allowbreak{} and d\textgreater g
proves all F3,\allowbreak{}F4,\allowbreak{}F5 labels LOW. With epsilon=\allowbreak{}delta-14,\allowbreak{} the
leading/stars/pairs give d\textgreater=\allowbreak{}107+\allowbreak{}g+\allowbreak{}2X,\allowbreak{}
X\_\allowbreak{}v,\allowbreak{}g\_\allowbreak{}v\textgreater=\allowbreak{}38,\allowbreak{} and h\_\allowbreak{}s\textgreater=\allowbreak{}15.

The exact32-type upper-tail classification can therefore reproduce the
HIGH21,\allowbreak{}LOW29,\allowbreak{}touched23 ledger and e=\allowbreak{}27 before normality. Two places
previously using C5\_\allowbreak{}v\textgreater0 are replaced by actual core-hit
lower C5\_\allowbreak{}v\textgreater=\allowbreak{}1-R\_\allowbreak{}v: C5\_\allowbreak{}s+\allowbreak{}s3+\allowbreak{}h\_\allowbreak{}s\textgreater26 follows from
e\_\allowbreak{}s+\allowbreak{}X\_\allowbreak{}s+\allowbreak{}g\_\allowbreak{}s+\allowbreak{}h\_\allowbreak{}s+\allowbreak{}1\textgreater26,\allowbreak{} and
C5\_\allowbreak{}t+\allowbreak{}M\_\allowbreak{}t+\allowbreak{}2X\_\allowbreak{}t+\allowbreak{}g\_\allowbreak{}t\textgreater16 follows from
e\_\allowbreak{}t+\allowbreak{}2X\_\allowbreak{}t+\allowbreak{}g\_\allowbreak{}t+\allowbreak{}1\textgreater16. The tail argument only needs
h\_\allowbreak{}s\textgreater2,\allowbreak{} not29,\allowbreak{} to exclude d\_\allowbreak{}s+\allowbreak{}1,\allowbreak{}d\_\allowbreak{}s+\allowbreak{}2. Once e=\allowbreak{}27,\allowbreak{} the BAD
native-loss floor73-max(r\_\allowbreak{}s,\allowbreak{}r\_\allowbreak{}t) exceeds epsilon=\allowbreak{}delta-27 uniformly,\allowbreak{}
so all retained centres are normal.

At e27,\allowbreak{} the target caps are S\_\allowbreak{}s\textless=\allowbreak{}K+\allowbreak{}546-e and
S\_\allowbreak{}t\textless=\allowbreak{}K+\allowbreak{}623-e. If B=\allowbreak{}alpha\_\allowbreak{}t\textgreater=\allowbreak{}a or
D=\allowbreak{}C5\_\allowbreak{}t\textgreater=\allowbreak{}a,\allowbreak{} the normal H\_\allowbreak{}t endpoint and lower endpoint of
the other centre give M+\allowbreak{}2X+\allowbreak{}G+\allowbreak{}h\_\allowbreak{}s\textless=\allowbreak{}R+\allowbreak{}567-e+\allowbreak{}endpointloss. Rebate
those unused endpoint losses N in both the four-star and triple/pair gap
bounds to get A+\allowbreak{}6f+\allowbreak{}N\textless=\allowbreak{}407+\allowbreak{}delta-3e. With adaptive
f=\allowbreak{}max(137-R,\allowbreak{}65-(tau-xi)/2) one has2tau+\allowbreak{}6f\textgreater=\allowbreak{}398; A\textless=\allowbreak{}60
then forces2I+\allowbreak{}N+\allowbreak{}3e\textless=\allowbreak{}69,\allowbreak{} impossible at e27. Thus B,\allowbreak{}D\textless a.
The nine further LOW labels give e34. The fixed-only prefix argument,\allowbreak{}
now with X\_\allowbreak{}t\textgreater=\allowbreak{}38,\allowbreak{} gives C=\allowbreak{}C5\_\allowbreak{}s\textless b and the four
further LOW labels give e38.

Five additional distinct labels are LOW:

\begin{lstlisting}
C25_s=C-D-2b-g_t<0,       C25_t=D-C-2a-g_s<0,
C35_t=C25_t-h_t<0,
C45_s=C35_s-g_s<0,        C45_t=C35_t-g_t<0.
\end{lstlisting}

Normality gives C,\allowbreak{}D\textgreater=\allowbreak{}1,\allowbreak{} while the proved bounds
C\textless b,\allowbreak{}D\textless a give the two strict signs. C25\_\allowbreak{}s,\allowbreak{}C25\_\allowbreak{}t
occupy their one previously untouched baseline group; the other three
labels lie in groups already occupied respectively by V3\_\allowbreak{}t,\allowbreak{}V8\_\allowbreak{}s,\allowbreak{}V8\_\allowbreak{}t.
Thus the ledger has HIGH21,\allowbreak{}LOW47,\allowbreak{}touched26 and genuine compatible
credit42. This counts the eight O HIGH labels once. The labeled
incidence is recomputed against the original26 baseline groups by the
checker.

For completeness the actual endpoint/group transport is as follows. Pack
B=\allowbreak{}alpha\_\allowbreak{}t and D=\allowbreak{}C5\_\allowbreak{}t in their two-resource group Gamma. The retained
H\_\allowbreak{}s,\allowbreak{}U\_\allowbreak{}s,\allowbreak{}V\_\allowbreak{}s form a three-resource group Psi. The two g\_\allowbreak{}t pairs
U\_\allowbreak{}t,\allowbreak{}V\_\allowbreak{}t and C6\_\allowbreak{}s,\allowbreak{}C7\_\allowbreak{}s form Omega. The separate C=\allowbreak{}C5\_\allowbreak{}s and H\_\allowbreak{}t
endpoint losses are unused by these groups. All these are
original77-point native resources. The source Q chosen below is common
to their packing and the stars. If its size is q\textgreater=\allowbreak{}61 and its
overlap is f,\allowbreak{} then

\begin{lstlisting}
2B+D>=3(73-r_t)+f-3Gamma,
G+h_s>=q+f-Psi-Omega/2.
\end{lstlisting}

Use the two normal H upper endpoints,\allowbreak{} the normal C lower endpoint,\allowbreak{} and
the exact H centre identities. Rebate the distinct native losses in the
group bounds and in the separately valid weighted-star bound
\texttt{2X\_\allowbreak{}\allowbreak{}s+\allowbreak{}4X\_\allowbreak{}\allowbreak{}t>=\allowbreak{}6f-2epsilon}. Their combination
gives

\begin{lstlisting}
2M_t+2q+11f<=2R+827+2r_t-delta+5epsilon.
\end{lstlisting}

In particular q\textgreater=\allowbreak{}61 permits the conservative constant122.
Substitute M\_\allowbreak{}t\textgreater=\allowbreak{}r\_\allowbreak{}t+\allowbreak{}xi\_\allowbreak{}t+\allowbreak{}A+\allowbreak{}16,\allowbreak{} epsilon=\allowbreak{}delta-42,\allowbreak{} and
delta\textless=\allowbreak{}60-I-tau. The necessary scalar inequality becomes

\begin{lstlisting}
4I+4tau+11F-2xi_s+210<=929.
\end{lstlisting}

For each sign choose independently whether to retain all core values or
delete its largest j=\allowbreak{}1,\allowbreak{}2,\allowbreak{}3,\allowbreak{}4 magnitudes. Define
\texttt{phi\_\allowbreak{}\allowbreak{}v=\allowbreak{}min(xi\_\allowbreak{}\allowbreak{}v,\allowbreak{}min\_\allowbreak{}\allowbreak{}j(j+\allowbreak{}floor((tau\_\allowbreak{}\allowbreak{}v-xi\_\allowbreak{}\allowbreak{}v)/\allowbreak{}j)))}
and \texttt{F=\allowbreak{}69-phi\_\allowbreak{}\allowbreak{}s-phi\_\allowbreak{}\allowbreak{}t}. The minimizing choices
give one actual common Q of size at least61 with overlap at least F at
zero displacement; all displacements subtract their integral gap. Hence
this is an actual source selection,\allowbreak{} not a combination of different
budgets. Credit42 and A\textless=\allowbreak{}60 restrict tau\textless=\allowbreak{}18. Exact
enumeration of every nonnegative integer tau\_\allowbreak{}s,\allowbreak{}tau\_\allowbreak{}t,\allowbreak{}xi\_\allowbreak{}s,\allowbreak{}xi\_\allowbreak{}t with
xi\_\allowbreak{}v\textless=\allowbreak{}tau\_\allowbreak{}v and tau\textless=\allowbreak{}18 has minimum

\begin{lstlisting}
4tau+11F-2xi_s+210=935>929.
\end{lstlisting}

Dropping the nonnegative4I only enlarges the domain,\allowbreak{} so this finite
certificate excludes every size and every residual mass. The independent
derivation,\allowbreak{} labeled ledger,\allowbreak{} all7315 relaxed scalar states,\allowbreak{} and exact
checker are in
\texttt{../\allowbreak{}foundations/\allowbreak{}z5\_\allowbreak{}\allowbreak{}alpha\_\allowbreak{}\allowbreak{}final85.md} and
\texttt{../\allowbreak{}foundations/\allowbreak{}z5\_\allowbreak{}\allowbreak{}alpha\_\allowbreak{}\allowbreak{}final85.py}. The
scalar checker verifies the displayed finite consequence; the geometric
prerequisites above remain written actual-resource arguments.

\hypertarget{z5-a1-and-branch-b}{%
\subsection{Z5 A1 and branch B}\label{z5-a1-and-branch-b}}

The activity/outer-integrality classifier is the literal Appendix E
classifier following equation(z5-crosschain). It is valid with the fresh
base reductions above. For split beta the whole paired unit alone is
charged: its hit count N obeys77-delta\textless=\allowbreak{}N\textless=\allowbreak{}77 and
beta\_\allowbreak{}s+\allowbreak{}beta\_\allowbreak{}t\textgreater=\allowbreak{}N-M. There is no selected-component
capacity. For sole beta,\allowbreak{} normalize s to its actual owner; its base holes
are O\_\allowbreak{}s,\allowbreak{}H\_\allowbreak{}s,\allowbreak{}O\_\allowbreak{}t. When pair1 is empty,\allowbreak{} the final extra hole is exactly
alpha\_\allowbreak{}s,\allowbreak{}H\_\allowbreak{}t,\allowbreak{}alpha\_\allowbreak{}t,\allowbreak{}C5\_\allowbreak{}s,\allowbreak{}C5\_\allowbreak{}t,\allowbreak{}or U\_\allowbreak{}s. The independent fresh
exclusions for all split,\allowbreak{} retained-pair1,\allowbreak{} empty-pair1 and branch-B
alternatives are in \texttt{zero85\_\allowbreak{}\allowbreak{}z5\_\allowbreak{}\allowbreak{}a1\_\allowbreak{}\allowbreak{}b.md},\allowbreak{} with
exact arithmetic in
\texttt{check\_\allowbreak{}\allowbreak{}zero85\_\allowbreak{}\allowbreak{}z5\_\allowbreak{}\allowbreak{}a1\_\allowbreak{}\allowbreak{}b.py}. The final C5\_\allowbreak{}t
case\textquotesingle s cutoff and32-type tail received the independent
receipt \texttt{../\allowbreak{}scalar/\allowbreak{}z5\_\allowbreak{}\allowbreak{}c5t\_\allowbreak{}\allowbreak{}cutoff\_\allowbreak{}\allowbreak{}audit.md},\allowbreak{}
checked by
\texttt{../\allowbreak{}scalar/\allowbreak{}check\_\allowbreak{}\allowbreak{}z5\_\allowbreak{}\allowbreak{}c5t\_\allowbreak{}\allowbreak{}cutoff\_\allowbreak{}\allowbreak{}audit.py}.

\clearpage
% Conditional local exclusion. Requires amsmath,amssymb.
\subsection{The cardinality-85 ownerless $Z_4$ family}

Assume the actual activity classification has left exactly the four
missing whole units $\{P,B,O_s,O_t\}$ in branch A. Both paired units
are empty; no component is assigned an owner. The fourteen retained
singletons are $H_v,U_v,V_v,C_{5v},C_{6v},C_{7v},\alpha_v$ on the two
physical signs. Use the cardinality-$85$ structural source
\[
 |D|=77,\quad r_s+r_t=68,\quad
 I=(r_s-34)^2,\quad \xi_v=R_v-r_v,\quad \xi=\xi_s+\xi_t,
 \quad A=I+\tau+\delta.
\]
The source results used here, without changing the meaning of $\delta$,
are
\begin{equation}\label{eq:85-ownerless-source}
 A\le137,\qquad L=e_s+e_t=M-R\ge A+20,
 \qquad e_v\ge\max(4,A-69),\quad 0\le\xi\le\tau.
\end{equation}
The deliberately weaker last depth bound is enough; the sharper
one-maximum bound $e_v\ge A-61$ is unnecessary for this proof.

\paragraph{The fixed credit and the two $H$ gaps.}
The same actual gap-label inventories as in the ownerless $Z_4$ proof
give fixed credit $\eta=18$. Here are the fresh numerical checks needed
to use those inventories at cardinality $85$. The four-edge physical
arm gives
\[
 2X_v+X_o+g_v\ge252-\tau-\delta\ge115+I.
\]
Consequently the literal active-center identities imply
\[
 d_v\ge M+116+I,\qquad X+g_v\ge58,\qquad h_v\ge X_v+1.
\]
The raw gap first gives $h_v\le M+56<d_v$; the complete ten-label
inventory improves this to $h_v\le M+11$. Its only possible fixed
labels are
\[
 C_{33},C_{34},C_{44},C_{15},C_{26},C_{27},C_{28},U_2,U_3,V_2.
\]
Here indexed $U_i,V_i$ are fixed labels, not ordinary resources.
The paired centers are LOW. Per sign the sixteen LOW labels
\[
 C_{55},C_{56},C_{57},C_{58},\quad V_3,V_4,V_5,V_6,V_7,V_8,
 \quad U_5,U_6,U_7,\quad F_3,F_4,F_5
\]
occupy 25 of the 26 baseline groups over both signs. Three additional
LOW witnesses are supplied by the reflected pairs $V_2,C_{35},C_{45}$
in already occupied groups. In particular
$V_2^s+V_2^t\le20-3X-G<0$, since
$3X+G\ge230+2I$. The untouched $C_{25}$ group supplies another
LOW witness. Together with the eight HIGH $O$ labels this is 44
misses against the baseline 26. Thus
\[
 \varepsilon=\delta-18\ge0.
\]

For the truncated core gap below $H_v$, the original counts $37$ and
$30$ are unchanged. The two strict eliminations used to pass from
$37$ to $30$ remain valid: freshly
\[
 a-(\alpha_v+M_v)\ge79+I>0,\qquad
 a-U_3\ge X+g_o+h_o-36>0.
\]
Therefore
\begin{equation}\label{eq:85-ownerless-H}
 h_v\le R+31\quad(H_v\text{ has a core hit}),\qquad
 h_v\le R_v+M_o+30\quad(H_v\text{ has no core hit}).
\end{equation}
In the second case its actual native upper clipping loss is at least
$73$: the 69 core values and four positive source peels are all above
the target. This uses the complete target $[1,S_v]$, not its hull.

\paragraph{A source obtained by deleting just the two maxima.}
Put $D_0=\tau-\xi\ge0$. After the maximum on each core side is
deleted, its remaining rank deviations sum to $\tau_v-\xi_v$.
Hence the resulting signed source $Q$ has 67 values and span at most
$66+D_0$. At every integral displacement $q$,
\[
 |Q\cap(Q+q)|\ge67-D_0-q.
\]
Write $f=67-D_0$. Even when this lower bound is negative it is a valid
intersection bound, so no positivity assumption on $f$ is needed.
Each leading-to-$V$ edge contains the actual copies
$\lambda_v+M_v+Q$ and $\lambda_v+v_4+Q$ at gap $X_v$;
the reflected leading-to-$C_7$ edge has the same gap.
Designate the two full leadings as roots, with summed native loss
$\ell\ge0$. These stars avoid both alphas and both $H$ resources.

Let $\Gamma_v$ be the native alpha deficiency, $k_v$ its lower clip,
$N=\Gamma_s+\Gamma_t\ge k_s+k_t=k$, and let $B$ denote the sum of
the two actual rank errors. The fresh complete-target endpoints give
\[
 \alpha_s+\alpha_t\ge78-k-B.
\]
For a light clip $k_v<73$ its error is at most $\xi_v$.
At most one clip is heavy, because $2\cdot73>\varepsilon$.
The rooted stars give separately valid inequalities
\begin{equation}\label{eq:85-ownerless-stars}
 2X_v\ge2f+N+\ell-\varepsilon,\qquad
 2X\ge4f+N+\ell-\varepsilon.
\end{equation}
If an $H$ has no core hit, its unused native upper clip $u\ge73$
rebates the second row:
\begin{equation}\label{eq:85-ownerless-stars-H}
 N+u\le\varepsilon,\qquad
 2X\ge4f+N+u+\ell-\varepsilon.
\end{equation}
These are deficiencies of distinct original retained resources.

\paragraph{Both $H$ have core hits; both alpha clips are light.}
Now $B\le\xi$ and $h_s+h_t\le2R+62$. The identity
$\alpha_s+\alpha_t=h_s+h_t-M-X$ gives
$L+X\le R-16+k+B$. Insert $L\ge A+20$ and
\eqref{eq:85-ownerless-stars}, then use
$N\ge k$, $k\le\varepsilon$, $B\le\xi$. This yields
\[
 4f+2I+2\tau+2\eta+\ell\le64+4\xi.
\]
Since $f=67-\tau+\xi$ and $\eta=18$, it follows that
$2\tau\ge240+2I+\ell$. But
$\tau\le137-I-\delta\le119-I$. Thus
$2+4I+\ell\le0$, a contradiction.

\paragraph{Both $H$ have core hits; one alpha clip is heavy.}
Say $k_v\ge73$. The opposite alpha is light. Its endpoint and
$h_o\le R+31$ imply
$e_o+X_o\le26+\xi_v+k_o+\zeta_o$, where
$\zeta_o\le\xi_o$. Using the individual star,
$N\ge k_v+k_o$ and $k_o\le\varepsilon-k_v$ gives
\[
 2f+2e_o+2k_v+\ell\le52+2\xi+2\varepsilon.
\]
Consequently $2e_o+264+\ell\le2(\tau+\delta)
\le274-2I$, so $e_o\le5$. On the other hand
$\delta\ge18+73=91$, and the source floor gives
$e_o\ge A-69\ge22$. This is impossible; the heavy endpoint error
has not been bounded or discarded.

\paragraph{Exactly one $H$ has no core hit.}
Its clip $u\ge73$ leaves $N\le\varepsilon-u\le46$, so both alphas
are light and $B\le\xi$. If $v$ is the core-hit side, the two caps
in \eqref{eq:85-ownerless-H} give
$h_s+h_t\le2R+e_o+61$, hence
$e_v+X\le R-17+k+B$. Use
\eqref{eq:85-ownerless-stars-H} and $k\le\varepsilon-u$ to get
\[
 2e_v+4f+2u+\ell+34\le2R+2B+2\varepsilon.
\]
Thus $2e_v+202+2u+\ell\le4\tau+2\delta$.
Since $\delta\ge18+u\ge91$ and $A\le137$,
the right side is at most $366-4I$. Hence $e_v\le9$,
whereas $e_v\ge A-69\ge22$.

\paragraph{Neither $H$ has a core hit.}
The two distinct native deficiencies total at least $146>137\ge\delta$.
All four $H$ statuses are therefore impossible. This excludes the
conditional cardinality-$85$ ownerless $Z_4$ family without assigning
an owner or capacity to either empty pair.

\clearpage
\hypertarget{complete-exclusion-of-the-positive-hole-size85-stratum-p2}{%
\section{Complete exclusion of the positive-hole SIZE85 stratum
P2}\label{complete-exclusion-of-the-positive-hole-size85-stratum-p2}}

This is a fresh argument for j=\allowbreak{}2 and 1\textless=\allowbreak{}z\textless=\allowbreak{}14 at
cardinality85. It assumes the proved SIZE85 integral primitive
structural theorem,\allowbreak{} the common covering and source envelopes,\allowbreak{}
native-loss/forest lemmas,\allowbreak{} and the sharp raw-gap counts. All occurrences
on opposite signs remain distinct labelled occurrences. No canonical
endpoint model or old SIZE86 activity census is used.

\hypertarget{source-parameters}{%
\subsection{1. Source parameters}\label{source-parameters}}

Set d=\allowbreak{}77,\allowbreak{} \textbar E\textbar=\allowbreak{}69,\allowbreak{} r\_\allowbreak{}s+\allowbreak{}r\_\allowbreak{}t=\allowbreak{}68,\allowbreak{} I=\allowbreak{}(r\_\allowbreak{}s-34)²,\allowbreak{} G144=\allowbreak{}144.
Let rho=\allowbreak{}rho\_\allowbreak{}s+\allowbreak{}rho\_\allowbreak{}t,\allowbreak{} mu=\allowbreak{}mu\_\allowbreak{}s+\allowbreak{}mu\_\allowbreak{}t,\allowbreak{}
tau=\allowbreak{}sum(P\_\allowbreak{}s)+\allowbreak{}sum(P\_\allowbreak{}t)-T\_\allowbreak{}rs-T\_\allowbreak{}rt,\allowbreak{} H=\allowbreak{}h\_\allowbreak{}s+\allowbreak{}h\_\allowbreak{}t=\allowbreak{}tau-mu. Write
C\_\allowbreak{}v=\allowbreak{}sum(P\_\allowbreak{}v),\allowbreak{} R\_\allowbreak{}v=\allowbreak{}max(P\_\allowbreak{}v),\allowbreak{} R=\allowbreak{}R\_\allowbreak{}s+\allowbreak{}R\_\allowbreak{}t,\allowbreak{} M\_\allowbreak{}v=\allowbreak{}max(D\_\allowbreak{}v),\allowbreak{}
M=\allowbreak{}M\_\allowbreak{}s+\allowbreak{}M\_\allowbreak{}t,\allowbreak{} K=\allowbreak{}\textbar2D\textbar,\allowbreak{} and S=\allowbreak{}sum(D\_\allowbreak{}s)+\allowbreak{}sum(D\_\allowbreak{}t). The source
peel magnitudes are R\_\allowbreak{}v+\allowbreak{}a\_\allowbreak{}v,\allowbreak{}R\_\allowbreak{}v+\allowbreak{}b\_\allowbreak{}v,\allowbreak{}R\_\allowbreak{}v+\allowbreak{}c\_\allowbreak{}v,\allowbreak{}R\_\allowbreak{}v+\allowbreak{}e\_\allowbreak{}v,\allowbreak{} with
1\textless=\allowbreak{}a\_\allowbreak{}v\textless b\_\allowbreak{}v\textless c\_\allowbreak{}v\textless e\_\allowbreak{}v. Put

\begin{lstlisting}
U=a_s+b_s+a_t+b_t, C=c_s+c_t,
J=(e_s-c_s)+(e_t-c_t), A=I+tau+delta, m0=-126-w.
\end{lstlisting}

The source-only identities and inequalities are

\begin{lstlisting}
S=2K+1316+w-delta,
2K=m0+A+4R+U+2C+J,
J+U<=458+w-A, J>=A-144-w, U>=6, J>=2,
A<=298+w, K-M>=R-286-w+A,
K-delta>=16-w+I+tau.
\end{lstlisting}

The quantitative layers,\allowbreak{} for U0 the first four peeled masses,\allowbreak{} V0 the
last four,\allowbreak{} C0=\allowbreak{}C\_\allowbreak{}s+\allowbreak{}C\_\allowbreak{}t=\allowbreak{}1190+\allowbreak{}I+\allowbreak{}tau,\allowbreak{} and C1=\allowbreak{}C0+\allowbreak{}U0,\allowbreak{} give

\begin{lstlisting}
C0-U0>=320+2I+delta-2z,
C1-V0>=446+2I+mu+rho+delta+2z.
\end{lstlisting}

Initially w\textless=\allowbreak{}4z\textless=\allowbreak{}56. The structural budget gives
delta\textless=\allowbreak{}302-I-rho+\allowbreak{}2z. If the smaller side were at most17,\allowbreak{} then
I\textgreater=\allowbreak{}289 and rho\textgreater=\allowbreak{}49,\allowbreak{} giving delta\textless=\allowbreak{}-8.
Hence 18\textless=\allowbreak{}r\_\allowbreak{}v\textless=\allowbreak{}50; throughout this range
rho\textgreater=\allowbreak{}49 and mu\textgreater=\allowbreak{}62. Non-AP primitivity gives
R\_\allowbreak{}v\textgreater=\allowbreak{}r\_\allowbreak{}v+\allowbreak{}1 and tau\_\allowbreak{}v\textgreater=\allowbreak{}r\_\allowbreak{}v-3.

\hypertarget{common-strip-localization-using-the-actual-four-inner-occurrences}{%
\subsection{2. Common-strip localization using the actual four inner
occurrences}\label{common-strip-localization-using-the-actual-four-inner-occurrences}}

The source envelope gives H\textless=\allowbreak{}236+\allowbreak{}w-I-delta\textless=\allowbreak{}292-I-delta.
For every integer r,\allowbreak{} T\_\allowbreak{}r+\allowbreak{}(r-34)²\textgreater=\allowbreak{}397; hence each
C\_\allowbreak{}v\textgreater H. The common-strip lemma fills {[}H,\allowbreak{}C0-H{]} in the
absolute core support. Its width L=\allowbreak{}C0-2H satisfies

\begin{lstlisting}
L>=1016+2I+delta-w>=960.
\end{lstlisting}

Sort the four inner occurrences as
t1\textless=\allowbreak{}t2\textless=\allowbreak{}t3\textless=\allowbreak{}t4. All sixteen subset choices are
required. The largest consecutive gap among their subset sums is at most
max\_\allowbreak{}i(t\_\allowbreak{}i-sum\_\allowbreak{}(h\textless i)t\_\allowbreak{}h). The first three terms are at most
t3\textless=\allowbreak{}U0/2. The layer and source inequalities give

\begin{lstlisting}
U0<=1168+w-2I-2delta+2z<=1252,
\end{lstlisting}

so these three terms are at most626\textless L.

For the last term,\allowbreak{} t4 is necessarily one side\textquotesingle s second
occurrence R\_\allowbreak{}v+\allowbreak{}b\_\allowbreak{}v; the other three contain R\_\allowbreak{}v+\allowbreak{}a\_\allowbreak{}v and both
opposite occurrences. Thus

\begin{lstlisting}
t4-t1-t2-t3=b_v-a_v-2R_o-a_o-b_o
             <=U-8-2R_o.
\end{lstlisting}

Crucially,\allowbreak{} use the offset sum U rather than the much larger layer mass
U0. Since J+\allowbreak{}U\textless=\allowbreak{}458+\allowbreak{}w-A and J\textgreater=\allowbreak{}2,\allowbreak{} U\textless=\allowbreak{}512.
Since R\_\allowbreak{}o\textgreater=\allowbreak{}19,\allowbreak{} this last gap is at most466\textless L.
Therefore every translate of the core central interval by an actual
inner subset sum overlaps or abuts the next,\allowbreak{} filling {[}H,\allowbreak{}C1-H{]}.

For the outer occurrences,\allowbreak{} selecting none or all already suffices
because

\begin{lstlisting}
C1-V0-2H>=-26+4I+mu+rho+3delta+2z-2w
         >=85-6z>=1.
\end{lstlisting}

This fills {[}H,\allowbreak{}S-H{]}. Complement symmetry in the same fixed boundary
strip now gives z,\allowbreak{}w even,\allowbreak{} w\textless=\allowbreak{}z\textless=\allowbreak{}14,\allowbreak{} and both targets
have eta\_\allowbreak{}v=\allowbreak{}w/2 holes. The old coarse all-sixteen inner estimate is not
needed.

After localization,\allowbreak{} H\textless=\allowbreak{}250-I-delta and
R\textless=\allowbreak{}72+\allowbreak{}rho+\allowbreak{}H\textless=\allowbreak{}72+\allowbreak{}tau\textless=\allowbreak{}384-I-delta. Thus
T\_\allowbreak{}r-R\textgreater=\allowbreak{}13+\allowbreak{}delta. The first absolute-core gap f is below both
R\_\allowbreak{}v: if one side had all its entries below f,\allowbreak{} its total would be at
most f-1\textless R,\allowbreak{} contradicting C\_\allowbreak{}v\textgreater=\allowbreak{}T\_\allowbreak{}r\textgreater R.
All peeled occurrences exceed f,\allowbreak{} so f survives and w\textgreater=\allowbreak{}2.

Every integral ordinary center in {[}1-M\_\allowbreak{}v,\allowbreak{}S\_\allowbreak{}v{]} is active. To prove
this without an endpoint assumption,\allowbreak{} use R\_\allowbreak{}o\textless=\allowbreak{}r\_\allowbreak{}o+\allowbreak{}tau\_\allowbreak{}o and
tau\_\allowbreak{}v\textgreater=\allowbreak{}r\_\allowbreak{}v-3:

\begin{lstlisting}
S_v-H-R_o>=T_r+7r+2I-594-2w+2delta>=39+2delta.
\end{lstlisting}

The last integer quadratic has minimum39 at r26. A target hole u lies
within H of its upper endpoint,\allowbreak{} so all u-P\_\allowbreak{}o are positive hull
integers; at most eta\_\allowbreak{}v\textless=\allowbreak{}7 are missing and r\_\allowbreak{}o\textgreater=\allowbreak{}18
guarantees a target hit. For a nonpositive center u\textgreater=\allowbreak{}1-M\_\allowbreak{}v,\allowbreak{}
the value u+\allowbreak{}M\_\allowbreak{}v lies in the complete lower target interval,\allowbreak{} since
S\_\allowbreak{}v-M\_\allowbreak{}v\textgreater=\allowbreak{}C\_\allowbreak{}v\textgreater H. This proves the activity
claim.

Finally K-delta\textgreater=\allowbreak{}16-w+\allowbreak{}I+\allowbreak{}tau\textgreater0 makes both leading
translates active and their shifts integral. A nonintegral v3 or v4
would force respectively \{own2\_\allowbreak{}v,\allowbreak{}C6\_\allowbreak{}o\} or \{own3\_\allowbreak{}v,\allowbreak{}C7\_\allowbreak{}o\} to be
the two missing singleton units. With only these holes H\_\allowbreak{}v and alpha\_\allowbreak{}v
remain active,\allowbreak{} contradicting the integrality of
H\_\allowbreak{}v-alpha\_\allowbreak{}v=\allowbreak{}lambda\_\allowbreak{}v+\allowbreak{}v3+\allowbreak{}v4. Two nonintegral lower coordinates would
require four holes. A nonintegral v2 would make O\_\allowbreak{}v,\allowbreak{} H\_\allowbreak{}v,\allowbreak{}alpha\_\allowbreak{}v
three missing singletons. All outer coordinates are integral.

\hypertarget{the-first-gap-profiles-and-two-cap-exclusion}{%
\subsection{3. The first-gap profiles and two-cap
exclusion}\label{the-first-gap-profiles-and-two-cap-exclusion}}

Write t\_\allowbreak{}v=\allowbreak{}d\_\allowbreak{}v-M,\allowbreak{} where d\_\allowbreak{}v=\allowbreak{}v1-v2,\allowbreak{} and define T\_\allowbreak{}v to mean that O\_\allowbreak{}v
is active or both beta\_\allowbreak{}v,\allowbreak{}H\_\allowbreak{}v are active. Sharp raw gaps give

\begin{lstlisting}
T_v => t_v<=22+w/2<=29,
beta_v and O_v active => 2t_v<=27+w/2<=34.
\end{lstlisting}

Both cap predicates cannot hold. The four-star estimate is
2X+\allowbreak{}4R+\allowbreak{}delta\textgreater=\allowbreak{}548,\allowbreak{} and the high-tail count gives
d\_\allowbreak{}s+\allowbreak{}d\_\allowbreak{}t+\allowbreak{}G+\allowbreak{}delta-2K\textgreater=\allowbreak{}642. With Y=\allowbreak{}2M-2K+\allowbreak{}44+\allowbreak{}w+\allowbreak{}delta,\allowbreak{} two caps
and the leading counts imply

\begin{lstlisting}
2Y>=1190-4R-delta,
Y<=616+3w-2R-2I-2tau-delta.
\end{lstlisting}

Consequently 4I+\allowbreak{}4tau+\allowbreak{}delta\textless=\allowbreak{}42+\allowbreak{}6w,\allowbreak{} contrary to
tau\textgreater=\allowbreak{}62+\allowbreak{}2z and w\textless=\allowbreak{}z. At least one O singleton is
missing,\allowbreak{} and at least one cap predicate holds.

Let k\_\allowbreak{}v count core entries below the first gap f and k=\allowbreak{}k\_\allowbreak{}s+\allowbreak{}k\_\allowbreak{}t. The
greedy first-gap rule gives their total f-1 and interval completeness
below f. The common rank lower bound is

\begin{lstlisting}
I+tau>=69f-1-(f+34.5)k+k²/4
             +(r_s-34-(k_s-k_t)/2)².
\end{lstlisting}

Since f\textless R/2\textless=\allowbreak{}192 and k(k+\allowbreak{}2)/4\textless=\allowbreak{}f-1,\allowbreak{}
k\textless=\allowbreak{}26. For f\textgreater=\allowbreak{}7,\allowbreak{}k\textless=\allowbreak{}4 the displayed lower
bound is at least320\textgreater312. For k\textgreater=\allowbreak{}5 it is at least
(-k³+\allowbreak{}68k²-4k+\allowbreak{}272)/4\textgreater=\allowbreak{}1827/4\textgreater312. Thus
f\textless=\allowbreak{}6.

The complete low multisets are empty,\allowbreak{} \{1\},\allowbreak{} \{1,\allowbreak{}1\},\allowbreak{} \{1,\allowbreak{}2\},\allowbreak{}
\{1,\allowbreak{}1,\allowbreak{}2\},\allowbreak{} and at f6 either \{1,\allowbreak{}1,\allowbreak{}3\} or \{1,\allowbreak{}2,\allowbreak{}2\}. These are fixed by
the total f-1,\allowbreak{} interval completeness and multiplicity at most two. Put b
for the mass excess of the coordinatewise minimal side profiles,\allowbreak{} B for
their two maxima,\allowbreak{} and sigma=\allowbreak{}tau-b\textgreater=\allowbreak{}0. If one distinguished
side has r=\allowbreak{}34+\allowbreak{}x,\allowbreak{} their constants are:

\begin{longtable}[]{@{}rlrr@{}}
\toprule\noalign{}
f & I+\allowbreak{}b & B & minimum tau \\
\midrule\noalign{}
\endhead
\bottomrule\noalign{}
\endlastfoot
1 & I+\allowbreak{}68 & 70 & 68 \\
2 & 101+\allowbreak{}x(x-1) & 71 & 85 \\
3 & I+\allowbreak{}132 & 72 & 132 \\
4 split & I+\allowbreak{}199 & 74 & 199 \\
4 together & 199+\allowbreak{}(x-1)² & 74 & 168 \\
5 & 228+\allowbreak{}x(x-1) & 75 & 212 \\
6 both low multisets & 294+\allowbreak{}x(x-1) & 77 & 278 \\
\end{longtable}

In particular f\textgreater=\allowbreak{}2 implies tau\textgreater=\allowbreak{}85,\allowbreak{} and
I+\allowbreak{}b-2B\textgreater=\allowbreak{}-41 in every such row.

\hypertarget{no-c7-singleton-is-empty}{%
\subsection{4. No C7 singleton is
empty}\label{no-c7-singleton-is-empty}}

If C7\_\allowbreak{}v were empty,\allowbreak{} it and at least one missing O would be the two
holes,\allowbreak{} so C6\_\allowbreak{}v is active. Integral activity localization makes C7\_\allowbreak{}v
HIGH,\allowbreak{} hence lambda\_\allowbreak{}v\textgreater S\_\allowbreak{}v+\allowbreak{}M\_\allowbreak{}o. In the leading resource the
entirely HIGH set lambda\_\allowbreak{}v+\allowbreak{}(D union(M\_\allowbreak{}v+\allowbreak{}D)) has at least
d+\allowbreak{}(r\_\allowbreak{}v+\allowbreak{}4)=\allowbreak{}r\_\allowbreak{}v+\allowbreak{}81 points. At least three of the four higher native
singletons V\_\allowbreak{}v,\allowbreak{}U\_\allowbreak{}v,\allowbreak{}H\_\allowbreak{}v,\allowbreak{}O\_\allowbreak{}v are retained,\allowbreak{} and each loses r\_\allowbreak{}v+\allowbreak{}5
points. Thus delta\textgreater=\allowbreak{}4r\_\allowbreak{}v+\allowbreak{}96.

If f\textgreater=\allowbreak{}2,\allowbreak{}
A\textgreater=\allowbreak{}I+\allowbreak{}4r\_\allowbreak{}v+\allowbreak{}96+\allowbreak{}tau\textgreater=\allowbreak{}132+\allowbreak{}96+\allowbreak{}85=\allowbreak{}313\textgreater312,\allowbreak{}
impossible. Here I+\allowbreak{}4r=\allowbreak{}(r-32)²+\allowbreak{}132.

At f1 write tau=\allowbreak{}68+\allowbreak{}sigma. The source bound yields

\begin{lstlisting}
(r_v-32)²+sigma<=2+w<=16.
\end{lstlisting}

Thus 28\textless=\allowbreak{}r\_\allowbreak{}v\textless=\allowbreak{}36,\allowbreak{} the opposite size is32..40,\allowbreak{} and
sigma\textless=\allowbreak{}16. Each side contains the consecutive prefix
2,\allowbreak{}...,\allowbreak{}r\_\allowbreak{}v-sigma\_\allowbreak{}v+\allowbreak{}1,\allowbreak{} at least 2,\allowbreak{}...,\allowbreak{}13; every remaining coefficient is
at most r\_\allowbreak{}v+\allowbreak{}1+\allowbreak{}sigma\_\allowbreak{}v\textless=\allowbreak{}57. The prefix support {[}0,\allowbreak{}L{]}
minus\{1,\allowbreak{}L-1\},\allowbreak{} with L\textgreater=\allowbreak{}90,\allowbreak{} keeps exactly those two reflected
holes under every remaining coefficient. Therefore both side hole counts
are2,\allowbreak{} the absolute core has exactly two holes,\allowbreak{} z=\allowbreak{}2,\allowbreak{} and the persisted
gap1 gives w=\allowbreak{}2.

If the missing O were opposite C7\_\allowbreak{}v,\allowbreak{} all four higher v singletons would
be retained and delta\textgreater=\allowbreak{}5r\_\allowbreak{}v+\allowbreak{}101. Then
I+\allowbreak{}5r\_\allowbreak{}v+\allowbreak{}sigma\textless=\allowbreak{}129+\allowbreak{}w\textless=\allowbreak{}131,\allowbreak{} whereas the integer minimum
of I+\allowbreak{}5r is164. Therefore O\_\allowbreak{}v is the other hole. It is HIGH,\allowbreak{} and its
four fixed labels O\_\allowbreak{}v+\allowbreak{}v\_\allowbreak{}i are HIGH; they supply four additional fixed
credits beyond the baseline nonpositive witnesses. Hence
delta\textgreater=\allowbreak{}4r\_\allowbreak{}v+\allowbreak{}100. With w2 the source gives
(r\_\allowbreak{}v-32)²+\allowbreak{}sigma\textless=\allowbreak{}0. Thus r\_\allowbreak{}v32,\allowbreak{}I4,\allowbreak{}sigma0,\allowbreak{}delta228,\allowbreak{}A300.

Now m0=\allowbreak{}-128,\allowbreak{}T=\allowbreak{}m0+\allowbreak{}A=\allowbreak{}172. Source envelopes give
J+\allowbreak{}U\textless=\allowbreak{}160,\allowbreak{}J\textgreater=\allowbreak{}154,\allowbreak{}U\textgreater=\allowbreak{}6,\allowbreak{} so U6,\allowbreak{}J154. Both
signs have (a,\allowbreak{}b)=\allowbreak{}(1,\allowbreak{}2),\allowbreak{} for which the four-offset fringe counts
F\_\allowbreak{}s,\allowbreak{}F\_\allowbreak{}t are at most5 per sign. These F values are distinct from the
baseline core-maximum sum B in the profile table. The full fringe
comparison gives

\begin{lstlisting}
T+U<=J+2+2(F_s+F_t),
\end{lstlisting}

and would require178\textless=\allowbreak{}154+\allowbreak{}2+\allowbreak{}20=\allowbreak{}176. This excludes the final
equality. Therefore every C7 singleton is retained.

\hypertarget{rooted-forests-reduce-to-f1-and-exclude-simultaneous-betao-activity}{%
\subsection{5. Rooted forests reduce to f1 and exclude simultaneous
beta/O
activity}\label{rooted-forests-reduce-to-f1-and-exclude-simultaneous-betao-activity}}

Trim the four largest core entries on each sign,\allowbreak{} leaving an
actual61-point set Q. Its rank-certified overlap constant q satisfies

\begin{lstlisting}
4q=516-5B-sigma+R.
\end{lstlisting}

On a cap side v use the positive arm L\_\allowbreak{}v-V\_\allowbreak{}v-U\_\allowbreak{}v and the negative arm
L\_\allowbreak{}o-C7\_\allowbreak{}o-C6\_\allowbreak{}o. If V\_\allowbreak{}v is retained,\allowbreak{} C7\_\allowbreak{}o is also retained and at
most one of U\_\allowbreak{}v,\allowbreak{}C6\_\allowbreak{}o can be empty. Thus choose one full two-edge arm
and the first edge of the other; every selected resource is retained.
Designate L\_\allowbreak{}v as a root. The same original deficiency budget then gives

\begin{lstlisting}
t_v>=K-M+3q-delta
   >=101-w+I+b-2B+sigma/4.
\end{lstlisting}

For every f\textgreater=\allowbreak{}2 this is at
least60-w\textgreater=\allowbreak{}46\textgreater29,\allowbreak{} contradicting the cap.

If V\_\allowbreak{}v is empty,\allowbreak{} the other hole must be O\_\allowbreak{}o: choosing O\_\allowbreak{}v instead
would leave both O\_\allowbreak{}o and beta\_\allowbreak{}v,\allowbreak{}H\_\allowbreak{}v active,\allowbreak{} giving two caps. Both H
singletons are retained. Branch B would retain both beta singletons and
again give two caps; in A the retained beta pair is sole on v. Hence
beta\_\allowbreak{}v,\allowbreak{}O\_\allowbreak{}v are active and the sharp double-gap upper bound is
t\_\allowbreak{}v\textless=\allowbreak{}17. The empty V\_\allowbreak{}v is HIGH,\allowbreak{} since leading activity implies
V\_\allowbreak{}v\textgreater1-M\_\allowbreak{}v. Both H\_\allowbreak{}v,\allowbreak{}O\_\allowbreak{}v are retained HIGH resources and
each loses r\_\allowbreak{}v+\allowbreak{}5 native points. Use edges L\_\allowbreak{}v-U\_\allowbreak{}v and L\_\allowbreak{}o-C7\_\allowbreak{}o,\allowbreak{}
root L\_\allowbreak{}v and include those two unused losses. This gives

\begin{lstlisting}
t_v>=-18-w+I+b-B+2r_v+sigma/2.
\end{lstlisting}

The f1 integer minimum is47-w\textgreater=\allowbreak{}33; every f\textgreater=\allowbreak{}2 is
at least48-w\textgreater=\allowbreak{}34. Both contradict17. Thus every cap side has
V\_\allowbreak{}v retained,\allowbreak{} and f=\allowbreak{}1.

At f1 the rooted lower bound is t\_\allowbreak{}v\textgreater=\allowbreak{}29-w+\allowbreak{}I+\allowbreak{}sigma/4. If
beta\_\allowbreak{}v and O\_\allowbreak{}v were simultaneously active,\allowbreak{} the double gap would force

\begin{lstlisting}
4I+sigma<=5w-62<=8.
\end{lstlisting}

For w\textless=\allowbreak{}12 this is already impossible. For w14,\allowbreak{} I\textless=\allowbreak{}2 and
sigma\textless=\allowbreak{}8; the consecutive core prefixes and small final
coefficients then preserve exactly the two side holes as in the
preceding dense-profile argument. This gives z=\allowbreak{}w2,\allowbreak{} which contradicts the
same inequality. No side has both beta\_\allowbreak{}v and O\_\allowbreak{}v active.

\hypertarget{exhaustive-activity-forms-and-the-terminal-contradiction}{%
\subsection{6. Exhaustive activity forms and the terminal
contradiction}\label{exhaustive-activity-forms-and-the-terminal-contradiction}}

If both O singletons are missing,\allowbreak{} B is impossible because both beta,\allowbreak{}H
caps would hold; in A the retained beta pair has a sole actual owner. If
exactly O\_\allowbreak{}v is missing,\allowbreak{} O\_\allowbreak{}o is active,\allowbreak{} hence beta\_\allowbreak{}o is inactive. B
would make beta\_\allowbreak{}o the other hole and leave beta\_\allowbreak{}v,\allowbreak{}H\_\allowbreak{}v active,\allowbreak{} giving
two caps. In A,\allowbreak{} either H\_\allowbreak{}v is the other hole or the whole beta pair is
empty; every other choice would again leave both cap sides active. Thus
exactly the following A whole-hole forms remain:

\begin{lstlisting}
{O_v,beta pair}, {O_s,O_t}, {O_v,H_v}.
\end{lstlisting}

All unlisted whole units are retained. In the first and third forms
choose the side u opposite the missing O\_\allowbreak{}v; in the middle choose the
sole beta owner. That side has the cap predicate. The retained singleton
centers alpha\_\allowbreak{}o,\allowbreak{}C5\_\allowbreak{}u and their exact identity imply
t\_\allowbreak{}u\textgreater=\allowbreak{}2X+\allowbreak{}G+\allowbreak{}2. Hence

\begin{lstlisting}
2X+G<=20+w/2.
\end{lstlisting}

At f1 write sigma\_\allowbreak{}v=\allowbreak{}sum\_\allowbreak{}i(x\_\allowbreak{}vi-i-1),\allowbreak{} xi\_\allowbreak{}v=\allowbreak{}R\_\allowbreak{}v-r\_\allowbreak{}v-1\textgreater=\allowbreak{}0,\allowbreak{}
and sigma=\allowbreak{}sigma\_\allowbreak{}s+\allowbreak{}sigma\_\allowbreak{}t,\allowbreak{}xi=\allowbreak{}xi\_\allowbreak{}s+\allowbreak{}xi\_\allowbreak{}t. Source mass gives

\begin{lstlisting}
I+sigma+delta<=230+w.
\end{lstlisting}

Trim six largest core entries per sign,\allowbreak{} leaving57 actual points. Its
span is at most58+\allowbreak{}(sigma-xi)/6,\allowbreak{} so the certified intersection lower
bound at gapq is55-(sigma-xi)/6-q. All U\_\allowbreak{}v,\allowbreak{}V\_\allowbreak{}v,\allowbreak{}C7\_\allowbreak{}v are retained in
every form. On each target the retained tree L\_\allowbreak{}v-V\_\allowbreak{}v-U\_\allowbreak{}v with edge
L\_\allowbreak{}v-C7\_\allowbreak{}v has three gaps X\_\allowbreak{}v,\allowbreak{}g\_\allowbreak{}v,\allowbreak{}X\_\allowbreak{}o. Both targets together give,\allowbreak{} in
one budget,\allowbreak{}

\begin{lstlisting}
2X+G>=330-sigma+xi-delta.
\end{lstlisting}

Combining these three inequalities forces

\begin{lstlisting}
80+I+xi<=3w/2<=21,
\end{lstlisting}

a contradiction. This excludes the whole P2 stratum,\allowbreak{} including all
positions of empty endpoints and all paired-unit ownership choices.

\hypertarget{exact-finite-arithmetic-replay}{%
\subsection{7. Exact finite arithmetic
replay}\label{exact-finite-arithmetic-replay}}

\texttt{check85\_\allowbreak{}\allowbreak{}p2.py} checks the range,\allowbreak{} width,\allowbreak{}
activity quadratics,\allowbreak{} first-gap polynomial and profile minima,\allowbreak{} C7
equality,\allowbreak{} rooted-forest scalar margins,\allowbreak{} and terminal contradiction. It
is an arithmetic checker of these written lemmas,\allowbreak{} not an enumeration of
actual source sets or activity descriptors.

\clearpage
\hypertarget{fresh-n85-positive-case-prefix-lemmas}{%
\section{Fresh n=\allowbreak{}85 positive-case prefix
lemmas}\label{fresh-n85-positive-case-prefix-lemmas}}

This file contains the common positive-case notation,\allowbreak{} the P3
localization and activity classification,\allowbreak{} and its first-gap4 exclusion.
The completed first-gap1/2/3 P3 terminals are in p3\_\allowbreak{}terminal.md. The
complete P4 source-seed argument,\allowbreak{} classifier,\allowbreak{} and terminals are in
p4\_\allowbreak{}fresh.md. There are no open transport branches in these files.

The new structure and source results are proved in
inverse/proof85\_\allowbreak{}inverse.md and foundations/source\_\allowbreak{}triple\_\allowbreak{}fringe85.md
(paths relative to progress85). Throughout,\allowbreak{} the source has77 points,\allowbreak{} its
core E has69 labelled points including zero,\allowbreak{} and the two
positive/negative core sizes satisfy a+\allowbreak{}b=\allowbreak{}68. Its triangular baseline
is1190 and I=\allowbreak{}(a-34)\^{}2. Put A=\allowbreak{}I+\allowbreak{}tau+\allowbreak{}delta,\allowbreak{} m0=\allowbreak{}77j-280-w,\allowbreak{} T=\allowbreak{}m0+\allowbreak{}A,\allowbreak{}
R=\allowbreak{}R\_\allowbreak{}s+\allowbreak{}R\_\allowbreak{}t,\allowbreak{} M=\allowbreak{}M\_\allowbreak{}s+\allowbreak{}M\_\allowbreak{}t,\allowbreak{} and L=\allowbreak{}M-R. The first two peel offsets on both
signs total U; the third offsets total C and the last-minus-third
offsets total J,\allowbreak{} so L=\allowbreak{}C+\allowbreak{}J.

The exact source identity and the source envelopes are

\begin{itemize}
\tightlist
\item
  S=\allowbreak{}2K+\allowbreak{}1470-77j+\allowbreak{}w-delta.
\item
  2K=\allowbreak{}m0+\allowbreak{}A+\allowbreak{}4R+\allowbreak{}U+\allowbreak{}2C+\allowbreak{}J.
\item
  K\textless=\allowbreak{}2R+\allowbreak{}C+\allowbreak{}166,\allowbreak{} J+\allowbreak{}U\textless=\allowbreak{}332-m0-A.
\item
  J\textgreater=\allowbreak{}m0+\allowbreak{}A-18,\allowbreak{} L\textgreater=\allowbreak{}m0+\allowbreak{}A-8.
\item
  K-delta\textgreater=\allowbreak{}77j-142-w+\allowbreak{}I+\allowbreak{}tau.
\item
  K-M\textgreater=\allowbreak{}R+\allowbreak{}77j-440-w+\allowbreak{}A.
\item
  T\textless=\allowbreak{}165,\allowbreak{} equivalently A\textless=\allowbreak{}445-77j+\allowbreak{}w.
\end{itemize}

The deliberately weaker envelope A\textless=\allowbreak{}452-77j+\allowbreak{}w is used below
where it already suffices; it follows from the same source bounds and is
not a provisional assertion. Positive-core mu\textgreater=\allowbreak{}62. Use
rho\textgreater=\allowbreak{}49 initially,\allowbreak{} and rho\textgreater=\allowbreak{}58 once both sizes
are\textgreater=\allowbreak{}21. The new inverse identities give

\begin{itemize}
\tightlist
\item
  rho+\allowbreak{}kappa+\allowbreak{}delta\textless=\allowbreak{}456-77j+\allowbreak{}3z-I.
\item
  U0-C0+\allowbreak{}delta-2z\textless=\allowbreak{}-166-2I-77j.
\item
  V0-C1+\allowbreak{}delta+\allowbreak{}2z\textless=\allowbreak{}-292-2I-77j-mu-rho.
\end{itemize}

\hypertarget{p3-localization-before-a-first-gap-profile}{%
\subsection{P3 localization before a first-gap
profile}\label{p3-localization-before-a-first-gap-profile}}

Here j=\allowbreak{}3. Initially 1\textless=\allowbreak{}z\textless=\allowbreak{}14,\allowbreak{}w\textless=\allowbreak{}4z. The weaker
source bound gives A\textless=\allowbreak{}221+\allowbreak{}w\textless=\allowbreak{}277 and
H=\allowbreak{}tau-mu\textless=\allowbreak{}215-I-delta. The integer inequality
T\_\allowbreak{}r+\allowbreak{}(r-34)\^{}2\textgreater=\allowbreak{}397 implies that each core side mass
exceeds H,\allowbreak{} so the core central strip exists.

Its width L0=\allowbreak{}C0-2H is at least1037+\allowbreak{}2I+\allowbreak{}delta. The total U0 of the first
two peeled magnitudes on each sign is at most1098-2I-2delta,\allowbreak{} using
C0\textless=\allowbreak{}1467 and the inner row. For their four sorted weights,\allowbreak{} the
first three gap terms are\textless=\allowbreak{}U0/2\textless=\allowbreak{}549. The other three
weights at the final gap total at least72+\allowbreak{}r\_\allowbreak{}min. Since
A\textgreater=\allowbreak{}I+\allowbreak{}mu,\allowbreak{} r\_\allowbreak{}min\textgreater=\allowbreak{}20. Thus

\begin{lstlisting}
L0-(t4-t1-t2-t3)>=1037-1098+2(72+20)=123.
\end{lstlisting}

All sixteen inner subset translates overlap. At the outer peel the
none/all translates overlap because

\begin{lstlisting}
C1-V0-2H>=204+4I+3delta+2z>0,
\end{lstlisting}

using mu\textgreater=\allowbreak{}62,\allowbreak{}rho\textgreater=\allowbreak{}49. Consequently w,\allowbreak{}z are even,\allowbreak{}
0\textless=\allowbreak{}w\textless=\allowbreak{}z\textless=\allowbreak{}14,\allowbreak{} and the two targets have equal hole
counts and localized boundary holes.

\hypertarget{p3-finite-dense-seed-and-boundary-activity}{%
\subsection{P3 finite dense seed and boundary
activity}\label{p3-finite-dense-seed-and-boundary-activity}}

After localization tau\textless=\allowbreak{}235,\allowbreak{} C0\textless=\allowbreak{}1425,\allowbreak{} and
r\_\allowbreak{}min\textgreater=\allowbreak{}21. Six absent values in{[}1,\allowbreak{}17{]} force
tau\textgreater=\allowbreak{}276,\allowbreak{} so at least twelve distinct values occur there. The
exact finite lemma checked by dense\_\allowbreak{}seed\_\allowbreak{}check.py is:

Every twelve-element subset B of{[}1,\allowbreak{}17{]},\allowbreak{} with total T\_\allowbreak{}B,\allowbreak{} satisfies
{[}11,\allowbreak{}T\_\allowbreak{}B-11{]} subset FS(B),\allowbreak{} and its lower half has at most five
holes.

The checker covers all C(17,\allowbreak{}12)=\allowbreak{}6188 subsets by integer bitset addition.
To adjoin the remaining labelled core occurrences,\allowbreak{} each magnitude
is\textless=\allowbreak{}47+\allowbreak{}235=\allowbreak{}282. For n\textless=\allowbreak{}33 selected occurrences,\allowbreak{} a failed
join x\textgreater T\_\allowbreak{}B-21 implies x\textgreater=\allowbreak{}58; at least35
remaining entries would give C0\textgreater=\allowbreak{}78+\allowbreak{}35*58\textgreater1425.
For n\textgreater=\allowbreak{}34 the selected total is\textgreater=\allowbreak{}306 and
x\textless=\allowbreak{}282\textless=\allowbreak{}T\_\allowbreak{}B-21. Every join succeeds,\allowbreak{} and hence
z\textless=\allowbreak{}10. The weaker source bound then gives A\textless=\allowbreak{}231,\allowbreak{}
H\textless=\allowbreak{}169-I-delta,\allowbreak{} tau\textgreater=\allowbreak{}62+\allowbreak{}2z,\allowbreak{} and
delta\textless=\allowbreak{}159-z-I.

Five low holes among{[}1,\allowbreak{}10{]} omit five singleton values of
sum\textless=\allowbreak{}40 and force tau\textgreater=\allowbreak{}5*68+\allowbreak{}30-80=\allowbreak{}290. Therefore
z\textless=\allowbreak{}8. For the terminal first gaps m\textless=\allowbreak{}3,\allowbreak{} four low holes
have total\textless=\allowbreak{}m+\allowbreak{}8+\allowbreak{}9+\allowbreak{}10\textless=\allowbreak{}30 and force tau\textgreater=\allowbreak{}232.
The weaker consequence T\textless=\allowbreak{}167 of the proved T\textless=\allowbreak{}165 gives
A\textless=\allowbreak{}216+\allowbreak{}w\textless=\allowbreak{}224,\allowbreak{} excluding this possibility. Thus at
m\textless=\allowbreak{}3,\allowbreak{} z,\allowbreak{}w\textless=\allowbreak{}6.

The source row K-delta\textgreater=\allowbreak{}89-w+\allowbreak{}I+\allowbreak{}tau proves both leadings
active. From H\textless=\allowbreak{}169-I-delta and
R\textless=\allowbreak{}68+\allowbreak{}tau\textless=\allowbreak{}299-I-delta,\allowbreak{}

\begin{lstlisting}
S_v-H-R_o>=T_r+5r-468+2I+2delta>=141+2delta.
\end{lstlisting}

Also T\_\allowbreak{}r-R\textgreater=\allowbreak{}98+\allowbreak{}delta. A target-hole center has at least
r\_\allowbreak{}o\textgreater=\allowbreak{}21 distinct positive images under opposite core
entries,\allowbreak{} all within its hull,\allowbreak{} whereas it has at most w/2\textless=\allowbreak{}5
target holes. Nonholes use zero; a center in{[}1-M\_\allowbreak{}v,\allowbreak{}0{]} uses M\_\allowbreak{}v and
the complete lower target interval. Every integral center
in{[}1-M\_\allowbreak{}v,\allowbreak{}S\_\allowbreak{}v{]} is therefore active,\allowbreak{} and an empty integral center
above an active lower center is HIGH\textgreater=\allowbreak{}S\_\allowbreak{}v+\allowbreak{}1.

At the first failed greedy inequality,\allowbreak{} the entries below the gap m total
m-1 and generate{[}0,\allowbreak{}m-1{]},\allowbreak{} while every remaining entry exceeds m. If
R\_\allowbreak{}v\textless=\allowbreak{}m,\allowbreak{} that entire side would lie in the prefix,\allowbreak{} giving
C\_\allowbreak{}v\textless=\allowbreak{}m-1\textless R\textless T\_\allowbreak{}r,\allowbreak{} impossible. Thus
m\textless min(R\_\allowbreak{}s,\allowbreak{}R\_\allowbreak{}t); it survives both peels,\allowbreak{} w\textgreater=\allowbreak{}2,\allowbreak{} and
the target holes are exactly S\_\allowbreak{}v-h for the lower source holes h.

\hypertarget{p3-outer-integrality-and-exhaustive-activity-forms}{%
\subsection{P3 outer integrality and exhaustive activity
forms}\label{p3-outer-integrality-and-exhaustive-activity-forms}}

Write U\_\allowbreak{}v=\allowbreak{}own2\_\allowbreak{}v,\allowbreak{} V\_\allowbreak{}v=\allowbreak{}own3\_\allowbreak{}v,\allowbreak{} C\_\allowbreak{}iv=\allowbreak{}lambda\_\allowbreak{}v-o\_\allowbreak{}(i-3)
for4\textless=\allowbreak{}i\textless=\allowbreak{}7,\allowbreak{} X\_\allowbreak{}v=\allowbreak{}v4-M\_\allowbreak{}v,\allowbreak{} X=\allowbreak{}X\_\allowbreak{}s+\allowbreak{}X\_\allowbreak{}t,\allowbreak{} G=\allowbreak{}g\_\allowbreak{}s+\allowbreak{}g\_\allowbreak{}t. The
active leadings make lambda\_\allowbreak{}v integral. Each lower coordinate v3 has
singleton witnesses U\_\allowbreak{}v,\allowbreak{}C6\_\allowbreak{}o,\allowbreak{} and v4 has V\_\allowbreak{}v,\allowbreak{}C7\_\allowbreak{}o; the eight
witnesses are distinct. Two nonintegral lower coordinates require four
holes. A sole nonintegral lower coordinate requires its two witnesses
and one of H\_\allowbreak{}v,\allowbreak{}alpha\_\allowbreak{}v missing. If alpha\_\allowbreak{}v is missing,\allowbreak{} retained
O\_\allowbreak{}v,\allowbreak{}H\_\allowbreak{}v force both top coordinates integral,\allowbreak{} a contradiction. If H\_\allowbreak{}v
is missing,\allowbreak{} retained O\_\allowbreak{}v,\allowbreak{}alpha\_\allowbreak{}v and C5\_\allowbreak{}o in A or beta\_\allowbreak{}v in B force
v2 and then the lower coordinate integral. A nonintegral top pair would
make O\_\allowbreak{}v,\allowbreak{}H\_\allowbreak{}v,\allowbreak{}alpha\_\allowbreak{}v and C5\_\allowbreak{}o in A or beta\_\allowbreak{}v in B missing. All
eight outer magnitudes are integral.

The branch-independent retained-resource forest rows are

\begin{lstlisting}
2X_v+X_o+g_v>=252-tau-delta,
2X_v+X_o+G>=305-tau-delta,
2X+4R+delta>=548.
\end{lstlisting}

The two candidate arms are leading\_\allowbreak{}v-\/-V\_\allowbreak{}v-\/-U\_\allowbreak{}v and
leading\_\allowbreak{}o-\/-C7\_\allowbreak{}o-\/-C6\_\allowbreak{}o,\allowbreak{} both with gaps X\_\allowbreak{}v,\allowbreak{}g\_\allowbreak{}v. At most one
vertex of one arm is empty. Restrict core translates using trim3 for the
first row and trim4 for the second. Positive and negative restricted
arms on one target are disjoint since M\textgreater R. Inside a target
every vertex at a point is retained; outside it at most one is empty.
Thus edge counts are bounded by retained-resource union deficiency,\allowbreak{} with
each full leading capacity counted once. Four full-core stars give the
third row.

The following five assertions exhaust the activity restrictions:

\begin{enumerate}
\def\labelenumi{\arabic{enumi}.}
\tightlist
\item
  O\_\allowbreak{}v,\allowbreak{}C6\_\allowbreak{}v cannot both be active.
\item
  O\_\allowbreak{}s,\allowbreak{}O\_\allowbreak{}t cannot both be active.
\item
  beta\_\allowbreak{}v,\allowbreak{}H\_\allowbreak{}v,\allowbreak{}C6\_\allowbreak{}v cannot all be active.
\item
  Active C5\_\allowbreak{}v implies active C6\_\allowbreak{}v.
\item
  beta\_\allowbreak{}v,\allowbreak{}O\_\allowbreak{}v,\allowbreak{}C7\_\allowbreak{}v cannot all be active.
\end{enumerate}

For1 and3,\allowbreak{} actual C6\_\allowbreak{}v activity gives
d\_\allowbreak{}v-M\textgreater=\allowbreak{}306-tau-delta\textgreater=\allowbreak{}85-w+\allowbreak{}I; O\_\allowbreak{}v activity or
beta\_\allowbreak{}v,\allowbreak{}H\_\allowbreak{}v activity gives a raw cap\textless=\allowbreak{}27,\allowbreak{} impossible. For5 the
lower bound is253-tau-delta\textgreater=\allowbreak{}32-w+\allowbreak{}I,\allowbreak{} exceeding the double
cap\textless=\allowbreak{}16. These conservative raw caps follow from the proved
O22+\allowbreak{}eta,\allowbreak{} beta/H21+\allowbreak{}eta,\allowbreak{} and double27/2+\allowbreak{}eta/2 counts with w\textless=\allowbreak{}10.
For2,\allowbreak{} with Y=\allowbreak{}2M-2K+\allowbreak{}44+\allowbreak{}w+\allowbreak{}delta,\allowbreak{} the high-tail count565 and star count548
yield189+\allowbreak{}4I+\allowbreak{}4tau+\allowbreak{}delta\textless=\allowbreak{}6w\textless=\allowbreak{}60,\allowbreak{} impossible. For4,\allowbreak{} active
C5\_\allowbreak{}v and empty C6\_\allowbreak{}v make C6\_\allowbreak{}v HIGH. The leading then loses all77
points of lambda\_\allowbreak{}v+\allowbreak{}D\_\allowbreak{}v,\allowbreak{} and at least three retained higher singletons
among C7,\allowbreak{}V,\allowbreak{}U,\allowbreak{}H,\allowbreak{}O each lose r\_\allowbreak{}v+\allowbreak{}5. This forces
delta\textgreater=\allowbreak{}92+\allowbreak{}3r\_\allowbreak{}v,\allowbreak{} but delta\textless=\allowbreak{}169-I and
min\_\allowbreak{}r(I+\allowbreak{}3r)=\allowbreak{}100 give a contradiction.

The literal owner identities,\allowbreak{} with W=\allowbreak{}v3+\allowbreak{}v4+\allowbreak{}o3+\allowbreak{}o4\textgreater2M,\allowbreak{} are

\begin{lstlisting}
C4_s+C4_t=-(s2+s3+s4+t2+t3+t4),
beta_s+beta_t+C5_s+C5_t=-2W,
C4_v+beta_v+C5_o=-2v3-2v4-o2-o3-o4,
C4_v+C5_o=-(beta_s+beta_t)-2W-d_o.
\end{lstlisting}

Comparison with the actual lower activity endpoint1-M\_\allowbreak{}v proves that
retained pair1 has a sole owner,\allowbreak{} the displayed triple cannot be active,\allowbreak{}
and the two retained B pairs have matching sole owners.

To exhaust whole-unit holes,\allowbreak{} at least one O is missing by2. If both are
missing,\allowbreak{} neither C6 can be the third A hole,\allowbreak{} because4 would also delete
its C5 singleton. Both C6 are retained. An empty beta pair gives A0.
Otherwise3 forces H missing on the active beta\textquotesingle s side;
the opposite beta is inactive and owner grammar puts pair1 on the
opposite side,\allowbreak{} giving A1. In B,\allowbreak{} both O missing would require one further
hole in each disjoint triple\{beta\_\allowbreak{}v,\allowbreak{}H\_\allowbreak{}v,\allowbreak{}C6\_\allowbreak{}v\},\allowbreak{} exceeding three
holes.

If just O\_\allowbreak{}o is missing,\allowbreak{} active O\_\allowbreak{}v forces empty C6\_\allowbreak{}v and C5\_\allowbreak{}v by1
and4. In A these are all three holes: the beta owner is v and pair1
owner o; active beta\_\allowbreak{}v,\allowbreak{}O\_\allowbreak{}v,\allowbreak{}C7\_\allowbreak{}v contradict5. In B the holes
are\{O\_\allowbreak{}v,\allowbreak{}C6\_\allowbreak{}o,\allowbreak{}Z\} with Z in\{beta\_\allowbreak{}v,\allowbreak{}H\_\allowbreak{}v,\allowbreak{}C6\_\allowbreak{}v\}; pair owners are v,\allowbreak{}
and4 excludes Z=\allowbreak{}C6\_\allowbreak{}v. In either remaining choice beta\_\allowbreak{}o,\allowbreak{}O\_\allowbreak{}o,\allowbreak{}C7\_\allowbreak{}o are
active,\allowbreak{} contradict5. Exactly these two forms remain:

\begin{longtable}[]{@{}lll@{}}
\toprule\noalign{}
form & whole-unit holes & actual owners \\
\midrule\noalign{}
\endhead
\bottomrule\noalign{}
\endlastfoot
A0 & beta pair,\allowbreak{}O\_\allowbreak{}s,\allowbreak{}O\_\allowbreak{}t & pair1 sole owner s or t \\
A1 & O\_\allowbreak{}s,\allowbreak{}O\_\allowbreak{}t,\allowbreak{}H\_\allowbreak{}v & beta sole owner v,\allowbreak{} pair1 sole owner o \\
\end{longtable}

All unlisted whole units are retained. The fixed HIGH/LOW label counts
give credits8 for A0 and11 for A1; the terminal proofs retain their
exact native/grouped deficiencies.

\hypertarget{p3-first-gap-profiles-and-the-gap4-exclusion}{%
\subsection{P3 first-gap profiles and the gap4
exclusion}\label{p3-first-gap-profiles-and-the-gap4-exclusion}}

The first-gap5 low multiset must be\{1,\allowbreak{}1,\allowbreak{}2\}. If k is the side
containing1,\allowbreak{}2,\allowbreak{} its profiles are(1,\allowbreak{}2,\allowbreak{}6,\allowbreak{}...,\allowbreak{}k+\allowbreak{}3) and(1,\allowbreak{}6,\allowbreak{}...,\allowbreak{}q+\allowbreak{}4),\allowbreak{} q=\allowbreak{}68-k.
Therefore

\begin{lstlisting}
I+tau>=228+(k-34)(k-35)>=228>A<=224.
\end{lstlisting}

For m6,\allowbreak{} the total count below m is\textless=\allowbreak{}3,\allowbreak{} giving
tau\textgreater=\allowbreak{}4(68-3)=\allowbreak{}260. The dense seed puts every lower hole
in{[}1,\allowbreak{}10{]},\allowbreak{} so the first gap satisfies m\textless=\allowbreak{}10. Thus
T\_\allowbreak{}(k\_\allowbreak{}v)\textless=\allowbreak{}m-1\textless=\allowbreak{}9 gives k\_\allowbreak{}v\textless=\allowbreak{}3.
For7\textless=\allowbreak{}m\textless=\allowbreak{}10,\allowbreak{} each side floor (m-k\_\allowbreak{}v)(r\_\allowbreak{}v-k\_\allowbreak{}v) is at
least4(r\_\allowbreak{}v-3),\allowbreak{} yielding tau\textgreater=\allowbreak{}4(68-6)=\allowbreak{}248,\allowbreak{} again impossible.
Hence m\textless=\allowbreak{}4.

At m4 the low multiset is\{1,\allowbreak{}2\},\allowbreak{} either split or together. The minimal
I+\allowbreak{}tau is199+\allowbreak{}chi+\allowbreak{}sigma,\allowbreak{} with chi=\allowbreak{}I if split and chi=\allowbreak{}(k-35)\^{}2 if k is
the together side. A value x in{[}5,\allowbreak{}12{]} absent from both cores adds at
least76-2x\textgreater=\allowbreak{}52 to the baseline,\allowbreak{} exceeding even the weaker
true-excess cap32. Hence the actual seed\{1,\allowbreak{}2,\allowbreak{}5,\allowbreak{}...,\allowbreak{}12\} is present; it
has exactly the hole4 and its reflection. Every remaining magnitude
is\textless=\allowbreak{}83. While selected total is below92,\allowbreak{} a failed join
x\textgreater T\_\allowbreak{}B-9 implies x\textgreater=\allowbreak{}63 and at least54 remaining
occurrences would force mass\textgreater1425; once T\_\allowbreak{}B\textgreater=\allowbreak{}92,\allowbreak{}
x\textless=\allowbreak{}83\textless=\allowbreak{}T\_\allowbreak{}B-9. All joins succeed,\allowbreak{} so z=\allowbreak{}w=\allowbreak{}2. The weaker
T\textless=\allowbreak{}167 cap gives chi+\allowbreak{}sigma+\allowbreak{}delta\textless=\allowbreak{}19; below we retain
the still weaker24,\allowbreak{} since it already suffices. Charged endpoints are
normal because delta\textless=\allowbreak{}24\textless r\_\allowbreak{}min+\allowbreak{}5.

In A0 let s be sole pair1 owner. The identity
H\_\allowbreak{}s=\allowbreak{}P+\allowbreak{}alpha\_\allowbreak{}s+\allowbreak{}alpha\_\allowbreak{}t+\allowbreak{}d\_\allowbreak{}t+\allowbreak{}2V+\allowbreak{}G gives
S\_\allowbreak{}s\textgreater=\allowbreak{}d\_\allowbreak{}t+\allowbreak{}2V+\allowbreak{}G+\allowbreak{}155-E. The opposite leading and
K\textless=\allowbreak{}M+\allowbreak{}R+\allowbreak{}21 give

\begin{lstlisting}
S_t+M+(2X+G)+(2X_t+g_t)<=R+1105+w-e.
\end{lstlisting}

The six- and three-edge forests use f=\allowbreak{}(146-sigma+\allowbreak{}R)/4 and yield

\begin{lstlisting}
S_t+M+5R/4-9sigma/4-2delta+3e<=1553/2+w.
\end{lstlisting}

The physical baseline tau\_\allowbreak{}t\textgreater=\allowbreak{}2r\_\allowbreak{}t-4,\allowbreak{}
R\_\allowbreak{}t\textgreater=\allowbreak{}r\_\allowbreak{}t+\allowbreak{}2 gives S\_\allowbreak{}t+\allowbreak{}M\textgreater=\allowbreak{}T\_\allowbreak{}r+\allowbreak{}6r+\allowbreak{}R+\allowbreak{}L+\allowbreak{}14. With
R\textgreater=\allowbreak{}74,\allowbreak{} L\textgreater=\allowbreak{}142-w+\allowbreak{}chi+\allowbreak{}sigma+\allowbreak{}delta,\allowbreak{} and
chi+\allowbreak{}sigma+\allowbreak{}delta\textless=\allowbreak{}24,\allowbreak{}

\begin{lstlisting}
T_r+6r+9chi/4+3e<=488.
\end{lstlisting}

As a function of physical r,\allowbreak{} chi is(r-34)\^{}2 in the split placement
and(r-35)\^{}2 or(r-33)\^{}2 in the together placements. All three
integer quadratic minima exceed600,\allowbreak{} contradicting488.

In A1 let s be beta owner and t pair1 owner. Normal endpoints at
beta\_\allowbreak{}s,\allowbreak{}P\_\allowbreak{}t,\allowbreak{}V\_\allowbreak{}t give
S\_\allowbreak{}s+\allowbreak{}3g\_\allowbreak{}s+\allowbreak{}4X\_\allowbreak{}s+\allowbreak{}2s4+\allowbreak{}t4\textless=\allowbreak{}1089+\allowbreak{}w+\allowbreak{}r\_\allowbreak{}s+\allowbreak{}delta-2e. Source mass,\allowbreak{} the
ordinary gap pair and stars imply

\begin{lstlisting}
T_r+6r+chi+sigma/4+283/2-2w+2e<=21R/4+9epsilon/2.
\end{lstlisting}

The profile radial estimate is2R+\allowbreak{}chi+\allowbreak{}delta\textless=\allowbreak{}178+\allowbreak{}w,\allowbreak{} derived from
K\textless=\allowbreak{}C+\allowbreak{}2B+\allowbreak{}174+\allowbreak{}sigma with B74 and J\textgreater=\allowbreak{}A-67-w. Using w2
and chi+\allowbreak{}sigma+\allowbreak{}delta\textless=\allowbreak{}24 yields

\begin{lstlisting}
T_r+6r+11chi/2+17sigma/8+13e/2<=380.
\end{lstlisting}

Each physical orientation has a larger global integer quadratic minimum.
Thus m4 is excluded in both forms. Together with the complete
first-gap1/2/3 proof in p3\_\allowbreak{}terminal.md,\allowbreak{} this excludes P3.

\hypertarget{p4-proof-reference}{%
\subsection{P4 proof reference}\label{p4-proof-reference}}

p4\_\allowbreak{}fresh.md contains the full seven-seed availability and propagation
proof,\allowbreak{} all six resulting profiles,\allowbreak{} the exact CORE8 classifier,\allowbreak{} and every
terminal card. check\_\allowbreak{}p4\_\allowbreak{}cards.py checks the stated exact seed
supports,\allowbreak{} propagation mass bounds,\allowbreak{} fixed-label incidence ledger,\allowbreak{}
adaptive scalar combinations,\allowbreak{} integer quadratics,\allowbreak{} and terminal
arithmetic.

\clearpage
\hypertarget{n85-p3-terminal-exclusions-using-the-fresh-universal-source-cap}{%
\section{n=\allowbreak{}85 P3 terminal exclusions using the fresh universal source
cap}\label{n85-p3-terminal-exclusions-using-the-fresh-universal-source-cap}}

This proof uses the new structure and the weaker consequence
T=\allowbreak{}m0+\allowbreak{}A\textless=\allowbreak{}167 of the proved universal source theorem
T\textless=\allowbreak{}165,\allowbreak{} where m0=\allowbreak{}-49-w and A=\allowbreak{}I+\allowbreak{}tau+\allowbreak{}delta. Thus

\begin{lstlisting}
A<=216+w.
\end{lstlisting}

The source is77 points,\allowbreak{} the core69 labelled points including zero,\allowbreak{} and
a+\allowbreak{}b=\allowbreak{}68. The baseline I is(a-34)\^{}2. Source mass is S=\allowbreak{}2K+\allowbreak{}1239+\allowbreak{}w-delta,\allowbreak{}
and full-source central-fringe bound K\textless=\allowbreak{}M+\allowbreak{}R+\allowbreak{}21. The exact source
offset bound is L\textgreater=\allowbreak{}A-57-w.

The localization/classification in
\texttt{fresh\_\allowbreak{}\allowbreak{}positive\_\allowbreak{}\allowbreak{}notes.md} supplies the two
forms A0 and A1. The dense seed lemma first gives z\textless=\allowbreak{}10. Five
low holes among{[}1,\allowbreak{}10{]} would avoid five values of sum at most40,\allowbreak{}
forcing tau\textgreater=\allowbreak{}5*68+\allowbreak{}30-80=\allowbreak{}290,\allowbreak{} impossible. Hence z\textless=\allowbreak{}8.
For a first gap m\textless=\allowbreak{}3,\allowbreak{} four low holes include m and have total at
most m+\allowbreak{}8+\allowbreak{}9+\allowbreak{}10\textless=\allowbreak{}30. They force tau\textgreater=\allowbreak{}4*68+\allowbreak{}20-60=\allowbreak{}232,\allowbreak{}
whereas A\textless=\allowbreak{}224. Therefore z\textless=\allowbreak{}6,\allowbreak{} w\textless=\allowbreak{}6.

At m=\allowbreak{}5,\allowbreak{} the profile I+\allowbreak{}tau\textgreater=\allowbreak{}228 contradicts A\textless=\allowbreak{}224. At
m\textgreater=\allowbreak{}6 the generic low-rank floor also contradicts the cap. At
m=\allowbreak{}4 the seed lemma gives z=\allowbreak{}w=\allowbreak{}2 and sigma+\allowbreak{}delta+\allowbreak{}chi\textless=\allowbreak{}19 (the
former notes used the weaker24); the separate normal-endpoint m=\allowbreak{}4
exclusions in the notes apply unchanged. It remains to treat m=\allowbreak{}1,\allowbreak{}2,\allowbreak{}3.

\hypertarget{profile-and-endpoint-constants}{%
\subsection{Profile and endpoint
constants}\label{profile-and-endpoint-constants}}

In the following table sigma is excess over the displayed minimal
profile and B the sum of baseline maxima. Q61 is obtained by deleting
four largest core magnitudes on each side.

\begin{longtable}[]{@{}
  >{\raggedright\arraybackslash}p{(\columnwidth - 10\tabcolsep) * \real{0.1667}}
  >{\raggedright\arraybackslash}p{(\columnwidth - 10\tabcolsep) * \real{0.1667}}
  >{\raggedright\arraybackslash}p{(\columnwidth - 10\tabcolsep) * \real{0.1667}}
  >{\raggedright\arraybackslash}p{(\columnwidth - 10\tabcolsep) * \real{0.1667}}
  >{\raggedright\arraybackslash}p{(\columnwidth - 10\tabcolsep) * \real{0.1667}}
  >{\raggedright\arraybackslash}p{(\columnwidth - 10\tabcolsep) * \real{0.1667}}@{}}
\toprule\noalign{}
\begin{minipage}[b]{\linewidth}\raggedright
m
\end{minipage} & \begin{minipage}[b]{\linewidth}\raggedright
I+\allowbreak{}tau
\end{minipage} & \begin{minipage}[b]{\linewidth}\raggedright
B
\end{minipage} & \begin{minipage}[b]{\linewidth}\raggedright
4f
\end{minipage} & \begin{minipage}[b]{\linewidth}\raggedright
L lower bound
\end{minipage} & \begin{minipage}[b]{\linewidth}\raggedright
mass cap
\end{minipage} \\
\midrule\noalign{}
\endhead
\bottomrule\noalign{}
\endlastfoot
1 & 68+\allowbreak{}I+\allowbreak{}sigma & 70 & 166-sigma+\allowbreak{}R & 11-w+\allowbreak{}I+\allowbreak{}sigma+\allowbreak{}delta &
sigma+\allowbreak{}delta+\allowbreak{}I\textless=\allowbreak{}148+\allowbreak{}w \\
2 & 101+\allowbreak{}chi+\allowbreak{}sigma & 71 & 161-sigma+\allowbreak{}R & 44-w+\allowbreak{}chi+\allowbreak{}sigma+\allowbreak{}delta &
sigma+\allowbreak{}delta+\allowbreak{}chi\textless=\allowbreak{}115+\allowbreak{}w \\
3 & 132+\allowbreak{}I+\allowbreak{}sigma & 72 & 156-sigma+\allowbreak{}R & 75-w+\allowbreak{}I+\allowbreak{}sigma+\allowbreak{}delta &
sigma+\allowbreak{}delta+\allowbreak{}I\textless=\allowbreak{}84+\allowbreak{}w \\
\end{longtable}

For m=\allowbreak{}2,\allowbreak{} k is the physical unique-1 owner\textquotesingle s size and
chi=\allowbreak{}(k-34)(k-35)\textgreater=\allowbreak{}0. It need not equal the paired-resource
owner. If that physical owner is side s then k=\allowbreak{}r\_\allowbreak{}s; otherwise
k=\allowbreak{}68-r\_\allowbreak{}s.

For a native retained ordinary/sole-pair center Z on target v the
nominal endpoint rows are

\begin{lstlisting}
Z>=73-r_v-l_Z-A_Z,
Z<=S_v+r_o-72+h_Z+D_Z.
\end{lstlisting}

Normal endpoints use zero error. Write c for the selected lower or upper
clipping count. An outside lower endpoint has c=\allowbreak{}73-r\_\allowbreak{}v+\allowbreak{}j,\allowbreak{} with the next
positive core rank j+\allowbreak{}1 if c\textless73; reflection gives the same upper
argument. If its baseline displacement is d,\allowbreak{} then

\begin{lstlisting}
(73-c)d<=sigma_v.
\end{lstlisting}

The baseline endpoint error is at most b+\allowbreak{}d,\allowbreak{} where b=\allowbreak{}1 at m1 and b=\allowbreak{}2 at
m2,\allowbreak{}m3. For a light c\textless73,\allowbreak{} the licensed sigma+\allowbreak{}epsilon bound gives
d\textless=\allowbreak{}c. For c\textless=\allowbreak{}71 this follows from
c(74-c)\textgreater=\allowbreak{}213 on the relevant interval26..71,\allowbreak{} and
sigma+\allowbreak{}epsilon\textless=\allowbreak{}146; smaller actual side lower bounds only
strengthen it. For c=\allowbreak{}72:

\begin{itemize}
\tightlist
\item
  A1 uses credit11,\allowbreak{} so sigma+\allowbreak{}epsilon\textless=\allowbreak{}148+\allowbreak{}6-11=\allowbreak{}143 at m1; the
  other profiles have still smaller caps.
\item
  A0 m1 with z2 has sigma+\allowbreak{}epsilon\textless=\allowbreak{}142. With z4 it
  is\textless=\allowbreak{}144. With z6,\allowbreak{} the two additional low holes imply
  tau\textgreater=\allowbreak{}176 and delta\textless=\allowbreak{}46,\allowbreak{} so c72 cannot occur. A0
  m2,\allowbreak{}m3 again have smaller caps.
\end{itemize}

Thus every light error is at most b+\allowbreak{}c\textless=\allowbreak{}b+\allowbreak{}Gamma\_\allowbreak{}Z.

For a heavy c\textgreater=\allowbreak{}73,\allowbreak{} at most one relevant selected resource is
heavy. In A1 epsilon\textless=\allowbreak{}143 at m1,\allowbreak{} and \textless=\allowbreak{}110,\allowbreak{}79 at m2,\allowbreak{}m3,\allowbreak{}
so two73 clips are impossible. In A0,\allowbreak{} m2/m3 also have
epsilon\textless146. At m1,\allowbreak{} z2 gives epsilon\textless=\allowbreak{}142; z4 implies
tau\textgreater=\allowbreak{}120,\allowbreak{} delta\textless=\allowbreak{}100,\allowbreak{} epsilon\textless=\allowbreak{}92; z6 gives
epsilon\textless=\allowbreak{}38. Thus again two73 clips are impossible.

The heavy next source magnitude is peeled. Its endpoint error is at most
b+\allowbreak{}sigma\_\allowbreak{}v+\allowbreak{}e\_\allowbreak{}v,\allowbreak{} where e\_\allowbreak{}v is the fourth peeled offset on the physical
source side providing that next magnitude. This deliberately loose bound
is used in A0. For A1,\allowbreak{} activity bounds the error by
M\_\allowbreak{}v-r\_\allowbreak{}v+\allowbreak{}72-c\textless=\allowbreak{}M\_\allowbreak{}v-r\_\allowbreak{}v-1. The profile maxima are r\_\allowbreak{}v+\allowbreak{}1 at m1
and at most r\_\allowbreak{}v+\allowbreak{}2 at m2,\allowbreak{}m3,\allowbreak{} so the sharper heavy error is at most
e\_\allowbreak{}v+\allowbreak{}sigma\_\allowbreak{}v+\allowbreak{}b-1. Together with the other two light errors,\allowbreak{} this is the
stated aggregate3b-1 below. The native deficiency of a heavy resource is
at least73,\allowbreak{} and delta\textgreater=\allowbreak{}e+\allowbreak{}c.

\hypertarget{a0-physical-target-capacity-replaces-the-old-total-mass-transport}{%
\subsection{A0: physical-target capacity replaces the old total-mass
transport}\label{a0-physical-target-capacity-replaces-the-old-total-mass-transport}}

Let s be the actual sole pair1 owner,\allowbreak{} and e=\allowbreak{}8,\allowbreak{} epsilon=\allowbreak{}delta-e. The
original unit holes are beta-pair,\allowbreak{}O\_\allowbreak{}s,\allowbreak{}O\_\allowbreak{}t. Target t receives exactly
its leading and seven retained ordinary resources H,\allowbreak{}U,\allowbreak{}V,\allowbreak{}C5,\allowbreak{}C6,\allowbreak{}C7,\allowbreak{}alpha,\allowbreak{}
with capacity K+\allowbreak{}539. Their grouping gives

\begin{lstlisting}
|F_t|>=K+539-epsilon,
|F_s|<=K+700-e,
S_s<=K+700-e+w/2.
\end{lstlisting}

The literal identity is

\begin{lstlisting}
H_s=P+2alpha_s+alpha_t+C5_t+4M+4X+2G.
\end{lstlisting}

Use endpoint resources P,\allowbreak{}alpha\_\allowbreak{}s,\allowbreak{}alpha\_\allowbreak{}t,\allowbreak{}C5\_\allowbreak{}t,\allowbreak{}H\_\allowbreak{}s with
clipping/error weights1,\allowbreak{}2,\allowbreak{}1,\allowbreak{}1,\allowbreak{}1. Their nominal lower/upper constants
give

\begin{lstlisting}
S_s>=233+4M+4X+2G-E_clip-B_err.
\end{lstlisting}

On each target v the retained tree is leading\_\allowbreak{}v-\/-V\_\allowbreak{}v-\/-U\_\allowbreak{}v,\allowbreak{} with
the extra edge leading\_\allowbreak{}v-\/-C7\_\allowbreak{}v. Its three gaps are X\_\allowbreak{}v,\allowbreak{}g\_\allowbreak{}v,\allowbreak{}X\_\allowbreak{}o.
For the positive edge restrict the full leading to lambda\_\allowbreak{}v+\allowbreak{}M\_\allowbreak{}v+\allowbreak{}Q\_\allowbreak{}v;
for the negative edge use lambda\_\allowbreak{}v-M\_\allowbreak{}o+\allowbreak{}Q\_\allowbreak{}v,\allowbreak{} where Q\_\allowbreak{}s=\allowbreak{}Q61 and
Q\_\allowbreak{}t=\allowbreak{}-Q61. Ordinary vertices are restricted to their matching Q\_\allowbreak{}v
translates. Each edge contributes at least f minus its gap. Both
restricted copies remain subsets of one full leading capacity K. Across
the two targets these six edges have total gap2X+\allowbreak{}G,\allowbreak{} so their grouped
deficiency satisfies D\_\allowbreak{}F\textgreater=\allowbreak{}6f-(2X+\allowbreak{}G). They are disjoint from
the five endpoint resources,\allowbreak{} and E\_\allowbreak{}clip+\allowbreak{}2D\_\allowbreak{}F\textless=\allowbreak{}2epsilon.
Combining the last three displays with K\textless=\allowbreak{}M+\allowbreak{}R+\allowbreak{}21 yields the
master inequality

\begin{lstlisting}
3M+12f+3e<=R+488+w/2+2delta+B_err.
\end{lstlisting}

All target holes are already in these native/grouped deficiencies.

Let l,\allowbreak{}D,\allowbreak{}Mcap,\allowbreak{}chi0,\allowbreak{}b be the following constants:

\begin{longtable}[]{@{}llllll@{}}
\toprule\noalign{}
m & l & D=\allowbreak{}3(4f constant)-488 & Mcap & chi0 & b \\
\midrule\noalign{}
\endhead
\bottomrule\noalign{}
\endlastfoot
1 & 11 & 10 & 148 & I & 1 \\
2 & 44 & -5 & 115 & chi & 2 \\
3 & 75 & -20 & 84 & I & 2 \\
\end{longtable}

The master uses M\textgreater=\allowbreak{}R+\allowbreak{}L,\allowbreak{} L\textgreater=\allowbreak{}l-w+\allowbreak{}chi0+\allowbreak{}sigma+\allowbreak{}delta,\allowbreak{}
and R\textgreater=\allowbreak{}70,\allowbreak{}71,\allowbreak{}72 respectively.

If every endpoint is light,\allowbreak{} B\_\allowbreak{}err\textless=\allowbreak{}6b+\allowbreak{}2epsilon. Direct
substitution forces

\begin{lstlisting}
5R+4chi0+sigma+5e <= Mcap+6b-3l-D+(9/2)w.
\end{lstlisting}

The maximum right sides are138,\allowbreak{}27,\allowbreak{}-82 respectively; the left side is at
least350+\allowbreak{}40,\allowbreak{}355+\allowbreak{}40,\allowbreak{}360+\allowbreak{}40. Thus no all-light case survives.

If there is one heavy endpoint of weight q in\{1,\allowbreak{}2\},\allowbreak{} then

\begin{lstlisting}
B_err<=6b+2epsilon-qc+q sigma+q e_v.
\end{lstlisting}

Cancel q e\_\allowbreak{}v against qL and use the mass cap. The exact terminal row is

\begin{lstlisting}
5R+(3+q)chi0+(q-1)delta+5e+qc+(3-q)l+D-6b
    <=2q Mcap+(7/2+q)w.
\end{lstlisting}

For m1 a heavy forces w\textless=\allowbreak{}4. For q1,\allowbreak{} left\textgreater=\allowbreak{}489 and
right\textless=\allowbreak{}314. For q2,\allowbreak{} delta\textgreater=\allowbreak{}e+\allowbreak{}c\textgreater=\allowbreak{}81 gives
left\textgreater=\allowbreak{}632 and right\textless=\allowbreak{}614. For m2 the q1/q2 left
floors are539,\allowbreak{}649 and right ceilings257,\allowbreak{}493. For m3 they are591,\allowbreak{}670 and
ceilings195,\allowbreak{}369. All are strict contradictions. Hence all A0 first gaps
are excluded.

\hypertarget{a1-the-same-ten-edge-forest-applies-to-all-three-profiles}{%
\subsection{A1: the same ten-edge forest applies to all three
profiles}\label{a1-the-same-ten-edge-forest-applies-to-all-three-profiles}}

Let s be physical beta owner,\allowbreak{} t actual pair1 owner,\allowbreak{} r=\allowbreak{}r\_\allowbreak{}s,\allowbreak{} e=\allowbreak{}11 and
epsilon=\allowbreak{}delta-e. Reserve b\_\allowbreak{}*=\allowbreak{}Gamma\_\allowbreak{}beta\_\allowbreak{}s+\allowbreak{}Gamma\_\allowbreak{}P\_\allowbreak{}t and
h\_\allowbreak{}*=\allowbreak{}Gamma\_\allowbreak{}H\_\allowbreak{}t.

The ten-edge forest consists of these retained arms:

\begin{lstlisting}
on s: leading_s--V_s--U_s;
      leading_s--C7_s--C6_s--C5_s;
on t: leading_t--V_t--U_t--H_t;
      leading_t--C7_t--C6_t.
\end{lstlisting}

Its literal edge inventory is:

\begin{longtable}[]{@{}lll@{}}
\toprule\noalign{}
target & edge & gap \\
\midrule\noalign{}
\endhead
\bottomrule\noalign{}
\endlastfoot
s & leading\_\allowbreak{}s,\allowbreak{}V\_\allowbreak{}s & X\_\allowbreak{}s \\
s & V\_\allowbreak{}s,\allowbreak{}U\_\allowbreak{}s & g\_\allowbreak{}s \\
s & leading\_\allowbreak{}s,\allowbreak{}C7\_\allowbreak{}s & X\_\allowbreak{}t \\
s & C7\_\allowbreak{}s,\allowbreak{}C6\_\allowbreak{}s & g\_\allowbreak{}t \\
s & C6\_\allowbreak{}s,\allowbreak{}C5\_\allowbreak{}s & h\_\allowbreak{}t \\
t & leading\_\allowbreak{}t,\allowbreak{}V\_\allowbreak{}t & X\_\allowbreak{}t \\
t & V\_\allowbreak{}t,\allowbreak{}U\_\allowbreak{}t & g\_\allowbreak{}t \\
t & U\_\allowbreak{}t,\allowbreak{}H\_\allowbreak{}t & h\_\allowbreak{}t \\
t & leading\_\allowbreak{}t,\allowbreak{}C7\_\allowbreak{}t & X\_\allowbreak{}s \\
t & C7\_\allowbreak{}t,\allowbreak{}C6\_\allowbreak{}t & g\_\allowbreak{}s \\
\end{longtable}

Their gaps total2X\_\allowbreak{}s+\allowbreak{}2X\_\allowbreak{}t+\allowbreak{}2g\_\allowbreak{}s+\allowbreak{}2g\_\allowbreak{}t+\allowbreak{}2h\_\allowbreak{}t. Restricted positive and
negative hub copies are lambda\_\allowbreak{}v+\allowbreak{}M\_\allowbreak{}v+\allowbreak{}Q\_\allowbreak{}v and lambda\_\allowbreak{}v-M\_\allowbreak{}o+\allowbreak{}Q\_\allowbreak{}v;
each stays inside the same full leading vertex. Ordinary vertices use
Q\_\allowbreak{}v translates. Omitting both paired resources leaves deficit at most
epsilon-b\_\allowbreak{}*,\allowbreak{} and hence

\begin{lstlisting}
X_s+X_t+g_s+g_t+h_t>=5f-(epsilon-b_*)/2.
\end{lstlisting}

The four-edge subforest uses leading\_\allowbreak{}s-\/-V\_\allowbreak{}s-\/-U\_\allowbreak{}s and
leading\_\allowbreak{}t-\/-C7\_\allowbreak{}t-\/-C6\_\allowbreak{}t. It omits both pairs and H\_\allowbreak{}t,\allowbreak{} so

\begin{lstlisting}
X_s+g_s>=2f-(epsilon-b_*-h_*)/2.
\end{lstlisting}

The two stars leading\_\allowbreak{}s-\/-V\_\allowbreak{}s and leading\_\allowbreak{}t-\/-C7\_\allowbreak{}t use the full
core E69,\allowbreak{} whose span is R,\allowbreak{} and give

\begin{lstlisting}
X_s>=137-R-(epsilon-b_*-h_*)/2.
\end{lstlisting}

Add the ten-edge row,\allowbreak{} twice the four-edge row and three times the
two-star row. The resulting positive scalar combination gives

\begin{lstlisting}
W=6X_s+X_t+3g_s+g_t+h_t
  >=9f+411-3R-3epsilon+3b_*+(5/2)h_*.
\end{lstlisting}

The center identity H\_\allowbreak{}t=\allowbreak{}2d\_\allowbreak{}s+\allowbreak{}beta\_\allowbreak{}s+\allowbreak{}P\_\allowbreak{}t+\allowbreak{}2s4+\allowbreak{}t4+\allowbreak{}g\_\allowbreak{}s+\allowbreak{}g\_\allowbreak{}t+\allowbreak{}h\_\allowbreak{}t and
its three endpoint rows,\allowbreak{} followed by the leading-s bound twice,\allowbreak{} give

\begin{lstlisting}
S_s+3g_s+4X_s+2s4+t4+g_t+h_t
  <=1089+w+r+delta-2e+E_err.
\end{lstlisting}

All-light E\_\allowbreak{}err\textless=\allowbreak{}3b+\allowbreak{}b\_\allowbreak{}*+\allowbreak{}h\_\allowbreak{}*. One-heavy
E\_\allowbreak{}err\textless=\allowbreak{}e\_\allowbreak{}u+\allowbreak{}sigma\_\allowbreak{}u+\allowbreak{}(3b-1)+\allowbreak{}G\_\allowbreak{}light,\allowbreak{} and

\begin{lstlisting}
3b_*+(5/2)h_*-G_light>=365/2,
L-e_u>=4, sigma_s-sigma_u>=-sigma.
\end{lstlisting}

The physical source lower rows are exactly Tr+\allowbreak{}6r+\allowbreak{}19 at m1,\allowbreak{} Tr+\allowbreak{}6r+\allowbreak{}18 for
m2,\allowbreak{}k=\allowbreak{}r,\allowbreak{} Tr+\allowbreak{}7r+\allowbreak{}24 for m2,\allowbreak{}k=\allowbreak{}68-r,\allowbreak{} and Tr+\allowbreak{}7r+\allowbreak{}22 at m3,\allowbreak{} plus
sigma\_\allowbreak{}s+\allowbreak{}R+\allowbreak{}L+\allowbreak{}6X\_\allowbreak{}s+\allowbreak{}X\_\allowbreak{}t. Hence the resulting scalar right constants
are571/2,\allowbreak{}1191/4,\allowbreak{}1167/4,\allowbreak{}305 respectively,\allowbreak{} in the shared form

\begin{lstlisting}
Tr+c_r r+sigma_s-9sigma/4+R/4+L+5e+3b_*+5h_*/2
  <=constant+w+4delta+E_err,
\end{lstlisting}

with c\_\allowbreak{}r=\allowbreak{}5,\allowbreak{}5,\allowbreak{}6,\allowbreak{}6 respectively.

\hypertarget{m1}{%
\subsubsection{m1}\label{m1}}

All-light substitution gives

\begin{lstlisting}
Tr+5r+4I+7sigma/4+sigma_s+R/4+5e<=1443/2+5w.
\end{lstlisting}

For integer r,\allowbreak{} Tr+\allowbreak{}5r+\allowbreak{}4(r-34)\^{}2\textgreater=\allowbreak{}679. R\textgreater=\allowbreak{}70
and5e55 give a left floor1503/2=\allowbreak{}751.5. For w\textless=\allowbreak{}4 the right
is\textless=\allowbreak{}741.5. If w6,\allowbreak{} three low holes include1 and force
tau\textgreater=\allowbreak{}176,\allowbreak{} so sigma\textgreater=\allowbreak{}108 and the left is strictly
larger than751.5,\allowbreak{} whereas the right is at most751.5.

One-heavy gives

\begin{lstlisting}
Tr+5r+4I+3sigma/4+R/4+5e<=693+5w<=723,
\end{lstlisting}

contradicting the same left floor751.5.

\hypertarget{m2-both-physical-unique-1-owners}{%
\subsubsection{m2: both physical unique-1
owners}\label{m2-both-physical-unique-1-owners}}

If k=\allowbreak{}r,\allowbreak{} all-light gives

\begin{lstlisting}
Tr+4r+I+3(r-34)(r-35)+7sigma/4+sigma_s+R/4+5e
  <=2283/4+5w<=2403/4.
\end{lstlisting}

Its quadratic part is at least661. If k=\allowbreak{}68-r the corresponding row is

\begin{lstlisting}
Tr+7r+I+3(r-34)(r-33)+7sigma/4+sigma_s+R/4+5e
  <=2531/4+5w<=2651/4,
\end{lstlisting}

whose quadratic part is at least723. Both contradict even without the
other positive terms.

One-heavy gives the two rows

\begin{lstlisting}
Tr+5r+4(r-34)(r-35)+3sigma/4+R/4+5e<=2305/4+5w<=2425/4,
Tr+6r+4(r-34)(r-33)+3sigma/4+R/4+5e<=2281/4+5w<=2401/4.
\end{lstlisting}

Their quadratic parts have minima695 and689 respectively.

\hypertarget{m3}{%
\subsubsection{m3}\label{m3}}

All-light gives

\begin{lstlisting}
Tr+6r+4I+7sigma/4+sigma_s+R/4+5e<=488+5w<=518.
\end{lstlisting}

One-heavy gives

\begin{lstlisting}
Tr+6r+4I+3sigma/4+R/4+5e<=919/2+5w<=979/2.
\end{lstlisting}

Here Tr+\allowbreak{}6r+\allowbreak{}4(r-34)\^{}2\textgreater=\allowbreak{}709. Thus neither alternative
survives. Together with the normal-endpoint m4 exclusions,\allowbreak{} A1 is
impossible.

Every use of repeated scalar inequalities is a positive linear
combination of separately valid budget consequences. Each forest vertex
keeps only its original full resource capacity.

\clearpage
\hypertarget{fresh-p4-proof-and-cards-for-n85}{%
\section{Fresh P4 proof and cards for
n85}\label{fresh-p4-proof-and-cards-for-n85}}

These calculations use the fresh structure,\allowbreak{} the proved sharp raw caps,\allowbreak{}
and the universal source theorem T=\allowbreak{}m0+\allowbreak{}A\textless=\allowbreak{}165. With j4,\allowbreak{}w2 one has
m0=\allowbreak{}26,\allowbreak{} A=\allowbreak{}I+\allowbreak{}tau+\allowbreak{}delta,\allowbreak{} I=\allowbreak{}(r\_\allowbreak{}s-34)\^{}2,\allowbreak{} and

\begin{lstlisting}
A<=139, S=2K+1164-delta,
2K=26+A+4R+U+2C+J,
L=M-R>=18+A,
K<=M+R+21.
\end{lstlisting}

This note includes the source-seed argument,\allowbreak{} all six resulting core-hole
profiles,\allowbreak{} the exhaustive classification,\allowbreak{} and the terminal cards using
their actual fresh constants. The shared positive-case prefix notation
is in \texttt{fresh\_\allowbreak{}\allowbreak{}positive\_\allowbreak{}\allowbreak{}notes.md}.

\hypertarget{source-localization-and-profile-reduction}{%
\subsection{Source localization and profile
reduction}\label{source-localization-and-profile-reduction}}

Originally z\textless=\allowbreak{}14,\allowbreak{}w\textless=\allowbreak{}4z. If the minimum side size
is\textless=\allowbreak{}20 then I\textgreater=\allowbreak{}196,\allowbreak{}rho\textgreater=\allowbreak{}49 and the
structural T\_\allowbreak{}floor48 imply delta\textless=\allowbreak{}244-I-rho\textless0. At
sizes21..23,\allowbreak{} rho\textgreater=\allowbreak{}58,\allowbreak{}I\textgreater=\allowbreak{}121 and the structural
floor144-2z\textgreater=\allowbreak{}116 give delta\textless=\allowbreak{}118-I\textless0. Thus
both sizes\textgreater=\allowbreak{}24.

The structural H bound is
H\textless=\allowbreak{}90+\allowbreak{}4z-I-delta\textless=\allowbreak{}146\textless T\_\allowbreak{}24. The fresh
inner/outer layer rows give central-strip join margins\textgreater=\allowbreak{}154
and456. Therefore w\textless=\allowbreak{}z,\allowbreak{} both even. In particular the loose old
source bound A\textless=\allowbreak{}144+\allowbreak{}w\textless=\allowbreak{}158 gives core mass\textless=\allowbreak{}1348
and core maxima\textless=\allowbreak{}201; these intentionally conservative bounds
suffice for the seven-seed argument in the notes. That argument gives
w\textless=\allowbreak{}z\textless=\allowbreak{}4 and the exact six profiles.

Here is that seed argument explicitly. The greedy first-gap rule leaves
total m-1 below m,\allowbreak{} with all later entries\textgreater m and no entry
equal to m. First gaps\textgreater=\allowbreak{}4 have minimal I+\allowbreak{}tau\textgreater=\allowbreak{}199
(m4),\allowbreak{}228 (m5),\allowbreak{} or tau\textgreater=\allowbreak{}248 (m\textgreater=\allowbreak{}6),\allowbreak{} contradicting
A\textless=\allowbreak{}158. At m1,\allowbreak{} two values u,\allowbreak{}v from\{2,\allowbreak{}...,\allowbreak{}6\} absent on both
signs force tau\textgreater=\allowbreak{}214-2u-2v\textgreater=\allowbreak{}192,\allowbreak{} so at most one is
absent. At m2,\allowbreak{} an absent q in\{3,\allowbreak{}4,\allowbreak{}5\} forces tau\textgreater=\allowbreak{}208-r-2q,\allowbreak{}
where r owns the unique1; the quadratic I+\allowbreak{}tau has
minimum174-2q\textgreater=\allowbreak{}164. At m3,\allowbreak{} an absent q in\{4,\allowbreak{}5,\allowbreak{}6\} forces
tau\textgreater=\allowbreak{}206-2q\textgreater=\allowbreak{}194. Thus the following actual
labelled seeds are available:

For the generic m\textgreater=\allowbreak{}6 bound,\allowbreak{} let k\_\allowbreak{}v count the low
occurrences on sign v. The core maximum\textless=\allowbreak{}201 gives
m\textless=\allowbreak{}201 and T\_\allowbreak{}(k\_\allowbreak{}v)\textless=\allowbreak{}m-1,\allowbreak{} hence k\_\allowbreak{}v\textless=\allowbreak{}19. At m6
the total low count is\textless=\allowbreak{}3,\allowbreak{} giving tau\textgreater=\allowbreak{}4(68-3)=\allowbreak{}260.
At m\textgreater=\allowbreak{}7,\allowbreak{} k\_\allowbreak{}v\textless=\allowbreak{}3 gives the side floor4(r\_\allowbreak{}v-3).
Otherwise m-k\_\allowbreak{}v\textgreater=\allowbreak{}k\_\allowbreak{}v+\allowbreak{}3; as r\_\allowbreak{}v\textgreater=\allowbreak{}24 and
k\_\allowbreak{}v\textless=\allowbreak{}19,\allowbreak{} the difference (k\_\allowbreak{}v+\allowbreak{}3)(r\_\allowbreak{}v-k\_\allowbreak{}v)-4(r\_\allowbreak{}v-3) is at
least k\_\allowbreak{}v+\allowbreak{}7\textgreater0. Summing again gives tau\textgreater=\allowbreak{}248. The
m4/m5 exact profile summations are displayed in the shared prefix
lemmas.

\begin{longtable}[]{@{}
  >{\raggedright\arraybackslash}p{(\columnwidth - 8\tabcolsep) * \real{0.1765}}
  >{\raggedright\arraybackslash}p{(\columnwidth - 8\tabcolsep) * \real{0.1765}}
  >{\raggedleft\arraybackslash}p{(\columnwidth - 8\tabcolsep) * \real{0.2353}}
  >{\raggedright\arraybackslash}p{(\columnwidth - 8\tabcolsep) * \real{0.1765}}
  >{\raggedleft\arraybackslash}p{(\columnwidth - 8\tabcolsep) * \real{0.2353}}@{}}
\toprule\noalign{}
\begin{minipage}[b]{\linewidth}\raggedright
condition
\end{minipage} & \begin{minipage}[b]{\linewidth}\raggedright
seed
\end{minipage} & \begin{minipage}[b]{\linewidth}\raggedleft
seed mass
\end{minipage} & \begin{minipage}[b]{\linewidth}\raggedright
lower holes
\end{minipage} & \begin{minipage}[b]{\linewidth}\raggedleft
minimum mass of any9 core occurrences
\end{minipage} \\
\midrule\noalign{}
\endhead
\bottomrule\noalign{}
\endlastfoot
m1,\allowbreak{}2 absent & 3,\allowbreak{}4,\allowbreak{}5,\allowbreak{}6 & 18 & 1,\allowbreak{}2 & 43 \\
m1,\allowbreak{}3 absent & 2,\allowbreak{}4,\allowbreak{}5,\allowbreak{}6 & 17 & 1,\allowbreak{}3 & 41 \\
m1,\allowbreak{}4 absent & 2,\allowbreak{}3,\allowbreak{}5,\allowbreak{}6 & 16 & 1,\allowbreak{}4 & 39 \\
m1,\allowbreak{}2,\allowbreak{}3,\allowbreak{}4,\allowbreak{}5 present & 2,\allowbreak{}3,\allowbreak{}4,\allowbreak{}5 & 14 & 1 & 34 \\
m1,\allowbreak{}2,\allowbreak{}3,\allowbreak{}4,\allowbreak{}6 present & 2,\allowbreak{}3,\allowbreak{}4,\allowbreak{}6 & 15 & 1 & 34 \\
m2 & 1,\allowbreak{}3,\allowbreak{}4,\allowbreak{}5 & 13 & 2 & 37 \\
m3 & 1,\allowbreak{}1,\allowbreak{}4,\allowbreak{}5,\allowbreak{}6 & 17 & 3 & 39 \\
\end{longtable}

Direct bitset addition,\allowbreak{} or a short enumeration of each
seed\textquotesingle s sums,\allowbreak{} gives the full support except the stated
lower holes and their reflections. The last column follows by taking at
most two occurrences of each smallest permitted magnitude. Adjoin the
remaining labelled core entries in increasing order. A failed
central-interval join at next entry x means x\textgreater=\allowbreak{}T-2h,\allowbreak{} where h
is the largest current lower seed-hole position. Any N remaining
magnitudes at least x,\allowbreak{} multiplicity\textless=\allowbreak{}2,\allowbreak{} have
total\textgreater=\allowbreak{}Nx+\allowbreak{}floor((N-1)\^{}2/4). For n\textless=\allowbreak{}8 selected
entries,\allowbreak{}N\textgreater=\allowbreak{}60; the seven resulting core-mass lower bounds
are1728,\allowbreak{}1547,\allowbreak{}1366,\allowbreak{}1604,\allowbreak{}1665,\allowbreak{}1423,\allowbreak{}1547.
For9\textless=\allowbreak{}n\textless=\allowbreak{}32,\allowbreak{}N\textgreater=\allowbreak{}36 and the last column
applies; the bounds are1753,\allowbreak{}1607,\allowbreak{}1461,\allowbreak{}1492,\allowbreak{}1492,\allowbreak{}1531,\allowbreak{}1533. All
exceed1348. For n\textgreater=\allowbreak{}33,\allowbreak{} T\textgreater=\allowbreak{}289 and
x\textless=\allowbreak{}201\textless T-2h-1. Thus all joins succeed and the absolute
core has at most four holes.

In the first three rows the profile baselines tau are136,\allowbreak{}134,\allowbreak{}132
respectively. For the missing4 row,\allowbreak{} missing2 on one sign would add r+\allowbreak{}1
above the baseline and force
I+\allowbreak{}tau\textgreater=\allowbreak{}167\textgreater A\textless=\allowbreak{}148. Both opposite2 labels
are therefore present,\allowbreak{} filling4; only1 remains as a lower hole. In the
other rows the sole displayed first gap persists. Every peeled magnitude
exceeds4,\allowbreak{} so all lower core holes persist to the source. The exact
(tau,\allowbreak{}z=\allowbreak{}w) pairs are therefore(68+\allowbreak{}sigma,\allowbreak{}2) for ordinarym1,\allowbreak{}
(135-r+\allowbreak{}sigma,\allowbreak{}2) for m2,\allowbreak{} (132+\allowbreak{}sigma,\allowbreak{}2) for m3,\allowbreak{} (136+\allowbreak{}sigma,\allowbreak{}4) for
missing1,\allowbreak{}2,\allowbreak{} (134+\allowbreak{}sigma,\allowbreak{}4) for missing1,\allowbreak{}3,\allowbreak{} and (132+\allowbreak{}sigma,\allowbreak{}2) for
missing1,\allowbreak{}4.

Every target-hole center is active: its hole offset is at most4 and
subtracting any of the at least25 opposite core magnitudes produces
distinct positive hull images well below the top two holes. Nonpositive
centers in{[}1-M\_\allowbreak{}v,\allowbreak{}0{]} use source M\_\allowbreak{}v. Hence every integral center
in{[}1-M\_\allowbreak{}v,\allowbreak{}S\_\allowbreak{}v{]} is active before any endpoint license.

Both leading resources are active by
K-delta\textgreater=\allowbreak{}166-w+\allowbreak{}I+\allowbreak{}tau\textgreater0. If both O were active,\allowbreak{}
their sharp caps give d\_\allowbreak{}s+\allowbreak{}d\_\allowbreak{}t\textless=\allowbreak{}2M+\allowbreak{}44+\allowbreak{}w. Set
Y=\allowbreak{}2M-2K+\allowbreak{}44+\allowbreak{}w+\allowbreak{}delta. The leading and high-tail counts imply
2Y\textgreater488,\allowbreak{} while K-M\textgreater=\allowbreak{}R-132-w+\allowbreak{}A and
R\textgreater=\allowbreak{}70,\allowbreak{}tau\textgreater=\allowbreak{}68,\allowbreak{}w\textless=\allowbreak{}4 give Y\textless=\allowbreak{}44,\allowbreak{}
contradiction.

If one O\_\allowbreak{}v were active,\allowbreak{} its sharp cap and leading row imply

\begin{lstlisting}
2X_v+g_v<=154+3w/2-R-I-tau.
\end{lstlisting}

An active lower-coordinate witness makes that coordinate integral. The
two full-core overlaps with the two leading resources (allowing an empty
mate with zero capacity) then give q\textgreater=\allowbreak{}137-R-delta/2. Since
q\textless=\allowbreak{}2X\_\allowbreak{}v+\allowbreak{}g\_\allowbreak{}v,\allowbreak{} this would force
I+\allowbreak{}tau\textless=\allowbreak{}17+\allowbreak{}3w/2+\allowbreak{}delta/2\textless=\allowbreak{}63,\allowbreak{} using the conservative
delta\textless=\allowbreak{}80-I. This contradicts tau\textgreater=\allowbreak{}68. Thus all four
lower-coordinate witnesses would be empty,\allowbreak{} forcing both O active,\allowbreak{}
already excluded. Consequently both O are empty.

Their two known holes leave only two others. A nonintegral lower
coordinate forces two witness holes and additionally H\_\allowbreak{}v or alpha\_\allowbreak{}v
nonintegral,\allowbreak{} impossible. After lower integrality,\allowbreak{} a nonintegral v2
forces H\_\allowbreak{}v,\allowbreak{}alpha\_\allowbreak{}v and also C5\_\allowbreak{}o in A or beta\_\allowbreak{}v in B missing. Thus
all outer coordinates are integral. O\_\allowbreak{}v\textgreater-M\_\allowbreak{}v by positive
leading activity,\allowbreak{} so each empty O is HIGH. Its four fixed labels
C11,\allowbreak{}C12,\allowbreak{}C13,\allowbreak{}C14 are HIGH and outside every baseline group. These eight
labels license credit8: delta\textgreater=\allowbreak{}8 and epsilon=\allowbreak{}delta-8.

The universal source cap A\textless=\allowbreak{}137+\allowbreak{}w now eliminates all extra
high-mass profiles immediately. First-gap3 and missing1,\allowbreak{}4 have
tau\textgreater=\allowbreak{}132,\allowbreak{}w2,\allowbreak{} hence delta\textless=\allowbreak{}7. Missing1,\allowbreak{}3 has
tau\textgreater=\allowbreak{}134,\allowbreak{}w4,\allowbreak{} delta\textless=\allowbreak{}7. Missing1,\allowbreak{}2 has
tau\textgreater=\allowbreak{}136,\allowbreak{}w4,\allowbreak{} delta\textless=\allowbreak{}5. Each contradicts credit8.

Only the ordinary first-gap1 profile and first-gap2 profile remain,\allowbreak{} with
exact target support F\_\allowbreak{}v=\allowbreak{}{[}1,\allowbreak{}S\_\allowbreak{}v{]} minus\{S\_\allowbreak{}v-m\}. At first gap1,\allowbreak{}
tau=\allowbreak{}68+\allowbreak{}sigma and I+\allowbreak{}sigma+\allowbreak{}delta\textless=\allowbreak{}71,\allowbreak{} so both physical sizes lie
in{[}26,\allowbreak{}42{]}. At first gap2,\allowbreak{} unique1 owner h satisfies h in{[}29,\allowbreak{}40{]}
and true profile excess sigma obeys I+\allowbreak{}sigma+\allowbreak{}delta\textless=\allowbreak{}h+\allowbreak{}4. In both
profiles delta\textless=\allowbreak{}71 and epsilon\textless=\allowbreak{}63.

\hypertarget{empty-center-rule-endpoints-and-core8-classification}{%
\subsection{Empty-center rule,\allowbreak{} endpoints and CORE8
classification}\label{empty-center-rule-endpoints-and-core8-classification}}

Source U obeys U\textless=\allowbreak{}350-2T,\allowbreak{} so U\textless=\allowbreak{}162. The first opposite
peeled magnitude is\textless=\allowbreak{}R\_\allowbreak{}o+\allowbreak{}U\textless=\allowbreak{}304. On both profiles
S\_\allowbreak{}v\textgreater=\allowbreak{}469 and R\textless=\allowbreak{}142. The points M\_\allowbreak{}v,\allowbreak{}0,\allowbreak{}-R\_\allowbreak{}o give
activity throughout{[}1-M\_\allowbreak{}v,\allowbreak{}S\_\allowbreak{}v+\allowbreak{}R\_\allowbreak{}o{]},\allowbreak{} except possibly at
S\_\allowbreak{}v+\allowbreak{}R\_\allowbreak{}o-m. There the first opposite peel produces a positive hit below
S\_\allowbreak{}v-m (at least163). Thus an empty integral center is
LOW\textless=\allowbreak{}-M\_\allowbreak{}v or HIGH\textgreater=\allowbreak{}S\_\allowbreak{}v+\allowbreak{}R\_\allowbreak{}o+\allowbreak{}1.

An empty own center cannot be LOW by leading activity. A retained
singleton above any HIGH center loses at least73 source points: all
positive points,\allowbreak{} zero,\allowbreak{} and all opposite core points. Such a loss exceeds
epsilon\textless=\allowbreak{}63. An empty cross center cannot be HIGH,\allowbreak{} since it
would put all77 points of lambda\_\allowbreak{}v+\allowbreak{}D\_\allowbreak{}v above the target. Cross
activity is upward closed
C4=\allowbreak{}\textgreater C5=\allowbreak{}\textgreater C6=\allowbreak{}\textgreater C7.

Before CORE8 retention,\allowbreak{} the target-independent separated-arm forests
give

\begin{lstlisting}
2X_v+X_o+g_v>=252-tau-delta,
2X_v+X_o+G>=305-tau-delta.
\end{lstlisting}

The arm choices involve at most one empty vertex at each point,\allowbreak{} and
positive/negative restricted supports are disjoint because
M\textgreater R. If beta\_\allowbreak{}v,\allowbreak{}H\_\allowbreak{}v are active,\allowbreak{} the sharp beta/H raw cap is
d\_\allowbreak{}v-M\textless=\allowbreak{}21+\allowbreak{}w/2\textless=\allowbreak{}22. Activity of C7\_\allowbreak{}v or C6\_\allowbreak{}v would
instead give d\_\allowbreak{}v-M\textgreater=\allowbreak{}253-tau-delta\textgreater=\allowbreak{}114
or306-tau-delta\textgreater=\allowbreak{}167. Thus neither triple beta\_\allowbreak{}v,\allowbreak{}H\_\allowbreak{}v,\allowbreak{}C7\_\allowbreak{}v
nor beta\_\allowbreak{}v,\allowbreak{}H\_\allowbreak{}v,\allowbreak{}C6\_\allowbreak{}v can be wholly active.

These rules reproduce the exhaustive whole-unit classification with the
fresh native-loss threshold73:

\begin{itemize}
\tightlist
\item
  A1 retained beta: two contributing betas would require one further
  hole in each\{H\_\allowbreak{}v,\allowbreak{}C6\_\allowbreak{}v\}; then both C5 singletons are active and the
  four-center identity beta\_\allowbreak{}s+\allowbreak{}beta\_\allowbreak{}t+\allowbreak{}C5\_\allowbreak{}s+\allowbreak{}C5\_\allowbreak{}t=\allowbreak{}-2W contradicts their
  activity lower bounds. Thus beta is sole on s. Active H\_\allowbreak{}s would force
  both C6\_\allowbreak{}s,\allowbreak{}C5\_\allowbreak{}s empty,\allowbreak{} leaving C7\_\allowbreak{}s active and contradicting the
  forbidden triple. Hence the known holes are O\_\allowbreak{}s,\allowbreak{}H\_\allowbreak{}s,\allowbreak{}O\_\allowbreak{}t,\allowbreak{} with one
  further Z. CORE8 is retained by the cross chain and the
  higher-retained73-loss rule. Z lies
  in\{H\_\allowbreak{}t,\allowbreak{}alpha\_\allowbreak{}s,\allowbreak{}alpha\_\allowbreak{}t,\allowbreak{}C5\_\allowbreak{}s,\allowbreak{}C5\_\allowbreak{}t,\allowbreak{}pair1\}. Pair1 is sole if
  retained because its two centers sum\textless-3M. If it is empty,\allowbreak{} use
  T5 below. If retained with opposite owner,\allowbreak{} use the opposite-owner
  card. If same owner s,\allowbreak{} the triple C4\_\allowbreak{}s,\allowbreak{}beta\_\allowbreak{}s,\allowbreak{}C5\_\allowbreak{}t is impossible,\allowbreak{}
  so Z=\allowbreak{}C5\_\allowbreak{}t and alpha\_\allowbreak{}s is retained; use same-owner card.
\item
  B: U,\allowbreak{}V must be retained by the higher-retained rule,\allowbreak{} and the forbidden
  triples plus cross chain force C6,\allowbreak{}C7 retained. The two further holes
  are one from each\{H\_\allowbreak{}s,\allowbreak{}beta\_\allowbreak{}s\},\allowbreak{}\{H\_\allowbreak{}t,\allowbreak{}beta\_\allowbreak{}t\}. Both alphas and
  both pairs remain. Empty beta is LOW. They cannot both be empty
  because beta\_\allowbreak{}s+\allowbreak{}beta\_\allowbreak{}t\textgreater0. Pair1 and C5 pair have sole
  matching owners by the literal center identities. Any active beta on
  the opposite owner is excluded by the opposite-owner card,\allowbreak{} and the
  remaining active beta on the common owner with retained alpha is
  excluded by the same-owner card.
\item
  A0 beta-pair empty: CORE8 is retained. Its one further hole Z belongs
  to\{H\_\allowbreak{}s,\allowbreak{}H\_\allowbreak{}t,\allowbreak{}alpha\_\allowbreak{}s,\allowbreak{}alpha\_\allowbreak{}t,\allowbreak{}C5\_\allowbreak{}s,\allowbreak{}C5\_\allowbreak{}t,\allowbreak{}pair1\}. If pair1 is
  retained it is sole; normalize owner s. Then C5\_\allowbreak{}s is retained and Z
  lies in\{H\_\allowbreak{}s,\allowbreak{}H\_\allowbreak{}t,\allowbreak{}alpha\_\allowbreak{}s,\allowbreak{}alpha\_\allowbreak{}t,\allowbreak{}C5\_\allowbreak{}t\}. At least one of
  H\_\allowbreak{}v,\allowbreak{}alpha\_\allowbreak{}v is retained above each empty beta,\allowbreak{} so both beta
  components are LOW. V6\_\allowbreak{}s,\allowbreak{}V6\_\allowbreak{}t are LOW and meet one baseline group.
  With the eight HIGH O labels the fixed credit is9.
\end{itemize}

First-gap1 charged endpoints have uniform allowance1:
I+\allowbreak{}sigma+\allowbreak{}epsilon\textless=\allowbreak{}63\textless73 makes the next core rank exact.
First-gap2 has delta\textless=\allowbreak{}38 and epsilon\textless=\allowbreak{}30,\allowbreak{} while an
outside-hull center clips at least33 source points,\allowbreak{} so charged endpoints
are normal. With a=\allowbreak{}1 for m1,\allowbreak{}a=\allowbreak{}0 for m2,\allowbreak{} both are

\begin{lstlisting}
Z>=73-r_v-l_Z-a,
Z<=S_v+r_o-72+h_Z+a.
\end{lstlisting}

\hypertarget{common-sourcerank-bounds-for-the-terminal-cards}{%
\subsection{Common source/rank bounds for the terminal
cards}\label{common-sourcerank-bounds-for-the-terminal-cards}}

Here sigma=\allowbreak{}tau-68 is the generic excess; at m1 each
sigma\_\allowbreak{}v\textgreater=\allowbreak{}0,\allowbreak{} at m2 sigma\_\allowbreak{}v\textgreater=\allowbreak{}-1 and
sigma\textgreater=\allowbreak{}27. Put eta=\allowbreak{}0 at m1 and eta=\allowbreak{}1 at m2. The penultimate
core estimate gives

\begin{lstlisting}
K<=C+312+sigma+eta,
2R+I+delta<=224+eta<=225.
\end{lstlisting}

Q61 overlap floor is f\_\allowbreak{}eta=\allowbreak{}(166-eta-sigma+\allowbreak{}R)/4. Full-core69 overlap
floor is137-R. Ordinary gaps and stars can use any actual common trimmed
core,\allowbreak{} with no change to capacities. At m1,\allowbreak{} writing xi=\allowbreak{}R-70,\allowbreak{} three
available overlap floors are

\begin{lstlisting}
F0=67-xi, F1=65-sigma+xi, F2=63-sigma/2+xi/2.
\end{lstlisting}

At m2 using the true profile excess sigma\textquotesingle{} and
xi\textquotesingle=\allowbreak{}R-71 they are

\begin{lstlisting}
F0=66-xi', F1=64-sigma'+xi', F2=62-sigma'/2+xi'/2.
\end{lstlisting}

The same full resources certify all these floors; each scalar inequality
may use the maximum F.

\hypertarget{opposite-owner-card}{%
\subsection{Opposite-owner card}\label{opposite-owner-card}}

Suppose beta\_\allowbreak{}s and P\_\allowbreak{}t=\allowbreak{}C4\_\allowbreak{}t are licensed capacity77 native resources
and CORE8 is retained. The identity V\_\allowbreak{}t=\allowbreak{}2d\_\allowbreak{}s+\allowbreak{}beta\_\allowbreak{}s+\allowbreak{}P\_\allowbreak{}t+\allowbreak{}2s4+\allowbreak{}t4+\allowbreak{}g\_\allowbreak{}s
and three endpoint rows followed by the leading-s inequality twice give

\begin{lstlisting}
S_s+3g_s+4X_s+2s4+t4<=1014+3a+r_s+delta-2e.
\end{lstlisting}

Two ordinary gap-g\_\allowbreak{}s pairs give3g\_\allowbreak{}s\textgreater=\allowbreak{}3f\_\allowbreak{}eta-3epsilon/2.
The stars give6X\_\allowbreak{}s+\allowbreak{}X\_\allowbreak{}t\textgreater=\allowbreak{}959-7R-3epsilon. Combining these
with the source mass,\allowbreak{} L\textgreater=\allowbreak{}86+\allowbreak{}I+\allowbreak{}sigma+\allowbreak{}delta and the radial
bound yields

\begin{lstlisting}
T_r+5r+29I/8+sigma/4+sigma_s+13e/2
  <=827/2+27eta/8+15delta/8+3a.
\end{lstlisting}

Use delta\textless=\allowbreak{}71-I,\allowbreak{} eta\textless=\allowbreak{}1,\allowbreak{}a\textless=\allowbreak{}1 and
sigma\_\allowbreak{}s\textgreater=\allowbreak{}-1. Then

\begin{lstlisting}
T_r+5r+11I/2+sigma/4+sigma_s+13e/2<=553.
\end{lstlisting}

The global integral quadratic minimum T\_\allowbreak{}r+\allowbreak{}5r+\allowbreak{}11(r-34)\^{}2/2 is1401/2.
Credit e\textgreater=\allowbreak{}8 and sigma\_\allowbreak{}s\textgreater=\allowbreak{}-1 give
left\textgreater=\allowbreak{}1503/2\textgreater553. Thus opposite owners are
impossible.

\hypertarget{same-owner-card-with-active-alpha}{%
\subsection{Same-owner card with active
alpha}\label{same-owner-card-with-active-alpha}}

Suppose alpha\_\allowbreak{}s,\allowbreak{}P\_\allowbreak{}s,\allowbreak{}beta\_\allowbreak{}s are licensed distinct native resources.
The literal identity and outer order give

\begin{lstlisting}
alpha_s-P_s-beta_s> d_t+2s4+t4+G.
\end{lstlisting}

Three endpoints and the distinct leading-t loss give

\begin{lstlisting}
S_t+g_t+2X_t+2s4+t4+G<=K+1082+3a-r_t-e.
\end{lstlisting}

Three disjoint ordinary pairs give
g\_\allowbreak{}s+\allowbreak{}2g\_\allowbreak{}t\textgreater=\allowbreak{}3f\_\allowbreak{}eta-epsilon,\allowbreak{} while stars
give2X\_\allowbreak{}s+\allowbreak{}3X\_\allowbreak{}t\textgreater=\allowbreak{}685-5R-3epsilon/2. Source and radial
substitution gives

\begin{lstlisting}
T_r+5r+25I/8+sigma/4+sigma_t+5delta/8+7e/2
  <=1349/2+23eta/8+3a.
\end{lstlisting}

At m1 the right is677.5,\allowbreak{} sigma\_\allowbreak{}t\textgreater=\allowbreak{}0; at m2 it is677.375,\allowbreak{}
sigma\_\allowbreak{}t\textgreater=\allowbreak{}-1. The quadratic part has exact integer
minimum5265/8=\allowbreak{}658.125. Credit8 makes the respective left
floors5489/8=\allowbreak{}686.125 and5481/8=\allowbreak{}685.125. Both exceed the right side.

\hypertarget{t5-sole-retained-beta-empty-pair1}{%
\subsection{T5: sole retained beta,\allowbreak{} empty
pair1}\label{t5-sole-retained-beta-empty-pair1}}

The holes are O\_\allowbreak{}s,\allowbreak{}H\_\allowbreak{}s,\allowbreak{}O\_\allowbreak{}t,\allowbreak{}pair1,\allowbreak{} beta sole owner s. Credit13 follows
by an exact label ledger. Leading
K-delta\textgreater=\allowbreak{}164+\allowbreak{}I+\allowbreak{}tau\textgreater R gives d\_\allowbreak{}v\textgreater g\_\allowbreak{}v,\allowbreak{}
so F3,\allowbreak{}F4,\allowbreak{}F5 are LOW. Empty beta\_\allowbreak{}t is LOW by retained alpha\_\allowbreak{}t. Both
empty C4 are LOW,\allowbreak{} giving C55,\allowbreak{}C56,\allowbreak{}C57,\allowbreak{}C58 LOW on each sign. Genuine
beta\_\allowbreak{}s,\allowbreak{}C5\_\allowbreak{}t imply d\_\allowbreak{}t\textgreater d\_\allowbreak{}s; C5\_\allowbreak{}s gives t2\textless d\_\allowbreak{}s,\allowbreak{}
so V3\_\allowbreak{}t,\allowbreak{}V4\_\allowbreak{}t LOW. There are20 LOW labels meeting18 baseline groups,\allowbreak{}
and11 HIGH labels from O\_\allowbreak{}s,\allowbreak{}H\_\allowbreak{}s,\allowbreak{}O\_\allowbreak{}t outside those groups. Thus
e=\allowbreak{}11+\allowbreak{}20-18=\allowbreak{}13,\allowbreak{} with no extra hole allowance.

Upper endpoints alpha\_\allowbreak{}s,\allowbreak{}H\_\allowbreak{}t,\allowbreak{} lower endpoints alpha\_\allowbreak{}t and twice
beta\_\allowbreak{}s,\allowbreak{} plus C5\_\allowbreak{}t,\allowbreak{} together with the leading-s inequality twice give

\begin{lstlisting}
2M+6X_s+2X_t+3g_s+g_t<=932+6a+delta-2e.
\end{lstlisting}

The weighted native losses are\textless=\allowbreak{}2epsilon and the six endpoint
allowances are exactly6a. Stars and pairs using overlap F
give6X\_\allowbreak{}s+\allowbreak{}2X\_\allowbreak{}t\textgreater=\allowbreak{}8F-3epsilon
and3g\_\allowbreak{}s+\allowbreak{}g\_\allowbreak{}t\textgreater=\allowbreak{}4F-3epsilon/2. Since
M\textgreater=\allowbreak{}R+\allowbreak{}86+\allowbreak{}I+\allowbreak{}sigma+\allowbreak{}delta and I+\allowbreak{}sigma+\allowbreak{}delta\textless=\allowbreak{}71,\allowbreak{}

\begin{lstlisting}
2R+11I/2+11sigma/2+12F<=924+6a.
\end{lstlisting}

At m1,\allowbreak{} R=\allowbreak{}70+\allowbreak{}xi and13F0/2+\allowbreak{}11F1/2
implies12F+\allowbreak{}2xi+\allowbreak{}11sigma/2\textgreater=\allowbreak{}793+\allowbreak{}xi,\allowbreak{} so the
left\textgreater=\allowbreak{}933\textgreater930. At m2,\allowbreak{} use generic sigma,\allowbreak{}xi=\allowbreak{}R-70,\allowbreak{}
F0=\allowbreak{}67-xi,\allowbreak{}F1=\allowbreak{}64-sigma+\allowbreak{}xi; the same combination gives a left
floor927.5+\allowbreak{}xi\textgreater=\allowbreak{}928.5\textgreater924. This excludes T5.

\hypertarget{a0-case1-pair1-empty}{%
\subsection{A0 case1: pair1 empty}\label{a0-case1-pair1-empty}}

All H,\allowbreak{}alpha,\allowbreak{}C5 and CORE8 are retained,\allowbreak{} with credit9. Pair-sum joint-loss
counting and the two H identities,\allowbreak{} using uniformly weakened
endpoints72-r\_\allowbreak{}v and S\_\allowbreak{}v+\allowbreak{}r\_\allowbreak{}o-71,\allowbreak{} give

\begin{lstlisting}
S>=4M+4X+2G+302+2F-3epsilon.
\end{lstlisting}

The source upper bound S\textless=\allowbreak{}2M+\allowbreak{}2R+\allowbreak{}1206-delta and the stars/pairs
with overlap F then give

\begin{lstlisting}
14F+2I+2tau<=814+3delta.
\end{lstlisting}

At m1,\allowbreak{} mass I+\allowbreak{}sigma+\allowbreak{}delta\textless=\allowbreak{}71 forces14F+\allowbreak{}5I+\allowbreak{}5sigma\textless=\allowbreak{}891.
But5F0+\allowbreak{}F1+\allowbreak{}8F2 gives14F+\allowbreak{}5sigma\textgreater=\allowbreak{}904,\allowbreak{} contradiction.

At m2 with true excess sigma\textquotesingle{} and unique1 owner h,\allowbreak{}
tau=\allowbreak{}135-h+\allowbreak{}sigma\textquotesingle{} and
I+\allowbreak{}sigma\textquotesingle+\allowbreak{}delta\textless=\allowbreak{}h+\allowbreak{}4.
Thus14F+\allowbreak{}5sigma\textquotesingle+\allowbreak{}5I\textless=\allowbreak{}556+\allowbreak{}5h. Yet5F0+\allowbreak{}F1+\allowbreak{}8F2
gives14F+\allowbreak{}5sigma\textquotesingle\textgreater=\allowbreak{}890,\allowbreak{} forcing
I-h\textless=\allowbreak{}-334/5,\allowbreak{} while min{[}(h-34)\^{}2-h{]}=\allowbreak{}-34. Contradiction.

\hypertarget{a0-case2-retained-pair1-opposite-alpha-retained}{%
\subsection{A0 case2: retained pair1,\allowbreak{} opposite alpha
retained}\label{a0-case2-retained-pair1-opposite-alpha-retained}}

The conservative transport being used is

\begin{lstlisting}
T_k+4k+128-tau_s<=7R+3delta,
\end{lstlisting}

where k=\allowbreak{}r\_\allowbreak{}t is the side opposite the pair1 owner. Its ordinary
first-gap1 proof uses the sharper actual fringe bound
B\_\allowbreak{}s+\allowbreak{}B\_\allowbreak{}t\textless=\allowbreak{}min(20,\allowbreak{}183-T); the root\textquotesingle s
common-source certificate gives positive margins3,\allowbreak{}2,\allowbreak{}1 for
T\textless=\allowbreak{}163,\allowbreak{}T164,\allowbreak{}T165. This is a separate licensed source lemma,\allowbreak{} not
an assumed129 transport.

At m2,\allowbreak{} use R\textless=\allowbreak{}71+\allowbreak{}sigma\textquotesingle{} and the unique1-owner
mass cap. If h=\allowbreak{}k,\allowbreak{} tau\_\allowbreak{}s=\allowbreak{}136-2k+\allowbreak{}sigma\_\allowbreak{}s and
I+\allowbreak{}sigma\textquotesingle+\allowbreak{}delta\textless=\allowbreak{}k+\allowbreak{}4. The conservative row forces

\begin{lstlisting}
T_k-2k+8I+5delta<=537.
\end{lstlisting}

Its quadratic part is at least496,\allowbreak{} and delta\textgreater=\allowbreak{}9 makes the
left\textgreater=\allowbreak{}541. If h=\allowbreak{}68-k,\allowbreak{} tau\_\allowbreak{}s=\allowbreak{}67-k+\allowbreak{}sigma\_\allowbreak{}s and the cap is
I+\allowbreak{}sigma\textquotesingle+\allowbreak{}delta\textless=\allowbreak{}72-k. The resulting row is

\begin{lstlisting}
T_k+13k+8I+5delta<=1012.
\end{lstlisting}

Its quadratic minimum is971,\allowbreak{} and delta\textgreater=\allowbreak{}9 gives
left\textgreater=\allowbreak{}1016. Thus both physical unique1 owners are excluded
with margin4.

\hypertarget{a0-case3-retained-pair1-opposite-alpha-empty}{%
\subsection{A0 case3: retained pair1,\allowbreak{} opposite alpha
empty}\label{a0-case3-retained-pair1-opposite-alpha-empty}}

Every other H,\allowbreak{}alpha\_\allowbreak{}s,\allowbreak{}C5 and CORE8 resource is retained. The same seven
weighted endpoint uses as the old identity,\allowbreak{} now with the fresh weakened
constants,\allowbreak{} give

\begin{lstlisting}
S>=5s4+4t4+2g_s+3g_t+2h_t+298-r_s-E,
E<=2epsilon.
\end{lstlisting}

Hence

\begin{lstlisting}
2M+M_s+5X_s+4X_t+2g_s+3g_t+2h_t
  <=2R+908+r_s-delta+2epsilon.
\end{lstlisting}

Use f=\allowbreak{}(233+\allowbreak{}R-tau)/4,\allowbreak{} which is valid for both profiles. The disjoint
ordinary groups give2g\_\allowbreak{}s+\allowbreak{}3g\_\allowbreak{}t+\allowbreak{}2h\_\allowbreak{}t\textgreater=\allowbreak{}5f+\allowbreak{}122-2epsilon. Stars
give5X\_\allowbreak{}s+\allowbreak{}4X\_\allowbreak{}t\textgreater=\allowbreak{}1233-9R-5epsilon/2.
Also2M+\allowbreak{}M\_\allowbreak{}s\textgreater=\allowbreak{}2R+\allowbreak{}41+\allowbreak{}r\_\allowbreak{}s+\allowbreak{}2I+\allowbreak{}2tau+\allowbreak{}2delta. Thus

\begin{lstlisting}
488+2I+2tau+5f+3delta<=9R+13epsilon/2,
3351/4+2I+3tau/4<=31R/4+7delta/2.
\end{lstlisting}

Use2R+\allowbreak{}I+\allowbreak{}delta\textless=\allowbreak{}225. This forces

\begin{lstlisting}
-273+47I+6tau+3delta<=0.
\end{lstlisting}

Since tau\textgreater=\allowbreak{}68,\allowbreak{} the left is at least135. Contradiction.

All original P4 configurations are exhausted by these owner/empty-pair
cases and the A0 three-case split.

\clearpage
% Standalone local theorem and proof; requires amsmath,amssymb,amsthm.
% This is a conditional local exclusion, not a replacement for the structural
% theorem or for the actual activity/retention classification.
\subsection{An analytic exclusion of the first-gap-one A0 cross case}

Consider the normalized cardinality-$85$ source with $d=77$ and
$E=P_s\cup\{0\}\cup(-P_t)$, where $r=|P_s|$, $k=|P_t|$ and
$r+k=68$. Assume that the first missing positive subset sum of the
absolute core is $1$, that $z=w=2$, and that the A0 classification gives
the retained centers $P=C_{4s}$, $\alpha_t$, $C_{6t}$ and $U_s$ and all
eight CORE8 resources. The selected $P$ is an actual contributor of its
sole-owned paired unit. The fixed-label credit is $9$; write
$\varepsilon=\delta-9\ge0$.

The inputs used here are the proved source bound
\[
 T=26+I+\tau+\delta\le165,\qquad I=(k-34)^2,
\]
the nominal first-gap-one endpoint bounds
\[
 Z\ge r_o+4-\ell_Z,\qquad Z\le S_v-r_v-3+h_Z,
\]
and the actual native/grouped loss budgets. All subset copies used below
belong to the specified retained resources; no capacity is assigned to
an empty center.

Write the four increasing source-tail offsets as $a_v<b_v<c_v<e_v$.
Let $U=a_s+b_s+a_t+b_t$, $J=e_s-c_s+e_t-c_t$, and let $B_v$ count
the distinct sums of two offsets on sign $v$ that exceed $e_v$.
Set $B=B_s+B_t$ and let $p$ be the number of odd integers
$a_v+b_v$. The two source envelopes and the quartet fringe bound give
\[
 J+U\le332-T,\quad J\ge T-18+p,\quad
 2B\le U+16-p.
\]
Thus $2T+U+p\le350$ and
\begin{equation}\label{eq:85-cross-fringe}
 B\le\min(20,183-T).
\end{equation}
The full-source central interval gives
$K\le M+R+1+B$, where $M=M_s+M_t$ and $R=R_s+R_t$.
Consequently the exact original mass identity yields
\begin{equation}\label{eq:85-cross-source}
 S_s+S_t=2K+1164-\delta
 \le2M+2R+1166+2B-\delta.
\end{equation}

\paragraph{The endpoint transport, including every constant.}
Put $a=s_4=M_s+X_s$, $b=t_4=M_t+X_t$ and $G=g_s+g_t$.
The literal center identity is
\[
 U_s=P+\alpha_t+C_{6t}+2a+4b+2G.
\]
Three lower endpoints and the upper endpoint of $U_s$ imply
\[
 S_s\ge219-2k+2a+4b+2G-E_{\rm end},\qquad
 E_{\rm end}=\ell_P+\ell_{\alpha_t}+\ell_{C_{6t}}+h_{U_s}
 \le\varepsilon.
\]
Indeed their constant is
$(k+4)+2(r+4)+(r+3)=219-2k$.
The opposite source satisfies
$S_t\ge T_k+\tau_t+4R_t+10$, while
$R_t\ge k+1$ and $M_t\ge k+5$.
Combining with \eqref{eq:85-cross-source} therefore gives
\begin{equation}\label{eq:85-cross-before-overlap}
 T_k+4k+\tau_t+243+2X_s+4X_t+2G
 \le2R+1166+2B-\delta+E_{\rm end}.
\end{equation}
Here $243=219+10+14$, and $14$ comes from
$2M_t+4R_t-2k\ge4k+14$.

\paragraph{Admissible asymmetric core copies.}
Put
\[
 E_v=\tau_v-R_v+1\ge0,\quad
 \zeta=\tau_t-k\ge0,\quad D=E_s+E_t=\tau-R+2.
\]
Every core magnitude is at least its rank plus $1$.
After deleting the maximum, its remaining rank deviations sum to $E_v$.
Deleting the $j_v$ largest core entries therefore leaves maximum at most
$r_v-j_v+1+\lfloor E_v/j_v\rfloor$.
For the signed trimmed core $Q$, this proves
\[
 |Q\cap(Q+q)|\ge
 67-j_s-j_t-\lfloor E_s/j_s\rfloor-\lfloor E_t/j_t\rfloor-q.
\]
Choose $j_s=3$ and choose $j_t\in\{1,2,3\}$ minimizing
$j_t+\lfloor E_t/j_t\rfloor$. Define
\[
 \phi_s=3+\lfloor E_s/3\rfloor,\quad
 \phi_t=\min_{j=1,2,3}\bigl(j+\lfloor E_t/j\rfloor\bigr),\quad
 f=67-\phi_s-\phi_t,
 \quad F=\max(137-R,f).
\]
The full core gives the bound $137-R-q$; the trimmed core gives $f-q$.

For a leading-to-own3 edge, the leading resource contains
$\lambda_v+M_v+Q$ and the ordinary resource contains
$\lambda_v+v_4+Q$, with gap $X_v=v_4-M_v$.
The reflected leading-to-$C_7$ edge uses the reflected copies at the
same gap. Thus these are actual subsets of the retained resources.
The rooted forest count, separately for an individual two-edge star
and for the four-edge star, gives
\[
 2X_t\ge2F-\varepsilon,\qquad
 2X_s+2X_t\ge4F-\varepsilon,
\]
and hence $2X_s+4X_t\ge6F-2\varepsilon$.
The four disjoint CORE8 ordinary pairs similarly give
$2G\ge4f-\varepsilon$.
Substituting into \eqref{eq:85-cross-before-overlap} and using
$E_{\rm end}\le\varepsilon$ and $\varepsilon=\delta-9$ yields
\begin{equation}\label{eq:85-cross-raw}
 T_k+4k+\tau_t+6F+4f
 \le2R+887+2B+3\delta.
\end{equation}
In particular the conservative old-$B=20$ residual is
\[
 L=T_k+4k+\tau_t-927-2R-3\delta+6F+4f.
\]
The actual residual in \eqref{eq:85-cross-raw} is
$L+2(20-B)$; this is exactly the fringe correction, not an additional
credit against $\delta$.

\paragraph{A short integer certificate.}
Put $q=T_k+5k$. Since $T=26+I+\tau+\delta$,
\[
 L+3T=q+3I-855+\zeta+3(E_s+E_t)+R+6F+4f.
\]
The two inequalities $F\ge137-R$ and $F\ge f$ give
$6F\ge137-R+5f$, so
\[
 L+3T\ge q+3I-115
 +(3E_s-9\phi_s)+(\zeta+3E_t-9\phi_t).
\]
First, $3E_s-9\phi_s\ge-27$.
If $E_t=0$, the second parenthesis is at least $-9$.
If $E_t>0$, then $R_t-k-1\ge1$: equality would force the entire
$k$-element core to be $\{2,\ldots,k+1\}$, for which $E_t=0$.
Hence $\zeta=E_t+(R_t-k-1)\ge E_t+1$.
For $E_t\ge6$, the choice $j_t=3$ gives
\[
 \zeta+3E_t-9\phi_t\ge E_t-26\ge-20.
\]
For $E_t=1,2,3,4,5$, the respective bounds are
$-13,-18,-14,-19,-15$, by using $j_t=1,2,2,2,3$.
Thus the second parenthesis is always at least $-20$.
Finally, for every integer $k$, with $x=k-28$,
\[
 q+3I=T_k+5k+3(k-34)^2
 =654+\frac{7x^2-5x}{2}\ge654.
\]
We have proved $L\ge492-3T$. By \eqref{eq:85-cross-fringe},
\[
 L+2(20-B)\ge492-3T+2\max(0,T-163)>0
 \qquad(T\le165).
\]
For $T\le163$ the margin is at least $3$; at $T=164,165$ it is at
least $2,1$, respectively. This contradicts \eqref{eq:85-cross-raw}
and excludes the conditional A0 cross case with first gap $1$.

\clearpage
\hypertarget{complete-z5-a0-alpha-hole-finish-at-size-85}{%
\section{Complete Z5 A0 alpha-hole finish at size
85}\label{complete-z5-a0-alpha-hole-finish-at-size-85}}

This is the A0 family with both paired units empty and exact whole-unit
holes \texttt{\{pair1,\allowbreak{}pair2,\allowbreak{}O\_\allowbreak{}\allowbreak{}s,\allowbreak{}O\_\allowbreak{}\allowbreak{}t,\allowbreak{}alpha\_\allowbreak{}\allowbreak{}s\}}.
There is no assigned paired owner. All thirteen remaining ordinary
singleton units are retained. The fresh structural and retention
arguments have already given integral outer coordinates,\allowbreak{} leading
activity,\allowbreak{} \texttt{w=\allowbreak{}0},\allowbreak{} and the source notation

\[
 r_s+\allowbreak{}r_t=\allowbreak{}68,\allowbreak{}\quad R_v=\allowbreak{}r_v+\allowbreak{}\xi_v,\allowbreak{}\quad
 0\le\xi_v\le\tau_v,\allowbreak{}\quad \tau_s+\allowbreak{}\tau_t=\allowbreak{}\tau,\allowbreak{}
 \quad A=\allowbreak{}I+\allowbreak{}\tau+\allowbreak{}\delta\le60.
\]

Write \texttt{a=\allowbreak{}s\_\allowbreak{}\allowbreak{}4=\allowbreak{}M\_\allowbreak{}\allowbreak{}s+\allowbreak{}X\_\allowbreak{}\allowbreak{}s},\allowbreak{}
\texttt{b=\allowbreak{}t\_\allowbreak{}\allowbreak{}4=\allowbreak{}M\_\allowbreak{}\allowbreak{}t+\allowbreak{}X\_\allowbreak{}\allowbreak{}t},\allowbreak{}
\texttt{G=\allowbreak{}g\_\allowbreak{}\allowbreak{}s+\allowbreak{}g\_\allowbreak{}\allowbreak{}t},\allowbreak{}
\texttt{B=\allowbreak{}alpha\_\allowbreak{}\allowbreak{}t},\allowbreak{} \texttt{C=\allowbreak{}C5\_\allowbreak{}\allowbreak{}s},\allowbreak{}
\texttt{D=\allowbreak{}C5\_\allowbreak{}\allowbreak{}t}. The alpha hole is LOW,\allowbreak{} so
\texttt{h\_\allowbreak{}\allowbreak{}s<=\allowbreak{}X\_\allowbreak{}\allowbreak{}s}. The generic retained-core-hit
argument and complete-target rank bounds apply before normality; they do
not import the old Z5 normality bound. The refined source lemma gives
\texttt{e\_\allowbreak{}\allowbreak{}v>=\allowbreak{}A+\allowbreak{}16} on both signs.

\hypertarget{fixed-credit-and-the-order-in-which-normality-is-established}{%
\subsection{Fixed credit and the order in which normality is
established}\label{fixed-credit-and-the-order-in-which-normality-is-established}}

The following stages are the fresh tail/order deductions proved in the
zero-case companion. They are listed here with their actual labels and
incidence so the final loss budget can be checked.

Initially all retained singleton deficiencies are at most60,\allowbreak{} giving at
least17 target hits,\allowbreak{} more than the eight peeled source entries; each has
a core hit. All selected native clips are below73,\allowbreak{} so the generic rank
errors are bounded by the appropriate \texttt{xi\_\allowbreak{}\allowbreak{}v}.
Both paired components and the alpha hole are LOW.

The initial ledger gives \texttt{e=\allowbreak{}14}. Its source
support obtained by trimming four core ranks on both signs has
\texttt{q=\allowbreak{}61} and overlap at least
\texttt{f=\allowbreak{}61-(tau-xi\_\allowbreak{}\allowbreak{}s-xi\_\allowbreak{}\allowbreak{}t)/\allowbreak{}4}. The stars,\allowbreak{} g pairs,\allowbreak{}
leading,\allowbreak{} and the retained \texttt{\{H\_\allowbreak{}\allowbreak{}s,\allowbreak{}U\_\allowbreak{}\allowbreak{}s\}} pair
then give

\[
 X_v,\allowbreak{}g_v\ge38,\allowbreak{}\quad h_s\ge15,\allowbreak{}
 \quad d_v\ge107+\allowbreak{}g_v+\allowbreak{}2X_v.
\]

These inequalities and the exact32-type upper-tail inventory give
\texttt{S\_\allowbreak{}\allowbreak{}s<=\allowbreak{}H\_\allowbreak{}\allowbreak{}s+\allowbreak{}M\_\allowbreak{}\allowbreak{}s+\allowbreak{}2} and
\texttt{S\_\allowbreak{}\allowbreak{}t<=\allowbreak{}H\_\allowbreak{}\allowbreak{}t+\allowbreak{}M\_\allowbreak{}\allowbreak{}t+\allowbreak{}16},\allowbreak{} yielding the additional13
HIGH labels below. The two places requiring cross-centre lower endpoints
use the proved core-hit lower bound
\texttt{C5\_\allowbreak{}\allowbreak{}v>=\allowbreak{}1-R\_\allowbreak{}\allowbreak{}v}; they have strict margins larger
than26 and16. The consecutive-tail argument needs
\texttt{h\_\allowbreak{}\allowbreak{}s>2},\allowbreak{} supplied by15. Thus this step gives
\texttt{e=\allowbreak{}27} without assuming normality.

Now a BAD ordinary centre has deficiency at least
\texttt{73-max(r\_\allowbreak{}\allowbreak{}s,\allowbreak{}r\_\allowbreak{}\allowbreak{}t)}. Put
\texttt{k=\allowbreak{}|r\_\allowbreak{}\allowbreak{}s-34|},\allowbreak{} so
\texttt{I=\allowbreak{}k\^2}. Its excess over the available budget
\texttt{epsilon=\allowbreak{}delta-27} is at least

\[
 73-(34+\allowbreak{}k)-(60-k^2-\tau-27)
   =\allowbreak{}k^2-k+\allowbreak{}6+\allowbreak{}\tau>0.
\]

Hence every retained centre is normal. The native upper endpoints and
unused-loss rebates prove \texttt{B,\allowbreak{}D<a}; the
fixed-only prefix count proves \texttt{C<b}. All three
retained centres \texttt{B,\allowbreak{}C,\allowbreak{}D} are now positive.

The complete HIGH set is

\[
 \mathcal H=\allowbreak{}\{C_{11},\allowbreak{}C_{12},\allowbreak{}C_{13},\allowbreak{}C_{14}\}^{s,\allowbreak{}t}
 \cup\{C_{22},\allowbreak{}C_{23},\allowbreak{}C_{24},\allowbreak{}C_{16},\allowbreak{}C_{17},\allowbreak{}C_{18}\}^{s,\allowbreak{}t}
 \cup\{C_{33}^{s}\},\allowbreak{}\qquad|\mathcal H|=\allowbreak{}21.
\]

It meets no baseline group. The LOW stages and exact incidence are:

\begin{longtable}[]{@{}
  >{\raggedright\arraybackslash}p{(\columnwidth - 10\tabcolsep) * \real{0.1364}}
  >{\raggedright\arraybackslash}p{(\columnwidth - 10\tabcolsep) * \real{0.1364}}
  >{\raggedleft\arraybackslash}p{(\columnwidth - 10\tabcolsep) * \real{0.1818}}
  >{\raggedleft\arraybackslash}p{(\columnwidth - 10\tabcolsep) * \real{0.1818}}
  >{\raggedleft\arraybackslash}p{(\columnwidth - 10\tabcolsep) * \real{0.1818}}
  >{\raggedleft\arraybackslash}p{(\columnwidth - 10\tabcolsep) * \real{0.1818}}@{}}
\toprule\noalign{}
\begin{minipage}[b]{\linewidth}\raggedright
Stage
\end{minipage} & \begin{minipage}[b]{\linewidth}\raggedright
Added LOW labels
\end{minipage} & \begin{minipage}[b]{\linewidth}\raggedleft
\texttt{H}
\end{minipage} & \begin{minipage}[b]{\linewidth}\raggedleft
\texttt{L}
\end{minipage} & \begin{minipage}[b]{\linewidth}\raggedleft
Groups touched
\end{minipage} & \begin{minipage}[b]{\linewidth}\raggedleft
Credit
\end{minipage} \\
\midrule\noalign{}
\endhead
\bottomrule\noalign{}
\endlastfoot
Initial & \texttt{C55,\allowbreak{}C56,\allowbreak{}C57,\allowbreak{}C58,\allowbreak{}V5,\allowbreak{}V6,\allowbreak{}V7,\allowbreak{}V8} on both
signs; \texttt{U5,\allowbreak{}U6,\allowbreak{}U7,\allowbreak{}U8,\allowbreak{}V2,\allowbreak{}V3,\allowbreak{}V4} on
\texttt{s}; \texttt{F3,\allowbreak{}F4,\allowbreak{}F5} on both
signs & 8 & 29 & 23 & 14 \\
Tail & No new LOW labels;13 HIGH additions & 21 & 29 & 23 & 27 \\
\texttt{B,\allowbreak{}D<a} &
\texttt{U5,\allowbreak{}U6,\allowbreak{}U7,\allowbreak{}U8,\allowbreak{}V3,\allowbreak{}V4,\allowbreak{}C66,\allowbreak{}C67,\allowbreak{}C68} on
\texttt{t} & 21 & 38 & 25 & 34 \\
\texttt{C<b} &
\texttt{C66,\allowbreak{}C67,\allowbreak{}C68,\allowbreak{}C35} on \texttt{s}
& 21 & 42 & 25 & 38 \\
Two cross labels & \texttt{C25} on both signs & 21 & 44
& 26 & 39 \\
Remaining three & \texttt{C35\_\allowbreak{}\allowbreak{}t,\allowbreak{}C45\_\allowbreak{}\allowbreak{}s,\allowbreak{}C45\_\allowbreak{}\allowbreak{}t} & 21 &
47 & 26 & 42 \\
\end{longtable}

For the last two stages the literal identities are

\[
 \begin{aligned}
 C25_s&=\allowbreak{}C-D-2b-g_t,\allowbreak{}&C25_t&=\allowbreak{}D-C-2a-g_s,\allowbreak{}\\
 C35_t&=\allowbreak{}C25_t-h_t,\allowbreak{}&C45_s&=\allowbreak{}C35_s-g_s,\allowbreak{}&C45_t&=\allowbreak{}C35_t-g_t.
 \end{aligned}
\]

All are negative because \texttt{1<=\allowbreak{}C<b},\allowbreak{}
\texttt{1<=\allowbreak{}D<a},\allowbreak{} and the outer gaps are positive. The
\texttt{C25\_\allowbreak{}\allowbreak{}s,\allowbreak{}C25\_\allowbreak{}\allowbreak{}t} group is the sole previously
untouched unstarred baseline row1,\allowbreak{} so two new labels cost one new
touched group. \texttt{C35\_\allowbreak{}\allowbreak{}t} is in the reflected row9
already touched by \texttt{V3\_\allowbreak{}\allowbreak{}t}; the two
\texttt{C45} labels are in the two row12 groups already
touched by \texttt{V8\_\allowbreak{}\allowbreak{}s,\allowbreak{}V8\_\allowbreak{}\allowbreak{}t}. Thus the final
compatible fixed credit is \textbf{42},\allowbreak{} and all native losses below
share \texttt{epsilon=\allowbreak{}delta-42>=\allowbreak{}0}.

\hypertarget{one-actual-adaptive-core-support}{%
\subsection{One actual adaptive core
support}\label{one-actual-adaptive-core-support}}

The final credit and \texttt{A<=\allowbreak{}60} imply
\texttt{I+\allowbreak{}tau<=\allowbreak{}18}. Let
\texttt{E\_\allowbreak{}\allowbreak{}v=\allowbreak{}tau\_\allowbreak{}\allowbreak{}v-xi\_\allowbreak{}\allowbreak{}v}. For
\texttt{j\_\allowbreak{}\allowbreak{}v} in \texttt{0,\allowbreak{}1,\allowbreak{}2,\allowbreak{}3,\allowbreak{}4},\allowbreak{}
remove that many largest positive core coefficients on sign
\texttt{v}. For zero removals use the excess
\texttt{xi\_\allowbreak{}\allowbreak{}v}. For positive removals,\allowbreak{} monotonicity of
the core deviations gives the retained maximum at most

\[
 r_v-j_v+\allowbreak{}\lfloor E_v/j_v\rfloor.
\]

Choose on each sign a minimizing removal count for

\[
 \phi_v=\allowbreak{}\min\left(\xi_v,\allowbreak{}
       \min_{1\le j\le4}\{j+\allowbreak{}\lfloor E_v/j\rfloor\}\right).
\]

This determines one \textbf{actual} subset \texttt{Q}
of the core,\allowbreak{} containing zero. It has
\texttt{q=\allowbreak{}69-j\_\allowbreak{}\allowbreak{}s-j\_\allowbreak{}\allowbreak{}t>=\allowbreak{}61}. Its integer displacement
overlap is at least \texttt{F-u},\allowbreak{} where

\[
 F=\allowbreak{}69-\phi_s-\phi_t.
\]

Both physical targets use this same subset and its reflection. All group
inequalities below use this chosen support; no overlaps of different
subsets are added or charged to the same budget. The triple bound keeps
\texttt{q>=\allowbreak{}61} explicitly,\allowbreak{} so replacing the old overlap
by \texttt{F} does not discard its source cardinality
condition.

\hypertarget{disjoint-groups-and-exact-transport}{%
\subsection{Disjoint groups and exact
transport}\label{disjoint-groups-and-exact-transport}}

Use the following original retained singleton groups,\allowbreak{} and no paired
capacity:

\begin{longtable}[]{@{}lll@{}}
\toprule\noalign{}
Group & Resources & Deficiency \\
\midrule\noalign{}
\endhead
\bottomrule\noalign{}
\endlastfoot
Pair & \texttt{B,\allowbreak{}D} &
\texttt{Gamma} \\
Triple & \texttt{H\_\allowbreak{}\allowbreak{}s,\allowbreak{}U\_\allowbreak{}\allowbreak{}s,\allowbreak{}V\_\allowbreak{}\allowbreak{}s} &
\texttt{Psi} \\
Two ordinary pairs & \texttt{U\_\allowbreak{}\allowbreak{}t,\allowbreak{}V\_\allowbreak{}\allowbreak{}t} and
\texttt{C6\_\allowbreak{}\allowbreak{}s,\allowbreak{}C7\_\allowbreak{}\allowbreak{}s} & Summed
\texttt{Omega} \\
Singleton & \texttt{C=\allowbreak{}C5\_\allowbreak{}\allowbreak{}s} &
\texttt{Gamma\_\allowbreak{}\allowbreak{}C} \\
Singleton & \texttt{H\_\allowbreak{}\allowbreak{}t} &
\texttt{Gamma\_\allowbreak{}\allowbreak{}Ht} \\
\end{longtable}

They contain eleven distinct retained singleton units. Their actual
union deficiencies sum to at most \texttt{epsilon}. The
joint counting lemma gives

\[
 2B+\allowbreak{}D\ge3(73-r_t)+\allowbreak{}F-3\Gamma,\allowbreak{}
 \quad g_s+\allowbreak{}h_s\ge q+\allowbreak{}F-\Psi,\allowbreak{}
 \quad2g_t\ge2F-\Omega.
\]

The first row combines
\texttt{B+\allowbreak{}D>=\allowbreak{}2(73-r\_\allowbreak{}\allowbreak{}t)+\allowbreak{}F-2Gamma} with the normal
native lower bound \texttt{B>=\allowbreak{}73-r\_\allowbreak{}\allowbreak{}t-Gamma}; it uses
the same \texttt{Gamma} throughout. Also
\texttt{l\_\allowbreak{}\allowbreak{}C<=\allowbreak{}Gamma\_\allowbreak{}\allowbreak{}C},\allowbreak{}
\texttt{u\_\allowbreak{}\allowbreak{}Ht<=\allowbreak{}Gamma\_\allowbreak{}\allowbreak{}Ht},\allowbreak{} and
\texttt{u\_\allowbreak{}\allowbreak{}Hs<=\allowbreak{}Psi}.

The literal centre identities and normal upper endpoints are

\[
 H_s=\allowbreak{}a+\allowbreak{}2b+\allowbreak{}G+\allowbreak{}h_s+\allowbreak{}B+\allowbreak{}C,\allowbreak{}\quad
 H_t=\allowbreak{}a+\allowbreak{}2b+\allowbreak{}G+\allowbreak{}h_s+\allowbreak{}B+\allowbreak{}D,\allowbreak{}
 \quad H_v\le S_v+\allowbreak{}r_o-72+\allowbreak{}u_{H_v}.
\]

Use the exact source mass \texttt{S=\allowbreak{}2K+\allowbreak{}1085-delta} and
\texttt{K<=\allowbreak{}M+\allowbreak{}R+\allowbreak{}21}. Adding the upper endpoints and
subtracting the packed lower endpoints gives precisely

\[
 2M_t+\allowbreak{}2X_s+\allowbreak{}4X_t+\allowbreak{}2(G+\allowbreak{}h_s)+\allowbreak{}F
 \le2R+\allowbreak{}827+\allowbreak{}2r_t-\delta+\allowbreak{}3\Gamma+\allowbreak{}l_C+\allowbreak{}u_{H_s}+\allowbreak{}u_{H_t}.
\]

Insert the two gap inequalities. The loss terms are bounded by
\texttt{3Gamma+\allowbreak{}3Psi+\allowbreak{}Omega+\allowbreak{}Gamma\_\allowbreak{}\allowbreak{}C+\allowbreak{}Gamma\_\allowbreak{}\allowbreak{}Ht<=\allowbreak{}3epsilon},\allowbreak{}
so

\[
 2M_t+\allowbreak{}2X_s+\allowbreak{}4X_t+\allowbreak{}2q+\allowbreak{}5F
 \le2R+\allowbreak{}827+\allowbreak{}2r_t-\delta+\allowbreak{}3\epsilon.
\]

Separately valid weighted stars on the same support give
\texttt{2X\_\allowbreak{}\allowbreak{}s+\allowbreak{}4X\_\allowbreak{}\allowbreak{}t>=\allowbreak{}6F-2epsilon}. Combining these
scalar inequalities does not assert that the star grouping is disjoint
from the displayed eleven-unit grouping. Finally use
\texttt{q>=\allowbreak{}61} to obtain

\[
 2M_t+\allowbreak{}122+\allowbreak{}11F\le2R+\allowbreak{}827+\allowbreak{}2r_t-\delta+\allowbreak{}5\epsilon.        \tag{*}
\]

\hypertarget{exact-finite-scalar-contradiction}{%
\subsection{Exact finite scalar
contradiction}\label{exact-finite-scalar-contradiction}}

Since \texttt{M\_\allowbreak{}\allowbreak{}t=\allowbreak{}r\_\allowbreak{}\allowbreak{}t+\allowbreak{}xi\_\allowbreak{}\allowbreak{}t+\allowbreak{}e\_\allowbreak{}\allowbreak{}t} and
\texttt{e\_\allowbreak{}\allowbreak{}t>=\allowbreak{}A+\allowbreak{}16},\allowbreak{} (*) would require

\[
 2I+\allowbreak{}2\tau-2\delta-2\xi_s+\allowbreak{}11F+\allowbreak{}5\cdot42-809\le0.
\]

Using \texttt{delta<=\allowbreak{}60-I-tau},\allowbreak{} its left side is at
least

\[
 4I+\allowbreak{}4\tau-2\xi_s+\allowbreak{}11F-719.
\]

The nonnegative \texttt{4I} may be omitted. The
complete remaining finite domain is
\texttt{0<=\allowbreak{}tau<=\allowbreak{}18},\allowbreak{}
\texttt{tau\_\allowbreak{}\allowbreak{}s+\allowbreak{}tau\_\allowbreak{}\allowbreak{}t=\allowbreak{}tau},\allowbreak{}
\texttt{0<=\allowbreak{}xi\_\allowbreak{}\allowbreak{}s<=\allowbreak{}tau\_\allowbreak{}\allowbreak{}s},\allowbreak{}
\texttt{0<=\allowbreak{}xi\_\allowbreak{}\allowbreak{}t<=\allowbreak{}tau\_\allowbreak{}\allowbreak{}t},\allowbreak{} all integers. It has7315
tuples. The exact checker finds

\[
 \min(4\tau-2\xi_s+\allowbreak{}11F-719)=\allowbreak{}6,\allowbreak{}
\]

attained in the relaxed domain at
\texttt{(tau,\allowbreak{}tau\_\allowbreak{}\allowbreak{}s,\allowbreak{}tau\_\allowbreak{}\allowbreak{}t,\allowbreak{}xi\_\allowbreak{}\allowbreak{}s,\allowbreak{}xi\_\allowbreak{}\allowbreak{}t,\allowbreak{}F)} equal to
\texttt{(16,\allowbreak{}11,\allowbreak{}5,\allowbreak{}5,\allowbreak{}3,\allowbreak{}61)}. Relaxing unrealizable core
parameters only enlarges the domain,\allowbreak{} so a positive minimum on this
domain suffices. This contradicts (*),\allowbreak{} and excludes the final A0
alpha-hole family.

The companion \texttt{z5\_\allowbreak{}\allowbreak{}alpha\_\allowbreak{}\allowbreak{}final85.py}
reconstructs every baseline group,\allowbreak{} checks all staged incidence counts
and literal new-label identities,\allowbreak{} selects one actual trimming
prescription per tuple,\allowbreak{} and checks this exact finite minimum. It does
not substitute arithmetic for the preceding actual-cover,\allowbreak{} normality,\allowbreak{} and
joint-deficiency arguments.

\clearpage
\hypertarget{fresh-size-85-z5-all-a1-and-b-alternatives}{%
\section{Fresh size-85 Z5: all A1 and B
alternatives}\label{fresh-size-85-z5-all-a1-and-b-alternatives}}

This proof uses the fresh structure and the actual-retention
classification in \texttt{zero85\_\allowbreak{}\allowbreak{}z4\_\allowbreak{}\allowbreak{}z5\_\allowbreak{}\allowbreak{}cards.md}.
Its source has 77 points,\allowbreak{} core 69 points,\allowbreak{} physical positive side sizes
\texttt{r\_\allowbreak{}\allowbreak{}s+\allowbreak{}r\_\allowbreak{}\allowbreak{}t=\allowbreak{}68},\allowbreak{} and no target holes. It does not
transplant the old Z5 normal-endpoint assumption.

Write

\begin{lstlisting}
I=(r_s-34)^2, A=I+tau+delta<=60,
R_v=r_v+xi_v, xi=xi_s+xi_t<=tau, R=68+xi,
M_v=R_v+e_v, L=e_s+e_t, M=R+L,
a=s4=M_s+X_s, b=t4=M_t+X_t, G=g_s+g_t.
\end{lstlisting}

The source rows used below are

\begin{lstlisting}
S_s+S_t=2K+1085-delta, K<=M+R+21,
L>=99+A, e_v>=A+16.
\end{lstlisting}

Several deliberately weaker computations use only
\texttt{e\_\allowbreak{}\allowbreak{}v>=\allowbreak{}A+\allowbreak{}8}. On each physical side,\allowbreak{}

\begin{lstlisting}
S_v>=T_(r_v)+tau_v+4R_v+6+e_v.
\end{lstlisting}

All retained ordinary resources and licensed sole pair components have a
core hit,\allowbreak{} since their deficiencies are at most 60,\allowbreak{} less than 69. With
lower and upper clipping counts l\_\allowbreak{}Z,\allowbreak{}u\_\allowbreak{}Z and native deficiency
Gamma\_\allowbreak{}Z,\allowbreak{}

\begin{lstlisting}
Z>=73-r_v-l_Z-xi_v,
Z<=S_v+r_o-72+u_Z+xi_o,
Gamma_Z>=l_Z+u_Z.
\end{lstlisting}

These rows include centers outside the target hull. Monotonicity of
\texttt{x\_\allowbreak{}\allowbreak{}i-i} bounds every core rank displacement by
xi\_\allowbreak{}v. No assertion that the centers themselves are inside the target is
made.

The eight O HIGH labels give fixed credit e\textgreater=\allowbreak{}8. Put
epsilon=\allowbreak{}delta-e. All compatible group deficiencies and unused native
losses share this one budget. The four extreme leaves are retained. Each
of the ordinary gap pairs contains at most one empty translate; their
intersections with an empty translate are off target and are charged
only to its retained partner\textquotesingle s actual loss.

Delete the four largest core magnitudes on each side. The Q61 support
gives

\begin{lstlisting}
f=61-(tau-xi)/4,
c_s X_s+c_t X_t >= (c_s+c_t)f-max(c_s,c_t)*epsilon/2,
c_s g_s+c_t g_t >= (c_s+c_t)f-max(c_s,c_t)*epsilon/2.
\end{lstlisting}

The second row uses the two disjoint ordinary pairs for each physical
gap. These star and ordinary-group inequalities are separately valid
scalar budget consequences. When endpoint resources are omitted from a
grouping,\allowbreak{} their native deficiencies can be rebated from epsilon. The
full-core star also gives
\texttt{2X\_\allowbreak{}\allowbreak{}v>=\allowbreak{}274-2R-epsilon}.

\hypertarget{four-reusable-transport-exclusions}{%
\subsection{Four reusable transport
exclusions}\label{four-reusable-transport-exclusions}}

\hypertarget{high-alpha_s}{%
\subsubsection{HIGH alpha\_\allowbreak{}s}\label{high-alpha_s}}

Suppose alpha\_\allowbreak{}s is HIGH,\allowbreak{} while H\_\allowbreak{}t,\allowbreak{}alpha\_\allowbreak{}t,\allowbreak{}C5\_\allowbreak{}t are retained. Since
alpha\_\allowbreak{}s\textgreater=\allowbreak{}S\_\allowbreak{}s+\allowbreak{}R\_\allowbreak{}t+\allowbreak{}1,\allowbreak{} the identity

\begin{lstlisting}
alpha_s+C5_t=d_t-2(a+b)-G
\end{lstlisting}

and the endpoints of C5\_\allowbreak{}t and H\_\allowbreak{}t-alpha\_\allowbreak{}t=\allowbreak{}d\_\allowbreak{}t give

\begin{lstlisting}
2S_s+R_t+2X+G<=2R+976+r_t+xi+xi_t.
\end{lstlisting}

The three endpoint losses share epsilon. Using
\texttt{S\_\allowbreak{}\allowbreak{}s>=\allowbreak{}T\_\allowbreak{}\allowbreak{}r+\allowbreak{}tau\_\allowbreak{}\allowbreak{}s+\allowbreak{}4R\_\allowbreak{}\allowbreak{}s+\allowbreak{}10},\allowbreak{}
\texttt{2X>=\allowbreak{}4f-epsilon},\allowbreak{} and discarding G\textgreater=\allowbreak{}0
gives

\begin{lstlisting}
2T_r+8r+I+2tau_s+6xi_s-2xi_t+e<=928.
\end{lstlisting}

Because xi\_\allowbreak{}t\textless=\allowbreak{}tau\_\allowbreak{}t and tau+\allowbreak{}delta\textless=\allowbreak{}60-I,\allowbreak{}

\begin{lstlisting}
2T_r+8r+3I+2delta+e<=1048.
\end{lstlisting}

For every integer r,\allowbreak{} \texttt{2T\_\allowbreak{}\allowbreak{}r+\allowbreak{}8r+\allowbreak{}3(r-34)\^2>=\allowbreak{}1092}
(minimum at r=\allowbreak{}24). Thus this HIGH-alpha alternative is impossible.

\hypertarget{both-pairs-retained-opposite-owners}{%
\subsubsection{Both pairs retained,\allowbreak{} opposite
owners}\label{both-pairs-retained-opposite-owners}}

Let beta\_\allowbreak{}s and P\_\allowbreak{}t be the actual sole owners. Use their lower
endpoints,\allowbreak{} the upper endpoint of V\_\allowbreak{}t,\allowbreak{} and the leading-s inequality
twice. The literal identity is

\begin{lstlisting}
V_t=2d_s+beta_s+P_t+2a+b+g_s.
\end{lstlisting}

It gives

\begin{lstlisting}
S_s+3g_s+4X_s+2a+b
   <=935+r_s+delta-2e+xi+xi_s.
\end{lstlisting}

The source lower row and the separately valid bounds
\texttt{6X\_\allowbreak{}\allowbreak{}s+\allowbreak{}X\_\allowbreak{}\allowbreak{}t>=\allowbreak{}7f-3epsilon} and
\texttt{3g\_\allowbreak{}\allowbreak{}s>=\allowbreak{}3f-3epsilon/\allowbreak{}2} yield,\allowbreak{} using the weaker
e\_\allowbreak{}v\textgreater=\allowbreak{}A+\allowbreak{}8,\allowbreak{}

\begin{lstlisting}
T_r+4r+3I+tau_s+tau/2+4xi_s+5xi/2
   <=136+5delta/2-13e/2.
\end{lstlisting}

The source mass cap therefore implies

\begin{lstlisting}
T_r+4r+11I/2+tau_s+3tau+4xi_s+5xi/2
   <=286-13e/2.
\end{lstlisting}

Its left side is at least \texttt{T\_\allowbreak{}\allowbreak{}r+\allowbreak{}4r+\allowbreak{}3I>=\allowbreak{}626},\allowbreak{}
whereas its right side is at most 234. Opposite owners are impossible.

\hypertarget{both-pairs-retained-same-owner-alpha_s-retained}{%
\subsubsection{Both pairs retained,\allowbreak{} same owner,\allowbreak{} alpha\_\allowbreak{}s
retained}\label{both-pairs-retained-same-owner-alpha_s-retained}}

Let both sole owners be s and put r=\allowbreak{}r\_\allowbreak{}t. Since
alpha\_\allowbreak{}t=\allowbreak{}h\_\allowbreak{}t-b\textgreater-b,\allowbreak{} the identity

\begin{lstlisting}
alpha_s-P_s=beta_s+beta_t+2d_t+2(a+b)+G
\end{lstlisting}

gives the strict lower relation
\texttt{alpha\_\allowbreak{}\allowbreak{}s-P\_\allowbreak{}\allowbreak{}s>beta\_\allowbreak{}\allowbreak{}s+\allowbreak{}d\_\allowbreak{}\allowbreak{}t+\allowbreak{}2a+\allowbreak{}b+\allowbreak{}G}. Use the
endpoints of the three retained s-resources and one leading-t
inequality. The transport is

\begin{lstlisting}
S_t+M_s+2X_s+3X_t+g_s+2g_t
   <=R+1023-r-e+2xi_s+xi_t.
\end{lstlisting}

Here the source row is
\texttt{S\_\allowbreak{}\allowbreak{}t+\allowbreak{}M\_\allowbreak{}\allowbreak{}s>=\allowbreak{}T\_\allowbreak{}\allowbreak{}r+\allowbreak{}3r+\allowbreak{}tau\_\allowbreak{}\allowbreak{}t+\allowbreak{}4xi\_\allowbreak{}\allowbreak{}t+\allowbreak{}xi\_\allowbreak{}\allowbreak{}s+\allowbreak{}173+\allowbreak{}A}.
Together with \texttt{2X\_\allowbreak{}\allowbreak{}s+\allowbreak{}3X\_\allowbreak{}\allowbreak{}t>=\allowbreak{}5f-3epsilon/\allowbreak{}2},\allowbreak{}
\texttt{g\_\allowbreak{}\allowbreak{}s+\allowbreak{}2g\_\allowbreak{}\allowbreak{}t>=\allowbreak{}3f-epsilon},\allowbreak{} and
e\_\allowbreak{}v\textgreater=\allowbreak{}A+\allowbreak{}8 imply

\begin{lstlisting}
T_r+4r+5I/2+tau_t+tau/2+4xi_t<=520-7e/2.
\end{lstlisting}

The left side is at least
\texttt{T\_\allowbreak{}\allowbreak{}r+\allowbreak{}4r+\allowbreak{}5(r-34)\^2/\allowbreak{}2>=\allowbreak{}608},\allowbreak{} while the right
side is at most 492. This alternative is impossible.

\hypertarget{the-branch-independent-four-endpoint-exclusion}{%
\subsubsection{The branch-independent four-endpoint
exclusion}\label{the-branch-independent-four-endpoint-exclusion}}

Suppose pair1 has sole owner s,\allowbreak{} alpha\_\allowbreak{}t is retained,\allowbreak{} and all eight
lower own/cross singletons are retained. The literal identity

\begin{lstlisting}
U_s=P_s+alpha_t+C6_t+2a+4b+2G
\end{lstlisting}

uses four distinct endpoints. Their losses share epsilon,\allowbreak{} giving

\begin{lstlisting}
S_t+2M_t+2X_s+4X_t+2G
   <=904+2r_t+2R-e+xi_s+3xi_t.
\end{lstlisting}

Use \texttt{S\_\allowbreak{}\allowbreak{}t+\allowbreak{}2M\_\allowbreak{}\allowbreak{}t>=\allowbreak{}T\_\allowbreak{}\allowbreak{}r+\allowbreak{}tau\_\allowbreak{}\allowbreak{}t+\allowbreak{}6R\_\allowbreak{}\allowbreak{}t+\allowbreak{}18},\allowbreak{}
\texttt{2X\_\allowbreak{}\allowbreak{}s+\allowbreak{}4X\_\allowbreak{}\allowbreak{}t>=\allowbreak{}6f-2epsilon},\allowbreak{} and
\texttt{2G>=\allowbreak{}4f-epsilon}. Then

\begin{lstlisting}
T_r+4r+3I+tau_t+tau/2-xi_s/2+7xi_t/2<=592-4e.
\end{lstlisting}

Since xi\_\allowbreak{}s\textless=\allowbreak{}tau\_\allowbreak{}s,\allowbreak{} its remaining excess terms are nonnegative.
Again the left side is at least 626 and the right side at most 560. This
exclusion uses no beta capacity,\allowbreak{} endpoint,\allowbreak{} or ownership.

\hypertarget{a1-with-pair1-empty}{%
\subsection{A1 with pair1 empty}\label{a1-with-pair1-empty}}

The sole beta owner is s. The base holes are O\_\allowbreak{}s,\allowbreak{}H\_\allowbreak{}s,\allowbreak{}O\_\allowbreak{}t,\allowbreak{}pair1 and
the remaining hole eta is exactly one of

\begin{lstlisting}
alpha_s, H_t, alpha_t, C5_s, C5_t, U_s.
\end{lstlisting}

For eta=\allowbreak{}alpha\_\allowbreak{}s,\allowbreak{} emptiness above an active beta makes alpha\_\allowbreak{}s HIGH;
the HIGH-alpha exclusion applies. Explicitly beta\_\allowbreak{}s has a core hit,\allowbreak{} so
beta\_\allowbreak{}s\textgreater=\allowbreak{}1-R\_\allowbreak{}s. The leading inequality and A\textless=\allowbreak{}60
imply \texttt{d\_\allowbreak{}\allowbreak{}s-R\_\allowbreak{}\allowbreak{}s>=\allowbreak{}93-r\_\allowbreak{}\allowbreak{}s+\allowbreak{}I>0},\allowbreak{} hence
alpha\_\allowbreak{}s=\allowbreak{}d\_\allowbreak{}s+\allowbreak{}beta\_\allowbreak{}s\textgreater0. Since
\texttt{S\_\allowbreak{}\allowbreak{}s>R\_\allowbreak{}\allowbreak{}t},\allowbreak{} an empty alpha\_\allowbreak{}s cannot have
\texttt{S\_\allowbreak{}\allowbreak{}s<alpha\_\allowbreak{}\allowbreak{}s<=\allowbreak{}S\_\allowbreak{}\allowbreak{}s+\allowbreak{}R\_\allowbreak{}\allowbreak{}t}: its image at source
point -R\_\allowbreak{}t would be a target hit. Thus
alpha\_\allowbreak{}s\textgreater=\allowbreak{}S\_\allowbreak{}s+\allowbreak{}R\_\allowbreak{}t+\allowbreak{}1,\allowbreak{} as the exclusion requires.

For eta=\allowbreak{}C5\_\allowbreak{}s or U\_\allowbreak{}s,\allowbreak{} the retained endpoints are
alpha\_\allowbreak{}s,\allowbreak{}beta\_\allowbreak{}s,\allowbreak{}H\_\allowbreak{}t,\allowbreak{}alpha\_\allowbreak{}t,\allowbreak{}C5\_\allowbreak{}t. The identities
alpha\_\allowbreak{}s-beta\_\allowbreak{}s=\allowbreak{}d\_\allowbreak{}s,\allowbreak{} H\_\allowbreak{}t-alpha\_\allowbreak{}t=\allowbreak{}d\_\allowbreak{}t,\allowbreak{} and
beta\_\allowbreak{}s+\allowbreak{}C5\_\allowbreak{}t=\allowbreak{}d\_\allowbreak{}t-d\_\allowbreak{}s-2(a+\allowbreak{}b)-G,\allowbreak{} followed by two leading-s bounds,\allowbreak{} give

\begin{lstlisting}
2M+6X_s+2X_t+3g_s+g_t
   <=853+delta-2e+3xi.
\end{lstlisting}

All endpoint and leading loss coefficients are at most two on distinct
resources. The weighted star and gap bounds give
\texttt{12f-9epsilon/\allowbreak{}2} for the gap part,\allowbreak{} so the source
cap implies

\begin{lstlisting}
3+11I/2+5tau/2+13e/2<=0.
\end{lstlisting}

This is impossible. The ordinary gap pairs may contain the empty U\_\allowbreak{}s;
its intersection is paid solely by the retained V\_\allowbreak{}s loss.

For eta=\allowbreak{}H\_\allowbreak{}t,\allowbreak{} use lower alpha\_\allowbreak{}t,\allowbreak{}C5\_\allowbreak{}s,\allowbreak{}beta\_\allowbreak{}s,\allowbreak{}C5\_\allowbreak{}t and upper
alpha\_\allowbreak{}s,\allowbreak{}V\_\allowbreak{}t. Their identities give

\begin{lstlisting}
4M_s+3M_t+6X_s+5X_t+3g_s+2g_t
   <=2R+744+r_s+delta-2e+4xi+xi_s.
\end{lstlisting}

The source lower bound is
\texttt{4M\_\allowbreak{}\allowbreak{}s+\allowbreak{}3M\_\allowbreak{}\allowbreak{}t>=\allowbreak{}3R+\allowbreak{}305+\allowbreak{}r\_\allowbreak{}\allowbreak{}s+\allowbreak{}xi\_\allowbreak{}\allowbreak{}s+\allowbreak{}4A},\allowbreak{} using
e\_\allowbreak{}v\textgreater=\allowbreak{}A+\allowbreak{}8. The gap part is at least
\texttt{16f-9epsilon/\allowbreak{}2}. Consequently

\begin{lstlisting}
537+R+4I-3delta/2+13e/2<=0,
\end{lstlisting}

contradicting R\textgreater=\allowbreak{}68 and delta\textless=\allowbreak{}60.

For eta=\allowbreak{}alpha\_\allowbreak{}t,\allowbreak{} use upper alpha\_\allowbreak{}s,\allowbreak{}H\_\allowbreak{}t and lower beta\_\allowbreak{}s,\allowbreak{}C5\_\allowbreak{}t,\allowbreak{}C5\_\allowbreak{}s.
The latter beta and C5\_\allowbreak{}s clipping coefficients are two. The exact
transport is

\begin{lstlisting}
4M_s+M_t+6X_s+3X_t+3G+3h_t
   <=2R+754+3r_s+delta-2e+5xi_s+2xi_t.
\end{lstlisting}

Group the two g\_\allowbreak{}s pairs with total deficiency Omega and the triple
\{H\_\allowbreak{}t,\allowbreak{}U\_\allowbreak{}t,\allowbreak{}V\_\allowbreak{}t\} with deficiency Psi. The triple has
\texttt{g\_\allowbreak{}\allowbreak{}t+\allowbreak{}h\_\allowbreak{}\allowbreak{}t>=\allowbreak{}61+\allowbreak{}f-Psi}; hence
\texttt{3G+\allowbreak{}3h\_\allowbreak{}\allowbreak{}t>=\allowbreak{}6f+\allowbreak{}183-3epsilon}. Separately the
stars give \texttt{6X\_\allowbreak{}\allowbreak{}s+\allowbreak{}3X\_\allowbreak{}\allowbreak{}t>=\allowbreak{}9f-3epsilon}. The
physical source lower row
\texttt{4M\_\allowbreak{}\allowbreak{}s+\allowbreak{}M\_\allowbreak{}\allowbreak{}t>=\allowbreak{}R+\allowbreak{}123+\allowbreak{}3r\_\allowbreak{}\allowbreak{}s+\allowbreak{}3xi\_\allowbreak{}\allowbreak{}s+\allowbreak{}4A} then implies

\begin{lstlisting}
219+7I+13tau/4+3xi/4+8e<=0,
\end{lstlisting}

again impossible. No LOW-alpha endpoint assumption is needed.

\hypertarget{the-final-a1-hole-etac5_t}{%
\subsection{The final A1 hole
eta=\allowbreak{}C5\_\allowbreak{}t}\label{the-final-a1-hole-etac5_t}}

The exact holes are O\_\allowbreak{}s,\allowbreak{}H\_\allowbreak{}s,\allowbreak{}O\_\allowbreak{}t,\allowbreak{}pair1,\allowbreak{}C5\_\allowbreak{}t. All eight lower
own/cross singletons,\allowbreak{} alpha\_\allowbreak{}s,\allowbreak{}alpha\_\allowbreak{}t,\allowbreak{}H\_\allowbreak{}t,\allowbreak{}C5\_\allowbreak{}s and sole beta\_\allowbreak{}s are
retained.

Initially the eight O labels and C22\_\allowbreak{}s,\allowbreak{}C23\_\allowbreak{}s,\allowbreak{}C24\_\allowbreak{}s are HIGH. The LOW
labels are C55,\allowbreak{}C56,\allowbreak{}C57,\allowbreak{}C58 on both signs and C66\_\allowbreak{}t,\allowbreak{}C67\_\allowbreak{}t,\allowbreak{}C68\_\allowbreak{}t. There
are 11 HIGH,\allowbreak{} 11 LOW,\allowbreak{} and 9 baseline groups touched by LOW,\allowbreak{} giving
compatible fixed credit e=\allowbreak{}13.

\hypertarget{the-beta-cutoff-with-actual-physical-target-capacity}{%
\subsubsection{The beta cutoff,\allowbreak{} with actual physical-target
capacity}\label{the-beta-cutoff-with-actual-physical-target-capacity}}

We first prove \texttt{beta\_\allowbreak{}\allowbreak{}s+\allowbreak{}s3<d\_\allowbreak{}\allowbreak{}s-X\_\allowbreak{}\allowbreak{}s}. Suppose
the reverse. Since
\texttt{d\_\allowbreak{}\allowbreak{}s=\allowbreak{}2(a+\allowbreak{}b)+\allowbreak{}G+\allowbreak{}C5\_\allowbreak{}\allowbreak{}s+\allowbreak{}alpha\_\allowbreak{}\allowbreak{}t},\allowbreak{} it forces

\begin{lstlisting}
alpha_s>=3M_s+4M_t+2X_s+4X_t+g_s+2g_t
             +2C5_s+2alpha_t.
\end{lstlisting}

Target t has its leading and exactly six retained ordinary resources
H\_\allowbreak{}t,\allowbreak{}alpha\_\allowbreak{}t,\allowbreak{}U\_\allowbreak{}t,\allowbreak{}V\_\allowbreak{}t,\allowbreak{}C6\_\allowbreak{}t,\allowbreak{}C7\_\allowbreak{}t. Their total capacity is K+\allowbreak{}462.
Their grouped deficiency is at most epsilon,\allowbreak{} so the exact two-target
mass gives \texttt{S\_\allowbreak{}\allowbreak{}s<=\allowbreak{}K+\allowbreak{}623-e}. The endpoint
inequality above therefore gives

\begin{lstlisting}
2M_s+3M_t+2X_s+4X_t+g_s+2g_t
   <=R+416+r_t-e+E+2xi_s+3xi_t,
\end{lstlisting}

where
\texttt{E=\allowbreak{}2l\_\allowbreak{}\allowbreak{}(C5\_\allowbreak{}\allowbreak{}s)+\allowbreak{}2l\_\allowbreak{}\allowbreak{}(alpha\_\allowbreak{}\allowbreak{}t)+\allowbreak{}u\_\allowbreak{}\allowbreak{}(alpha\_\allowbreak{}\allowbreak{}s)}.
The source gives \texttt{2M\_\allowbreak{}\allowbreak{}s+\allowbreak{}3M\_\allowbreak{}\allowbreak{}t>=\allowbreak{}2R+\allowbreak{}R\_\allowbreak{}\allowbreak{}t+\allowbreak{}214+\allowbreak{}3A}.
The weighted stars give
\texttt{2X\_\allowbreak{}\allowbreak{}s+\allowbreak{}4X\_\allowbreak{}\allowbreak{}t>=\allowbreak{}6f-2epsilon}. If Omega is the
weighted ordinary-gap deficiency,\allowbreak{} then
\texttt{g\_\allowbreak{}\allowbreak{}s+\allowbreak{}2g\_\allowbreak{}\allowbreak{}t>=\allowbreak{}3f-Omega} and
\texttt{E+\allowbreak{}Omega<=\allowbreak{}2epsilon},\allowbreak{} because those ordinary
groups avoid the three endpoints. Thus

\begin{lstlisting}
3A+9f<=134-e+4epsilon+xi.
\end{lstlisting}

The source cap turns this into

\begin{lstlisting}
355+4I+7tau/4+5xi/4+5e<=0,
\end{lstlisting}

an immediate contradiction. This proves the beta cutoff.

\hypertarget{the-complete-actual-upper-tail-inventory}{%
\subsubsection{The complete actual upper-tail
inventory}\label{the-complete-actual-upper-tail-inventory}}

Put t0=\allowbreak{}d\_\allowbreak{}s-X\_\allowbreak{}s\textgreater0 and add the four s-outer magnitudes to a
target value u\textgreater t0. A four-term representation with at most
one top-two occurrence is at most s1+\allowbreak{}3s3. With exactly two and not both
s1,\allowbreak{} it is at most s1+\allowbreak{}s2+\allowbreak{}2s3. Both lie below the required total,\allowbreak{} since
\texttt{d\_\allowbreak{}\allowbreak{}s-g\_\allowbreak{}\allowbreak{}s-X\_\allowbreak{}\allowbreak{}s>=\allowbreak{}K-native\_\allowbreak{}\allowbreak{}leading\_\allowbreak{}\allowbreak{}loss+\allowbreak{}X\_\allowbreak{}\allowbreak{}s>0}.
If both are s1,\allowbreak{} both remaining entries must be s3 or s4; otherwise the
total is at most the threshold. This gives three types C33,\allowbreak{}C34,\allowbreak{}C44.
Exactly three top-two occurrences give the four source families
O\_\allowbreak{}s+\allowbreak{}D\_\allowbreak{}s,\allowbreak{}H\_\allowbreak{}s+\allowbreak{}D\_\allowbreak{}s,\allowbreak{}alpha\_\allowbreak{}s+\allowbreak{}D\_\allowbreak{}s,\allowbreak{}beta\_\allowbreak{}s+\allowbreak{}D\_\allowbreak{}s,\allowbreak{} plus 24 free-outer
types. Exactly four give five further fixed types.

The complete 32 non-source types,\allowbreak{} in the original fixed inventory,\allowbreak{} are

\begin{lstlisting}
C13,C14,C15,C16,C17,C18,
C23,C24,C25,C26,C27,C28,
U3,U4,U5,U6,U7,U8,
V3,V4,V5,V6,V7,V8,
V2,U2,C22,C12,C11,C44,C34,C33,
\end{lstlisting}

all on target s. O\_\allowbreak{}s and H\_\allowbreak{}s source families are empty. The cutoff
removes beta\_\allowbreak{}s+\allowbreak{}D\_\allowbreak{}s and the six V3..V8 fixed types. Thus only the
retained alpha\_\allowbreak{}s source family and at most 26 fixed types cover the
tail. Writing Gamma\_\allowbreak{}alpha for its native deficiency,\allowbreak{}

\begin{lstlisting}
S_s<=d_s-X_s+103-Gamma_alpha.
\end{lstlisting}

The inequality is automatic when S\_\allowbreak{}s\textless=\allowbreak{}t0.

\hypertarget{four-extra-high-types-and-the-sharper-tail}{%
\subsubsection{Four extra HIGH types and the sharper
tail}\label{four-extra-high-types-and-the-sharper-tail}}

The core hit of beta\_\allowbreak{}s gives beta\_\allowbreak{}s\textgreater=\allowbreak{}1-R\_\allowbreak{}s. Since
\texttt{h\_\allowbreak{}\allowbreak{}s=\allowbreak{}d\_\allowbreak{}\allowbreak{}s+\allowbreak{}a+\allowbreak{}beta\_\allowbreak{}\allowbreak{}s},\allowbreak{} the preceding tail bound
and the full-core star give

\begin{lstlisting}
h_s-S_s>=e_s+1+2X_s-103+Gamma_alpha
         >=52+I+tau-2xi+e
         >=-8+2I+delta+e>0.
\end{lstlisting}

Thus U4\_\allowbreak{}s=\allowbreak{}h\_\allowbreak{}s and U3\_\allowbreak{}s=\allowbreak{}h\_\allowbreak{}s+\allowbreak{}g\_\allowbreak{}s are HIGH. Also
\texttt{h\_\allowbreak{}\allowbreak{}s-d\_\allowbreak{}\allowbreak{}s=\allowbreak{}a+\allowbreak{}beta\_\allowbreak{}\allowbreak{}s>=\allowbreak{}e\_\allowbreak{}\allowbreak{}s+\allowbreak{}X\_\allowbreak{}\allowbreak{}s+\allowbreak{}1>0},\allowbreak{} so
V2\_\allowbreak{}s=\allowbreak{}2h\_\allowbreak{}s+\allowbreak{}g\_\allowbreak{}s-d\_\allowbreak{}s and U2\_\allowbreak{}s=\allowbreak{}2h\_\allowbreak{}s+\allowbreak{}g\_\allowbreak{}s are HIGH. These four distinct
labels are disjoint from the previous HIGH ledger,\allowbreak{} which now has 15
HIGH,\allowbreak{} 11 LOW,\allowbreak{} touched 9: the credit is e=\allowbreak{}17.

All four labels belong to the 26 remaining types in the displayed tail
inventory. Removing them leaves at most 22 fixed types,\allowbreak{} giving

\begin{lstlisting}
S_s<=d_s-X_s+99-Gamma_alpha.
\end{lstlisting}

The generic endpoint rows of beta\_\allowbreak{}s and alpha\_\allowbreak{}s now imply

\begin{lstlisting}
X_s<=22+l_beta-l_alpha+xi.
\end{lstlisting}

\hypertarget{the-final-trim-one-stars-with-native-rebates}{%
\subsubsection{The final trim-one stars,\allowbreak{} with native
rebates}\label{the-final-trim-one-stars-with-native-rebates}}

Remove the largest core magnitude on each sign. The actual Q67 support
has span at most \texttt{66+\allowbreak{}tau-xi},\allowbreak{} so

\begin{lstlisting}
|Q67 intersect (Q67+q)|>=f1-q, f1=67-tau+xi.
\end{lstlisting}

The two physical-s stars \{L\_\allowbreak{}s,\allowbreak{}V\_\allowbreak{}s\} and \{L\_\allowbreak{}t,\allowbreak{}C7\_\allowbreak{}t\} are retained
and avoid beta\_\allowbreak{}s,\allowbreak{}alpha\_\allowbreak{}s. Their original-budget inequality is

\begin{lstlisting}
2X_s>=2f1-epsilon+Gamma_beta+Gamma_alpha.
\end{lstlisting}

Combine this with the upper bound on X\_\allowbreak{}s and use
\texttt{Gamma\_\allowbreak{}\allowbreak{}beta>=\allowbreak{}l\_\allowbreak{}\allowbreak{}beta},\allowbreak{}
\texttt{Gamma\_\allowbreak{}\allowbreak{}alpha>=\allowbreak{}l\_\allowbreak{}\allowbreak{}alpha},\allowbreak{} and
\texttt{Gamma\_\allowbreak{}\allowbreak{}beta+\allowbreak{}Gamma\_\allowbreak{}\allowbreak{}alpha<=\allowbreak{}epsilon}. It yields

\begin{lstlisting}
2f1<=44+2xi+2epsilon,
tau+delta>=45+e=62.
\end{lstlisting}

This contradicts \texttt{tau+\allowbreak{}delta<=\allowbreak{}60-I}. All ranks,\allowbreak{}
outside-hull endpoints,\allowbreak{} and possible numerical coincidences are
included; no normality assumption is needed.

\hypertarget{routing-the-retained-pair1-a1-alternatives}{%
\subsection{Routing the retained-pair1 A1
alternatives}\label{routing-the-retained-pair1-a1-alternatives}}

First exclude the initially split beta case using only its whole paired
capacity. Both beta components active force the holes
\texttt{O\_\allowbreak{}\allowbreak{}s,\allowbreak{}O\_\allowbreak{}\allowbreak{}t,\allowbreak{}H\_\allowbreak{}\allowbreak{}s,\allowbreak{}H\_\allowbreak{}\allowbreak{}t,\allowbreak{}C5\_\allowbreak{}\allowbreak{}t},\allowbreak{} with
pair1\textquotesingle s actual sole owner s and every lower own/cross
resource and both alphas retained. The whole-pair total hit count N
satisfies

\begin{lstlisting}
77-delta<=N<=77, beta_s+beta_t>=N-M.
\end{lstlisting}

Use the endpoints of P\_\allowbreak{}s,\allowbreak{}alpha\_\allowbreak{}s and the identity
\texttt{alpha\_\allowbreak{}\allowbreak{}s-P\_\allowbreak{}\allowbreak{}s=\allowbreak{}beta\_\allowbreak{}\allowbreak{}s+\allowbreak{}beta\_\allowbreak{}\allowbreak{}t+\allowbreak{}2d\_\allowbreak{}\allowbreak{}t+\allowbreak{}2(a+\allowbreak{}b)+\allowbreak{}G}.
Two leading-t inequalities then give

\begin{lstlisting}
S_t+M+2X_s+6X_t+g_s+3g_t<=931+2delta+xi.
\end{lstlisting}

The native endpoint and leading losses have coefficient at most two; no
selected beta component receives its own capacity or deficiency. The
weaker physical-source row
\texttt{S\_\allowbreak{}\allowbreak{}t+\allowbreak{}M>=\allowbreak{}T\_\allowbreak{}\allowbreak{}r+\allowbreak{}4r+\allowbreak{}R+\allowbreak{}109+\allowbreak{}A+\allowbreak{}tau\_\allowbreak{}\allowbreak{}t+\allowbreak{}4xi\_\allowbreak{}\allowbreak{}t},\allowbreak{}
together with \texttt{2X\_\allowbreak{}\allowbreak{}s+\allowbreak{}6X\_\allowbreak{}\allowbreak{}t>=\allowbreak{}8f-3epsilon} and
\texttt{g\_\allowbreak{}\allowbreak{}s+\allowbreak{}3g\_\allowbreak{}\allowbreak{}t>=\allowbreak{}4f-3epsilon/\allowbreak{}2},\allowbreak{} implies

\begin{lstlisting}
T_r+4r+13I/2+7tau/2+tau_t+3xi+4xi_t<=316.
\end{lstlisting}

But \texttt{T\_\allowbreak{}\allowbreak{}r+\allowbreak{}4r+\allowbreak{}13(r-34)\^2/\allowbreak{}2>=\allowbreak{}1357/\allowbreak{}2} for all
integers r. This excludes split beta before any sole-beta native loss is
licensed.

If pair1 is retained,\allowbreak{} opposite owners are excluded above. Same owners
with alpha\_\allowbreak{}s retained are excluded above. If alpha\_\allowbreak{}s is empty,\allowbreak{} the
owner identity forces C5\_\allowbreak{}t empty; the exact remaining holes are
O\_\allowbreak{}s,\allowbreak{}H\_\allowbreak{}s,\allowbreak{}O\_\allowbreak{}t,\allowbreak{}alpha\_\allowbreak{}s,\allowbreak{}C5\_\allowbreak{}t. Pair1,\allowbreak{}alpha\_\allowbreak{}t and all eight lower
resources are retained,\allowbreak{} so the branch-independent four-endpoint
exclusion applies. Together with the six pair1-empty placements this
exhausts A1.

\hypertarget{branch-b}{%
\subsection{Branch B}\label{branch-b}}

In B,\allowbreak{} beta\_\allowbreak{}s,\allowbreak{}beta\_\allowbreak{}t are singleton units,\allowbreak{} and C5\_\allowbreak{}s,\allowbreak{}C5\_\allowbreak{}t form a
retained paired unit with actual sole owner s. Pair1,\allowbreak{} when retained,\allowbreak{} has
the same owner. The classification proves integrality and extreme-leaf
retention before any native endpoint is used.

If alpha\_\allowbreak{}s is retained,\allowbreak{} branch B\textquotesingle s strict sign gives

\begin{lstlisting}
alpha_s-C5_s> d_t+2(a+b)+G.
\end{lstlisting}

Use its two endpoints and one leading-t inequality. The transport is

\begin{lstlisting}
S_t+M+2X_s+4X_t+g_s+2g_t
   <R+1029-e+xi.
\end{lstlisting}

The physical source lower row and the star/gap bounds,\allowbreak{} with the weaker
e\_\allowbreak{}v\textgreater=\allowbreak{}A+\allowbreak{}8,\allowbreak{} imply

\begin{lstlisting}
T_r+4r+3I+tau_t+3tau/4+4xi_t+5xi/4<427-4e.
\end{lstlisting}

The left side is at least 626 and the right side at most 395.

If alpha\_\allowbreak{}s is empty,\allowbreak{} it is the only extra hole outside the mandatory
beta/H choices. Consequently pair1,\allowbreak{}alpha\_\allowbreak{}t and all eight lower
resources are retained,\allowbreak{} and pair1 has sole owner s. The
branch-independent four-endpoint exclusion applies. It borrows no
beta-pair capacity from A. These alternatives exhaust B.

The exact arithmetic checker is
\texttt{check\_\allowbreak{}\allowbreak{}zero85\_\allowbreak{}\allowbreak{}z5\_\allowbreak{}\allowbreak{}a1\_\allowbreak{}\allowbreak{}b.py}. It checks the
displayed integral minima and terminal margins; the physical-target,\allowbreak{}
inventory,\allowbreak{} endpoint,\allowbreak{} and budget deductions are the written proof above.
The independent final-family audit is
\texttt{../\allowbreak{}scalar/\allowbreak{}z5\_\allowbreak{}\allowbreak{}c5t\_\allowbreak{}\allowbreak{}cutoff\_\allowbreak{}\allowbreak{}audit.md},\allowbreak{} with
standalone checker
\texttt{../\allowbreak{}scalar/\allowbreak{}check\_\allowbreak{}\allowbreak{}z5\_\allowbreak{}\allowbreak{}c5t\_\allowbreak{}\allowbreak{}cutoff\_\allowbreak{}\allowbreak{}audit.py}. It
independently confirms the physical-target cutoff,\allowbreak{} all four additional
HIGH types,\allowbreak{} the 26-to-22 tail pruning and the generic-endpoint trim-one
contradiction.

\clearpage
\hypertarget{independent-z5-c5_t-final-cutoff-and-tail-audit}{%
\section{Independent Z5 C5\_\allowbreak{}t-final cutoff and tail
audit}\label{independent-z5-c5_t-final-cutoff-and-tail-audit}}

This is a conditional local argument using the cardinality-85
structural/activity conclusions. Exact holes are
O\_\allowbreak{}s,\allowbreak{}H\_\allowbreak{}s,\allowbreak{}O\_\allowbreak{}t,\allowbreak{}pair1,\allowbreak{}C5\_\allowbreak{}t. The retained beta pair has sole actual
contributor beta\_\allowbreak{}s. Every resource below is its original full native
resource; neither empty pair nor an inactive component is assigned
capacity.

Notation: \textbar D\textbar=\allowbreak{}77,\allowbreak{} r\_\allowbreak{}s+\allowbreak{}r\_\allowbreak{}t=\allowbreak{}68,\allowbreak{} R=\allowbreak{}68+\allowbreak{}xi,\allowbreak{}
A=\allowbreak{}I+\allowbreak{}tau+\allowbreak{}delta\textless=\allowbreak{}60,\allowbreak{} epsilon=\allowbreak{}delta-e. Source bounds used are
K\textless=\allowbreak{}M+\allowbreak{}R+\allowbreak{}21,\allowbreak{} M\textgreater=\allowbreak{}R+\allowbreak{}99+\allowbreak{}A,\allowbreak{} and M\_\allowbreak{}t\textgreater=\allowbreak{}R\_\allowbreak{}t+\allowbreak{}16+\allowbreak{}A.
The fixed ledger initially gives e=\allowbreak{}13. Generic endpoints have
lower73-r\_\allowbreak{}v-l-xi\_\allowbreak{}v and upperS\_\allowbreak{}v+\allowbreak{}r\_\allowbreak{}o-72+\allowbreak{}u+\allowbreak{}xi\_\allowbreak{}o. All selected clips
are light because delta\textless=\allowbreak{}60\textless73.

\hypertarget{a-direct-physical-union-proof-of-the-beta-cutoff}{%
\subsection{A direct physical-union proof of the beta
cutoff}\label{a-direct-physical-union-proof-of-the-beta-cutoff}}

The exact physical t side has six retained native singletons
H\_\allowbreak{}t,\allowbreak{}U\_\allowbreak{}t,\allowbreak{}V\_\allowbreak{}t,\allowbreak{}C6\_\allowbreak{}t,\allowbreak{}C7\_\allowbreak{}t,\allowbreak{}alpha\_\allowbreak{}t and its full leading. Grouping them
gives actual deficiency at most epsilon,\allowbreak{} so

\texttt{S\_\allowbreak{}\allowbreak{}t >=\allowbreak{} K+\allowbreak{}6*77-epsilon =\allowbreak{} K+\allowbreak{}462-epsilon}.

From the unchanged mass identity S\_\allowbreak{}s+\allowbreak{}S\_\allowbreak{}t=\allowbreak{}2K+\allowbreak{}1085-delta,\allowbreak{}

\texttt{S\_\allowbreak{}\allowbreak{}s <=\allowbreak{} K+\allowbreak{}623-e}.

Write z=\allowbreak{}C5\_\allowbreak{}s,\allowbreak{}q=\allowbreak{}alpha\_\allowbreak{}t and beta=\allowbreak{}beta\_\allowbreak{}s. Suppose
beta+\allowbreak{}s3\textgreater=\allowbreak{}d\_\allowbreak{}s-X\_\allowbreak{}s. The literal identity d\_\allowbreak{}s=\allowbreak{}2(a+\allowbreak{}b)+\allowbreak{}G+\allowbreak{}z+\allowbreak{}q
then gives

\texttt{beta>=\allowbreak{}M\_\allowbreak{}\allowbreak{}s+\allowbreak{}2M\_\allowbreak{}\allowbreak{}t+\allowbreak{}2X\_\allowbreak{}\allowbreak{}t+\allowbreak{}g\_\allowbreak{}\allowbreak{}t+\allowbreak{}z+\allowbreak{}q},\allowbreak{}

\texttt{alpha\_\allowbreak{}\allowbreak{}s=\allowbreak{}d\_\allowbreak{}\allowbreak{}s+\allowbreak{}beta>=\allowbreak{}3M\_\allowbreak{}\allowbreak{}s+\allowbreak{}4M\_\allowbreak{}\allowbreak{}t+\allowbreak{}2X\_\allowbreak{}\allowbreak{}s+\allowbreak{}4X\_\allowbreak{}\allowbreak{}t+\allowbreak{}g\_\allowbreak{}\allowbreak{}s+\allowbreak{}2g\_\allowbreak{}\allowbreak{}t+\allowbreak{}2z+\allowbreak{}2q}.

The two lower endpoints of z,\allowbreak{}q have summed nominal constant

\texttt{(73-r\_\allowbreak{}\allowbreak{}s)+\allowbreak{}(73-r\_\allowbreak{}\allowbreak{}t)=\allowbreak{}78}.

Using their weight2,\allowbreak{} the upper alpha\_\allowbreak{}s endpoint,\allowbreak{} the physical cap,\allowbreak{} and
K\textless=\allowbreak{}M+\allowbreak{}R+\allowbreak{}21 yields

\texttt{2M\_\allowbreak{}\allowbreak{}s+\allowbreak{}3M\_\allowbreak{}\allowbreak{}t+\allowbreak{}2X\_\allowbreak{}\allowbreak{}s+\allowbreak{}4X\_\allowbreak{}\allowbreak{}t+\allowbreak{}g\_\allowbreak{}\allowbreak{}s+\allowbreak{}2g\_\allowbreak{}\allowbreak{}t <=\allowbreak{} R+\allowbreak{}416+\allowbreak{}r\_\allowbreak{}\allowbreak{}t-e+\allowbreak{}E+\allowbreak{}2xi\_\allowbreak{}\allowbreak{}s+\allowbreak{}3xi\_\allowbreak{}\allowbreak{}t},\allowbreak{}

where \textbf{E=\allowbreak{}2l\_\allowbreak{}z+\allowbreak{}2l\_\allowbreak{}q+\allowbreak{}u\_\allowbreak{}alpha\_\allowbreak{}s}. The constant is
exactly623+\allowbreak{}21-72-2*78=\allowbreak{}416. The endpoints z,\allowbreak{}q,\allowbreak{}alpha\_\allowbreak{}s are distinct and
avoid all CORE8 ordinary pairs.

Let Omega be the combined deficiencies of the four CORE8 pairs. Since
the endpoint weights are at most2 and those endpoint resources are
disjoint from the ordinary pairs,\allowbreak{}

\texttt{E+\allowbreak{}Omega<=\allowbreak{}2epsilon}.

Select the actual four-trim core Q61,\allowbreak{} with overlap f=\allowbreak{}61-(tau-xi)/4.
Independently,\allowbreak{} the retained extreme-star forests give

\texttt{2X\_\allowbreak{}\allowbreak{}s+\allowbreak{}4X\_\allowbreak{}\allowbreak{}t>=\allowbreak{}6f-2epsilon}.

The four ordinary pairs give 2g\_\allowbreak{}s\textgreater=\allowbreak{}2f-Omega\_\allowbreak{}s and
2g\_\allowbreak{}t\textgreater=\allowbreak{}2f-Omega\_\allowbreak{}t; hence

\texttt{g\_\allowbreak{}\allowbreak{}s+\allowbreak{}2g\_\allowbreak{}\allowbreak{}t>=\allowbreak{}3f-Omega}.

The source lower bound is

\texttt{2M\_\allowbreak{}\allowbreak{}s+\allowbreak{}3M\_\allowbreak{}\allowbreak{}t=\allowbreak{}2M+\allowbreak{}M\_\allowbreak{}\allowbreak{}t>=\allowbreak{}2R+\allowbreak{}R\_\allowbreak{}\allowbreak{}t+\allowbreak{}214+\allowbreak{}3A}.

Combine these rows and R\_\allowbreak{}t=\allowbreak{}r\_\allowbreak{}t+\allowbreak{}xi\_\allowbreak{}t to obtain

\texttt{3A+\allowbreak{}9f <=\allowbreak{}134-e+\allowbreak{}4epsilon+\allowbreak{}xi}.

This deliberately permits the ordinary gap groups and stars separate
copies of the same deficit. It makes no claim that they are disjoint.
Substitution gives

\texttt{415+\allowbreak{}3I+\allowbreak{}3tau/\allowbreak{}4+\allowbreak{}5xi/\allowbreak{}4-delta+\allowbreak{}5e<=\allowbreak{}0}.

But delta\textless=\allowbreak{}60-I-tau,\allowbreak{} so

\texttt{355+\allowbreak{}4I+\allowbreak{}7tau/\allowbreak{}4+\allowbreak{}5xi/\allowbreak{}4+\allowbreak{}5e<=\allowbreak{}0}.

Its left side is at least420 at e=\allowbreak{}13. Thus beta+\allowbreak{}s3\textless d\_\allowbreak{}s-X\_\allowbreak{}s,\allowbreak{}
strictly,\allowbreak{} for every position of all selected endpoints. No normality
assumption is used.

\hypertarget{independent-audit-of-the-four-label-tail-pruning}{%
\subsection{Independent audit of the four-label tail
pruning}\label{independent-audit-of-the-four-label-tail-pruning}}

The actual 32-type upper-tail inventory in the original report
explicitly contains U2\_\allowbreak{}s,\allowbreak{}U3\_\allowbreak{}s,\allowbreak{}U4\_\allowbreak{}s,\allowbreak{}V2\_\allowbreak{}s. The six removed beta fixed
types are V3\_\allowbreak{}s,\allowbreak{}...,\allowbreak{}V8\_\allowbreak{}s,\allowbreak{} so these four labels remain distinct
inventory types after the beta cutoff.

The cutoff and the initial 32-type classification leave only the
alpha\_\allowbreak{}s source family and at most26 fixed types in (d\_\allowbreak{}s-X\_\allowbreak{}s,\allowbreak{}S\_\allowbreak{}s{]}.
Its source capacity is77,\allowbreak{} so

\texttt{S\_\allowbreak{}\allowbreak{}s<=\allowbreak{}d\_\allowbreak{}\allowbreak{}s-X\_\allowbreak{}\allowbreak{}s+\allowbreak{}103-Gamma\_\allowbreak{}\allowbreak{}alpha}.

Since beta\_\allowbreak{}s has a core hit,\allowbreak{} beta\_\allowbreak{}s\textgreater=\allowbreak{}1-R\_\allowbreak{}s. Thus
h\_\allowbreak{}s=\allowbreak{}beta\_\allowbreak{}s+\allowbreak{}d\_\allowbreak{}s+\allowbreak{}M\_\allowbreak{}s+\allowbreak{}X\_\allowbreak{}s satisfies

\texttt{h\_\allowbreak{}\allowbreak{}s-S\_\allowbreak{}\allowbreak{}s>=\allowbreak{}e\_\allowbreak{}\allowbreak{}s+\allowbreak{}1+\allowbreak{}2X\_\allowbreak{}\allowbreak{}s-103+\allowbreak{}Gamma\_\allowbreak{}\allowbreak{}alpha}.

The full-core two-edge physical-s stars
give2X\_\allowbreak{}s\textgreater=\allowbreak{}274-2R-epsilon. With e\_\allowbreak{}s\textgreater=\allowbreak{}A+\allowbreak{}16 and
xi\textless=\allowbreak{}tau,\allowbreak{}

\texttt{h\_\allowbreak{}\allowbreak{}s-S\_\allowbreak{}\allowbreak{}s>=\allowbreak{}52+\allowbreak{}I+\allowbreak{}tau-2xi+\allowbreak{}e>=\allowbreak{}-8+\allowbreak{}2I+\allowbreak{}delta+\allowbreak{}e>0}.

Therefore U4\_\allowbreak{}s=\allowbreak{}h\_\allowbreak{}s and U3\_\allowbreak{}s=\allowbreak{}h\_\allowbreak{}s+\allowbreak{}g\_\allowbreak{}s are HIGH. Also

\texttt{V2\_\allowbreak{}\allowbreak{}s=\allowbreak{}beta\_\allowbreak{}\allowbreak{}s+\allowbreak{}s2>=\allowbreak{}h\_\allowbreak{}\allowbreak{}s+\allowbreak{}e\_\allowbreak{}\allowbreak{}s+\allowbreak{}X\_\allowbreak{}\allowbreak{}s+\allowbreak{}g\_\allowbreak{}\allowbreak{}s+\allowbreak{}1>h\_\allowbreak{}\allowbreak{}s},\allowbreak{}

so V2\_\allowbreak{}s is HIGH. Actual beta activity implies U2\_\allowbreak{}s\textgreater H\_\allowbreak{}s,\allowbreak{}
so U2\_\allowbreak{}s is HIGH independently. These four additional HIGH labels are
distinct from the eleven selected LOW labels. The LOW set still touches
nine baseline groups,\allowbreak{} so the ledger becomes HIGH15,\allowbreak{} LOW11,\allowbreak{} touched9 and
raises e13 to e17.

Pruning these four HIGH types leaves at most22 fixed types in the same
tail. Hence

\texttt{S\_\allowbreak{}\allowbreak{}s<=\allowbreak{}d\_\allowbreak{}\allowbreak{}s-X\_\allowbreak{}\allowbreak{}s+\allowbreak{}99-Gamma\_\allowbreak{}\allowbreak{}alpha}.

Generic complete-target endpoints for alpha\_\allowbreak{}s=\allowbreak{}d\_\allowbreak{}s+\allowbreak{}beta\_\allowbreak{}s,\allowbreak{} with
l\_\allowbreak{}alpha+\allowbreak{}u\_\allowbreak{}alpha\textless=\allowbreak{}Gamma\_\allowbreak{}alpha,\allowbreak{} now give

\texttt{X\_\allowbreak{}\allowbreak{}s<=\allowbreak{}22+\allowbreak{}xi+\allowbreak{}l\_\allowbreak{}\allowbreak{}beta-l\_\allowbreak{}\allowbreak{}alpha}.

Choose the actual one-trim core Q67,\allowbreak{} with overlap f1=\allowbreak{}67-tau+\allowbreak{}xi. The two
physical-s stars are disjoint from the native beta and alpha resources,\allowbreak{}
so

\texttt{2X\_\allowbreak{}\allowbreak{}s>=\allowbreak{}2f1+\allowbreak{}Gamma\_\allowbreak{}\allowbreak{}beta+\allowbreak{}Gamma\_\allowbreak{}\allowbreak{}alpha-epsilon}.

Combining,\allowbreak{} then using l\_\allowbreak{}beta\textless=\allowbreak{}Gamma\_\allowbreak{}beta and
Gamma\_\allowbreak{}beta\textless=\allowbreak{}epsilon,\allowbreak{} gives

\texttt{2f1<=\allowbreak{}44+\allowbreak{}2xi+\allowbreak{}2epsilon}.

Therefore tau+\allowbreak{}delta\textgreater=\allowbreak{}45+\allowbreak{}e=\allowbreak{}62,\allowbreak{} contradicting A\textless=\allowbreak{}60.
This closes the conditional final family without a normality bootstrap.

Scope: this independently checks the local transport,\allowbreak{} resource-disjoint
loss use,\allowbreak{} constants,\allowbreak{} and listed tail types. It relies on the
structural/source/activity conclusions and the exhaustive 32-type
representation classification.

\hypertarget{final-independent-skim-of-the-other-a1b-transports}{%
\subsection{Final independent skim of the other A1/B
transports}\label{final-independent-skim-of-the-other-a1b-transports}}

The newly written \texttt{../\allowbreak{}zero/\allowbreak{}zero85\_\allowbreak{}\allowbreak{}z5\_\allowbreak{}\allowbreak{}a1\_\allowbreak{}\allowbreak{}b.md}
was independently reconstructed against its literal center identities
and generic source endpoints. Opposite owners,\allowbreak{} the pair1-empty
C5\_\allowbreak{}s/U\_\allowbreak{}s placements,\allowbreak{} H\_\allowbreak{}t placement,\allowbreak{} alpha\_\allowbreak{}t placement,\allowbreak{} and branch B
with retained alpha\_\allowbreak{}s have correct nominal constants,\allowbreak{} source
decompositions,\allowbreak{} and separately valid star/gap budgets.
HIGH-alpha\textquotesingle s displayed scalar bounds are conservative;
retaining its displayed +\allowbreak{}10 source-mass term actually strengthens928
to908,\allowbreak{} so its contradiction remains valid.

Two nonblocking transcription errors were found and corrected by the
source author before packaging:

* Same-owner alpha\_\allowbreak{}s-retained transport leaves \textbf{M\_\allowbreak{}s},\allowbreak{} not M,\allowbreak{}
after K\textless=\allowbreak{}M+\allowbreak{}R+\allowbreak{}21. Its corrected scalar row is
\texttt{T\_\allowbreak{}\allowbreak{}r+\allowbreak{}4r+\allowbreak{}(5/\allowbreak{}2)I+\allowbreak{}tau\_\allowbreak{}\allowbreak{}t+\allowbreak{}tau/\allowbreak{}2+\allowbreak{}4xi\_\allowbreak{}\allowbreak{}t<=\allowbreak{}520-7e/\allowbreak{}2}.
The integer quadratic minimum608 at r28 exceeds492 at e8,\allowbreak{} margin116. *
The branch-independent four-endpoint row has radial terms
\textbf{-xi\_\allowbreak{}s/2+\allowbreak{}7xi\_\allowbreak{}t/2},\allowbreak{} replacing the originally stronger positive
terms. Its correction still closes:
\texttt{tau/\allowbreak{}2-xi\_\allowbreak{}\allowbreak{}s/\allowbreak{}2>=\allowbreak{}tau\_\allowbreak{}\allowbreak{}t/\allowbreak{}2>=\allowbreak{}0},\allowbreak{} so the left side
remains at least Q3\textgreater=\allowbreak{}626 versus upper560.

The updated proof/checker therefore receives an independent
\textbf{PASS},\allowbreak{} conditional on the actual retention grammar and
structural/source prerequisites. Its routed A1 and B alternatives use
only retained native resources and already licensed sole pair
components; the branch-independent card assigns no beta capacity.

\clearpage
\fontsize{9.5}{11.4}\selectfont\setlength{\parskip}{0.15em}
\hypertarget{independent-final-audit-of-the-z4-retention-cards}{%
\section{Independent final audit of the Z4 retention
cards}\label{independent-final-audit-of-the-z4-retention-cards}}

Verdict: \textbf{PASS after the Card S light-row sign correction already
applied by the zero-hole author.} This is an analytical audit of actual
native resources,\allowbreak{} endpoint constants,\allowbreak{} and loss budgets,\allowbreak{} in addition to
the scalar replay. It is conditional on the proved cardinality-85
structural/activity classification and source lemmas.

Common fresh parameters are \textbar D\textbar=\allowbreak{}77,\allowbreak{} r\_\allowbreak{}s+\allowbreak{}r\_\allowbreak{}t=\allowbreak{}68,\allowbreak{}
S=\allowbreak{}2K+\allowbreak{}1162-delta,\allowbreak{} R=\allowbreak{}68+\allowbreak{}xi,\allowbreak{} A=\allowbreak{}I+\allowbreak{}tau+\allowbreak{}delta,\allowbreak{} epsilon=\allowbreak{}delta-e,\allowbreak{}
K\textless=\allowbreak{}M+\allowbreak{}R+\allowbreak{}21,\allowbreak{} L\textgreater=\allowbreak{}A+\allowbreak{}20 and A\textless=\allowbreak{}137. The card
calculations deliberately retain the weaker valid cap143. Generic
endpoints have lower73-r\_\allowbreak{}v and upperS\_\allowbreak{}v+\allowbreak{}r\_\allowbreak{}o-72,\allowbreak{} with their actual
clipping/rank errors. Choose the actual four-trim Q61 with
f=\allowbreak{}61-(tau-xi)/4.

\hypertarget{physical-target-capacities}{%
\subsection{Physical target
capacities}\label{physical-target-capacities}}

Every physical target bound uses a group of its original retained full
leading and ordinary resources. It is a union-deficiency bound and
requires no disjointness of their target images.

\begin{longtable}[]{@{}
  >{\raggedright\arraybackslash}p{(\columnwidth - 6\tabcolsep) * \real{0.2500}}
  >{\raggedright\arraybackslash}p{(\columnwidth - 6\tabcolsep) * \real{0.2500}}
  >{\raggedright\arraybackslash}p{(\columnwidth - 6\tabcolsep) * \real{0.2500}}
  >{\raggedright\arraybackslash}p{(\columnwidth - 6\tabcolsep) * \real{0.2500}}@{}}
\toprule\noalign{}
\begin{minipage}[b]{\linewidth}\raggedright
Card
\end{minipage} & \begin{minipage}[b]{\linewidth}\raggedright
Opposite physical t ordinary native resources
\end{minipage} & \begin{minipage}[b]{\linewidth}\raggedright
Capacity
\end{minipage} & \begin{minipage}[b]{\linewidth}\raggedright
Consequent cap on S\_\allowbreak{}s
\end{minipage} \\
\midrule\noalign{}
\endhead
\bottomrule\noalign{}
\endlastfoot
E-H & H\_\allowbreak{}t,\allowbreak{}U\_\allowbreak{}t,\allowbreak{}V\_\allowbreak{}t,\allowbreak{}C5\_\allowbreak{}t,\allowbreak{}C6\_\allowbreak{}t,\allowbreak{}C7\_\allowbreak{}t,\allowbreak{}alpha\_\allowbreak{}t & K+\allowbreak{}539 & K+\allowbreak{}623-e \\
Owner alpha\_\allowbreak{}s missing & Same seven & K+\allowbreak{}539 & K+\allowbreak{}623-e \\
Opposite alpha\_\allowbreak{}t missing & H\_\allowbreak{}t,\allowbreak{}U\_\allowbreak{}t,\allowbreak{}V\_\allowbreak{}t,\allowbreak{}C5\_\allowbreak{}t,\allowbreak{}C6\_\allowbreak{}t,\allowbreak{}C7\_\allowbreak{}t & K+\allowbreak{}462 &
K+\allowbreak{}700-e \\
E-C5 & H\_\allowbreak{}t,\allowbreak{}U\_\allowbreak{}t,\allowbreak{}V\_\allowbreak{}t,\allowbreak{}C6\_\allowbreak{}t,\allowbreak{}C7\_\allowbreak{}t,\allowbreak{}alpha\_\allowbreak{}t & K+\allowbreak{}462 & K+\allowbreak{}700-e \\
S & H\_\allowbreak{}t,\allowbreak{}U\_\allowbreak{}t,\allowbreak{}V\_\allowbreak{}t,\allowbreak{}C6\_\allowbreak{}t,\allowbreak{}C7\_\allowbreak{}t,\allowbreak{}alpha\_\allowbreak{}t & K+\allowbreak{}462 & K+\allowbreak{}700-e \\
\end{longtable}

In Card S,\allowbreak{} C5\_\allowbreak{}t is inactive by the same-owner identity. In branch A
beta has actual sole owner s; in branch B beta\_\allowbreak{}t is the fourth whole
hole and the second paired capacity belongs to its actual s owner. Thus
no unused paired component is included in the physical t count. All
fixed credit is already in epsilon.

\hypertarget{card-f}{%
\subsection{Card F}\label{card-f}}

The endpoint identity in the first two applications is

\texttt{H\_\allowbreak{}\allowbreak{}s=\allowbreak{}P\_\allowbreak{}\allowbreak{}s+\allowbreak{}2alpha\_\allowbreak{}\allowbreak{}s+\allowbreak{}alpha\_\allowbreak{}\allowbreak{}t+\allowbreak{}C5\_\allowbreak{}\allowbreak{}t+\allowbreak{}4(M+\allowbreak{}X)+\allowbreak{}2G}.

Its exact fresh nominal constant is

\texttt{3(73-r\_\allowbreak{}\allowbreak{}s)+\allowbreak{}2(73-r\_\allowbreak{}\allowbreak{}t)+\allowbreak{}(72-r\_\allowbreak{}\allowbreak{}t)=\allowbreak{}233}.

The third application has the reflected identity with sole native
beta\_\allowbreak{}s and no use of the empty pair1. The endpoint constant is
again233. The six-edge forest has W=\allowbreak{}2X+\allowbreak{}G,\allowbreak{} D\_\allowbreak{}F\textgreater=\allowbreak{}6f-W and
N+\allowbreak{}D\_\allowbreak{}F\textless=\allowbreak{}epsilon. Its weighted clipping E\textless=\allowbreak{}2N,\allowbreak{} so
W*2-E\textgreater=\allowbreak{}12f-2epsilon. Endpoints and retained forest vertices
are distinct.

If the one U/V vertex is empty,\allowbreak{} the actual separated positive/negative
restricted zones (M\textgreater R) license the same charge to retained
full resources,\allowbreak{} including an empty middle V,\allowbreak{} with zero capacity assigned
to that vertex. This is the literal one-empty-vertex argument,\allowbreak{} not
restoration of its capacity.

Adding the opposite source mass T\_\allowbreak{}r+\allowbreak{}tau\_\allowbreak{}res+\allowbreak{}4R\_\allowbreak{}r+\allowbreak{}10 to the endpoint
side and using the global source bound gives

\texttt{T\_\allowbreak{}\allowbreak{}r+\allowbreak{}4r+\allowbreak{}tau\_\allowbreak{}\allowbreak{}res+\allowbreak{}2L-3tau+\allowbreak{}2e<=\allowbreak{}229+\allowbreak{}delta+\allowbreak{}B},\allowbreak{}

after dropping nonnegative radial terms. All-light B\textless=\allowbreak{}2epsilon
yields the stated Q3 row332. A unique heavy endpoint of weight q=\allowbreak{}1,\allowbreak{}2
gives B\textless=\allowbreak{}2epsilon-qc+\allowbreak{}q*tau+\allowbreak{}q*e\_\allowbreak{}v and e\_\allowbreak{}v\textless=\allowbreak{}L; the
stated upper332+\allowbreak{}306q follows exactly. The minima626/649/665 and native
c\textgreater=\allowbreak{}73,\allowbreak{}delta\textgreater=\allowbreak{}8+\allowbreak{}c give all claimed strict
contradictions.

\hypertarget{empty-beta-e-cards}{%
\subsection{Empty-beta E cards}\label{empty-beta-e-cards}}

The following constants were reconstructed directly from the literal
endpoint identities,\allowbreak{} the physical capacities above,\allowbreak{} and
K\textless=\allowbreak{}M+\allowbreak{}R+\allowbreak{}21:

\begin{longtable}[]{@{}
  >{\raggedright\arraybackslash}p{(\columnwidth - 6\tabcolsep) * \real{0.2143}}
  >{\raggedleft\arraybackslash}p{(\columnwidth - 6\tabcolsep) * \real{0.2857}}
  >{\raggedright\arraybackslash}p{(\columnwidth - 6\tabcolsep) * \real{0.2143}}
  >{\raggedleft\arraybackslash}p{(\columnwidth - 6\tabcolsep) * \real{0.2857}}@{}}
\toprule\noalign{}
\begin{minipage}[b]{\linewidth}\raggedright
Card
\end{minipage} & \begin{minipage}[b]{\linewidth}\raggedleft
Fresh nominal endpoint constant
\end{minipage} & \begin{minipage}[b]{\linewidth}\raggedright
Body after source magnitudes
\end{minipage} & \begin{minipage}[b]{\linewidth}\raggedleft
Final master constant
\end{minipage} \\
\midrule\noalign{}
\endhead
\bottomrule\noalign{}
\endlastfoot
E-H & 228-r\_\allowbreak{}t & 3a+\allowbreak{}4b+\allowbreak{}g\_\allowbreak{}s+\allowbreak{}2g\_\allowbreak{}t & 416 \\
Owner alpha\_\allowbreak{}s missing & 225-2r\_\allowbreak{}t & 2a+\allowbreak{}4b+\allowbreak{}2G,\allowbreak{} using raw
alpha\_\allowbreak{}s\textgreater=\allowbreak{}1-a & 419 \\
Opposite alpha\_\allowbreak{}t missing & 161+\allowbreak{}r\_\allowbreak{}t & 4a+\allowbreak{}3b+\allowbreak{}2G,\allowbreak{} using raw
alpha\_\allowbreak{}t\textgreater=\allowbreak{}1-b & 492 \\
E-C5 & 228-r\_\allowbreak{}t & 3a+\allowbreak{}4b+\allowbreak{}2G & 493 \\
\end{longtable}

For example E-H gives S\_\allowbreak{}s\textgreater=\allowbreak{}228-r\_\allowbreak{}t+\allowbreak{}3M\_\allowbreak{}s+\allowbreak{}4M\_\allowbreak{}t+\allowbreak{}W-E-B.
Physical cap\textless=\allowbreak{}M+\allowbreak{}R+\allowbreak{}644-e removes M and leaves2M\_\allowbreak{}s+\allowbreak{}3M\_\allowbreak{}t.
Using2M\_\allowbreak{}s+\allowbreak{}3M\_\allowbreak{}t=\allowbreak{}2R+\allowbreak{}2L+\allowbreak{}r\_\allowbreak{}t+\allowbreak{}xi\_\allowbreak{}t+\allowbreak{}e\_\allowbreak{}t then produces constant416,\allowbreak{}
without any residual source-mass estimate.

The forests and native endpoint losses have precisely these budgets:

\begin{longtable}[]{@{}
  >{\raggedright\arraybackslash}p{(\columnwidth - 4\tabcolsep) * \real{0.3333}}
  >{\raggedright\arraybackslash}p{(\columnwidth - 4\tabcolsep) * \real{0.3333}}
  >{\raggedright\arraybackslash}p{(\columnwidth - 4\tabcolsep) * \real{0.3333}}@{}}
\toprule\noalign{}
\begin{minipage}[b]{\linewidth}\raggedright
Card
\end{minipage} & \begin{minipage}[b]{\linewidth}\raggedright
Actual forest edges
\end{minipage} & \begin{minipage}[b]{\linewidth}\raggedright
Endpoint clipping/forest budget
\end{minipage} \\
\midrule\noalign{}
\endhead
\bottomrule\noalign{}
\endlastfoot
E-H & Six edges plus four stars; stars contain V\_\allowbreak{}s only &
E+\allowbreak{}D6+\allowbreak{}Dstar\textless=\allowbreak{}2epsilon-2N+\allowbreak{}h \\
Owner alpha missing & Eight double-arm edges plus two opposite stars;
all four endpoints avoided & E+\allowbreak{}D8+\allowbreak{}D2\textless=\allowbreak{}2epsilon-N \\
Opposite alpha missing & Six edges plus five edges; all weighted
endpoints avoided & E+\allowbreak{}D6+\allowbreak{}D5\textless=\allowbreak{}2epsilon \\
E-C5 & Seven edges plus four edges; all five endpoints avoided &
E+\allowbreak{}D7+\allowbreak{}D4\textless=\allowbreak{}2epsilon-N \\
\end{longtable}

These are separate valid inequalities with the same original deficit.
Overlapping forest systems are not declared one disjoint group. The E-H
upper endpoint\textquotesingle s overlap with the star system is exactly
the +\allowbreak{}h correction in Jstar=\allowbreak{}E+\allowbreak{}B-2N+\allowbreak{}h. Its light lower terms are
nonpositive because b\textless=\allowbreak{}clip; the upper term is at most its
actual rank error. A heavy lower term contributes at most
xi\_\allowbreak{}v+\allowbreak{}e\_\allowbreak{}v-4-c.

All published E scalar finishes are valid. Two clarifications are
useful:

* For owner-alpha heavy on t,\allowbreak{} the exact cancellation gives RHS-9/2; the
published RHS-1/2 is merely weaker and remains contradictory. * In
opposite-alpha,\allowbreak{} the only weight2 heavy endpoint is alpha\_\allowbreak{}s,\allowbreak{} whose error
source is s. Thus2L+\allowbreak{}e\_\allowbreak{}s-2e\_\allowbreak{}s=\allowbreak{}L+\allowbreak{}e\_\allowbreak{}t\textgreater=\allowbreak{}A+\allowbreak{}24. For weight1,\allowbreak{}
either physical source leaves at least A+\allowbreak{}24. The bound would not hold
for an invented weight2/source-t endpoint,\allowbreak{} but that endpoint is absent
from the literal list.

\hypertarget{cards-o-and-s}{%
\subsection{Cards O and S}\label{cards-o-and-s}}

The external/source reviewer independently rederived Card
O\textquotesingle s full endpoint master1006,\allowbreak{} its native rebates and all
four endpoint-heavy flags. The nominal constant in the V\_\allowbreak{}t identity
is150; the weighted leading rebate retains
E\_\allowbreak{}clip+\allowbreak{}2ell\_\allowbreak{}s\textless=\allowbreak{}2epsilon-(B+\allowbreak{}H). The actual CORE8 forest
vertices and endpoints are retained and disjoint exactly where rebates
are used. The scalar minima Q4\textgreater=\allowbreak{}649
and4tau+\allowbreak{}10f\textgreater=\allowbreak{}1990/3 give the displayed strict margins. No
empty component is used.

For Card S I independently reconstructed nominal constant155 and the
physical-union master

\texttt{2M+\allowbreak{}9f+\allowbreak{}4e+\allowbreak{}2N<=\allowbreak{}545+\allowbreak{}3delta+\allowbreak{}B}.

The leading t and the six/three-edge forests avoid all four endpoints.
Hence E+\allowbreak{}ell\_\allowbreak{}t+\allowbreak{}D6+\allowbreak{}D3\textless=\allowbreak{}3epsilon-2N,\allowbreak{} exactly as needed. The heavy
scalar row is correct.

The all-light row originally printed a wrong sign/coefficient on delta.
Its correct form is

\texttt{(9/\allowbreak{}4)I-(3/\allowbreak{}4)delta+\allowbreak{}4e+\allowbreak{}2N<=\allowbreak{}-577/\allowbreak{}4}.

It still strictly contradicts e11,\allowbreak{}delta\textless=\allowbreak{}143: its left side is
at least-253/4,\allowbreak{} a margin81. The source author has applied this
correction in both the proof notes and aggregate checker.

All exact A-E/A-H/B retention rows dispatched to these cards are
therefore excluded,\allowbreak{} and the separately audited ownerless certificate
completes Z4. This receipt checks local analytical applicability; it
does not treat scalar replay as a configuration classifier.

The added two-beta-hole B pre-dispatch row was also independently
checked. Insert the fresh source masses,\allowbreak{} K-M\textgreater=\allowbreak{}R-132+\allowbreak{}A,\allowbreak{}
M\_\allowbreak{}v\textgreater=\allowbreak{}r\_\allowbreak{}v+\allowbreak{}4,\allowbreak{} and the retained three-arm
row2X\_\allowbreak{}v+\allowbreak{}g\_\allowbreak{}v\textgreater=\allowbreak{}189-3tau/4-delta into its displayed
master1161. With R\textgreater=\allowbreak{}68,\allowbreak{}A\textless=\allowbreak{}143 the exact final row is
\texttt{2T\_\allowbreak{}\allowbreak{}r+\allowbreak{}9r+\allowbreak{}4I+\allowbreak{}7tau/\allowbreak{}4+\allowbreak{}2tau\_\allowbreak{}\allowbreak{}v<=\allowbreak{}1063}. Its global
integer minimum1192 at r26 gives margin129. The owner-side HIGH beta
alternative contradicts2d\_\allowbreak{}v\textless=\allowbreak{}M\_\allowbreak{}o-1 versus actual pair1
activity d\_\allowbreak{}v\textgreater M. Thus this extra B form does not leave an
undispatched case.
\end{document}